% !TeX encoding = UTF-8

\documentclass[fontset=none]{exam-zh}
\usepackage{mathtools}
\usepackage{tcolorbox}
% 431 项目统一文档主题。
% 各根入口应先声明 \documentclass{exam-zh}，再输入本文件。

% 项目统一中文字体：避免依赖某台电脑上的 STSong/STHeiti 等系统字体。
\setCJKmainfont{SourceHanSerifCN-Regular.ttf}[
    Path = fonts/,
    BoldFont = SourceHanSerifCN-Bold.ttf,
    ItalicFont = AR-PL-KaitiM-GB.ttf,
    BoldItalicFont = AR-PL-KaitiM-GB.ttf,
    BoldItalicFeatures = {FakeBold = 2},
]
\setCJKsansfont{SourceHanSansCN-Regular.ttf}[
    Path = fonts/,
    BoldFont = SourceHanSansCN-Bold.ttf,
]
\setCJKmonofont{SourceHanSansCN-Regular.ttf}[
    Path = fonts/,
]

% fontset=none 时 CTeX 不创建这些常用字体族命令，显式补齐。
\setCJKfamilyfont{zhsong}{SourceHanSerifCN-Regular.ttf}[Path = fonts/]
\setCJKfamilyfont{zhhei}{SourceHanSansCN-Bold.ttf}[Path = fonts/]
\setCJKfamilyfont{zhkai}{AR-PL-KaitiM-GB.ttf}[Path = fonts/]
\setCJKfamilyfont{zhfs}{SourceHanSerifCN-Regular.ttf}[Path = fonts/]
\providecommand{\songti}{\CJKfamily{zhsong}}
\providecommand{\heiti}{\CJKfamily{zhhei}}
\providecommand{\kaishu}{\CJKfamily{zhkai}}
\providecommand{\fangsong}{\CJKfamily{zhfs}}


\usepackage{booktabs}
\usepackage{multirow}
\usepackage{threeparttable}
\usepackage{threeparttablex}
\usepackage{tikz}
\usetikzlibrary{calc}
\usetikzlibrary{arrows.meta}
\usetikzlibrary{decorations.pathmorphing}
\usetikzlibrary{decorations.pathreplacing}
\usetikzlibrary{positioning}

\tikzset{
    mindmap_root/.style={rectangle, rounded corners, draw=black, fill=gray!10, thick, align=center, minimum height=0.8cm, font=\bfseries\small},
    mindmap_level1/.style={rectangle, rounded corners, draw=blue!60, fill=blue!5, thin, align=center, minimum height=0.6cm, font=\small},
    mindmap_level2/.style={rectangle, draw=none, fill=none, align=center, font=\footnotesize, text width=3cm},
    arrow/.style={->, >=stealth, thick, gray}
}

\usepackage{tabularx}
\usepackage{xcolor}
\usepackage{makecell}
\usepackage{extarrows}
\usepackage{longtable}
\usepackage{float}
\usepackage{enumitem}
\usepackage{perpage}
\usepackage{scalerel}
\usepackage{stackengine}
\MakePerPage{footnote}
\stackMath
\tcbuselibrary{breakable}
\tcbset{breakable}

% 适用于较长估计变量名的宽帽号。
\newcommand{\reallywidehat}[1]{%
  \savestack{\tmpbox}{\stretchto{%
    \scaleto{%
      \scalerel*[\widthof{\ensuremath{#1}}]{\kern-.6pt\bigwedge\kern-.6pt}{%
        \rule[-\textheight/2]{1ex}{\textheight}%
      }%
    }{\textheight}%
  }{0.5ex}}%
  \stackon[1pt]{#1}{\tmpbox}%
}

% 列表全局设置
\setlist{nosep}
% 列表内部保持正文节奏；整个列表块与上下内容各保留半行间距。
\setlist[itemize]{
    leftmargin=*, labelsep=0.5em,
    itemsep=0pt, parsep=0pt,
    topsep=0.5\baselineskip, partopsep=0pt,
    after={}
}
\setlist[itemize,1]{label=$\bullet$}
\setlist[itemize,2]{label=\small$\bullet$}
\setlist[itemize,3]{label=\scriptsize$\bullet$}
\setlist[enumerate]{
    leftmargin=*, labelsep=0.5em,
    itemsep=0pt, parsep=0pt,
    topsep=0.5\baselineskip, partopsep=0pt,
    after={}
}
\setlist[description]{leftmargin=\parindent, labelindent=0pt}

% 章节公共样式。需要每节另起一页的资料在入口中单独覆盖 section/break。
\IfFileExists{ctex.sty}{
    \ctexset{
        section = {
            format = \Large\bfseries\raggedright,
            beforeskip = 1.5ex plus 0.5ex minus 0.2ex,
            afterskip = 1.5ex plus 0.2ex,
        },
        subsection = {
            format = \large\bfseries\raggedright,
            beforeskip = 1.0ex plus 0.5ex minus 0.2ex,
            afterskip = 1.0ex plus 0.2ex,
        },
        subsubsection = {
            format = \normalsize\bfseries\raggedright,
            beforeskip = 0.5ex plus 0.5ex minus 0.2ex,
            afterskip = 0.5ex plus 0.2ex,
        },
        paragraph = {
            runin = false,
            number = \arabic{paragraph}.,
            format = \normalsize\bfseries\raggedright,
            indent = 2em,
            beforeskip = 0.5ex plus 0.3ex minus 0.1ex,
            afterskip = 0.3ex plus 0.2ex,
        }
    }
}{}

% 按 chapter 计数器自动生成“第 X 讲”。
% exam-zh 在 \chapter 完成时推进 chapter，因此标题使用当前值加一。
\ExplSyntaxOn
\NewDocumentCommand{\lecture}{m}
  {
    \exp_args:Nx \chapter
      {第\zhnumber{\int_eval:n {\value{chapter} + 1}}讲：\exp_not:n {#1}}
  }
\ExplSyntaxOff

% 431 资料通用的 exam-zh 设置；可选参数为居中页眉，必选参数为右侧页眉。
\newcommand{\SetupFourThreeOneDocument}[2][]{
    \examsetup{
        page={
            size = a4paper,
            show-head = true,
            head-content = {
                \fancyhf{}
                \fancyhead[C]{#1}
                \fancyhead[EL]{\leftmark}
                \fancyhead[OR]{#2}
            },
            foot-content = {\hfill 第;页 / 共;页 \hfill}
        },
        style/student-version-cleanaux = true,
        fillin/type = paren,
        question = {
            show-points = auto,
            show-answer = true,
        },
        problem = {
            points-separate-par = false,
        },
        paren = {
            text-color = blue,
            show-paren = true,
            show-answer = true,
        },
        solution = {
            show-solution = show-stay,
            label-indentation = true,
            text-color = blue,
            blank-type = hide,
        },
    }

    \ExamPrintAnswerSet{
        question/show-answer = false,
        paren/show-answer = false,
        solution/show-solution = hide,
        paren/text-color = black,
        solution/text-color = black,
    }
}

\renewcommand{\chaptermark}[1]{\markboth{#1}{}}
\renewcommand{\headrulewidth}{1pt}

\newcommand{\true}{$\checkmark$}
\newcommand{\false}{$\times$}

\usepackage{hyperref}
\usepackage{tocloft}

\renewcommand{\contentsname}{\hspace*{\fill}目录\hspace*{\fill}}
\setlength{\cftbeforetoctitleskip}{2em}
\setlength{\cftaftertoctitleskip}{2em}
\setlength{\cftchapindent}{0em}
\setlength{\cftsecindent}{2.0em}
\setlength{\cftsecnumwidth}{2em}
\setlength{\cftsubsecindent}{4.0em}
\setlength{\cftsubsecnumwidth}{2.0em}
\setlength{\cftsubsubsecindent}{6.0em}
\setlength{\cftsubsubsecnumwidth}{2.0em}

\setlength{\intextsep}{1em}

\hypersetup{
    colorlinks = true,
    linkcolor = blue,
    urlcolor = blue,
    citecolor = blue,
    pdfborder = {0 0 0},
}

\newcommand{\kref}[1]{\ref{#1} \textsuperscript{(p.\pageref{#1})}}

\usepackage[group-separator={,}]{siunitx}

\tcbuselibrary{breakable}
\tcbset{breakable}
\usetikzlibrary{intersections}

% 完整版：latexmk 默认，或 -usepretex -pretex='\\def\\WJHandoutKnowledge{}'
% 仅习题：-usepretex -pretex='\\def\\WJExercisesOnly{}' -jobname=金融经济学_王江_exercises_only
\ifdefined\WJExercisesOnly
\else
  \ifdefined\WJHandoutKnowledge
  \else
    \def\WJHandoutKnowledge{}
  \fi
\fi

\ifdefined\WJExercisesOnly
  \SetupFourThreeOneDocument{金融经济学·王江·习题}
\else
  \SetupFourThreeOneDocument{金融经济学·王江·讲义及习题}
\fi

\title{金融经济学}
\ifdefined\WJExercisesOnly
  \subject{王江《金融经济学》课后习题}
\else
  \subject{王江《金融经济学》讲义及习题解答}
\fi
\author{Bilibili@园丁小区詹姆斯}

% 原书题号含星号（如 2.2*），不用 exam-zh 自动序号。
\examsetup{
    problem = {
        label = {},
        label-align = left,
        show-points = false,
    }
}

\newcommand{\problemno}[1]{\textbf{#1}\quad}

% 王江教材式事件树：0 期节点水平接到分叉，竖线只连接各分枝、不出头。
\makeatletter
\newif\ifWJ@tree@stem
\pgfkeys{
    /WJ tree/.is family,
    /WJ tree/left/.initial={},
    /WJ tree/stem/.code=\WJ@tree@stemtrue,
}
\NewDocumentCommand{\WJEventTree}{O{} m}{%
    \begingroup
    \WJ@tree@stemfalse
    \pgfqkeys{/WJ tree}{#1}%
    \begin{tikzpicture}[x=1em,y=1em,baseline=(current bounding box.center)]
        \def\WJ@n{0}%
        \foreach \WJ@lab [count=\WJ@i] in {#2} {\xdef\WJ@n{\WJ@i}}%
        \pgfmathsetmacro{\WJ@dy}{1.35}%
        \pgfmathsetmacro{\WJ@ymax}{\WJ@dy*(\WJ@n-1)/2}%
        \pgfmathsetmacro{\WJ@tk}{0.7}%
        \draw[line width=0.4pt,line cap=butt,line join=miter]
            (0,\WJ@ymax) -- (0,-\WJ@ymax);
        \foreach \WJ@lab [count=\WJ@i] in {#2} {%
            \pgfmathsetmacro{\WJ@y}{\WJ@ymax-\WJ@dy*(\WJ@i-1)}%
            \draw[line width=0.4pt,line cap=butt,line join=miter]
                (0,\WJ@y) -- (\WJ@tk,\WJ@y);
            \node[anchor=base west,inner sep=0.28em,yshift=0.05em]
                at (\WJ@tk,\WJ@y) {$\WJ@lab$};
        }%
        \pgfkeysgetvalue{/WJ tree/left}\WJ@left
        \ifx\WJ@left\empty\else
            \WJ@tree@stemtrue
            \node[anchor=east,inner sep=0.22em] at (-0.85,0) {\WJ@left};
        \fi
        \ifWJ@tree@stem
            \draw[line width=0.4pt,line cap=butt]
                (-0.85,0) -- (0,0);
        \fi
    \end{tikzpicture}%
    \endgroup
}
\makeatother

\begin{document}

\FourThreeOneFrontMatter

\ifdefined\FourThreeOnePrintInterior
\else
\MakeFourThreeOneCover
    {431 FINANCE \quad/\quad FINANCIAL ECONOMICS}
    {金融经济学}
    {\ifdefined\WJExercisesOnly 王江《金融经济学》课后习题\else 王江《金融经济学》讲义及习题解答\fi}
    {从偏好与状态价格出发，完成资产定价的推演}
    {Bilibili@园丁小区詹姆斯}
    {PREFERENCE \quad EQUILIBRIUM \quad PRICING}

\null
\thispagestyle{empty}
\fi

\MakeFourThreeOneContents[2]

\cleardoublepage

% 完整讲义包含引论；原书解答从第 2 章开始，仅习题版保持原编号。
\ifdefined\WJHandoutKnowledge
  \setcounter{chapter}{0}
  % !TeX root = ../../金融经济学_王江_讲义及习题.tex
% 知识点：references/金融经济学/金融经济学-王江_261013.pdf，印刷页 1–5（PDF 第 16–20 页）
% 补充：references/金融经济学/金融经济学_第1章 引论.pptx，第 7–46 页
% 本章无原书习题；PPT 的教学组织、考核安排与联系方式不纳入正文。

\lecture{引论}

本章先说明金融经济学研究什么，再区分均衡定价与无套利定价，最后用一个两状态期权例子展示如何从明确假设推导可检验的价格约束。教材的引论给出全书路线；PPT 进一步讨论学科发展、数学公理化方法与模型的适用边界。

\begin{tcolorbox}[knowledge,title=1.1 金融经济学的研究对象]
\textbf{金融学研究不确定环境下资源的跨期配置。}今天与未来的资源不同，未来不同状态下的资源也不同；金融交易使参与者能够交换对这些资源的要求权。金融市场既是金融要求权（financial claim）的交易场所，也是形成金融资产价格、引导资源配置的机制。

资源配置有两个相关层面：\textbf{微观层面}，个人、公司等如何利用既有资源满足需求；\textbf{整体层面}，不同参与者之间如何转移资源与分担风险。个人在既定约束下最优，并不自动意味着全社会配置有效，后者还取决于市场结构与交易条件。

\textbf{教材的三个核心问题：}
\begin{enumerate}
 \item 参与者如何作出金融决策，尤其是交易与组合选择？
 \item 这些选择如何共同决定市场行为与金融要求权的价格？
 \item 价格如何影响资源的实际配置，市场能否实现有效配置？
\end{enumerate}
金融经济学用经济学的一般原理分析这些问题，为投资、金融市场与机构、金融工程以及公司财务提供理论基础。金融资产的突出特点是\textbf{面向未来的支付要求权}；其价格取决于未来支付、信息、风险及参与者的风险态度。金融学与经济学并非两套互斥的方法，二者共享选择与均衡分析，金融学尤其重视跨期、不确定性与无套利约束。
\end{tcolorbox}

\begin{tcolorbox}[knowledge,title=1.2 均衡定价与无套利定价]
\textbf{均衡定价：}从经济的基本要素出发，包括资源禀赋、偏好、未来风险和交易机会，在价格给定时求每个人的最优选择，再寻找使市场出清的价格。\textbf{最优化与出清须同时成立}；仅仅所有人各自没有意愿继续调整，并不足以代替出清条件。

\textbf{无套利定价：}给定若干已有证券的价格与支付，约束其他相关证券的价格。如果一个支付能由已有证券组合逐状态复制，则其价格必须等于复制成本；否则买入便宜的一方、卖出昂贵的一方即可套利。\textbf{复制是相对定价的核心技术。}

\noindent\begin{tabularx}{\linewidth}{@{}lXX@{}}
\toprule
 & 均衡定价 & 无套利定价\\
\midrule
给定条件 & 偏好、禀赋、风险、市场结构 & 已有证券价格、支付与可行交易\\
基本要求 & 每人最优且市场出清 & 不存在无风险的免费净收益\\
回答问题 & 基础价格与配置如何共同形成 & 相关支付的价格如何保持一致\\
局限 & 结论依赖经济假设与均衡条件 & 不完全市场中未必确定唯一价格\\
\bottomrule
\end{tabularx}

PPT 将二者分别称作“绝对定价”与“相对定价”。这里“绝对”不是脱离计价单位的数值：均衡价格仍须选定计价物。无套利也不意味着偏好完全不影响价格，而是\textbf{偏好对已有价格的影响已被包含在这些价格中}，复制定价时无须再次估计。

在本书的无摩擦与适当单调性条件下，均衡必无套利；反过来，无套利通常只给出价格一致性的必要约束，并不保证市场已经出清。第四章给出正式定义与基本定理，第五章展开期权定价。
\end{tcolorbox}

\begin{tcolorbox}[knowledge,title=1.3 基本框架与全书路线]
教材采用一个简化的\textbf{静态、两时点证券市场框架}：以个人参与者为原型，研究外生给定支付的证券，先把市场结构视为给定。静态并不是没有未来或没有风险，而是不展开中间日期不断接收信息、调仓交易的过程。

\textbf{暂时抽掉的两个关键因素是信息不对称与交易成本。}现实金融要求权可能依赖行为、难以完全写入合约，交易也可能通过中介而非公开证券市场进行；机构与市场结构还会随技术、激励和制度变化。本书先建立基准，不表示这些因素不重要，更不表示基准结论在存在摩擦时可以直接照搬。

\noindent\begin{tabularx}{\linewidth}{@{}lX@{}}
\toprule
章节 & 依次解决的问题\\
\midrule
第 2–3 章 & 建立经济环境、偏好与市场模型；在 Arrow--Debreu 市场分析选择、均衡与配置效率。\\
第 4–5 章 & 从一般市场结构推出无套利约束，再用于期权定价与市场完全化。\\
第 6–9 章 & 期望效用、风险厌恶与投资组合选择，进一步分析参与者行为。\\
第 10–11 章 & 完全和不完全市场中的均衡及资产定价。\\
第 12–14 章 & 以均值—方差偏好及因子结构获得 CAPM、APT 等更具体模型。\\
第 15 章 & 将框架用于公司财务，包括价值最大化、所有权与管理权分离、MM 定理。\\
\bottomrule
\end{tabularx}

教材结尾再讨论信息不对称、交易成本等被简化因素，附录补充证明和数学工具。\textbf{建模的顺序是先给出一般基准，再增加具体结构}，而不是把每章的公式看成相互独立的经验规则。
\end{tcolorbox}

\begin{tcolorbox}[knowledge,title=1.4 学科发展的主要线索（PPT 补充）]
PPT 用若干历史节点说明工具如何形成：Bernoulli 的不确定性下效用分析（1738）、Bachelier 的价格随机运动研究（1900）、Fisher 的跨期选择与分离思想（1930）、Keynes 的流动性偏好（1936），以及 von Neumann--Morgenstern 的期望效用公理体系（1944）。这些分别回应了\textbf{如何比较风险选择、如何描述价格变化、如何处理跨期配置、如何分析货币利率和如何建立一致的决策标准}。

此后，投资组合理论、一般均衡理论、无套利与期权定价逐步形成现代微观金融学的分析框架。Arrow--Debreu 的一般均衡研究与 Debreu 的《价值理论》强调用集合、偏好和凸性来明确可行选择与价格体系。组合选择、CAPM、期权定价、MM 定理与 APT 则把这些思想用于不同的金融问题。

PPT 概括的三个变化是：分析工具更加定量化，分析基础更加重视个体决策，理论逐步成为设计与评估金融安排的工具。\textbf{历史联系不应替代具体假设：}例如鞅说的是条件期望关系，并不推出“最可能的未来价格等于现在价格”；后者混淆了均值与众数，也忽略了计价单位和所用概率测度。
\end{tcolorbox}

\begin{tcolorbox}[knowledge,title=1.5 数学公理化：假设、推导与经验检验]
公理化方法先定义对象、可行选择和行为假设，再按逻辑推导结论。其作用是把“为什么成立”与“在哪些条件下成立”写清楚，使增加、放松某项假设时，能够识别结论随之发生的变化。

PPT 围绕几种方法论立场展开讨论，以下保留争论的要点，而不把评价性表述当作数学定理：
\begin{itemize}
 \item \textbf{Keynes 式的提醒：}符号复杂不等于解释深刻，推导仍依赖初始假设，不能因形式精确而忽略经济含义。
 \item \textbf{Debreu 强调的严格性：}明确对象和推理，能发现不一致、复用已证结果，并为研究新问题提供基础。
 \item \textbf{Popper 强调的可检验与可证伪：}理论应推出可能受经验反驳的结论；从假设演绎出来不等于已经得到数据支持。
 \item \textbf{Friedman 强调的预测表现：}评价理论要重视其解释与预测能力；但这不能被理解为所有不现实的假设都无须关注。假设决定模型可能在哪些环境下失效。
\end{itemize}
演绎推理回答“若假设成立，结论是否成立”；统计分析、计量和案例证据帮助回答“这些机制在观察对象中是否起作用”。样本中观察到的规律仍有外推与识别问题，不能被当作无条件的逻辑必然。模型、数据图表与案例可互相补充，不能用一种证据完全替代另一种。

\textbf{CAPM 的假设清单示例。}PPT 列出相同单期投资期限、资产可分且可交易、允许卖空、均值—方差选择、对收益均值与协方差的同质预期、同一利率下自由借贷、无税费与价格接受者等条件。列出它们的目的，是把定价结论与其成立环境对应起来；现实应用须检查约束、估计误差和模型偏离，而非只记一条收益公式。

\textbf{实证与规范。}实证分析研究“是什么、为何如此”，规范分析还须引入“应当如何”的价值标准。经济学研究和投资实践都可能使用这两类分析，不能按职业作绝对划分。配置有效与否也不同于分配公平与否，第二章将给出 Pareto 标准。
\end{tcolorbox}

\begin{tcolorbox}[example,title=PPT 推理例子：公共信息与共同知识]
将 PPT 的夫妻推理题写成等价的红标问题：$50$ 人头上各有一个红标，每人看得见别人的标记、看不见自己的。所有人都知道彼此严格按逻辑推理；每天统一报告，能确定自己有红标的人当天报告，所有人都能观察前一天的报告。起初没有公共宣告，后来一位可靠的主持人\textbf{公开宣布“至少一人有红标”}，所有人都知道所有人听到了这一宣告。

\textbf{结论：}从宣告后的第一天计，若实际有 $n$ 个红标，这 $n$ 人将在第 $n$ 天一起报告；因此本题为第 $50$ 天。

\textbf{归纳推导：}只有一个红标时，该人看不到其他红标，由宣告立即确定自己有标，在第一天报告。若有两个，每人看到一个；第一天无人报告，就排除了“只有对方一个”的可能，第二天确定自己也有标。一般地，假设 $n-1$ 个红标会在第 $n-1$ 天报告，实际有 $n$ 个时，每名红标者看到 $n-1$ 个；若自己无标，他们应按上述规则报告。观察到第 $n-1$ 天仍无人报告后，才能排除该可能，在第 $n$ 天报告。

当实际人数很多时，每个人在宣告前早已看到至少一个红标；宣告新增的是\textbf{所有人知道、知道别人知道、并如此递推的共同知识}，从而启动上述归纳。相同的一阶事实不等于相同的高阶知识结构；信息公开、观察规则和理性假设都是结论的一部分。
\end{tcolorbox}

\begin{tcolorbox}[knowledge,title=1.6 抽象时必须保留哪些差异]
PPT 用“外延与内涵”提醒我们：把证券仅描述为支付，会抽掉法律、流动性、交易约束等特征。应精确区分：\textbf{集合具有相同成员意味着集合相等，但并不意味着集合中任意两个成员具有相同性能、可以互换。}

无套利复制要求\textbf{同一日期、每个状态的净支付相同}，同时能按所用价格执行买卖；仅有相同的收益分布或期望收益还不够。例如两个等概率状态下，支付 $[1,0]^{\mathsf T}$ 与 $[0,1]^{\mathsf T}$ 的分布相同，却发生在不同状态，其价格可分别为 $\phi_1$ 与 $\phi_2$，不必相等。

期权与复制组合在特定模型下可以有相同支付，但现实中可能有不同的融资、保证金、税费或对手方安排。此时应扩充模型中的成本与约束，而非直接否定复制逻辑。反之，模型的逻辑自洽也不能替代适用条件的核对。
\end{tcolorbox}

\begin{tcolorbox}[knowledge,title=1.7 两状态期权：三种定价方法的同一逻辑]
沿用 PPT 的符号，股票现价为 $S_0>0$，未来为 $aS_0$ 或 $bS_0$，设 $a>b$。执行价为 $K$ 的欧式看涨期权支付为
\[
 C_a=[aS_0-K]_+,\qquad C_b=[bS_0-K]_+.
\]
其中 $[x]_+=\max\{x,0\}$。设一期无风险总收益率为 $R=1+r_F>0$，交易无摩擦、可卖空且两状态均可能发生。股票与债券无套利要求 $b<R<a$：否则股票与按无风险利率融资的组合会有零成本、非负非零的净支付。

\textbf{A. 复制组合。}持有 $z$ 股股票与一笔当前投入 $y$ 的无风险账户，未来支付为 $zaS_0+Ry$、$zbS_0+Ry$。逐状态匹配期权支付：
\[
 zaS_0+Ry=C_a,\qquad zbS_0+Ry=C_b.
\]
两式相减并代回，得到
\[
 z=\frac{C_a-C_b}{S_0(a-b)},\qquad
 y=\frac{aC_b-bC_a}{R(a-b)},\qquad C_0=zS_0+y.
\]
期权与复制组合未来相同，所以初始成本相同；若报价偏高就卖出期权、买入复制组合，偏低则反向操作。

\textbf{B. 风险对冲。}当 $C_a\ne C_b$ 时，卖空一股股票、买入 $x$ 份期权，取
\[
 x=\frac{S_0(a-b)}{C_a-C_b}=\frac1z,
\]
使两状态支付相同。无风险组合必须按 $R$ 增长，故
\[
 R(-S_0+xC_0)=-aS_0+xC_a=-bS_0+xC_b.
\]
解得与复制法相同的价格。若 $C_a=C_b$，不能用此处含零分母的对冲比，而应直接把期权视为确定支付，价格为 $C_a/R$。

\textbf{C. 风险中性定价。}定义 $q=(R-b)/(a-b)$，则 $0<q<1$，并有 $qa+(1-q)b=R$。将复制价格整理为
\[
 \boxed{C_0=\frac{qC_a+(1-q)C_b}{R}.}
\]
这不是假设投资者真的风险中性，也不是估计真实上涨概率，而是将已有价格约束改写为另一种期望。按无风险账户计价后，价格在该测度下满足鞅关系。

\textbf{与 PPT 的公式对应：}第 43–45 页公式没有写利息，等价于本例取 $R=1$，此时须有 $b<1<a$，$q=(1-b)/(a-b)$，且
\[
 C_0=\frac{C_a-C_b+aC_b-bC_a}{a-b}.
\]
若利率非零，就必须保留贴现因子；不能把未经规范化的价格直接说成等于未来价格的期望。
\end{tcolorbox}

\begin{tcolorbox}[knowledge,title=本章小结]
\begin{enumerate}
 \item 金融经济学研究不确定条件下的金融决策、价格形成与资源配置，核心是跨期配置和风险分担。
 \item 均衡定价从偏好、禀赋与市场结构推出最优化且出清的价格；无套利定价从已有价格约束相关支付的价格。
 \item 数学推导须写清假设，经济解释须与公式一致，经验应用还须检验模型与数据的对应关系。
 \item 共同知识不同于人人分别知道；逐状态相同支付也不同于相同概率分布。
 \item 两状态期权的复制、对冲与风险中性方法给出同一价格；PPT 的无贴现公式是零利率特例。
\end{enumerate}
\end{tcolorbox}

\else
  \setcounter{chapter}{1}
\fi
% !TeX root = ../../金融经济学_王江_讲义及习题.tex
% 知识点：references/金融经济学/金融经济学-王江_261013.pdf，第 2 章
% 补充：references/金融经济学/金融经济学_第2章 基本框架.pptx（第 1–44 页，含 WMF 图片公式）
% 习题：references/金融经济学/史上最全最完整的金融经济学答案(王江版).pdf，页 2–6（印刷页 1–5）

\lecture{基本框架}

\ifdefined\WJHandoutKnowledge

\begin{tcolorbox}[knowledge,title=2.1 经济环境]
本章先描述经济的三个基本要素：环境、参与者和证券市场，再讨论个人选择、市场均衡与配置效率。环境的两个关键维度是\textbf{时间与风险}：0 期交易和消费，1 期自然状态实现，再按该状态兑现支付。不同日期、不同状态的消费是不同的或有商品。
\begin{enumerate}

\item 形式上可以用一棵状态树来表示上面描述的经济环境。用 $\omega$ 标识状态，假设总共有 $\Omega$ 个状态，那么表示所有状态可能性的集合 $\Omega$ 称为\textbf{状态空间}。记状态发生的概率为 $p(\omega)$，概率测度记为 $P=\{p_{\omega}, \omega\in\Omega\}$，其中 $p_\omega\ge0$、$\sum_\omega p_\omega=1$。同一符号 $\Omega$ 在教材中兼指状态集合与状态数，须由上下文区分。

\item 假设经济中只有一种不可储存的商品，且本模型没有参与者可选择的生产技术。因此不能直接把今天的商品储存到明天；证券交易转移的是对各日期、各状态资源的要求权，不改变给定状态下的总资源。
\item 用事件树表示未来状态及其概率：
\begin{center}
\WJEventTree[left={$\Omega$},stem]{
    \omega=1 \text{，概率为 }p_1,
    \omega=2 \text{，概率为 }p_2,
    \vdots,
    \omega=\Omega \text{，概率为 }p_{\Omega}
}
\end{center}
\item \textbf{例 2.1（Lucas 树经济）：}一棵桃树今年产出 $100$，明年好天气产出 $200$、坏天气产出 $50$，两状态等概率，此后不再产出。状态编号与产出数量应分开：$\omega=a,b$ 是标签，$Y_1(a)=200$、$Y_1(b)=50$ 才是支付。
\begin{center}
\WJEventTree[left={$100$},stem]{200\text{，状态 }a\text{，概率 }1/2,50\text{，状态 }b\text{，概率 }1/2}
\end{center}
\end{enumerate}
\end{tcolorbox}


\begin{tcolorbox}[knowledge,title=2.2 经济参与者]
    把经济参与者描述为：为自己切身利益而行动的个体。一个参与者的经济特征包括：
    \begin{enumerate}
        \item 他控制的经济资源（任何有经济价值，可以用于满足经济需求的东西，可以是有形 or 无形的）。参与者的经济资源包括
        \begin{itemize}
            \item 所占有的实物商品/资本品
            \item 信息
            \item 生产技术
        \end{itemize}
        \item 他的经济需求。
    \end{enumerate}

\tcbline

参与者的经济资源：
\begin{enumerate}
    \item 禀赋：即参与者初始占有的资源，是模型开始时归他所有的商品或资本品（不要求来自出生时）。给定两个时期，且第二时期有多个状态，那么他的禀赋可表示为：
    \[
    e_k=\begin{bmatrix}e_{k,0} \\ e_{k,1}\end{bmatrix}
    =\begin{bmatrix}e_{k,0} \\ \begin{bmatrix}e_{k,11} \\ \vdots \\ e_{k,1\Omega}\end{bmatrix}
    \end{bmatrix}
    \quad k=1,\cdots,K
    \]

    在数学上，可以把参与者的禀赋看成是 $1+\Omega$ 维实空间 $\mathrm{R}^{1+\Omega}$ 中的一个元素，且假设参与者的禀赋非负，那么禀赋 $e_k \in \mathrm{R}^{1+\Omega}_{+}$。
\begin{tcolorbox}[warning,title=向量的“大于”]
    对于向量 $a=\begin{bmatrix}a_1 \\ a_2 \\ \vdots \\ a_n\end{bmatrix}$，有以下定义：
    \begin{enumerate}
        \item $a \ge 0$，即 a 的所有分量都非负，即 $a_i \ge 0, i=1,\cdots,n$。
        \item $a > 0$，即 a 的所有分量都非负，且至少有一个分量严格大于 0，即 $\exists i \in \{1,\cdots,n\}$ 使得 $a_i > 0$。
        \item $a \gg 0$，即 a 的所有分量都严格大于零，即 $a_i > 0, i \in \{1,\cdots,n\}$。
    \end{enumerate}
\end{tcolorbox}
    \item 信息：信息具备经济价值。参与者的信息可以分为两种类型：公开信息和私有信息。如果参与者仍然拥有私有信息，信息就是不对称的。信息不对称会导致很多问题，例如道德风险和逆向选择。但在本书中，将假设经济者对于经济，特别是其未来状态发生的可能性，拥有相同的信息。
    \item 生产技术：指可以把商品转化为更多商品，或者把现在的商品转化为未来的商品。在本书中，假设所有参与者不拥有任何生产技术。
\end{enumerate}

\tcbline

参与者的经济需求：
参与者的经济需求基于他的不同时期和未来不同状态下的消费。对于参与者 k，记他在第零期的消费是$c_{k,0}$，第一期的消费是 $c_{k,\omega}$，其中 $\omega \in \Omega$。消费的完整表述是：
\[
c_k \equiv \begin{bmatrix}c_{k,0} \\ c_{k,1}\end{bmatrix}=
\begin{bmatrix}c_{k,0} \\ \begin{bmatrix}c_{k,11} \\ \vdots \\ c_{k,1\omega} \\ \vdots \\ c_{k,1\Omega}\end{bmatrix}
\end{bmatrix} \quad k=1,\cdots,K
\]


\begin{enumerate}
\item 消费集：消费者可能的消费选择叫做消费计划 $c=\begin{bmatrix}
    c_0 \\ c_1
\end{bmatrix}$。所有消费计划的集合叫做消费集 $C$。和禀赋一样，消费计划属于 $\mathrm{R}^{1+\Omega}$ 中的一个元素，且假设消费计划非负，那么消费计划 $c \in \mathrm{R}^{1+\Omega}_{+}$，消费集就是 $C=\mathrm{R}^{1+\Omega}_{+}$。

\begin{tcolorbox}[knowledge]
\textbf{定义 2.1} 设 A 是 $R_n$ 的一个集合，对于 $\forall a,b \in A$ 和 $\alpha \in [0,1]$，如果 $a\alpha+(1-\alpha)b$ 也在A 中，则称集合 A 是凸的（convex）。

\textbf{定义 2.2} 对于 $\mathrm R^n$ 的集合 $A$，如果其中任意\textbf{收敛}序列 $a_i\to a$ 的极限也在 $A$ 中，则称集合 A 是闭的（closed）。

\textbf{假设 1} 消费集 $C=\mathrm{R}^{1+\Omega}_{+}$ 是 $\mathrm{R}^{1+\Omega}$ 中的一个闭凸子集。

\textbf{直接验证：}若 $a,b\ge0$、$\alpha\in[0,1]$，则每个分量 $\alpha a_i+(1-\alpha)b_i\ge0$，故 $C$ 凸。若 $c^{(n)}\in C$ 且 $c^{(n)}\to c$，则 $c_i=\lim_{n\to\infty}c_i^{(n)}\ge0$，故 $c\in C$，所以 $C$ 闭。“闭”不要求集合中每个序列都有极限，只要求已有极限不落到集合外。
\end{tcolorbox}

\item 偏好：偏好是参与者对所有可能性计划的排序这样的排序定义了他的经济需求。

\begin{tcolorbox}[knowledge]
\textbf{定义 2.3} 偏好是C 上的一个二元关系，表示为 $\succeq$，且满足以下三个条件：
\begin{itemize}
    \item 完备性：对于 $\forall a,b \in C$，有 $a\succeq b$ 或 $b\succeq a$。
    \item 传递性：对于 $\forall a,b,c \in C$，如果 $a\succeq b$ 且 $b\succeq c$，则 $a\succeq c$。
    \item 自反性：对于 $\forall a \in C$，有 $a\succeq a$。
\end{itemize}
\end{tcolorbox}

记 $a\succ b$ 表示 $a\succeq b$ 且不是 $b\succeq a$（严格偏好）；$a\sim b$ 表示二者相互弱偏好（无差异）。这是对整段消费计划的排序，不能与向量逐分量比较 $a>b$ 混为一谈。

\item 偏好的基本假设：\textcolor{red}{\textbf{参与者喜多厌少}：如果两个消费计划 a 和 b 之间，在所有时期和状态下，a 都不比 b 提供的消费少，且在有些状态下提供更多的消费，那么 a 优于 b。}
    \begin{tcolorbox}[knowledge]
        \textcolor{red}{\textbf{公理 1 不满足性：}} 如果 $a>b$，那么 $a\succ b$（严格偏好）。

        \textbf{公理 2 连续性：} $\forall c \in C$，集合 $\{a\in C: a\succeq c\}$ 和 $\{b\in C: b\preceq c\}$ 都是闭集。

        \textbf{公理 3 凸性：} $\forall a,b \in C$ 以及 $\alpha \in [0,1]$，如果 $a \succeq b$ ，那么有 $\alpha a+(1-\alpha)b\succeq b$ 。
    \end{tcolorbox}

    \item 效用函数：从消费计划到实数的映射叫做效用函数。
    
    \begin{tcolorbox}[knowledge]
        \textbf{定义 2.4} 对应于偏好关系的效用函数 U 是从 C 到 R 的函数，使得 $\forall a,b \in C$，当且仅当 $a\succeq b$ 时，有 $U(a) \ge U(b)$。
    \end{tcolorbox}

    \begin{tcolorbox}[knowledge]
        \textbf{定理 2.1} 对于一个在闭凸消费集 C 上，有定义 2.3 定义的（完备性、传递性和自反性）、满足公理 2 的偏好（连续性），存在一个连续的效用函数 U 使得 $\forall a,b \in C, a\succeq b$，当且仅当 $U(a) \ge U(b)$。
    \end{tcolorbox}

    \begin{tcolorbox}[example]
        \begin{question}
            \textbf{例 2.3}\quad 在例 2.1 中描述的 Lucas 树经济中，消费空间是 $\mathrm{R}^{3}_{+}$，表示在今年和明年的两个可能状态下的消费；状态 $a$ 对应于好天气，而状态 $b$ 则对应于坏天气。假设经济中的某个参与者有如下的效用函数：
            \[
                U(c_0,c_{1a},c_{1b})=\log c_0+\tfrac12(\log c_{1a}+\log c_{1b})
            \]
            在如下三个消费计划中，他将会选择哪一个？
            \begin{center}
            \begin{tabular}{@{}c@{\quad}c@{\quad}c@{}}
                \textbf{计划 A} & \textbf{计划 B} & \textbf{计划 C} \\[0.35em]
                \WJEventTree[left={$1$},stem]{1,1} &
                \WJEventTree[left={$\tfrac12$},stem]{\tfrac32,\tfrac32} &
                \WJEventTree[left={$1$},stem]{\tfrac32,\tfrac12}
            \end{tabular}
            \end{center}
        \end{question}

        \begin{solution}
            对于 A、B 和 C，效用函数给出的相应的效用值分别为 $U_A=0$、$U_B=\log(3/4)\approx-0.288$、$U_C=\tfrac12\log(3/4)\approx-0.144$。显然，参与者对三个消费计划的排序是 $A>C>B$，因此他会选择 A。
        \end{solution}

        \begin{question}
            \textbf{例 2.4}\quad 让我们考虑同一经济中的另外一个参与者，他的效用函数是
            \[
                V(c_0,c_{1a},c_{1b})=c_0\sqrt{c_{1a}c_{1b}}
            \]
            他对上面三个消费计划的排序是什么？
        \end{question}
        \begin{solution}
            \[
                \begin{aligned}
                    V_A &= 1\cdot\sqrt{1\cdot 1}=1,\\
                    V_B &= \tfrac12\cdot\sqrt{\tfrac32\cdot\tfrac32}=\tfrac34,\\
                    V_C &= 1\cdot\sqrt{\tfrac32\cdot\tfrac12}=\tfrac{\sqrt3}{2}.
                \end{aligned}
            \]
            因此 $V_A>V_C>V_B$，参与者对三个消费计划的排序仍是 $A>C>B$。
        \end{solution}
    \end{tcolorbox}
\end{enumerate}

上面的例子说明了如下的结论：从定理 2.1 中得到的效用函数只是\textbf{序数的}（ordinal），即它只给出了消费计划的排序。效用函数的任意\textbf{严格递增}变换不会改变这种排序：若 $V=f(U)$、$f$ 严格递增，则 $V(a)\ge V(b)\iff U(a)\ge U(b)$。本例在正消费域内恰有
\[
 e^{U(c)}=e^{\log c_0+\frac12\log c_{1a}+\frac12\log c_{1b}}
 =c_0\sqrt{c_{1a}c_{1b}}=V(c),
\]
所以两种公式给出相同偏好。效用差额或倍数没有直接福利数量含义，也不能直接进行人际比较。对数效用只在正消费域定义，不能据此宣称它在包含零消费的整个 $\mathrm R^{1+\Omega}_+$ 上连续。

利用效用函数作为偏好的一种表述形式，我们可以把偏好的三个基本假设重新表示如下：
\begin{enumerate}
    \item \textbf{不满足性：}如果 $a>b$，那么 $U(a)>U(b)$；相应有 $a\ge b\Rightarrow U(a)\ge U(b)$。
    \item \textbf{连续性：}$\lim_{a_n \to a}U(a_n)=U(a)$。
    \item \textbf{凸性：}对任意 $a,b$ 和 $\alpha\in[0,1]$，
    \[
      U(\alpha a+(1-\alpha)b)\ge\min\{U(a),U(b)\}.
    \]
    这是\textbf{拟凹性}，即上等值集为凸集，并非 $U$ 作为函数的凸性。若选取凹效用表示，则更强的 Jensen 不等式 $U(\alpha a+(1-\alpha)b)\ge\alpha U(a)+(1-\alpha)U(b)$ 足以推出偏好凸。严格递增变换保留排序与拟凹性，却未必保留函数的凹性；不能把任一序数表示的曲率直接当作风险厌恶。
\end{enumerate}

本章习题里反复出现的是下面几种效用函数$u(x)$的形式。两状态等概率时，跨期效用写成
\[
    U(c)=u(c_0)+\tfrac12\bigl(u(c_{1a})+u(c_{1b})\bigr).
\]

表中只记函数形状，以及它如何对应不满足性和凸性。绝对风险厌恶、相对风险厌恶以及 CARA、CRRA 的比较留到第 6、7 章。

{\small
\renewcommand{\arraystretch}{1.25}
\noindent\begin{tabularx}{\linewidth}{@{}llX@{}}
\toprule
名称 & 效用函数 $u(x)$ & 要点 \\
\midrule
对数 & $\log x$ & $x>0$。严格递增、严格凹。幂效用去掉与消费无关的常数后，在 $\gamma\to1$ 时趋于对数。 \\
幂 & $\dfrac{x^{1-\gamma}}{1-\gamma}$ & $\gamma>0$、$\gamma\neq 1$，$x>0$。严格递增、严格凹。平方根效用即 $\gamma=\tfrac12$。常相对风险厌恶（CRRA）。 \\
负指数 & $-e^{-ax}$ & $a>0$。严格递增、严格凹。常绝对风险厌恶（CARA）。 \\
二次 & $x-\tfrac12 a x^{2}$ & $a=0$ 时线性，处处递增；$a>0$ 时仅当 $x<1/a$ 边际效用为正，不满足性只在此范围内成立。 \\
\bottomrule
\end{tabularx}}

下图只示意形状，纵轴各自平移和缩放。幂效用取 $\gamma=2$，负指数与二次取 $a=1$。二次图中的虚线是边际效用降为零的位置 $x=1/a$。
\begin{center}
\setlength{\tabcolsep}{2pt}
\footnotesize
\begin{tabular}{@{}cccc@{}}
\begin{tikzpicture}[baseline=(current bounding box.center), x=1cm, y=1cm]
    \path (-0.28,-0.42) rectangle (2.62,1.72);
    \draw[-{Stealth[length=1.5mm]}, thin] (0.22,0) -- (2.32,0) node[below left, inner sep=1pt] {$x$};
    \draw[-{Stealth[length=1.5mm]}, thin] (0.22,0) -- (0.22,1.48) node[above left, inner sep=1pt] {$u$};
    \draw[blue!65!black, thick, domain=0.32:2.15, samples=50, smooth] plot (\x,{0.95+0.62*ln(\x)});
\end{tikzpicture}
&
\begin{tikzpicture}[baseline=(current bounding box.center), x=1cm, y=1cm]
    \path (-0.28,-0.42) rectangle (2.62,1.72);
    \draw[-{Stealth[length=1.5mm]}, thin] (0.22,0) -- (2.32,0) node[below left, inner sep=1pt] {$x$};
    \draw[-{Stealth[length=1.5mm]}, thin] (0.22,0) -- (0.22,1.48) node[above left, inner sep=1pt] {$u$};
    \draw[blue!65!black, thick, domain=0.34:2.15, samples=50, smooth] plot (\x,{1.42-0.40/\x});
\end{tikzpicture}
&
\begin{tikzpicture}[baseline=(current bounding box.center), x=1cm, y=1cm]
    \path (-0.28,-0.42) rectangle (2.62,1.72);
    \draw[-{Stealth[length=1.5mm]}, thin] (0,0) -- (2.32,0) node[below left, inner sep=1pt] {$x$};
    \draw[-{Stealth[length=1.5mm]}, thin] (0,0) -- (0,1.48) node[above left, inner sep=1pt] {$u$};
    \draw[blue!65!black, thick, domain=0:2.15, samples=40, smooth] plot (\x,{1.32-1.12*exp(-\x)});
\end{tikzpicture}
&
\begin{tikzpicture}[baseline=(current bounding box.center), x=1cm, y=1cm]
    \path (-0.28,-0.42) rectangle (2.62,1.72);
    \draw[-{Stealth[length=1.5mm]}, thin] (0,0) -- (2.32,0) node[below left, inner sep=1pt] {$x$};
    \draw[-{Stealth[length=1.5mm]}, thin] (0,0) -- (0,1.48) node[above left, inner sep=1pt] {$u$};
    \draw[blue!65!black, thick, domain=0:2, samples=40, smooth] plot ({1.05*\x},{2.4*(\x-0.5*\x*\x)});
    \draw[dashed, black!45] (1.05,0) -- (1.05,1.2);
    \node[below, inner sep=1pt] at (1.05,0) {$1/a$};
\end{tikzpicture}
\\[0.15em]
对数 $\log x$ & 幂 $\dfrac{x^{1-\gamma}}{1-\gamma}$ & 负指数 $-e^{-ax}$ & 二次 $x-\tfrac12 ax^{2}$
\end{tabular}
\end{center}

\end{tcolorbox}



\begin{tcolorbox}[knowledge,title=2.3 证券市场]
资源的配置是通过金融市场中的交易来完成的，市场允许参与者在不同时期和状态上配置资源。

\tcbline

\begin{enumerate}
    \item \textbf{证券及其支付：}在我们所要构建的简单框架中，金融市场由一组\textbf{证券}（security）构成。一只证券即是一份\textbf{金融要求权}（financial claim），它在 1 期给其所有者带来\textbf{支付}（payoff）；支付的数量一般依赖于当时的经济状态。令 $x_\omega$ 为证券在状态 $\omega$ 时的支付，其中 $\omega\in\Omega$，则一只证券由它在各个可能状态下的支付 $\{x_\omega,\,\omega=1,\dots,\Omega\}$ 来定义。形式上可用一棵\textbf{支付树}表示：
    \begin{center}
    \WJEventTree[left={$0$},stem]{x_1,x_2,\vdots,x_\Omega}
    \end{center}
    定义 $\mathrm{R}^{\Omega}$ 为\textbf{支付空间}（payoff space），即所有可能支付的集合；其中的向量
    \[
        x=\begin{bmatrix}x_1 \\ x_2 \\ \vdots \\ x_\Omega\end{bmatrix}
    \]
    就定义了一只证券。在定义证券时，我们未考虑任何在 0 期的可能支付；到期以前（在我们的框架中就是 0 期）证券产生的支付通常叫做\textbf{股利}（dividend）。
    
    在 Lucas 树例子中，把代表树的所有权的证券叫做\textbf{股票}（stock），它代表对实物资产（树）所产生的支付的要求权；这样的证券也叫做\textbf{金融资产}（financial asset）或\textbf{资产}，其支付来自\textbf{标的}（underlying）实物资产。

    \textbf{无风险债券}（riskless bond）是这样一种证券，它在所有状态下的支付都相等；若各状态下的支付均为 $1$，则其支付向量为 $\begin{bmatrix}1 \\ \vdots \\ 1\end{bmatrix}$。
    
    \textbf{风险证券}（如股票）的支付则随状态变化，例如支付向量为 $\begin{bmatrix}x_1 \\ \vdots \\ x_\Omega\end{bmatrix}$。
    
    金融证券具有三个基本特征：
    \begin{enumerate}
        \item 支付只依赖于未来的经济状况；
        \item 支付是外生给定的，不受参与者行为的影响；
        \item 状态集合及其概率分布是所有参与者的共同知识。
    \end{enumerate}

    \item \textbf{市场结构：}假设市场上有 $N$ 只可交易证券，用 $n=1,\dots,N$ 来标识。记第 $n$ 只证券的支付向量为
    \[
        X_n=\begin{bmatrix}x_{1,n} \\ x_{2,n} \\ \vdots \\ x_{\Omega,n}\end{bmatrix},
    \]
    把所有证券的支付向量并排，就得到 $\Omega\times N$ 的\textbf{市场结构矩阵}（也称支付矩阵）
    \[
        X\equiv[X_{\cdot,1},\,\cdots\,X_{\cdot,n},\,\cdots\,X_{\cdot,N}]=
        \begin{bmatrix}
            x_{1,1}   & x_{1,2}   & \cdots & x_{1,N}   \\
            \vdots    & \vdots    & \ddots & \vdots    \\
            x_{\omega,1}   & x_{\omega,2}   & \cdots & x_{\omega,N}   \\
            \vdots    & \vdots    & \ddots & \vdots    \\
            x_{\Omega,1} & x_{\Omega,2} & \cdots & x_{\Omega,N}
        \end{bmatrix}\equiv
        \begin{bmatrix}
            X_{1,\cdot} \\ \vdots \\ X_{\omega,\cdot} \\ \vdots \\ X_{\Omega,\cdot}
        \end{bmatrix}
    \]
    矩阵 $X$ 的每一列是一只证券在各状态下的支付（列向量），每一行是一个状态下各证券的支付（行向量）。通常有 $\Omega>N$。

    \item \textbf{证券组合：}对每只证券的持有数量构成\textbf{证券组合}（security portfolio），简称\textbf{组合}（portfolio）。若 $\theta_n$ 为证券 $n$ 的持有数量，则\textbf{组合}由向量$\theta=\begin{bmatrix}\theta_1 \\ \theta_2 \\ \vdots \\ \theta_N\end{bmatrix}$ 定义；该组合在 1 期的支付向量为 $X\theta$。对每个参与者 $k$，记其交易前的\textbf{初始证券组合}为$\bar{\theta}_k=\begin{bmatrix}\bar{\theta}_{k,1} \\ \bar{\theta}_{k,2} \\ \vdots \\ \bar{\theta}_{k,N}\end{bmatrix}$，所有参与者初始证券组合之和称为\textbf{市场组合}（market portfolio）$\theta_M$，即
    \[
        \theta_M=\sum_{k=1}^{K}\bar{\theta}_k,
    \]
    它表示市场上所有可交易证券的总供给。

    \item \textbf{市场化：}若支付 $x\in\mathrm{R}^{\Omega}$ 能由某交易组合复制出来，即存在 $\theta\in\mathrm{R}^{N}$ 使得 $X\theta=x$（有解），则称 $x$ 是\textbf{市场化的}（marketed）。所有市场化支付的集合为
    \[
        M=\{X\theta:\,\theta\in\mathrm{R}^{N}\}.
    \]
    一般地，$M$ 是支付空间 $\mathrm{R}^{\Omega}$ 的一个子空间，其维数等于 $\mathrm{rank}(X)\le\min\{N,\Omega\}$，只有证券支付列线性无关时才等于 $N$。允许任意实数组合保证了这个集合是线性子空间；若有卖空或持仓限制，受限可达集合未必是子空间。

\begin{tcolorbox}[example,title=PPT 补充：股票与债券能生成哪些支付]
Lucas 树股票支付 $[200,50]^{\mathsf T}$，单位无风险债券支付 $[1,1]^{\mathsf T}$，所以
\[
 X=\begin{bmatrix}200&1\\50&1\end{bmatrix}.
\]
给定目标支付 $y=[y_a,y_b]^{\mathsf T}$，解 $X\theta=y$ 得
\[
 \theta_{\rm stock}=\frac{y_a-y_b}{150},\qquad
 \theta_{\rm bond}=\frac{4y_b-y_a}{3}.
\]
例如支付 $[1,0]^{\mathsf T}$ 可由 $1/150$ 份股票与 $-1/3$ 份债券复制。股票与债券共同生成全部两状态支付；仅有股票时，只能生成其支付列的倍数。复制成本由价格 $S^{\mathsf T}\theta$ 决定，不能只看发生概率。
\end{tcolorbox}

    \item \textbf{交易过程：}交易可以通过\textbf{集中化交易}（centralized exchange）或\textbf{柜台交易}（over-the-counter，OTC）完成；具体的交易安排称为\textbf{市场微观结构}（market micro-structure）。记 $S=[S_1,\ldots,S_N]^{\mathsf T}$ 为证券价格，$h_k(S)$ 为参与者的\textbf{交易后总持仓}，$\bar\theta_k$ 为初始持仓。总持仓出清条件为
\[
 \sum_k h_k(S)=\sum_k\bar\theta_k=\theta_M.
\]
若改用净交易量 $\theta_k(S)=h_k(S)-\bar\theta_k$，则完全等价地写成 $\sum_k\theta_k(S)=0$。不能把净买入需求之和与总供给直接相比。PPT 的开盘报价例子说明的是价格下的\textbf{意愿需求}：若老王拟净卖出 $200$ 股、老张拟净买入 $300$ 股，尚有 $100$ 股超额净需求，须由其他参与者或调整后的意愿补足；实际成交的买卖量总是相等。此处只是出清机制的教学示意，不是对交易所撮合程序或某只股票当前报价的描述。

    交易成本、头寸限制（如卖空约束）、个体交易对价格的影响以及税收等，都会妨碍交易，统称为\textbf{市场摩擦}。\textbf{无摩擦市场}满足以下假设：
    \begin{enumerate}
        \item 参与者可无成本地进入或退出证券市场；
        \item 没有交易成本；
        \item 对参与者的证券持有没有限制（允许卖空）；
        \item 单个参与者的交易不影响市场价格（价格接受者）；
        \item 没有税收。
    \end{enumerate}
    税费、买卖价差与融资约束会使可执行交易的净成本偏离简单标价；此处用无摩擦假设建立基准，不把 PPT 中列举的某一时期税率当作模型常数。摩擦可能解释机构与交易安排的作用，但不保证存在任何人都能无风险获利的机会。
\end{enumerate}
\end{tcolorbox}

\begin{tcolorbox}[knowledge,title=2.4 基本经济模型]
现在我们已经对将要分析的经济的基本模型有了完整的描述。在作进一步的讨论和分析以前，先作一个总结。

\begin{tcolorbox}[knowledge]
\textbf{定义 2.5}\quad 一个\textbf{经济}定义如下：
\begin{enumerate}
    \item 有两个时期 $0$ 和 $1$；在 1 期有 $\Omega$ 个可能状态，对这些状态有概率测度 $P$。经济中只有一种不可储存的商品。
    \item 经济中有 $K$ 个参与者（$k=1,\dots,K$）：
    \begin{enumerate}[label=(\alph*)]
        \item 每一参与者对未来状态发生的可能性都有\textbf{相同的信息}，由 $P$ 描述；
        \item 每一参与者具有\textbf{禀赋} $e_k\in\mathrm{R}^{1+\Omega}_{+}$；
        \item 每一参与者有定义于 $C=\mathrm{R}^{1+\Omega}_{+}$ 上、且满足公理 1$\sim$公理 3 的偏好（不满足性、连续性、凸性）。
    \end{enumerate}
    \item 有一个市场结构为 $X$ 的\textbf{无摩擦证券市场}。
\end{enumerate}
\end{tcolorbox}

\begin{tcolorbox}[knowledge]
\textbf{定义 2.6}\quad 在上面定义的经济中，如果所有参与者的 1 期禀赋都可表示为其初始证券组合的支付，即对每个 $k$ 有 
\[
e_{k,1}=X\bar{\theta}_k
\]
则我们称之为\textbf{证券市场经济}（security-market economy）。
\end{tcolorbox}

这要求\textbf{每个人的实际禀赋可市场化}，不要求所有想象得到的支付都可交易，因此与第三章的“完全市场”不同。例如两状态只有支付 $[1,1]^{\mathsf T}$ 的债券，若人人禀赋也是该向量的倍数，它仍是证券市场经济，但并不完全。

PPT 以农业收入与劳动收入说明：如果某人的未来禀赋不在 $X$ 的列空间内，就不能完全表示为现有证券组合。仅凭资产未直接挂牌不能判断它不可复制，仍须检查支付。证券化、农业风险合约等可能扩大可达支付集合，但不会自动消除所有收入风险或保证帕累托效率。

为简化起见，在以后的讨论中有时会局限于证券市场经济的情形。
\end{tcolorbox}

\begin{tcolorbox}[knowledge,title=2.5 市场均衡]
在证券市场中，资源配置通过以下两方面的相互作用来实现：
\begin{enumerate}
    \item 给定证券及其未来支付与当前价格，每个参与者选择最优证券持有量；
    \item 所有参与者的总需求影响证券价格；当需求等于供给时达到\textbf{均衡}（equilibrium），此时各参与者都选择了最优证券需求且市场出清。
\end{enumerate}

\tcbline

\begin{enumerate}
    \item \textbf{参与者各自优化：}记 $\theta=\begin{bmatrix}\theta_1 \\ \vdots \\ \theta_N\end{bmatrix}$ 为\textbf{净购买}的证券组合（相对初始持有量的交易量），$\theta_n>0$ 表示\textbf{净买入}，$\theta_n<0$ 表示\textbf{净卖出}；最终总持仓 $h_{k,n}$ 的正负才分别表示多头和空头。该组合在 0 期的市场价值（成本）为 $S^{\mathsf T}\theta$。对禀赋为 $e_k$ 的参与者，净购买 $\theta_k$ 后消费计划 $c_k=\begin{bmatrix}c_{k,0} \\ c_{k,1}\end{bmatrix}$ 为
    \[
        c_{k,0}=\underbrace{e_{k,0}-S^{\mathsf T}\theta_k}_{\text{第零期禀赋减去购买证券的支出}},\qquad
        c_{k,1}=\underbrace{e_{k,1}+X\theta_k}_{\text{第 1 期禀赋加上证券在第一期的支付}}.
    \]
    其中 $S^{\mathsf T}\theta_k>0$ 表示净支出、减少当前消费，$S^{\mathsf T}\theta_k<0$ 表示净收入。对购买正价格债券这一特例，前者是储蓄、后者是借款；对一般含风险资产的组合，应称交易支出或收入，不能只凭符号笼统叫储蓄或借贷。所有可达消费计划的集合称为\textbf{预算集}
    \[
        B(e_k,\{X,S\})=\Bigl\{c\ge 0\;\Big|\;
        c_{k,0}=e_{k,0}-S^{\mathsf T}\theta_k,\;
        c_{k,1}=e_{k,1}+X\theta_k,\;
        \theta_k\in\mathrm{R}^{N}\Bigr\}.
    \]
    参与者选择 $\theta_k$ 以最大化效用：
    \[
    \begin{aligned}
        &\max_{\theta_k} U_k(c_k) \\
        &\text{s.t.}\quad
        c_{k,0}=e_{k,0}-S^{\mathsf T}\theta_k,\;
        c_{k,1}=e_{k,1}+X\theta_k,\;
        c_{k,0},\,c_{k,1}\ge 0
    \end{aligned}
    \]
    等价地，
    \[
        \max_{c_k\in B(e_k,\{X,S\})} U_k(c_k).
    \]
    该问题的解给出参与者 $k$ 对证券的\textbf{需求函数} $\theta_k(e_k,S)$，它依赖于价格向量 $S$ 与禀赋 $e_k$。

    \item \textbf{市场出清：}为避免对初始持仓重复计数，可区分两种一致写法。令 $\varepsilon_{k,1}$ 为不含证券支付的直接禀赋，$h_k=\bar\theta_k+\theta_k$ 为总持仓，则
    \[
      c_{k,1}=\varepsilon_{k,1}+Xh_k
      =\underbrace{(\varepsilon_{k,1}+X\bar\theta_k)}_{e_{k,1}}+X\theta_k.
    \]
    本章使用右边写法：$e_{k,1}$ 已包括初始证券支付，$\theta_k$ 是净交易，不能再把 $X\bar\theta_k$ 加一次。用总持仓时出清为 $\sum_kh_k=\sum_k\bar\theta_k$，用净交易时出清为 $\sum_k\theta_k=0$。这只是改写账目，并没有把实际证券供给变成零。

    如此一来，\textbf{证券市场出清}要求有人净买 $\theta_k>0$，就要有人净卖 $\theta_k<0$。对所有参与者加总：净买入之和必须为 0，否则多出来的份额没人卖、或少掉的份额没人接，即净需求为零：
    \[
        \underbrace{\sum_{k=1}^{K}\theta_k(e_k,S)=0}_{\text{证券市场出清条件}}.
    \]
    满足该条件的 $S$ 即为\textbf{均衡价格}。由 $c_{k,0}=e_{k,0}-S^{\mathsf T}\theta_k(e_k,S)$，对 $k$ 求和得
    \[
        \underbrace{\underbrace{\sum_{k=1}^{K}c_{k,0}}_{\text{0 期商品总需求}}
        =\sum_{k=1}^{K}e_{k,0}-S^{\mathsf T}\Bigl(\sum_{k=1}^{K}\theta_k(e_k,S)\Bigr)
        =\underbrace{\sum_{k=1}^{K}e_{k,0}}_{\text{0 期总供给}}}_{\text{第0期商品市场的出清}}
    \]
    左边是第0 期商品的总需求，右边是第0期的总供给，表示的是第0期商品市场的出清。同理，对每个状态 $\omega=1,\dots,\Omega$ 有
    \[
        \sum_{k=1}^{K}c_{k,\omega}=\sum_{k=1}^{K}e_{k,\omega},
    \]
    即\textbf{商品市场出清}。合起来得到\textbf{逐日期、逐状态的资源可行性}（PPT 在此也称作 Walras 法则的出清含义）：
    \[
        \sum_{k=1}^{K}c_k=\sum_{k=1}^{K}e_k.
    \]
    

    求均衡时，\textbf{首先}对任意价格 $S$ 解出各参与者的最优证券组合 $\theta_k(e_k,S)$，\textbf{第二}用市场出清条件解出均衡价格 $S$。形式上，
    \[
        S=S\bigl(P;\,\{U_k,\,e_k,\,k=1,\dots,K\};\,X\bigr),
    \]
    这是模型内均衡价格对基本要素的依赖关系，不自动保证均衡存在或唯一，也不是现实中所有波动的完备解释。在生产技术、信息和制度均固定的本模型中分析比较静态，不能把任何政策变化都直接等同于偏好 $U$ 的变化。即证券价格取决于经济面临的风险 $P$、参与者的偏好 $U_k$ 与禀赋 $e_k$ 以及市场结构 $X$ 等\textbf{基本面}（fundamentals）或\textbf{本源}（primitives）：经济面临的风险，经济中参与者的偏好和禀赋以及证券市场的结构。
\end{enumerate}
\end{tcolorbox}

\begin{tcolorbox}[knowledge,title=2.6 最优性,breakable=false]
通过证券市场进行的资源配置是否已是\textbf{最好的可能配置}？参与者各自选择最优消费计划，并不意味着社会资源配置一定最优；若市场不完备，结果可能是\textbf{次优的}。

\begin{tcolorbox}[knowledge]
\textbf{定义 2.7}\quad 配置 $\{c_k,\,\forall k\}$ \textbf{Pareto 占优于}配置 $\{c_k',\,\forall k\}$，如果
\[
    \forall k,\quad U_k(c_k)\ge U_k(c_k'),
\]
且至少对一个参与者严格不等式成立。

这就是说，对比两套分配结果，如果新配置帕累托占优于旧配置，说明：新配置对所有参与者的效用都至少不比原来差，且至少对一个参与者的效用严格提高，且这种“好”不以“损人”为前提，这是一种\textbf{更具效率}的体现。
\end{tcolorbox}



\begin{tcolorbox}[knowledge]
\textbf{定义 2.8}\quad 给定经济的总供给 $\{e_k,\,\forall k\}$，配置 $\{c_k,\,\forall k\}$ 是\textbf{可行的}（feasible），如果
\[
    \underbrace{\sum_{k=1}^{K}c_k=\sum_{k=1}^{K}e_k}_{\text{说明所有状态下所有人的消费之和等于总禀赋之和，没有闲置}}.
\]
\end{tcolorbox}

\begin{tcolorbox}[knowledge]
\textbf{定义 2.9}\quad 若配置 $\{c_k,\,\forall k\}$ 可行，且不存在其他可行配置 Pareto 占优于它，则该配置是\textbf{Pareto 最优的}（Pareto optimal）。Pareto 最优即在不牺牲其他参与者福利、也不使用更多资源的前提下，无法改进任一参与者的福利（某人福利的提高，必定以损害他人的利益为前提）。给定偏好后，通常只有可行配置的一个子集是最优的。
\end{tcolorbox}

这里比较的是人的\textbf{整段消费计划的效用}，并不要求每个日期、每个状态下的消费都逐项增加。

Pareto 最优配置也叫做\textbf{有效}（efficient）配置。若证券市场使参与者达到 Pareto 最优配置，则称该市场在\textbf{资源配置意义上有效}。这与“价格反映信息”的有效市场假说是不同概念；效率也不等于平均分配或公平。
\end{tcolorbox}


\begin{tcolorbox}[example,title=例题：帕累托最优]
    \begin{question}
        \textbf{例 2.5}\quad 考虑一个两期、1 期有两种状态 $a$、$b$ 的经济。经济总禀赋：
        \begin{center}
        \WJEventTree[left={$100$},stem]{400,25}
        \end{center}
        有两个参与者，效用函数均为 $u(c_0,c_{1a},c_{1b})=\sqrt{c_0}+\tfrac12(\sqrt{c_{1a}}+\sqrt{c_{1b}})$（$c=\begin{bmatrix}c_0 \\ c_{1a} \\ c_{1b}\end{bmatrix}$）。树中每个节点括号内为 $(c_{1,\cdot},\,c_{2,\cdot})$：该节点上参与者 1、2 的消费量（不是整段消费计划列向量）。考虑三种配置：
        \begin{center}
        \begin{tabular}{@{}c@{\quad}c@{\quad}c@{}}
            \textbf{配置 A} & \textbf{配置 B} & \textbf{配置 C} \\[0.35em]
            \WJEventTree[left={$(100,0)$},stem]{ {(0,400)}, {(0,25)} } &
            \WJEventTree[left={$(50,50)$},stem]{ {(200,200)}, {(\tfrac{25}{2},\tfrac{25}{2})} } &
            \WJEventTree[left={$(100,0)$},stem]{ {(400,0)}, {(25,0)} }
        \end{tabular}
        \end{center}
        配置 A 中参与者 1 在 0 期消费全部、参与者 2 在 1 期消费全部；B 为两参与者在各期、各状态均分；C 中参与者 1 始终消费全部。哪种配置是 Pareto 最优的？（完整消费计划 $c_1$、$c_2$ 见解答。）
    \end{question}
    \begin{solution}
        与前文及习题 2.2 相同，参与者 $i$ 的\textbf{消费计划}是 $\mathrm{R}^{3}_{+}$ 中的列向量（两状态记为 $a,b$）：
        \[
            c_i=\begin{bmatrix}c_{i,0} \\ c_{i,1a} \\ c_{i,1b}\end{bmatrix},
            \qquad U_i(c_i)=\sqrt{c_{i,0}}+\tfrac12\bigl(\sqrt{c_{i,1a}}+\sqrt{c_{i,1b}}\bigr).
        \]
        题中事件树每个节点上的 $(\cdot,\cdot)$ 只是该节点上\textbf{参与者 1、2 的消费量}；整段计划仍用上面的 $c_1$、$c_2$ 表示。

        \textbf{可行性:}三个配置合起来在每一日期—状态下满足 $c_{1,\cdot}+c_{2,\cdot}$ 等于总禀赋（0 期 $100$，状态 $a$ 为 $400$，状态 $b$ 为 $25$），故 A、B、C 均为可行配置。

        \textbf{各配置的效用：}
        \begin{itemize}
            \item \textbf{配置 A：}
            \[
                c_1=\begin{bmatrix}100 \\ 0 \\ 0\end{bmatrix},\qquad
                c_2=\begin{bmatrix}0 \\ 400 \\ 25\end{bmatrix}.
            \]
            \[
                \begin{aligned}
                    U_1(c_1) &= \sqrt{100}+\tfrac12(0+0)=10,\qquad 
                    U_2(c_2) &= 0+\tfrac12(\sqrt{400}+\sqrt{25})
                    =\tfrac12(20+5)=12.5.
                \end{aligned}
            \]
            效用对为 $(U_1,U_2)=(10,\,12.5)$。

            \item \textbf{配置 B}（均分）：$c_1=c_2=
            \begin{bmatrix}50 \\ 200 \\ \tfrac{25}{2}\end{bmatrix}$。
            \[
                U_1(c_1)=U_2(c_2)=\sqrt{50}+\tfrac12\Bigl(\sqrt{200}+\sqrt{\tfrac{25}{2}}\Bigr)
                \approx 7.071+8.839=15.91.
            \]
            效用对为 $(15.91,\,15.91)$。

            \item \textbf{配置 C：}
            \[
                c_1=\begin{bmatrix}100 \\ 400 \\ 25\end{bmatrix},\qquad
                c_2=\begin{bmatrix}0 \\ 0 \\ 0\end{bmatrix}.
            \]
            \[
                \begin{aligned}
                    U_1(c_1) = \sqrt{100}+\tfrac12(\sqrt{400}+\sqrt{25})=10+12.5=22.5, \quad
                    U_2(c_2) = 0.
                \end{aligned}
            \]
            效用对为 $(22.5,\,0)$。
        \end{itemize}

        \textbf{Pareto 占优与 Pareto 最优。}
        \begin{itemize}
            \item 比较 A 与 B：$U_1(B)=15.91>10=U_1(A)$ 且 $U_2(B)=15.91>12.5=U_2(A)$，两人都严格更好，故\textbf{配置 B Pareto 占优于配置 A}。因此 A \textbf{不是} Pareto 最优。

            \item 比较 B 与 C：$U_1(C)=22.5>U_1(B)$，但 $U_2(C)=0<U_2(B)$，有人变差，故 C 不 Pareto 占优 B，B 也不 Pareto 占优 C。

            \item \textbf{配置 B 是 Pareto 最优的}：仅比较 A、B、C 还不够，须排除\textbf{所有}可行改善。对任一节点总资源 $E$ 和分配 $x,E-x$，凹性给出 $\sqrt x+\sqrt{E-x}\le2\sqrt{E/2}$。将 0 期与两状态分别加权相加，得任何可行配置的 $U_1+U_2$ 都不超过均分 B 的总效用。若有人严格改善、另一人不变差，总效用必上升，矛盾。因此 B 最优；这里加总效用只是在给定这两个效用函数后证明结论的数学工具，不是一般的人际福利比较标准。

            \item \textbf{配置 C 是 Pareto 最优的}：若要提高参与者 2 的效用（必须给 2 正的消费），只能从参与者 1 已占用的 $100$、$400$、$25$ 中让出，从而必使 $U_1$ 下降，无法在不损害 1 的情况下改进 2。
        \end{itemize}

        
    \end{solution}

    \begin{tcolorbox}[warning,title=注意]
        \textbf{结论：}A 非 Pareto 最优；\textbf{B 与 C 都是} Pareto 最优配置（同一总禀赋下不同的有效分配）。可以看出：
        \begin{itemize}
            \item 在同一可行集（同一总禀赋、满足 $\sum_k c_k=\sum_k e_k$）里，Pareto 最优配置往往不止一个。
            \item 更能看到，帕累托最优仅考虑了“效率”，但没考虑“公平”，例如配置 C 中参与者 2 的效用为0，但仍属于帕累托最优。
        \end{itemize}
        \end{tcolorbox}
\end{tcolorbox}


\begin{tcolorbox}[knowledge,title=2.7 本章小结]
本章完成了在市场经济中通过资源配置分析所使用的基本框架的构建。资源配置是通过证券市场完成的：交易证券的价格引导参与者通过交易达到他们希望的配置，而交易又最终决定了证券的市场价格。证券价格在引导个体参与者和平衡整个经济方面都起着关键作用。因此，我们关注证券价格如何决定，这也称为\textbf{资产定价}（asset pricing）。

\begin{enumerate}
    \item 从金融的角度看，一个经济的结构包括三个方面：（1）它所处的\textbf{自然环境}；（2）\textbf{经济的参与者}；（3）\textbf{金融市场}。

    \item \textbf{外部环境}包括两个要素：\textbf{时间}和\textbf{风险}。为刻画时间，假设有 0 和 1 两个时期；为刻画风险，假设 1 期有 $\Omega$ 个可能状态，用概率测度 $P$ 描述。经济中只有一种不可储存的商品。

    \item 经济中有 $K$ 个\textbf{参与者}。每个参与者的经济特征包括其经济资源和经济需求。具体假设为：
    \begin{enumerate}[label=(\alph*)]
        \item 每个参与者 $k$ 具有禀赋 $e_k\in\mathrm{R}^{1+\Omega}_{+}$；
        \item 每个参与者对未来状态发生的可能性拥有相同的信息，完全由概率测度 $P$ 描述；
        \item 尽管参与者给定其禀赋和信息后可能拥有生产技术，这里先假设他们不拥有任何生产技术；
        \item 每个参与者在 $C=\mathrm{R}^{1+\Omega}_{+}$ 上具有偏好，且满足不满足性、连续性和凸性公理；
        \item 给定这些偏好关系，可用效用函数来描述。
    \end{enumerate}

    \item \textbf{金融市场}由一组证券构成：
    \begin{enumerate}[label=(\alph*)]
        \item 它们的支付只依赖于未来实际的经济状态；
        \item 它们的支付是外生给定的，且不受经济参与者行为的影响；
        \item 参与者知道且一致认同证券支付的概率分布；
        \item 市场结构为 $X$，且市场中无摩擦。
    \end{enumerate}

    \item \textbf{市场均衡}必须满足两个条件：
    \begin{enumerate}[label=(\alph*)]
        \item 每个参与者各自最优化其证券需求量；
        \item \textbf{市场出清}，即交易证券的总需求等于总供给，此时商品市场也是出清的。
    \end{enumerate}

    \item 如果配置 $\{c_k,\,\forall k\}$ 是可行的，且不存在另外的可行配置 Pareto 占优于它，则该配置是 \textbf{Pareto 最优的}。如果证券市场允许参与者达到 Pareto 最优配置，就称它为\textbf{有效市场}。

    \item \textbf{对数效用下的支出份额。}设 \(U=\log c_0+\rho\log c_1\)，两期且 1 期只有一个状态，只有收益率为 \(r_F\) 的无风险证券。以 0 期消费计价的财富为 \(w=e_0+e_1/(1+r_F)\)。最优消费满足
    \[
        c_1^{*}=\rho(1+r_F)\,c_0^{*},
    \]
    因而 0 期消费占财富的份额只由 \(\rho\) 决定：
    \[
        c_0^{*}=\frac{w}{1+\rho},\qquad
        c_1^{*}=\frac{\rho(1+r_F)}{1+\rho}\,w.
    \]
    \(c_0^{*}/w=1/(1+\rho)\) 不随 \(r_F\) 改变；利率只通过 \(w\) 里未来禀赋的现值影响 \(c_0^{*}\)。这是对数效用（相对风险厌恶为 1）的结论。对幂效用 $u(c)=c^{1-\gamma}/(1-\gamma)$，一阶条件为 $c_1/c_0=[\rho(1+r_F)]^{1/\gamma}$，所以
    \[
      \frac{c_0^*}{w}=\frac1{1+\rho^{1/\gamma}(1+r_F)^{1/\gamma-1}}.
    \]
    幂效用的份额仍不依赖财富水平，但一般会随利率改变；对数情形 $\gamma=1$ 才消去上式的利率项。负指数效用一般也不具有同样的固定份额。
\end{enumerate}
\end{tcolorbox}
\fi

% !TeX root = ../../金融经济学_王江_讲义及习题.tex
% 02_基本框架 课后习题（由章节文件拆分）

\begin{problem}
\problemno{2.1} $U(c)$ 和 $V(c)$ 是两个效用函数，$c\in\mathrm{R}_{+}^{n}$，且 $V(x)=f(U(x))$，其中 $f(\cdot)$ 是一正单调函数。证明这两个效用函数表示了相同的偏好。
\end{problem}
\begin{solution}
定义 2.4：效用函数表示偏好 $\succeq$，是指对任意两个消费计划，$U(c_{1})\ge U(c_{2})$ 当且仅当 $c_{1}\succeq c_{2}$。两个效用函数表示相同偏好，就是它们给出完全一样的排序。正单调变换不改变不等号的方向，所以 $V=f(U)$ 与 $U$ 排序相同。

假设 $U(c)$ 表示的偏好关系为 $\succeq$，那么 $\forall\, c_{1},\,c_{2}\in\mathrm{R}_{+}^{N}$ 有
\[
    U(c_{1})\ge U(c_{2})\Leftrightarrow c_{1}\succeq c_{2}
\]
而 $f(\cdot)$ 是正单调函数，因而
\[
    V(c_{1})=f(U(c_{1}))\ge f(U(c_{2}))=V(c_{2})\Leftrightarrow U(c_{1})\ge U(c_{2})
\]
因此 $V(c_{1})\ge V(c_{2})\Leftrightarrow c_{1}\succeq c_{2}$，即 $V(c)$ 表示的偏好也是 $\succeq$。
\end{solution}

\begin{problem}
\problemno{2.2*} 在 1 期，经济有两个可能状态 $a$ 和 $b$，它们的发生概率相等：
\begin{center}
\WJEventTree{a,b}
\end{center}
考虑定义在消费计划 $c=\begin{bmatrix}c_{0} \\ c_{1a} \\ c_{1b}\end{bmatrix}$ 上的效用函数：
\[
    \begin{aligned}
        U(c) &= \log c_{0}+\tfrac12(\log c_{1a}+\log c_{1b})\\
        U(c) &= \tfrac{1}{1-\gamma}c_{0}^{1-\gamma}+\tfrac12\Bigl(\tfrac{1}{1-\gamma}c_{1a}^{1-\gamma}+\tfrac{1}{1-\gamma}c_{1b}^{1-\gamma}\Bigr)\\
        U(c) &= -e^{-ac_{0}}-\tfrac12\bigl(e^{-ac_{1a}}+e^{-ac_{1b}}\bigr)
    \end{aligned}
\]
证明它们满足：不满足性、连续性和凸性。
\end{problem}
\begin{solution}
三个效用函数都具有形式 $U(c)=u(c_{0})+\tfrac12\bigl(u(c_{1a})+u(c_{1b})\bigr)$。对数与幂效用要求底数为正，指数效用 $u(x)=-e^{-ax}$ 在 $a>0$ 时于整个实数轴上有定义。下面统一在严格正消费集
\[
    C=\bigl\{c\in\mathrm{R}^{3}:c_{0}>0,\;c_{1a}>0,\;c_{1b}>0\bigr\}
\]

上证明；并设幂效用中 $\gamma>0$、$\gamma\neq 1$，指数效用中 $a>0$。对应的期效用为
\[
    u(x)=\log x,\qquad
    u(x)=\frac{x^{1-\gamma}}{1-\gamma},\qquad
    u(x)=-e^{-ax},\qquad x>0.
\]

三者都满足 $u'(x)>0$、$u''(x)<0$，因而严格递增且严格凹：
\[
    (\log x)'=\frac1x,\quad (\log x)''=-\frac1{x^{2}},
\]
\[
    \Bigl(\frac{x^{1-\gamma}}{1-\gamma}\Bigr)'=x^{-\gamma},\quad
    \Bigl(\frac{x^{1-\gamma}}{1-\gamma}\Bigr)''=-\gamma x^{-\gamma-1},
\]
\[
    (-e^{-ax})'=ae^{-ax},\quad (-e^{-ax})''=-a^{2}e^{-ax}.
\]

\begin{enumerate}[label=(\alph*)]
    \item 不满足性（喜多厌少）。设 $c,c'\in C$ 且 $c\ge c'$。由 $u$ 在 $(0,\infty)$ 上递增，$u(c_{i})\ge u(c_{i}')$，故 $U(c)\ge U(c')$。
    \item 连续性。要证：若 $\{c^{(n)}\}\subset C$ 且 \(\lim_{n\to\infty}c^{(n)}=c\in C\)，则
    \[
        \lim_{n\to\infty}U\bigl(c^{(n)}\bigr)=U(c).
    \]

    按分量有 \(\lim_{n\to\infty}c_{i}^{(n)}=c_{i}>0\)。\(u\) 在 \((0,\infty)\) 上连续，因而在点 \(c_{i}\) 连续。由连续的序列定义，极限可以代入 \(u\)：
    \[
        \lim_{n\to\infty}u\bigl(c_{i}^{(n)}\bigr)=u\Bigl(\lim_{n\to\infty}c_{i}^{(n)}\Bigr)=u(c_{i}).
    \]
    $U$ 只是三项 $u$ 的和，极限穿过和式：
    \[
        \begin{aligned}
            \lim_{n\to\infty}U\bigl(c^{(n)}\bigr)
            &=\lim_{n\to\infty}u\bigl(c_{0}^{(n)}\bigr)
            +\tfrac12\lim_{n\to\infty}u\bigl(c_{1a}^{(n)}\bigr)
            +\tfrac12\lim_{n\to\infty}u\bigl(c_{1b}^{(n)}\bigr)\\
            &=u(c_{0})+\tfrac12\bigl(u(c_{1a})+u(c_{1b})\bigr)
            =U(c).
        \end{aligned}
    \]
    \item 凸性。$u$ 严格凹，故对 $\alpha\in(0,1)$ 与 $c,c'\in C$，
    \[
        u\bigl(\alpha c_{i}+(1-\alpha)c_{i}'\bigr)\ge\alpha u(c_{i})+(1-\alpha)u(c_{i}'),
    \]
    且凸组合仍在 $C$ 中。因此
    \[
        U\bigl(\alpha c+(1-\alpha)c'\bigr)\ge\alpha U(c)+(1-\alpha)U(c').
    \]

    若 $U(c)>U(c')$，则右端严格大于 $U(c')$，即 $U\bigl(\alpha c+(1-\alpha)c'\bigr)>U(c')$。
\end{enumerate}
\end{solution}

\begin{problem}
\problemno{2.3} $U(c)=c-\tfrac12 a c^{2}$ 是一可能的效用函数，其中 $c\in\mathrm{R}_{+}$，$a$ 是非负的系数。$U(c)$ 具有不满足性吗？如果不是，那么 $a$ 取什么值和/或 $c$ 在什么范围内时 $U(c)$ 具有不满足性？
\end{problem}
\begin{solution}
不一定，比如当 $a=1$ 时，$U(\tfrac12)=\tfrac38<U(1)=\tfrac12>U(3)=-1.5$。$U(c)$ 不具有不满足性。

当 $a=0$ 时，$U(c)=c$ 具有不满足性；

当 $a>0$ 时，当 $c\in[0,\tfrac{1}{a}]$ 时 $U(c)$ 具有不满足性。

\begin{center}
\footnotesize
\begin{tabular}{@{}ccc@{}}
\begin{tikzpicture}[baseline=(current bounding box.center), x=1cm, y=1cm]
    \path (-0.32,-1.55) rectangle (2.55,1.62);
    \draw[-{Stealth[length=1.5mm]}, thin] (0,0) -- (2.28,0) node[below left, inner sep=1pt] {$c$};
    \draw[-{Stealth[length=1.5mm]}, thin] (0,0) -- (0,1.48) node[above left, inner sep=1pt] {$U$};
    \draw[blue!65!black, thick, domain=0:2] plot (\x,{0.65*\x});
\end{tikzpicture}
&
\begin{tikzpicture}[baseline=(current bounding box.center), x=1cm, y=1cm]
    \path (-0.32,-1.55) rectangle (2.55,1.62);
    \draw[-{Stealth[length=1.5mm]}, thin] (0,0) -- (2.28,0) node[below left, inner sep=1pt] {$c$};
    \draw[-{Stealth[length=1.5mm]}, thin] (0,0) -- (0,1.48) node[above left, inner sep=1pt] {$U$};
    \draw[blue!65!black, thick, domain=0:1, samples=40, smooth] plot ({2*\x},{2.2*(\x-0.5*\x*\x)});
    \draw[dashed, black!45] (2,0) -- (2,1.1);
    \node[below, inner sep=1pt] at (2,0) {$1/a$};
\end{tikzpicture}
&
\begin{tikzpicture}[baseline=(current bounding box.center), x=1cm, y=1cm]
    \path (-0.32,-1.55) rectangle (2.55,1.62);
    \draw[-{Stealth[length=1.5mm]}, thin] (0,0) -- (2.35,0) node[above right, inner sep=1pt] {$c$};
    \draw[-{Stealth[length=1.5mm]}, thin] (0,-1.4) -- (0,0.85) node[above left, inner sep=1pt] {$U$};
    \draw[blue!65!black, thick, domain=0:3.05, samples=50, smooth] plot ({0.68*\x},{0.85*(\x-0.5*\x*\x)});
    \draw[dashed, black!45] (0.34,0) -- (0.34,0.319);
    \draw[dashed, black!45] (0.68,0) -- (0.68,0.425);
    \draw[dashed, black!45] (2.04,0) -- (2.04,-1.275);
    \fill (0.34,0.319) circle (1pt);
    \fill (0.68,0.425) circle (1pt);
    \fill (2.04,-1.275) circle (1pt);
    \draw (0.34,0.045) -- (0.34,-0.045) node[below, inner sep=1.5pt] {$\tfrac12$};
    \draw (0.68,0.045) -- (0.68,-0.045) node[below, inner sep=1.5pt] {$1$};
    \draw (2.04,0.045) -- (2.04,-0.045) node[below, inner sep=1.5pt] {$3$};
\end{tikzpicture}
\\[0.15em]
$a=0$，$U(c)=c$ & $a>0$，$c\in[0,1/a]$ & $a=1$，$U(\tfrac12)<U(1)>U(3)$
\end{tabular}
\end{center}
\end{solution}

\begin{problem}
\problemno{2.4} 考虑一个经济，它在 1 期有三个可能状态：$a$，$b$ 和 $c$：
\begin{center}
\WJEventTree{a,b,c}
\end{center}
证券市场包括证券 1 和 2，它们具有如下的支付向量：$X_{1}=\begin{bmatrix}1 \\ 1 \\ 1\end{bmatrix}$ 以及 $X_{2}=\begin{bmatrix}1 \\ 2 \\ 3\end{bmatrix}$。它们的价格分别为 $S_{1}$ 和 $S_{2}$。
\begin{enumerate}[label=(\alph*)]
    \item 描述这个经济的支付空间。
    \item 写出这个经济的市场结构矩阵 $X$。
    \item 考虑含有 $\theta_{1}$ 单位的证券 1 和 $\theta_{2}$ 单位的证券 2 的组合。写出这个组合的支付向量。这个组合的价格是多少？
    \item 假设这个市场中总共有 $K$ 个参与者。每个参与者的禀赋是 1 单位的证券 1 和 2 单位的证券 2。这时的市场组合是什么？市场组合的支付向量是什么？市场组合的总价值是多少？
    \item 写出市场化支付的集合。
    \item 如果市场不允许卖空。市场化支付的集合是什么？
    \item 现在引入新的证券 3，它的支付向量为 $X_{3}=\begin{bmatrix}0 \\ 0 \\ 1\end{bmatrix}$。写出新的市场结构矩阵。在这个市场结构下，市场化支付集合是什么？
\end{enumerate}
\end{problem}
\begin{solution}
三个状态 $a,b,c$ 对应 $\Omega=3$。记
\[
    X_1=\begin{bmatrix}1 \\ 1 \\ 1\end{bmatrix},\quad
    X_2=\begin{bmatrix}1 \\ 2 \\ 3\end{bmatrix},\quad
    \theta=\begin{bmatrix}\theta_1 \\ \theta_2\end{bmatrix}.
\]

组合在 1 期的支付向量为 $X_1\theta_1+X_2\theta_2$，也可写成矩阵乘积 $X\theta$，其中 $X=[X_1,\,X_2]$。

\begin{enumerate}[label=(\alph*)]
    \item \textbf{支付空间}：所有可能在三个状态下实现的支付向量的集合，即 $\mathrm{R}^{\Omega}=\mathrm{R}^{3}$（每个分量对应状态 $a,b,c$ 下的支付）。

    \item \textbf{市场结构矩阵}：两列分别为两只证券的支付向量，
    \[
        X=\begin{bmatrix}1&1\\1&2\\1&3\end{bmatrix}=[X_1,\,X_2].
    \]

    \item \textbf{组合支付与价格：}持有 $(\theta_1,\theta_2)$ 时，
    \[
        X\theta=\theta_1 X_1+\theta_2 X_2=
        \begin{bmatrix}\theta_1+\theta_2 \\ \theta_1+2\theta_2 \\ \theta_1+3\theta_2\end{bmatrix}.
    \]

    0 期购买成本（组合价格）为 $\theta_1 S_1+\theta_2 S_2$。

    \item \textbf{市场组合：}$K$ 个参与者每人禀赋为 $1$ 单位证券 1、$2$ 单位证券 2，加总得市场组合
    \[
        \theta_M=\begin{bmatrix}K \\ 2K\end{bmatrix}.
    \]

    其支付向量为
    \[
        X\theta_M=K X_1+2K X_2=
        \begin{bmatrix}3K \\ 5K \\ 7K\end{bmatrix}.
    \]

    总价值为 $K S_1+2K S_2$。

    \item \textbf{市场化支付（允许卖空）：}能由证券组合复制的 1 期支付集合为
    \[
        M=\bigl\{Y\in\mathrm{R}^{3}: Y=X\theta=[X_1,\,X_2]\begin{bmatrix}\theta_1 \\ \theta_2\end{bmatrix}=\theta_1 X_1+\theta_2 X_2,\;\theta_1,\theta_2\in\mathrm{R}\bigr\},
    \]

    即 $X_1,X_2$ 张成的二维子空间（只有两只证券，维度 $\le 2$）。

    \item \textbf{禁止卖空：}仅允许 $\theta_1,\theta_2\ge 0$，为 $X_1,X_2$ 的\textbf{非负}线性组合：
    \[
        M_{+}=\bigl\{Y\in\mathrm{R}^{3}: Y=X\theta=[X_1,\,X_2]\begin{bmatrix}\theta_1 \\ \theta_2\end{bmatrix}=\theta_1 X_1+\theta_2 X_2,\;\theta_1,\theta_2\in\mathrm{R}_{+}\bigr\}.
    \]

    \item \textbf{加入证券 3：}其支付向量为 $X_3=\begin{bmatrix}0 \\ 0 \\ 1\end{bmatrix}$，仅在状态 $c$ 支付 $1$。两只证券时的市场结构是 $X=[X_1,\,X_2]$。加入证券 3 后的市场结构另记为 $X'$，以免与 $X$ 相混：
    \[
        X'=\begin{bmatrix}1&1&0\\1&2&0\\1&3&1\end{bmatrix}=[X_1,\,X_2,\,X_3].
    \]

    三列线性无关，可表示任意 $Y\in\mathrm{R}^{3}$，故（允许卖空时）
    \[
        M=\bigl\{Y\in\mathrm{R}^{3}: Y=\theta_1 X_1+\theta_2 X_2+\theta_3 X_3,\;\theta_1,\theta_2,\theta_3\in\mathrm{R}\bigr\}=\mathrm{R}^{3}.
    \]
\end{enumerate}
\end{solution}

\begin{problem}
\problemno{2.5} 在练习 2.4 中定义的只存在证券 1 和 2 的经济中。考虑一个禀赋为 $\theta_{1}$ 单位的证券 1 和 $\theta_{2}$ 单位的证券 2 的参与者。写出他的预算集。
\end{problem}
\begin{solution}
无摩擦，允许卖空。预算集是花费不超过禀赋市值的那些 1 期支付：
\[
    B_{12}=\Bigl\{C\in\mathrm{R}_{+}^{3}:
    C=\alpha_{1}X_{1}+\alpha_{2}X_{2},\;
    \alpha_{1}S_{1}+\alpha_{2}S_{2}\le\theta_{1}S_{1}+\theta_{2}S_{2},\;
    \alpha_{1},\alpha_{2}\in\mathrm{R}\Bigr\}.
\]
\end{solution}

\begin{problem}
\problemno{2.6} 在上面的练习中引入练习 2.4 中定义的证券 3，它的价格为 $S_{3}$。这时，参与者的预算集是什么（他在证券 3 上的禀赋为 0）？证明由证券 1、2、3 构成的预算集包含仅由证券 1、2 构成的预算集。
\end{problem}
\begin{solution}
证券 3 上的禀赋为 $0$，可交易财富仍是 $\theta_{1}S_{1}+\theta_{2}S_{2}$。预算集为
\[
    B_{123}=\Bigl\{C\in\mathrm{R}_{+}^{3}:
    C=\alpha_{1}X_{1}+\alpha_{2}X_{2}+\alpha_{3}X_{3},\;
    \alpha_{1}S_{1}+\alpha_{2}S_{2}+\alpha_{3}S_{3}\le\theta_{1}S_{1}+\theta_{2}S_{2},\;
    \alpha_{1},\alpha_{2},\alpha_{3}\in\mathrm{R}\Bigr\}.
\]

练习 2.5 的预算集是
\[
    B_{12}=\Bigl\{C\in\mathrm{R}_{+}^{3}:
    C=\alpha_{1}X_{1}+\alpha_{2}X_{2},\;
    \alpha_{1}S_{1}+\alpha_{2}S_{2}\le\theta_{1}S_{1}+\theta_{2}S_{2},\;
    \alpha_{1},\alpha_{2}\in\mathrm{R}\Bigr\}.
\]

设 $C\in B_{12}$，对应头寸为 $\alpha_{1},\alpha_{2}$。取 $\alpha_{3}=0$，则
\[
    C=\alpha_{1}X_{1}+\alpha_{2}X_{2}+\alpha_{3}X_{3},
\]

并且
\[
    \alpha_{1}S_{1}+\alpha_{2}S_{2}+\alpha_{3}S_{3}=\alpha_{1}S_{1}+\alpha_{2}S_{2}\le\theta_{1}S_{1}+\theta_{2}S_{2}.
\]

故 $C\in B_{123}$，即 $B_{12}\subseteq B_{123}$。
\end{solution}

\begin{problem}
\problemno{2.7*} 考虑一个在 1 期只有一个可能状态的经济。（在这种情况下不存在不确定性。）参与者 1 的 0 期禀赋为 100 而 1 期禀赋为 1，即他的禀赋向量为 $\begin{bmatrix}100 \\ 1\end{bmatrix}$。他的偏好可以表示成如下形式：
\[
    U(c_{0},c_{1})=\log c_{0}+\rho\log c_{1}.
\]

系数 $\rho$ 为反映参与者在当前消费和未来消费之间相对偏好的参数。有一只证券，它的 0 期价格为 1、1 期支付为 $1+r_{F}$。这里，$r_{F}$ 是利率。
\begin{enumerate}[label=(\alph*)]
    \item 如果这个参与者不能在市场上进行交易，那么他的消费计划以及相应的效用 $U^{a}$ 是什么？
    \item 现在假设他可以在市场上进行交易。
    \begin{itemize}
        \item 他的预算集是什么？以当前消费为单位，他的总财富 $w$ 是多少？
        \item 写出参与者的优化问题。令 $c_{0}$ 为参与者的当前（即 0 期）最优消费，$s$ 为最优储蓄以及 $U^{b}$ 为在最优策略下得到的效用。求解他的最优消费/储蓄选择以及相应的效用。把 $U^{b}$ 表示成财富 $w$、利率 $r_{F}$ 和偏好系数 $\rho$ 的函数。
        \item 讨论参与者的最优选择如何依赖于利率 $r_{F}$ 和偏好系数 $\rho$。给出解释。
    \end{itemize}
    \item 证明 $U^{b}\ge U^{a}$。
    \item 令 $g$ 为参与者由于能够在证券市场上交易而获得的益处。它的定义为
    \[
        U^{b}(w-g)=U^{a}.
    \]

    计算 $g$。讨论 $g$ 如何依赖于 $\rho$？$g$ 如何依赖于 $r_{F}$？给出解释。
\end{enumerate}
\end{problem}
\begin{solution}
禀赋 $e_0=100$、$e_1=1$。无风险证券：0 期价格 $1$，1 期支付 $1+r_F$。记 $S$ 为 0 期\textbf{储蓄}（买入该证券的金额；$S<0$ 表示借贷）。

\begin{enumerate}[label=(\alph*)]
    \item \textbf{不能交易（自给自足）：}只能消费禀赋，
    \[
        c_0=100,\qquad c_1=1,\qquad
        U^{a}=\log(100)+\rho\log(1)=\log 100.
    \]

    \item \textbf{可以交易。}
    \begin{itemize}
        \item \textbf{预算集与财富 $w$。}买卖证券等价于在两期之间转移消费：0 期少消费 $S$ 用于储蓄，1 期增加 $S(1+r_F)$。故
        \[
            c_0=100-S,\qquad c_1=1+S(1+r_F),\qquad S\in\mathrm{R}.
        \]

        以 0 期消费为计价单位，把 1 期禀赋折现到 0 期（用市场利率 $r_F$）：
        \[
            w=100+\frac{1}{1+r_F}.
        \]

        这是参与者按市价出售全部禀赋后可在 0 期支配的\textbf{总财富}（现值）。任意满足上述两式的 $(c_0,c_1)$ 都在预算集内。

        \item \textbf{优化问题。}在预算集上最大化对数效用：
        \[
            \max_{S\in\mathrm{R}}\ \log(100-S)+\rho\log\bigl[1+S(1+r_F)\bigr].
        \]

        一阶条件（边际效用替代率 = 跨期相对价格）：
        \[
            \underbrace{\frac{1}{100-S}}_{\text{0 期边际效用}}
            =\rho\underbrace{\frac{1+r_F}{1+S(1+r_F)}}_{\text{1 期边际效用}\times\text{价格比}}.
        \]

        整理得 $\bigl[1+S(1+r_F)\bigr]=\rho(1+r_F)(100-S)$，即
        \[
            1+S(1+r_F)(1+\rho)=100\rho(1+r_F).
        \]
        
        其解记为最优储蓄
        \[
            S^{*}=\frac{100\rho(1+r_F)-1}{(1+\rho)(1+r_F)}.
        \]

        预算约束是前面的两式 \(c_0=100-S\)、\(c_1=1+S(1+r_F)\)，其中 \(c_0,c_1,S\) 仍是可行选择，不是最优值。把 \(S^{*}\) 代回去，得到最优消费 \(c_0^{*},c_1^{*}\)：
        \[
            \begin{aligned}
                c_0^{*}&=100-S^{*}
                =100-\frac{100\rho(1+r_F)-1}{(1+\rho)(1+r_F)}
                =\frac{100(1+r_F)+1}{(1+\rho)(1+r_F)},\\[0.4em]
                c_1^{*}&=1+S^{*}(1+r_F)
                =1+(1+r_F)\cdot\frac{100\rho(1+r_F)-1}{(1+\rho)(1+r_F)}
                =\frac{\rho\bigl[100(1+r_F)+1\bigr]}{1+\rho}.
            \end{aligned}
        \]

        财富 \(w=100+\dfrac{1}{1+r_F}=\dfrac{100(1+r_F)+1}{1+r_F}\) 的分子正是上式的 \(100(1+r_F)+1\)，故
        \[
            c_0^{*}=\frac{w}{1+\rho},\qquad
            c_1^{*}=\frac{\rho(1+r_F)}{1+\rho}\,w.
        \]

        两式相除，零期与一期的最优消费只差偏好权重和总收益：
        \[
            c_1^{*}=\rho(1+r_F)\,c_0^{*}.
        \]

        因此最优现期消费占财富 $\dfrac{1}{1+\rho}$，这个份额不随 \(r_F\) 改变；最优 1 期消费占 \(\dfrac{\rho(1+r_F)}{1+\rho}\)，是在该份额上再乘以 \(\rho(1+r_F)\)。这是对数效用下的固定支出份额。最优效用
        \[
            U^{b}=\log\Bigl(\frac{w}{1+\rho}\Bigr)+\rho\log\Bigl(\frac{\rho(1+r_F)}{1+\rho}\,w\Bigr).
        \]

        \item \textbf{比较静态。}上面已经求出最优消费量：
        \[
            c_0^{*}=\frac{100(1+r_F)+1}{(1+\rho)(1+r_F)}
            =\frac{1}{1+\rho}\Bigl(100+\frac{1}{1+r_F}\Bigr),\qquad
            c_1^{*}=\frac{\rho\bigl[100(1+r_F)+1\bigr]}{1+\rho}.
        \]
        \(r_F\) 只出现在 \(c_0^{*}\) 的 \(1/(1+r_F)\) 里，故 \(r_F\) 上升时 \(c_0^{*}\) 下降；\(\rho\) 在分母，\(\rho\) 上升时 \(c_0^{*}\) 也下降。
        
        \(c_1^{*}\) 的分子含 \(\rho(1+r_F)\)，分母是 \(1+\rho\)，两项都使 \(c_1^{*}\) 随 \(r_F\)、\(\rho\) 上升。
        
        直觉：\(r_F\) 升高提高储蓄回报，\(\rho\) 升高加重未来消费权重，都会促使少消费于 0 期、多留于 1 期。
    \end{itemize}

    \item \textbf{$U^{b}\ge U^{a}$。}不能交易时只能消费禀赋，\((c_0,c_1)=(100,1)\)，效用为 \(U^{a}=\log 100\)。
    
    能交易时储蓄 \(S\) 可任取，其中 \(S=0\) 仍可行，对应的消费还是 \((100,1)\)，效用仍是 \(\log 100=U^{a}\)。\(U^{b}\) 是所有可行储蓄里的最大效用，而 \(S=0\) 只是其中一个可行点，最大值不会低于这一点：
    \[
        U^{b}=\max_{S\in\mathrm{R}}\,U\bigl(100-S,\,1+S(1+r_F)\bigr)
        \ge U(100,1)=U^{a}.
    \]
    等号当且仅当最优储蓄恰好是 \(S^{*}=0\)，即 \(100\rho(1+r_F)=1\)。

    \item \textbf{交易益处 $g$。}定义 $U^{b}(w-g)=U^{a}$：$w-g$ 是与交易后最优效用相同、但用一条 0 期财富数字表示的\textbf{确定性等价财富}；$g=w-(w-g)$ 即为因能交易而多出来的等价 0 期资源。

    将 $U^{a}=\log 100$ 与上式 $U^{b}$ 联立，记 $w-g=\hat w$，则
    \[
        (1+\rho)\log\hat w+\rho\log\bigl(\rho(1+r_F)\bigr)-(1+\rho)\log(1+\rho)=\log 100.
    \]

    取指数并整理：
    \[
        \hat w^{1+\rho}\,\frac{\bigl[\rho(1+r_F)\bigr]^{\rho}}{(1+\rho)^{1+\rho}}=100
        \quad\Rightarrow\quad
        \hat w=\rho^{-\frac{\rho}{1+\rho}}(1+r_F)^{-\frac{\rho}{1+\rho}}
        100^{\frac{1}{1+\rho}}(1+\rho).
    \]

    故
    \[
        g=w-\hat w
        =w-\rho^{-\frac{\rho}{1+\rho}}(1+r_F)^{-\frac{\rho}{1+\rho}}
        100^{\frac{1}{1+\rho}}(1+\rho).
    \]
    $g$ 随 $\rho$、$r_F$ 增加而增加：二者都提高把消费从 0 期转到 1 期」的动机，能在市场上做跨期配置带来的等价收益更大。
\end{enumerate}
\end{solution}



% !TeX root = ../../金融经济学_王江_讲义及习题.tex
% 知识点：references/金融经济学/金融经济学-王江_261013.pdf，第 3 章
% 补充：references/金融经济学/金融经济学_第3章 Arrow-Debreu经济.pptx（第 1–30 页）
% 习题：references/金融经济学/史上最全最完整的金融经济学答案(王江版).pdf，页 8–12（印刷页 7–11）

\lecture{阿罗--德布鲁经济}

\ifdefined\WJHandoutKnowledge


本章运用第 2 章的框架分析\textcolor{red}{\textbf{参与者如何选择投资组合、证券价格如何影响这些选择以及市场如何决定这些价格}}。这一简单的市场结构也是理解资源配置效率的基准。

\begin{tcolorbox}[knowledge,title=3.1 Arrow-Debreu 证券市场]
下面引入一类最简单的证券：它们只在 1 期支付，且支付是已实现状态的函数。

\begin{tcolorbox}[knowledge]
\textbf{定义 3.1}\quad 状态 $\omega$ \textbf{或有要求权}（state-$\omega$ contingent claim）就是当状态为 $\omega$ 时支付为 $1$、而在其他状态下支付为 $0$ 的证券。
\end{tcolorbox}

因此，状态 $\omega$ 或有要求权具有如下支付分布（$x_{\omega'}$ 为状态 $\omega'$ 下的支付）：
\[
    x_{\omega'}=
    \begin{cases}
        1, & \omega'=\omega,\\
        0, & \omega'\neq\omega,
    \end{cases}
    \qquad\text{即}\qquad
    x(\omega')=\mathbf{1}_{\omega}(\omega'),
\]
其中 $\mathbf{1}_{\omega}(\cdot)$ 为示性函数：当 $\omega'=\omega$ 时取 $1$，否则取 $0$。

\begin{tcolorbox}[example]
\textbf{例 3.1}\quad 在例 2.1 的 Lucas 树经济中，1 期有两个可能状态，因而有两个或有要求权，其支付向量分别为
\[
    \text{状态 $a$ 或有要求权：}\begin{bmatrix}1 \\ 0\end{bmatrix},\qquad
    \text{状态 $b$ 或有要求权：}\begin{bmatrix}0 \\ 1\end{bmatrix}.
\]
\end{tcolorbox}

对每个状态都可以定义一个相应的或有要求权，共有 $\Omega$ 个。这些或有要求权称为\textbf{Arrow-Debreu 证券}；由它们的完全集合构成的证券市场称为\textbf{Arrow-Debreu 证券市场}。若按状态排列证券，支付矩阵（市场结构）为单位阵：
\[
    X^{A\text{-}D}=
    \begin{bmatrix}
        1 & \cdots & 0 & \cdots & 0 \\
        \vdots & \ddots & \vdots & \ddots & \vdots \\
        0 & \cdots & 1 & \cdots & 0 \\
        \vdots & \ddots & \vdots & \ddots & \vdots \\
        0 & \cdots & 0 & \cdots & 1
    \end{bmatrix}
    =I.
\]
\end{tcolorbox}

\begin{tcolorbox}[knowledge,title=3.2 状态价格]
记 $\phi_{\omega}$ 为 0 期购买、在状态 $\omega$ 才支付 $1$ 单位消费品的状态 $\omega$ 或有要求权的价格（以 0 期消费品计量）。\textcolor{red}{\textbf{因为状态 $\omega$ 或有要求权在 1 期只有当状态为 $\omega$ 时才支付 $1$ 单位的消费品}}，$\phi_{\omega}$ 也称为\textbf{状态价格}（state price）。所有状态价格组成向量
\[
    \phi=\begin{bmatrix}\phi_1 \\ \vdots \\ \phi_{\omega} \\ \vdots \\ \phi_{\Omega}\end{bmatrix}.
\]
在无套利（没有免费的午餐）下，状态价格严格为正：
\[
    \phi\gg 0.
\]
若 $\phi_\omega\le0$，买入一份该状态证券，无须正的初始支出，未来支付非负且在状态 $\omega$ 为正，形成套利（只考虑可能发生的状态）。状态价格不是发生概率，也不必加总为 $1$；它还反映时间价值和该状态下资源的稀缺性。在无摩擦、一价定律成立时，支付为 $x$ 的可复制证券价格为 $\phi^{\mathsf T}x$。
\end{tcolorbox}

\begin{tcolorbox}[knowledge,title=3.3 市场的完全性]
在 Arrow-Debreu 市场中，1 期的消费计划可写成 $x=\begin{bmatrix}c_{11} \\ \vdots \\ c_{1\omega} \\ \vdots \\ c_{1\Omega}\end{bmatrix}$。持有状态证券组合 $\theta=\begin{bmatrix}c_{11} \\ \vdots \\ c_{1\omega} \\ \vdots \\ c_{1\Omega}\end{bmatrix}$ 即可复制该消费计划：
\[
    \begin{aligned}
        \underbrace{X\theta}_{\substack{\text{组合在 1 期各状态下的总支付}\\\text{（一般：任意证券组合的收益）}}}
        =\underbrace{X^{A\text{-}D}\begin{bmatrix}c_{11} \\ \vdots \\ c_{1\Omega}\end{bmatrix}}_{\substack{\theta\text{：状态 }\omega\text{ 阿罗证券持有 }c_{1\omega}\text{ 份}\\\text{研究阿罗市场，此处 }X=X^{A\text{-}D}}}
        =\underbrace{\begin{bmatrix}c_{11} \\ \vdots \\ c_{1\Omega}\end{bmatrix}}_{X^{A\text{-}D}=I\text{，故状态 }\omega\text{ 下收入 }c_{1\omega}}
        =\underbrace{x}_{\text{欲实现的 1 期消费计划}}.
    \end{aligned}
\]
该组合在 0 期的成本为
\[
    \underbrace{\phi^{\mathsf T}\theta}_{\text{复制该计划在 0 期支付的总成本}}
    =\underbrace{\sum_{\omega\in\Omega}\phi_{\omega}c_{1\omega}}_{\text{各状态消费 }\times\text{ 对应状态价格}}.
\]

\textcolor{red}{是否 \textbf{完全}，取决于市场能否以 \textbf{有限成本}为任意 \textbf{有限消费计划} \textbf{融资}}。

\begin{tcolorbox}[knowledge]
\textbf{定义 3.2}\quad 如果市场中的任一有限消费计划都可以通过有限成本的可交易证券的组合来融资，那么我们就称这个证券市场是\textbf{完全的}（complete）。即，任给消费计划 $c=\begin{bmatrix}c_{11} \\ \vdots \\ c_{1\Omega}\end{bmatrix}$，存在证券组合 $\theta\in\mathrm R^N$ 使得 $X\theta=c$ 有解。
\end{tcolorbox}



\begin{tcolorbox}[example]
\textbf{例 3.2}\quad 在 Lucas 树经济中，若市场上只有一只无风险债券，其支付为 $\begin{bmatrix}1 \\ 1\end{bmatrix}$。参与者若想要消费计划 $\begin{bmatrix}2 \\ 1\end{bmatrix}$（在状态 $a$ 消费 $2$、在状态 $b$ 消费 $1$），无法仅靠该债券实现，市场\textbf{不完备}。
\end{tcolorbox}


\begin{enumerate}
    \item 在完全市场中，任意消费计划 $x=[c_{11},\dots,c_{1\Omega}]$ 都是状态证券的线性组合；投资者可通过买卖状态证券在不同状态间调整消费。
    \item 完全市场中，消费者有最高的灵活度，可以实现最优的风险配置。
    \item 在相同状态空间、无摩擦且允许自由组合时，完全市场都能生成全部状态支付；比较价格与均衡配置，还须保持对应的定价关系、偏好和禀赋，不能把支付空间等价理解为不同经济有相同价格。
    \item 完全不表示每个人能负担任何消费计划：支付能被复制，与个人财富足以支付复制成本，是两个问题。
    \item \textbf{金融创新}的价值之一，正是通过增加可交易证券使市场更完全，从而更好地实现风险分担与效率。
\end{enumerate}

\begin{tcolorbox}[example,title=PPT 补充：三组支付矩阵的完全性]
\[
 X_1=\begin{bmatrix}1&3\\2&4\end{bmatrix},\quad
 X_2=\begin{bmatrix}1&2\\2&4\end{bmatrix},\quad
 X_3=\begin{bmatrix}1&0&1\\0&1&1\\1&1&1\end{bmatrix}.
\]
$\det X_1=-2\ne0$，两状态的任意支付都可复制；$X_2$ 第二列是第一列的两倍，秩为 $1$，两只证券仍不能生成两状态全部支付。$\det X_3=-1\ne0$，所以第三个市场完全，虽没有直接列出全部阿罗证券，也能复制它们。一般判据为 $\mathrm{rank}(X)=\Omega$；两只证券面对三个状态时，秩最多为 $2$，必不完全。
\end{tcolorbox}

\textbf{为什么现实中仍有未被交易的风险？}PPT 以保险说明：\textbf{道德风险}是投保后的行为激励改变，\textbf{逆向选择}是投保前风险类型的信息不对称；可验证性不足、交易和执行成本也会使某些状态合约缺失。大量独立风险有利于分散，但不能据此断言系统性风险绝不能承保或交易；分散困难与支付不可交易是不同问题。

\end{tcolorbox}

\begin{tcolorbox}[knowledge,title=3.4 参与者的优化]
考虑单个参与者的效用优化问题。在 Arrow-Debreu 市场中，设参与者禀赋为 $e$、效用函数为 $U$。若 1 期禀赋以状态证券组合 $\bar{\theta}=\begin{bmatrix}e_{11} \\ \vdots \\ e_{1\Omega}\end{bmatrix}$ 持有（初始持有），则 1 期禀赋的支付为
\[
    X^{A\text{-}D}\bar{\theta}=\bar{\theta}=e_1=
    \begin{bmatrix}e_{11} \\ \vdots \\ e_{1\Omega}\end{bmatrix}.
\]

上式的 $\bar\theta=e_1$ 是\textbf{复制禀赋}的初始组合；交易后复制目标消费的总持仓应记作 $\theta=c_1$，两者不能混同。参与者的\textbf{财富}为禀赋在 0 期的市场价值：
\[
 w=e_0+\phi^{\mathsf T}e_1.
 \tag{3.2}
\]
也可以直接从净交易推预算：令 $z=\theta-\bar\theta=c_1-e_1$ 为状态证券净买入量，则
\[
 c_0+\phi^{\mathsf T}z\le e_0
 \iff c_0+\phi^{\mathsf T}c_1\le e_0+\phi^{\mathsf T}e_1=w.
 \tag{3.3}
\]
因此“先把禀赋换成财富再选择消费”是等价分析方法，不要求参与者在现实交易中先卖光全部禀赋。

记扩展价格向量 $\hat{\phi}=\begin{bmatrix}1 \\ \phi\end{bmatrix}$（0 期消费品价格为 $1$，即计价物），则预算约束可写为 $\hat{\phi}^{\mathsf T}c\le\hat{\phi}^{\mathsf T}e$，或 $\hat{\phi}^{\mathsf T}(e-c)\ge 0$。

参与者的优化问题为（在预算内最大化效用，且消费非负）
\begin{align}
    \max_{c\in C}\ &U(c) \tag{3.4}\\
    \mathrm{s.t.}\quad &\hat{\phi}^{\mathsf T}(e-c)\ge 0, \notag\\
    &c\ge 0. \notag
\end{align}

\begin{tcolorbox}[knowledge]
\textbf{定理 3.1}\quad 设 $C=\mathrm{R}^{1+\Omega}_{+}$，$\hat\phi\gg0$、$0\le w<\infty$，$U(\cdot)$ 在 $C$ 上连续，则问题 (3.4) 有解。
\end{tcolorbox}

\textbf{证明（补齐 PPT 的有界性步骤）：}预算集
\[
 B=\{c\ge0:\hat\phi^{\mathsf T}c\le w\}
\]
非空（含零消费），且每个分量均满足 $0\le c_i\le w/\phi_i\le w/\min_j\phi_j$，其中 $\phi_0=1$。所以可行消费不会无限增大。若可行序列 $c^{(n)}\to c$，取极限仍有 $c\ge0$、$\hat\phi^{\mathsf T}c\le w$，故极限仍可行。连续函数在这样的有限维有界闭集上取得最大值。此处先用不等式预算集；只有加入不满足性后，才可推出最优时预算取等号。

以下另假设 $U$ 可微且各消费分量的边际效用严格为正。\textbf{这是本章使用的强化单调性条件}；一般的不满足性本身并不等价于处处 $DU\gg0$。记 $\partial_0 U$ 为 0 期消费的边际效用，$\partial_{\omega}U$ 为状态 $\omega$ 下 1 期消费的边际效用，则
\[
    \partial_0 U>0,\qquad \partial_{\omega}U>0,\quad\forall\omega\in\Omega,
    \tag{3.5}
\]
边际效用向量为 $DU=\begin{bmatrix}\partial_0 U \\ \partial_1 U \\ \vdots \\ \partial_{\Omega}U\end{bmatrix}$，后续一阶分析采用 $DU\gg0$。

\begin{tcolorbox}[knowledge]
\textbf{定理 3.2}\quad 设 $w>0$、$U$ 可微且 $DU\gg 0$。则 (3.4) （参与者的效用优化问题，在预算内最大化效用，且消费非负）的解满足：存在 $\lambda>0$、$\mu_i\ge 0$（$i=0,1,\dots,\Omega$）使得
\begin{align}
    \partial_i U&=\lambda\phi_i-\mu_i,\quad i=0,1,\dots,\Omega\quad(\phi_0=1),\tag{3.6a}\\
    \hat{\phi}^{\mathsf T}(e-c)&=0,\tag{3.6b}\\
    \mu_i c_i&=0,\quad i=0,1,\dots,\Omega.\tag{3.6c}
\end{align}


定理 3.2 说的是：在问题 (3.4)（预算约束 $+$ 消费非负）下，若 $DU\gg0$，则最优消费满足 KKT 条件 (3.6a)–(3.6c)（不等式约束下的一阶必要条件；若 $U$ 凹且可行集凸，亦为充分条件；不能把“偏好凸”误写成“效用函数凸”）。含义简述如下：
\begin{itemize}
    \item \textbf{(3.6a)} \textbf{边际效用定价：}$\partial_i U=\lambda\phi_i-\mu_i$。内点（$c_i>0$）时 $\mu_i=0$，边际效用与扩展价格成比例；若某期/某状态消费为 $0$，则 $\partial_iU/\phi_i\le\lambda$；该分量的单位支出边际效用不高于正消费分量，非负约束阻止继续减少它。
    \item \textbf{(3.6b)} \textbf{预算紧绑：}按 $\hat{\phi}$ 计价时花光财富，$\hat{\phi}^{\mathsf T}c=\hat{\phi}^{\mathsf T}e$（在 $DU\gg 0$ 下不会故意少花）。
    \item \textbf{(3.6c)} \textbf{互补松弛：}$\mu_i c_i=0$。要么该分量消费为正、对应乘子为零；要么消费为零、允许 $\mu_i>0$（角点解）。
\end{itemize}
\end{tcolorbox}

对问题 (3.4) 构造拉格朗日函数（$\lambda$ 为预算约束 $\hat{\phi}^{\mathsf T}(e-c)\ge 0$ 的乘子，$\mu_i$ 为 $c_i\ge 0$ 的乘子）：
\[
    \mathcal{L}
    =U(c)+\lambda\,\hat{\phi}^{\mathsf T}(e-c)+\sum_{i=0}^{\Omega}\mu_i c_i.
\]
对 $c_i$ 求偏导并令其为零，即得 (3.6a)；与 (3.6b)、(3.6c) 联立即为定理 3.2。\textbf{符号校正：}对本处“最大化、$c_i\ge0$、$\mu_i\ge0$”的约定，拉格朗日函数应加 $\mu_i c_i$，故边际条件应为 $\partial_iU=\lambda\phi_i-\mu_i$；PPT 所列加号应相应校正。最优时至少一个消费分量为正，故由该分量的一阶条件可知 $\lambda>0$。对参与者 $k$ 将 $U,c,e$ 换为 $U_k,c_k,e_k$ 即可。

\textbf{角点检验：}取 $U=2c_0+c_1$、价格均为 $1$、财富为 $1$，最优消费为 $(1,0)$。此时 $\lambda=2$、$\mu_0=0$、$\mu_1=1$，满足 $\partial_1U=1=2-1$；若误用加号便无法用非负乘子表达这个最优角点。

对参与者 $k$，\textbf{内点解}（$c_k\gg 0$，故 $\mu_i=0$）时由 (3.6a) 可写成 $DU_k(c_k)=\lambda_k\hat{\phi}$，即
\begin{align}
    \underbrace{\frac{\partial U_k(c_k)}{\partial c_k(0)}=\lambda_k}_{\substack{\text{对第 0 期消费求偏导}\\(\phi_0=1)}},
    &\quad
    \underbrace{\frac{\partial U_k(c_k)}{\partial c_k(\omega)}=\lambda_k\phi_{\omega}}_{\substack{\text{对第 1 期各状态消费求偏导}}},\quad
    \forall\omega\in\Omega. \tag{3.7}
\end{align}

等价地，\textbf{跨期边际替代率}等于状态价格：
\[
 \frac{\partial_\omega U_k(c_k)}{\partial_0 U_k(c_k)}=\phi_\omega
 \quad\Longleftrightarrow\quad
 \partial_0U_k(c_k)=\frac1{\phi_\omega}\partial_\omega U_k(c_k).
 \tag{3.8}
\]
\textbf{用预算内的小调整理解：}将状态 $\omega$ 的消费增加 $\mathrm dc_\omega$，须减少当前消费 $\mathrm dc_0=-\phi_\omega\mathrm dc_\omega$，效用的一阶变化为
\[
 \mathrm dU_k=(\partial_\omega U_k-\phi_\omega\partial_0U_k)\mathrm dc_\omega.
\]
内点最优要求这个系数为零；若为正，增加该状态消费有利，若为负，反向调整有利。这是\textbf{边际}上的无差异，不表示任意幅度转移消费都保持效用不变。$\phi_\omega$ 衡量为未来该状态消费 $1$ 单位所放弃的当前消费量，$1/\phi_\omega$ 是反向交换的数量比，不能把它解释成“效用单位的交换率”。

\begin{tcolorbox}[example,title=PPT 补充：例 3.3 的对数消费规则]
对 $U=\log c_0+\tfrac12(\log c_{1a}+\log c_{1b})$，假设 $w>0$、$\phi_a,\phi_b>0$。消费趋近零时对数效用趋于负无穷，最优解在正消费区域。一阶条件给出
\[
 c_0=\frac1\lambda,\qquad c_{1a}=\frac1{2\lambda\phi_a},\qquad
 c_{1b}=\frac1{2\lambda\phi_b}.
\]
代入预算得 $1/\lambda+1/(2\lambda)+1/(2\lambda)=w$，于是
\[
 \lambda=\frac2w,\qquad c_0^*=\frac w2,\quad
 c_{1a}^*=\frac w{4\phi_a},\quad c_{1b}^*=\frac w{4\phi_b}.
\]
一半财富用于现在，每个未来状态分得四分之一的\textbf{支出}，不是相同的消费数量。价格对消费的反比关系是在财富 $w$ 给定下的结论；当初始禀赋含未来消费时，价格变化也会改变 $w$。PPT 第 18 页末项的分母应为 $\phi_b$。完整题目与解答保留在本章后面的例 3.3。
\end{tcolorbox}
\end{tcolorbox}


\begin{tcolorbox}[knowledge,title=3.5 市场均衡]
给定状态价格，Arrow-Debreu 经济中的\textbf{均衡}是一组状态价格 $\{\phi_{\omega},\,\omega\in\Omega\}$，使得：
\begin{enumerate}
    \item \textbf{给定价格，每个参与者最优化：}对参与者 $k$，
    \[
        c_k(e_k,\phi)=\operatorname*{Argmax}\,U_k(c_k)
    \]
    满足 $c_{k,0}+\phi^{\mathsf T}c_{k,1}=e_{k,0}+\phi^{\mathsf T}e_{k,1}$ 且 $c_k\ge 0$；
    \item \textbf{商品市场出清}：
    \[
        \sum_{k=1}^{K}c_{k,0}(e_k,\phi)=\sum_{k=1}^{K}e_{k,0},\qquad
        \sum_{k=1}^{K}c_{k,1}(e_k,\phi)=\sum_{k=1}^{K}e_{k,1}.
    \]
\end{enumerate}
令总超额需求为 $z=\sum_k(c_k-e_k)$。将花光预算的等式加总，得到 $\hat\phi^{\mathsf T}z=0$，这就是 Walras 法则。在价格严格为正时，若其中 $\Omega$ 个分量为零，剩余一个也必须为零。因此可以用 $1+\Omega$ 个出清条件中的 $\Omega$ 个求解价格，最后一式用于核验。出清按每个未来状态分别成立，不能仅检查概率加权后的总量。

\end{tcolorbox}

\begin{tcolorbox}[knowledge,title=瓦尔拉斯均衡]
\textbf{瓦尔拉斯均衡}（Walrasian equilibrium，竞争均衡）指一组价格与配置，使得：
\begin{enumerate}
    \item 每个参与者在给定价格下最大化效用（或厂商最大化利润），且为\textbf{价格接受者}；
    \item 每个市场（每种商品）上\textbf{总需求等于总供给}（市场出清）。
\end{enumerate}
在本章中，商品是 $1+\Omega$ 种\textbf{或有商品}：0 期消费品（计价物，价格 $1$）与 1 期各状态 $\omega$ 下的消费品（价格 $\phi_\omega$，即状态 $\omega$ 阿罗证券的价格）。上文两条均衡条件正是瓦尔拉斯均衡在 Arrow-Debreu 证券市场经济中的写法，故也称\textbf{Arrow-Debreu 均衡}。第 2.5 节中的 Walras 法则说明：在标准预算下，$1+\Omega$ 个商品市场里若 $1+\Omega-1$ 个已出清，则剩余一个自动出清；实物总禀赋 $\sum_k e_k$ 仅被重新分配，不会凭空增减。
\end{tcolorbox}

\begin{tcolorbox}[knowledge,title=3.5 市场均衡（存在性与边际条件）]
\begin{tcolorbox}[knowledge]
\textbf{定理 3.3}\quad 对定义 2.5 （两期、$\Omega$ 状态、单一不可储存商品、$K$ 名同质信息参与者、禀赋与满足公理 1–3 的偏好（不满足性、连续性、凸性）、无摩擦证券市场 $X$）中的经济，Arrow-Debreu 证券市场存在均衡。
\end{tcolorbox}

在可微且各参与者最优消费均为内点时，所有人面对相同价格，因而在任一状态 $\omega$ 上\textcolor{red}{\textbf{相对边际效用都相等}}：
\[
    \frac{\partial_{\omega}U_k(c_k)}{\partial_0 U_k(c_k)}
    =\phi_{\omega}
    =\frac{\partial_{\omega}U_{k'}(c_{k'})}{\partial_0 U_{k'}(c_{k'})},
    \quad\forall k,k',\,\omega\in\Omega.
    \tag{3.9}
\]
$\phi_{\omega}$ 可理解为：为在状态 $\omega$ 多获得 1 单位 1 期消费，0 期须放弃的消费量——市场上关于跨期、跨状态资源配置的\textbf{共识价格}。

\end{tcolorbox}

\begin{tcolorbox}[example,title=PPT 补充：例 3.4 从个人需求到均衡价格]
两人禀赋为 $(100;0,0)$、$(0;200,50)$，具有例 3.3 的相同对数效用。财富为 $w_1=100$、$w_2=200\phi_a+50\phi_b$。记 $W=w_1+w_2$，三项市场出清为
\[
 \frac W2=100,\qquad \frac W{4\phi_a}=200,\qquad
 \frac W{4\phi_b}=50.
\]
先得 $W=200$，再得 $\phi_a=1/4$、$\phi_b=1$，代回财富恰有 $w_1=w_2=100$。两人最优消费均为 $(50;100,25)$，逐状态加总正好等于总禀赋。两状态概率相同却价格不同：状态 $b$ 的总消费更少，同样一单位消费更稀缺，状态价格更高。均衡公式 $\phi_\omega=c_{k,0}/(2c_{k,\omega})$ 也给出相同结果。详细推导保留在后面的例 3.4。
\end{tcolorbox}



\begin{tcolorbox}[knowledge,title=3.5 市场均衡（续）]
\textbf{不完全市场为何未必有逐状态 MRS 相等？}若只能买卖一般支付矩阵 $X$ 的证券，内点、无交易限制时，最优条件只要求
\[
 X^{\mathsf T}m_k=S,\qquad
 (m_k)_\omega=\frac{\partial_\omega U_k}{\partial_0U_k}.
\]
所以 $X^{\mathsf T}(m_k-m_{k'})=0$。只有当 $X$ 行满秩，才能推出 $m_k=m_{k'}$；不完全市场可能留下未被交易消除的状态间估值差异。交易成本和持仓约束还可能把相应等式改成不等式。

\textbf{金融创新如何增加完全性？}关键是新支付列不在原有列空间中，而非名称或数量增加。PPT 中的巨灾债券、通胀挂钩债券和信用衍生品，是针对特定灾害、购买力或违约支付的例子；它们能否增加秩须逐一检查，不能仅凭引入一种产品就断言市场完全。
可行配置满足 $\sum_{k=1}^{K}c_k=\sum_{k=1}^{K}e_k$。对均衡配置，将各参与者预算约束在 $\hat{\phi}$ 下加总，得到
\[
    \sum_{k=1}^{K}\hat{\phi}^{\mathsf T}c_k=\sum_{k=1}^{K}\hat{\phi}^{\mathsf T}e_k.
    \tag{3.10}
\]

\begin{tcolorbox}[knowledge]
\textbf{定理 3.4}\quad 对存在 Arrow-Debreu 证券市场的经济（定义 2.5），均衡达到的资源配置是\textbf{Pareto 最优的}。
\end{tcolorbox}

\textbf{证明（按 PPT 的预算比较法）：}记均衡配置为 $\{c_k\}$，财富为 $w_k=\hat\phi^{\mathsf T}e_k=\hat\phi^{\mathsf T}c_k$。反设存在可行配置 $\{a_k\}$，所有人弱偏好 $a_k$，至少一人 $j$ 严格偏好。

先证弱改善者必须满足 $\hat\phi^{\mathsf T}a_k\ge w_k$。若其成本小于 $w_k$，剩余预算可购买一个各分量为正的小消费增量 $d$，在不满足性下 $a_k+d$ 严格优于 $a_k$，从而严格优于原最优消费 $c_k$，且仍在预算内，矛盾。对严格改善者 $j$，则必有 $\hat\phi^{\mathsf T}a_j>w_j$，否则 $a_j$ 本来就买得起，与 $c_j$ 最优矛盾。加总得
\[
 \hat\phi^{\mathsf T}\sum_k a_k>\sum_k w_k
 =\hat\phi^{\mathsf T}\sum_k e_k,
\]
却与可行性 $\sum_k a_k=\sum_k e_k$ 矛盾。故均衡配置 Pareto 最优。这里比较的是\textbf{偏好改善}，不要求每个人的每个消费分量都增加；并且需要完全市场使相关可行消费计划能按同一状态价格购买。

定理 3.4 即\textbf{福利经济学第一定理}：完全、无摩擦的证券市场使参与者达到\textcolor{red}{\textbf{资源在不同时期和不同状态下的 Pareto 最优配置}}，为证券市场的资源配置效率提供了理论基础（亚当·斯密“看不见的手”的现代表述）。

\begin{tcolorbox}[knowledge,title=福利经济学两大定理]
\textbf{第一定理}（定理 3.4）：在完全、无摩擦的竞争市场中，任一均衡配置都是 Pareto 最优的。

\textbf{第二定理}（福利经济学第二定理）：在消费集凸、偏好凸（满足公理 3）且市场完全（如 Arrow-Debreu 市场）等条件下，\textbf{任一} Pareto 最优配置都可以成为某个竞争均衡的配置——只需在交易开始前对禀赋作适当的\textbf{一次性总付转移}（lump-sum redistribution），使参与者的初始禀赋变为 $e_k'$，再令他们在状态价格 $\phi$ 下自由交易即可达到该配置。

第二定理的含义：在一定条件下，\textbf{效率}问题与\textbf{分配}问题可以分开处理——先确定有效的总资源配置，再通过不扭曲价格信号的转移支付来实现合意的福利分配。（现实中转移支付往往受信息、政治与再分配成本限制，故第二定理更多是规范基准。）
\end{tcolorbox}


\end{tcolorbox}

\begin{tcolorbox}[knowledge,title=本章小结]
\begin{enumerate}
    \item 状态 $\omega$ 或有要求权在状态 $\omega$ 支付 $1$、其他状态支付 $0$；由全部状态或有证券构成的市场为 Arrow-Debreu 证券市场。
    \item 任一有限消费计划可用有限成本证券组合融资时，市场完全；Arrow-Debreu 市场是完全市场。
    \item 均衡存在，且各参与者最优化、市场出清；在可微、内点条件下，均衡时相对边际效用（跨期边际替代率）相等。
    \item \textbf{福利经济学第一定理：}Arrow-Debreu 均衡配置 Pareto 最优。
    \item \textbf{福利经济学第二定理：}在凸性与完全市场下，任一 Pareto 最优配置可通过一次性总付转移与市场交易实现。
\end{enumerate}
\end{tcolorbox}
\clearpage
\fi

% !TeX root = ../../金融经济学_王江_讲义及习题.tex
% 03_Arrow-Debreu经济 课后习题（由章节文件拆分）

    \begin{problem}
        \problemno{例 3.3} 考虑“Lucas 树经济”，1 期有两个概率相等的状态 $a$ 和 $b$。假设 1 期的两个可能状态的状态价格分别为 $\phi_{a}$ 和 $\phi_{b}$。考虑一个参与者，他的禀赋为 $(e_{0};e_{1a};e_{1b})$。他具有如下对数形式的效用函数：
        \[
            U(c_{0},c_{1a},c_{1b})=\log c_{0}+\tfrac12(\log c_{1a}+\log c_{1b})
        \]
        
        他的最优消费/组合选择是什么？
        \end{problem}
        \begin{solution}
        给定状态价格和他的禀赋，他的总财富是 $w=e_{0}+\phi_{a}e_{1a}+\phi_{b}e_{1b}$。他的最优化问题是
        \[
            \begin{aligned}
                \max_{\{c_{0},c_{1a},c_{1b}\}}\ &\log c_{0}+\tfrac12(\log c_{1a}+\log c_{1b})\\
                \mathrm{s.t.}\quad &w-(c_{0}+\phi_{a}c_{1a}+\phi_{b}c_{1b})=0\\
                &c_{0},c_{1a},c_{1b}>0
            \end{aligned}
        \]

        假设 $w>0$、$\phi_a,\phi_b>0$。对数效用的消费域为 $c_0,c_{1a},c_{1b}>0$；最优解在内部，可直接使用等式预算的拉格朗日函数
        \[
            \mathcal{L}=\log c_{0}+\tfrac12(\log c_{1a}+\log c_{1b})+\lambda\bigl(w-c_0-\phi_a c_{1a}-\phi_b c_{1b}\bigr).
        \]
        一阶条件为
        \[
          \frac1{c_0}=\lambda,\qquad
          \frac1{2c_{1a}}=\lambda\phi_a,\qquad
          \frac1{2c_{1b}}=\lambda\phi_b,\qquad
          c_0+\phi_a c_{1a}+\phi_b c_{1b}=w.
        \]
        当消费趋近 $0$ 时，对数效用趋于负无穷、边际效用趋于正无穷，因此最优消费在每一时期、每一状态都严格为正。由上述一阶条件得到最优消费：
        \[
            c_{0}^{*}=\frac{1}{\lambda},\qquad
            c_{1a}^{*}=\frac{1}{\lambda}\frac{1}{2\phi_{a}},\qquad
            c_{1b}^{*}=\frac{1}{\lambda}\frac{1}{2\phi_{b}}
        \]

        把 $c^{*}$ 代入预算约束，我们可以得到 $\lambda$ 的解：
        \[
            \frac{1}{\lambda}+\frac{1}{2\lambda}+\frac{1}{2\lambda}=w\qquad\Rightarrow\qquad\lambda=\frac{2}{w}
        \]

        最后，我们有
        \[
            c_{0}^{*}=\frac12 w,\qquad
            c_{1a}^{*}=\frac14\frac{w}{\phi_{a}},\qquad
            c_{1b}^{*}=\frac14\frac{w}{\phi_{b}}
        \]

        可以看出，参与者把一半财富用作现在的消费，把另外一半财富作为未来的消费。在财富给定时，某一状态下的消费与对应状态价格成反比；若价格变化也改变禀赋市值，还须计入财富效应。这个结果很直观。状态价格高的状态下的消费更昂贵。结果，参与者在这些状态下选择较低的消费。
        \end{solution}

    \begin{problem}
        \problemno{例 3.4} 考虑一个经济，在 1 期有两个概率相等的状态 $a$ 和 $b$。经济中有参与者 1 和 2，他们具有的禀赋分别为：
        \begin{center}
        \WJEventTree[left={$e_{1}{:}100$}]{0,0}
        \qquad
        \WJEventTree[left={$e_{2}{:}0$}]{200,50}
        \end{center}
        两个参与者都具有如下形式的对数效用函数：
        \[
            U(c)=\log c_{0}+\tfrac12(\log c_{a}+\log c_{b})
        \]
        \end{problem}
        \begin{solution}
        在市场上存在一组完全的状态或有证券可以交易。因为有两个状态，因而只有两个状态或有证券。
        
        现在我们开始分析这个经济的均衡。从给定交易证券价格下参与者的最优化问题开始。记 $\phi=\begin{bmatrix}\phi_{a} \\ \phi_{b}\end{bmatrix}$ 为状态价格（向量），即两个状态或有证券的价格。我们可以定义每个参与者的财富为 $w=\hat{\phi}^{\top}e$，这里 $\hat{\phi}=\begin{bmatrix}1 \\ \phi\end{bmatrix}$，而 $e$ 是他的禀赋。这时，最优化问题变成了：
        \[
            \begin{aligned}
                \max_{c}\ &\log c_{0}+\tfrac12(\log c_{a}+\log c_{b})\\
                \mathrm{s.t.}\quad &c_{0}+\phi_{a}c_{a}+\phi_{b}c_{b}=w
            \end{aligned}
        \]

        由例 3.3，最优消费为
        \[
            c_{k,0}^{*}=\tfrac12 w_{k},\qquad
            c_{k,a}^{*}=\tfrac14\frac{w_{k}}{\phi_{a}},\qquad
            c_{k,b}^{*}=\tfrac14\frac{w_{k}}{\phi_{b}},\qquad k=1,2
        \]
        在状态价格 $\phi_{a},\phi_{b}$ 下，参与者按禀赋市值计的财富为 $w_{k}=e_{k,0}+\phi_{a}e_{k,a}+\phi_{b}e_{k,b}$（0 期消费品价格为 $1$）。代入禀赋 $(100;0,0)$ 与 $(0;200,50)$：
        \[
            w_{1}=100+0\cdot\phi_{a}+0\cdot\phi_{b}=100,\qquad
            w_{2}=0+200\phi_{a}+50\phi_{b}=200\phi_{a}+50\phi_{b}.
        \]

        \textbf{0 期商品市场出清}（两人 0 期消费之和等于经济 0 期总禀赋 $100$）。由例 3.3，$c_{k,0}^{*}=\tfrac12 w_{k}$，故
        \[
            c_{1,0}^{*}+c_{2,0}^{*}=\tfrac12 w_{1}+\tfrac12 w_{2}=\tfrac12(w_{1}+w_{2})=100
            \quad\Longrightarrow\quad w_{1}+w_{2}=200.
        \]

        把 $w_{1},w_{2}$ 代入 $w_{1}+w_{2}=200$：
        \[
            100+(200\phi_{a}+50\phi_{b})=200.
        \]

        移项得 $200\phi_{a}+50\phi_{b}=100$；两边同除以 $50$，得到与状态价格有关的第一个方程
        \[
            4\phi_{a}+\phi_{b}=2.
        \]

        \textbf{1 期·状态 $a$ 商品市场出清}（状态 $a$ 下总消费 $200$）。由 $c_{k,a}^{*}=\tfrac14 w_{k}/\phi_{a}$，
        \[
            c_{1,a}^{*}+c_{2,a}^{*}=\tfrac{w_{1}}{4\phi_{a}}+\tfrac{w_{2}}{4\phi_{a}}
            =\frac{w_{1}+w_{2}}{4\phi_{a}}=200
            \quad\Longrightarrow\quad w_{1}+w_{2}=800\phi_{a}.
        \]

        与 $w_{1}+w_{2}=200$ 联立得 $\phi_{a}=\dfrac14$，从而 $\phi_{b}=1$。状态 $b$ 出清 $\dfrac{w_{1}+w_{2}}{4\phi_{b}}=50$ 自动成立。        此时 $w_{1}=w_{2}=100$，最优消费为
        \[
            \begin{aligned}
                c_{1,0}^{*}=c_{2,0}^{*}&=50\\
                c_{1,a}^{*}=c_{2,a}^{*}&=100\\
                c_{1,b}^{*}=c_{2,b}^{*}&=25
            \end{aligned}
        \]
        \end{solution}

\begin{problem}
\problemno{3.1*} 考虑如下经济，在 1 期有两个可能状态 $a$ 和 $b$：
\begin{center}
\WJEventTree{a,b}
\end{center}
\begin{enumerate}[label=(\alph*)]
    \item 描述所有 Arrow-Debreu 证券的支付向量。记这些证券的价格向量为 $\phi$。
    \item 考虑一个拥有如下禀赋的参与者：
    \begin{center}
    \WJEventTree[left={$0$}]{2,1}
    \end{center}
    把他的禀赋表示成 Arrow-Debreu 证券的组合。
    \item 计算他的金融财富。写出他的预算集。
    \item 假设参与者的效用函数如下：
    \[
        U(c_{0},c_{1a},c_{1b})=-e^{-c_{0}}-\tfrac12\bigl(e^{-c_{1a}}+e^{-c_{1b}}\bigr).
    \]
    不考虑消费的非负约束，写出他的优化问题。求解他的最优消费选择。
    \item 讨论他的消费如何依赖于 Arrow-Debreu 证券的价格向量 $\phi$。
    \item 证明在某些价格下，他（在某些时期/状态下）的消费可能是负的。
\end{enumerate}
\end{problem}
\begin{solution}
\begin{enumerate}[label=(\alph*)]
    \item Arrow-Debreu 证券的支付向量是 $X_{a}=\begin{bmatrix}1 \\ 0\end{bmatrix}$，$X_{b}=\begin{bmatrix}0 \\ 1\end{bmatrix}$；
    \item $2X_{a}+X_{b}$；
    \item 0 期禀赋为 $0$，1 期禀赋是 $2$ 单位状态 $a$ 或有要求与 $1$ 单位状态 $b$ 或有要求，金融财富为 $w=2\phi_{a}+\phi_{b}$。在消费集 $\mathrm{R}_{+}^{3}$ 上，预算集是
    \[
        \bigl\{c\in\mathrm{R}_{+}^{3}:c_{0}+\phi_{a}c_{1a}+\phi_{b}c_{1b}=w\bigr\}.
    \]
    \item 本问不施加 $c\ge 0$，在 $\phi_{a}>0$、$\phi_{b}>0$ 时于 $c\in\mathrm{R}^{3}$ 上求解
    \[
        \begin{aligned}
            \max_{\{c_{0},c_{1a},c_{1b}\}}\ &\bigl\{-e^{-c_{0}}-\tfrac12\bigl(e^{-c_{1a}}+e^{-c_{1b}}\bigr)\bigr\}\\
            \mathrm{s.t.}\quad &c_{0}+\phi_{a}c_{1a}+\phi_{b}c_{1b}=w=2\phi_{a}+\phi_{b}
        \end{aligned}
    \]
    拉格朗日函数为
    \[
        \mathcal{L}=-e^{-c_{0}}-\tfrac12\bigl(e^{-c_{1a}}+e^{-c_{1b}}\bigr)+\lambda\bigl(w-c_{0}-\phi_{a}c_{1a}-\phi_{b}c_{1b}\bigr)
    \]
    对消费求导，一阶条件为
    \[
        e^{-c_{0}}=\lambda,\qquad
        \tfrac12 e^{-c_{1a}}=\lambda\phi_{a},\qquad
        \tfrac12 e^{-c_{1b}}=\lambda\phi_{b}
    \]
    从而最优消费为
    \[
        c_{0}^{*}=-\log\lambda,\qquad
        c_{1a}^{*}=-\log(2\lambda\phi_{a}),\qquad
        c_{1b}^{*}=-\log(2\lambda\phi_{b})
    \]
    代入预算约束：
    \[
        -\log\lambda-\phi_{a}\log(2\lambda\phi_{a})-\phi_{b}\log(2\lambda\phi_{b})=w
    \]
    即 $-(1+\phi_{a}+\phi_{b})\log\lambda-\phi_{a}\log(2\phi_{a})-\phi_{b}\log(2\phi_{b})=w$。记 $S=1+\phi_{a}+\phi_{b}$，则
    \[
        \log\lambda=-\frac{w+\phi_{a}\log(2\phi_{a})+\phi_{b}\log(2\phi_{b})}{S}
    \]
    从而 $c_{0}^{*}=-\log\lambda$，且 $c_{1a}^{*}=-\log\lambda-\log(2\phi_{a})$、$c_{1b}^{*}=-\log\lambda-\log(2\phi_{b})$。代入 $w=2\phi_{a}+\phi_{b}$ 得到
    \[
        \begin{aligned}
            c_{0}^{*}&=\frac{2\phi_{a}+\phi_{b}+\phi_{a}\log(2\phi_{a})+\phi_{b}\log(2\phi_{b})}{1+\phi_{a}+\phi_{b}}\\
            c_{1a}^{*}&=\frac{2\phi_{a}+\phi_{b}+\phi_{b}\log(2\phi_{b})-\log(2\phi_{a})-\phi_{b}\log(2\phi_{a})}{1+\phi_{a}+\phi_{b}}\\
            c_{1b}^{*}&=\frac{2\phi_{a}+\phi_{b}+\phi_{a}\log(2\phi_{a})-\log(2\phi_{b})-\phi_{a}\log(2\phi_{b})}{1+\phi_{a}+\phi_{b}}
        \end{aligned}
    \]
    \item 由一阶条件，$c_{1a}^{*}=c_{0}^{*}-\log(2\phi_{a})$，$c_{1b}^{*}=c_{0}^{*}-\log(2\phi_{b})$。记
    \[
        N=2\phi_{a}+\phi_{b}+\phi_{a}\log(2\phi_{a})+\phi_{b}\log(2\phi_{b}),\qquad
        D=1+\phi_{a}+\phi_{b},
    \]
    则 $c_{0}^{*}=N/D$。直接求导得
    \[
        \frac{\partial c_{0}^{*}}{\partial\phi_{a}}
        =\frac{3+\phi_{a}+\log(2\phi_{a})+\phi_{b}\bigl(2+\log(\phi_{a}/\phi_{b})\bigr)}{D^{2}},
        \qquad
        \frac{\partial c_{1a}^{*}}{\partial\phi_{a}}
        =\frac{\partial c_{0}^{*}}{\partial\phi_{a}}-\frac{1}{\phi_{a}},
    \]
    并且 $\partial c_{1b}^{*}/\partial\phi_{a}=\partial c_{0}^{*}/\partial\phi_{a}$。$\phi_{a}$ 上升时，禀赋价值 $w=2\phi_{a}+\phi_{b}$ 增加，同时状态 $a$ 的消费相对变贵；二者的相对强弱取决于价格，符号由分子决定。若 $2\phi_{a}\ge e^{-1}$ 且 $\phi_{a}\ge\phi_{b}$，分子为正，$c_{0}^{*}$ 与 $c_{1b}^{*}$ 上升；$c_{1a}^{*}$ 上升还要求 $\partial c_{0}^{*}/\partial\phi_{a}>1/\phi_{a}$。例如 $\phi_{a}=\phi_{b}=0.1$ 时分子为 $3.3+\log 0.2>0$，但导数约为 $1.17<10=1/\phi_{a}$，故 $c_{1a}^{*}$ 下降。若 $\phi_{a}=0.01$、$\phi_{b}=1$，分子为 $5.01+\log 0.02+\log 0.01<0$，$c_{0}^{*}$ 与 $c_{1b}^{*}$ 随 $\phi_{a}$ 上升而下降。

    对 $\phi_{b}$ 同样有
    \[
        \frac{\partial c_{0}^{*}}{\partial\phi_{b}}
        =\frac{2+\phi_{b}+\log(2\phi_{b})+\phi_{a}\log(\phi_{b}/\phi_{a})}{D^{2}},
        \qquad
        \frac{\partial c_{1b}^{*}}{\partial\phi_{b}}
        =\frac{\partial c_{0}^{*}}{\partial\phi_{b}}-\frac{1}{\phi_{b}},
    \]
    且 $\partial c_{1a}^{*}/\partial\phi_{b}=\partial c_{0}^{*}/\partial\phi_{b}$。若 $2\phi_{b}\ge e^{-1}$ 且 $\phi_{b}\ge\phi_{a}$，则 $c_{0}^{*}$ 与 $c_{1a}^{*}$ 不减。
    \item 当 $\phi_{a}=\phi_{b}=0.1$ 时，$w=0.3$，$D=1.2$，
    \[
        c_{0}^{*}=\frac{0.3+0.2\log(0.2)}{1.2}\approx -0.0182<0.
    \]
    该计划满足预算等式，但 $c_{0}^{*}<0$，因而不属于 (c) 的预算集。
\end{enumerate}
\end{solution}

\begin{problem}
\problemno{3.2} 考虑一个在 1 期只有一个可能状态的经济。（在这种情况下不存在不确定性。）参与者 1 的 0 期禀赋为 100 而 1 期禀赋为 1，即他的禀赋向量为 $\begin{bmatrix}100 \\ 1\end{bmatrix}$。他的偏好可以表示成如下形式：
\[
    U(c_{0},c_{1})=\log c_{0}+\rho\log c_{1}.
\]
有一只可交易证券，它的 0 期价格为 1、1 期支付为 $1+r_{F}$。这里，$r_{F}$ 是利率。
\begin{enumerate}[label=(\alph*)]
    \item 假设利率 $r_{F}$ 是给定的，导出参与者对证券的需求。
    \item 假设参与者 1 是经济中的唯一参与者。描述市场出清条件。
    \item 求解均衡利率。
    \item 均衡利率如何依赖于偏好参数 $\rho$？解释所得到的结论。
\end{enumerate}
\end{problem}
\begin{solution}
\begin{enumerate}[label=(\alph*)]
    \item 先不考虑消费的非负性，求解参与者的证券需求。此时参与者的优化问题变成了
    \[
        \max_{\theta}\bigl\{\log(100-\theta)+\rho\log\bigl(1+\theta(1+r_{F})\bigr)\bigr\}
    \]
    这与练习 2.7 是同一规划。一阶条件
    \[
        \frac{1}{100-\theta}=\rho\frac{1+r_{F}}{1+\theta(1+r_{F})}
    \]
    解得证券需求
    \[
        \theta^{*}=\frac{100\rho(1+r_{F})-1}{(1+\rho)(1+r_{F})}
    \]
    最优消费为
    \[
        c_{0}^{*}=\frac{100(1+r_{F})+1}{(1+\rho)(1+r_{F})},\quad
        c_{1}^{*}=\frac{\rho\bigl(100(1+r_{F})+1\bigr)}{1+\rho}
    \]
    由于最优消费都是非负的，因此上面的求解过程没有问题。
    \item 由于市场上只有一个参与者，因而市场出清条件是 $\theta^{*}=0$；
    \item 由上面的市场出清条件，我们可以求得均衡利率为 $r_{F}=\dfrac{1}{100\rho}-1$；
    \item $r_{F}$ 随着偏好系数的增加而减小。这是因为偏好系数增大，未来消费带来的边际效用相对增加，从而人们对能将财富在 0 期、1 期之间进行转换的证券需求增加，证券价格上升，利率下降。
\end{enumerate}
\end{solution}

\begin{problem}
\problemno{3.3*} 继续考虑练习 3.2 中定义的经济。现在假设还有另外一个参与者 2，他的效用函数和参与者 1 的一样但他的禀赋为 $\begin{bmatrix}1 \\ 100\end{bmatrix}$。
\begin{enumerate}[label=(\alph*)]
    \item 求解参与者 2 的最优消费/储蓄选择。
    \item 比较参与者 1 和 2 的最优消费/储蓄选择。
    \item 求解有两个参与者时的市场均衡。
    \item 计算均衡利率。它与前一情形（经济中只有参与者 1）下的均衡利率有何不同？
    \item 证明均衡配置是 Pareto 最优的。
\end{enumerate}
\end{problem}
\begin{solution}
\begin{enumerate}[label=(\alph*)]
    \item 先不考虑消费非负的约束，参与者 2 的优化问题就变成了
    \[
        \max_{\theta}\bigl\{\log(1-\theta)+\rho\log\bigl(100+\theta(1+r_{F})\bigr)\bigr\}
    \]
    一阶条件 $\dfrac{1}{1-\theta}=\rho\dfrac{1+r_{F}}{100+\theta(1+r_{F})}$ 给出证券需求
    \[
        \theta^{*}=\frac{\rho(1+r_{F})-100}{(1+\rho)(1+r_{F})}
    \]
    最优消费为
    \[
        c_{0}^{*}=\frac{(1+r_{F})+100}{(1+\rho)(1+r_{F})},\quad
        c_{1}^{*}=\frac{\rho\bigl((1+r_{F})+100\bigr)}{1+\rho}
    \]
    由于最优消费都是非负的，因此上面的求解过程没有问题。
    \item 两人的最优消费/储蓄选择形式上是完全一样的，差别在于两人的禀赋不一样；
    \item 此时的市场出清条件是 $\theta_{1}^{*}+\theta_{2}^{*}=0$，即
    \[
        \frac{100\rho(1+r_{F})-1}{(1+\rho)(1+r_{F})}+\frac{\rho(1+r_{F})-100}{(1+\rho)(1+r_{F})}=0
    \]
    分子之和为零：$101\rho(1+r_{F})=101$，故 $1+r_{F}=1/\rho$；
    \item 均衡利率是 $r_{F}=1/\rho-1$。只有参与者 1 时总禀赋为 $(100,1)$，均衡利率为 $1/(100\rho)-1$；现在总禀赋为 $(101,101)$，利率由跨期边际替代率 $\rho$ 单独决定，不再被极端的当期/未来禀赋扭曲；
    \item 配置 Pareto 最优，是指不存在另一种可行配置，使某人的效用严格提高而其他人都不下降。完全市场的均衡满足这一点（定理 3.4）；可检验的条件是两人的跨期边际替代率相等。边际效用之比为
    \[
        \frac{\partial U/\partial c_{1}}{\partial U/\partial c_{0}}=\rho\frac{c_{0}}{c_{1}}
    \]
    由两人的最优消费都有 $c_{0}^{*}/c_{1}^{*}=1/[\rho(1+r_{F})]$，因而两人的
    \[
        \rho\frac{c_{0}^{*}}{c_{1}^{*}}=\frac{1}{1+r_{F}}
    \]
    相等，均衡配置是 Pareto 最优的。
\end{enumerate}
\end{solution}

\begin{problem}
\problemno{3.4*} 在 1 期，经济有三个等可能的状态 $a$、$b$ 和 $c$。在证券市场上交易的 Arrow-Debreu 证券集合是完全的。经济中有两个参与者，参与者 1 和 2。他们的禀赋如下：
\begin{center}
\WJEventTree[left={$e_{1}{:}100$}]{0,0,0}
\qquad
\WJEventTree[left={$e_{2}{:}0$}]{200,100,50}
\end{center}
他们的偏好相同，形式如下：
\[
    U(c)=-\frac{1}{c_{0}}-\frac13\Bigl(\frac{1}{c_{1a}}+\frac{1}{c_{1b}}+\frac{1}{c_{1c}}\Bigr).
\]
\begin{enumerate}[label=(\alph*)]
    \item 求解每一参与者的最优消费/投资选择。
    \item 求解均衡状态价格。
    \item 证明对两个参与者的消费都有禀赋越高、消费越高的性质。
    \item 证明每一参与者的消费在总消费中所占的份额 $c_{k,1\omega}/e_{1\omega}$ 在所有状态下都是一样的，其中 $k=1,2$，而 $e_{1\omega}$ 是在状态 $\omega=a,b,c$ 下的 1 期总消费/禀赋。
\end{enumerate}
\end{problem}
\begin{solution}
\begin{enumerate}[label=(\alph*)]
    \item 求解最优消费选择时，先不考虑消费非负的约束。对于参与者 1，它的优化问题是：
    \[
        \begin{aligned}
            \max_{\{c_{0},c_{1}\}}\ &\Bigl\{-\frac{1}{c_{0}}-\frac13\Bigl(\frac{1}{c_{1a}}+\frac{1}{c_{1b}}+\frac{1}{c_{1c}}\Bigr)\Bigr\}\\
            \mathrm{s.t.}\quad &c_{0}+\phi_{a}c_{1a}+\phi_{b}c_{1b}+\phi_{c}c_{1c}=100
        \end{aligned}
    \]
    拉格朗日函数为
    \[
        \mathcal{L}=-\frac{1}{c_{0}}-\frac13\Bigl(\frac{1}{c_{1a}}+\frac{1}{c_{1b}}+\frac{1}{c_{1c}}\Bigr)+\lambda\bigl(100-c_{0}-\phi_{a}c_{1a}-\phi_{b}c_{1b}-\phi_{c}c_{1c}\bigr)
    \]
    一阶条件是边际效用与状态价格成比例：
    \[
        \frac{1}{c_{0}^{2}}=\lambda,\qquad
        \frac{1}{3c_{1\omega}^{2}}=\lambda\phi_{\omega},\quad\omega=a,b,c
    \]
    因而 $c_{1\omega}=c_{0}/\sqrt{3\phi_{\omega}}$。代入预算约束
    \[
        c_{0}\Bigl(1+\sqrt{\frac{\phi_{a}}{3}}+\sqrt{\frac{\phi_{b}}{3}}+\sqrt{\frac{\phi_{c}}{3}}\Bigr)=100
    \]
    得到参与者 1 的最优消费选择为
    \[
        \begin{aligned}
            c_{0}^{*}&=\frac{100}{1+\sqrt{\phi_{a}/3}+\sqrt{\phi_{b}/3}+\sqrt{\phi_{c}/3}}\\
            c_{1a}^{*}&=\frac{100}{\sqrt{3\phi_{a}}\bigl(1+\sqrt{\phi_{a}/3}+\sqrt{\phi_{b}/3}+\sqrt{\phi_{c}/3}\bigr)}\\
            c_{1b}^{*}&=\frac{100}{\sqrt{3\phi_{b}}\bigl(1+\sqrt{\phi_{a}/3}+\sqrt{\phi_{b}/3}+\sqrt{\phi_{c}/3}\bigr)}\\
            c_{1c}^{*}&=\frac{100}{\sqrt{3\phi_{c}}\bigl(1+\sqrt{\phi_{a}/3}+\sqrt{\phi_{b}/3}+\sqrt{\phi_{c}/3}\bigr)}
        \end{aligned}
    \]
    参与者 2 的财富为 $w_{2}=200\phi_{a}+100\phi_{b}+50\phi_{c}$，同一组一阶条件给出
    \[
        c_{0}^{*}=\frac{w_{2}}{1+\sqrt{\phi_{a}/3}+\sqrt{\phi_{b}/3}+\sqrt{\phi_{c}/3}},\qquad
        c_{1\omega}^{*}=\frac{c_{0}^{*}}{\sqrt{3\phi_{\omega}}}
    \]
    即
    \[
        \begin{aligned}
            c_{0}^{*}&=\frac{200\phi_{a}+100\phi_{b}+50\phi_{c}}{1+\sqrt{\phi_{a}/3}+\sqrt{\phi_{b}/3}+\sqrt{\phi_{c}/3}}\\
            c_{1a}^{*}&=\frac{200\phi_{a}+100\phi_{b}+50\phi_{c}}{\sqrt{3\phi_{a}}\bigl(1+\sqrt{\phi_{a}/3}+\sqrt{\phi_{b}/3}+\sqrt{\phi_{c}/3}\bigr)}\\
            c_{1b}^{*}&=\frac{200\phi_{a}+100\phi_{b}+50\phi_{c}}{\sqrt{3\phi_{b}}\bigl(1+\sqrt{\phi_{a}/3}+\sqrt{\phi_{b}/3}+\sqrt{\phi_{c}/3}\bigr)}\\
            c_{1c}^{*}&=\frac{200\phi_{a}+100\phi_{b}+50\phi_{c}}{\sqrt{3\phi_{c}}\bigl(1+\sqrt{\phi_{a}/3}+\sqrt{\phi_{b}/3}+\sqrt{\phi_{c}/3}\bigr)}
        \end{aligned}
    \]
    因为最优消费都是非负的，因而上面的求解过程是没有问题的；
    \item 市场出清条件是
    \[
        \begin{aligned}
            &\frac{100}{\sqrt{3\phi_{a}}\bigl(1+\sqrt{\phi_{a}/3}+\sqrt{\phi_{b}/3}+\sqrt{\phi_{c}/3}\bigr)}+\frac{200\phi_{a}+100\phi_{b}+50\phi_{c}}{\sqrt{3\phi_{a}}\bigl(1+\sqrt{\phi_{a}/3}+\sqrt{\phi_{b}/3}+\sqrt{\phi_{c}/3}\bigr)}=200\\
            &\frac{100}{\sqrt{3\phi_{b}}\bigl(1+\sqrt{\phi_{a}/3}+\sqrt{\phi_{b}/3}+\sqrt{\phi_{c}/3}\bigr)}+\frac{200\phi_{a}+100\phi_{b}+50\phi_{c}}{\sqrt{3\phi_{b}}\bigl(1+\sqrt{\phi_{a}/3}+\sqrt{\phi_{b}/3}+\sqrt{\phi_{c}/3}\bigr)}=100\\
            &\frac{100}{\sqrt{3\phi_{c}}\bigl(1+\sqrt{\phi_{a}/3}+\sqrt{\phi_{b}/3}+\sqrt{\phi_{c}/3}\bigr)}+\frac{200\phi_{a}+100\phi_{b}+50\phi_{c}}{\sqrt{3\phi_{c}}\bigl(1+\sqrt{\phi_{a}/3}+\sqrt{\phi_{b}/3}+\sqrt{\phi_{c}/3}\bigr)}=50
        \end{aligned}
    \]
    由 $c_{1\omega}=c_{0}/\sqrt{3\phi_{\omega}}$，0 期总禀赋为 100，故 $\sum_{k}c_{k,0}=100$。状态 $a$ 的出清化为
    \[
        \frac{\sum_{k}c_{k,0}}{\sqrt{3\phi_{a}}}=\frac{100}{\sqrt{3\phi_{a}}}=200
    \]
    即 $\sqrt{3\phi_{a}}=\dfrac12$，从而 $\phi_{a}=\dfrac{1}{12}$。同理 $\dfrac{100}{\sqrt{3\phi_{b}}}=100$ 给出 $\phi_{b}=\dfrac13$，$\dfrac{100}{\sqrt{3\phi_{c}}}=50$ 给出 $\phi_{c}=\dfrac43$。
    \item $\phi_{a}<\phi_{b}<\phi_{c}$，对于参与者 1、2 都有 $c_{1a}^{*}>c_{1b}^{*}>c_{1c}^{*}$，即总禀赋越高、消费越高；
    \item 因为对于消费者 1、2 都有
    \[
        \begin{aligned}
            \frac{c_{1a}^{*}}{c_{1b}^{*}}&=\sqrt{\frac{\phi_{b}}{\phi_{a}}}=2=\frac{200}{100}\\
            \frac{c_{1a}^{*}}{c_{1c}^{*}}&=\sqrt{\frac{\phi_{c}}{\phi_{a}}}=4=\frac{200}{50}
        \end{aligned}
    \]
    因此，每一参与者的消费在总消费中所占的份额在所有状态下都是相等的。
\end{enumerate}
\end{solution}



% !TeX root = ../../金融经济学_王江_讲义及习题.tex
% 知识点：references/金融经济学/金融经济学-王江_261013.pdf，印刷页 49–62
% 补充：references/金融经济学/金融经济学_第4章 套利和资产定价.pptx（第 1–36 页）
% 习题：references/金融经济学/史上最全最完整的金融经济学答案(王江版).pdf，页 14–18（印刷页 13–17）

\lecture{套利和资产定价}

\ifdefined\WJHandoutKnowledge

第 3 章在 Arrow--Debreu 市场结构下讨论均衡与效率；本章回到第 2 章的一般支付矩阵 $X$，讨论\textbf{证券价格的基本性质}、\textbf{定价原理}、\textbf{无套利}与\textbf{资产定价基本定理}。

\begin{tcolorbox}[knowledge,title=4.1 一般市场结构]
多数可交易证券在多个状态下都有正支付，称为\textbf{复合证券}（相对只在单一状态付 1 的阿罗证券而言）。设 1 期有 $\Omega$ 个状态，市场上有 $N$ 只证券（$n=1,\dots,N$），第 $n$ 只证券的支付向量为 $X_n=\begin{bmatrix}x_{1,n} \\ \vdots \\ x_{\Omega,n}\end{bmatrix}$，\textbf{市场结构矩阵}为（与第 2.3 节一致，行为状态、列为证券）：
\[
    X=[X_{\cdot,1},\,\cdots,\,X_{\cdot,N}]=
    \begin{bmatrix}
        x_{1,1}   & \cdots & x_{1,N}   \\
        \vdots    & \ddots & \vdots    \\
        x_{\omega,1} & \cdots & x_{\omega,N} \\
        \vdots    & \ddots & \vdots    \\
        x_{\Omega,1} & \cdots & x_{\Omega,N}
    \end{bmatrix}.
\]

\tcbline

\textbf{A.\ 冗余证券}\quad 若证券 $j$ 的支付列 $x_j$ 可由其余证券的\textbf{复制组合}（replicating portfolio）线性表出，即存在 $\theta_{\setminus j}$ 使得
\[
    x_j=X_{\setminus j}\theta_{\setminus j},
\]
则称 $j$ 为\textbf{支付意义上的冗余证券}。对任意原组合 $\theta$，
\[
 X\theta=X_{\setminus j}\theta_{\setminus j}+\theta_jx_j
 =X_{\setminus j}(\theta_{\setminus j}+\theta_j\theta^*_{\setminus j}),
\]
其中 $\theta^*_{\setminus j}$ 是复制 $x_j$ 的固定配方。这说明删去冗余列不改变可达支付集合。若市场无摩擦且满足一价定律，还必须有
\[
 S_j=S_{\setminus j}^{\mathsf T}\theta^*_{\setminus j},
\]
所以替换前后的成本也相同；若价格不相等，同支付的两种实现方式反而会产生套利。资源配置上的等价还要求复制交易不受限制，原有证券禀赋也按复制关系换算。存在交易成本、卖空约束等摩擦时，支付可复制不保证交易上可替代。

\begin{tcolorbox}[example,title=PPT 补充：第三只证券是否增加支付维度]
晴、雨两个状态下，证券 $A,B,C$ 的支付为
\[
 X=\begin{bmatrix}10&0&20\\0&10&10\end{bmatrix}.
\]
$C=2A+B$，故第三列虽不同于前两列，却没有增加独立支付维度；无套利价格必须满足 $S_C=2S_A+S_B$。例如原来持有 $aA+bB+cC$，可改为 $(a+2c)A+(b+c)B$，各状态支付与成本均相同。\textbf{证券种类数增加，不等于市场更完全。}
\end{tcolorbox}

\tcbline

\textbf{B.\ 证券市场的等价描述}\quad 无摩擦时，可用\textbf{线性无关}的 $N$ 只证券描述市场，$\mathrm{rank}(X)=N$（通常设 $N\le\Omega$）。任取 $N$ 个线性无关的组合向量 $\theta_1,\dots,\theta_N$，令 $H=[\theta_1,\cdots,\theta_N]$（$N\times N$ 可逆），则新支付矩阵 $X_\theta=XH$ 与原市场\textbf{等价}：原证券列空间不变，$X=X_\theta H^{-1}$，每只原证券都可由新组合复制。更一般地，原组合 $\theta$ 对应新组合 $H^{-1}\theta$，因为
\[
 X\theta=XH(H^{-1}\theta)=X_\theta(H^{-1}\theta).
\]
新证券的价格须同时变为 $S_\theta=H^{\mathsf T}S$，于是
\[
 S_\theta^{\mathsf T}H^{-1}\theta=S^{\mathsf T}\theta.
\]
支付和成本都不变，这才是同一市场的等价描述。可逆变换只是更换支付空间的一组基；若 $H$ 不可逆，可能丢掉支付维度。秩不变也可由 $X=XHH^{-1}$ 得到
\[
 \mathrm{rank}(X)\le\mathrm{rank}(XH)\le\mathrm{rank}(X),
 \qquad\mathrm{rank}(XH)=N.
\]

\tcbline

\textbf{C.\ 生成（span）与完全市场}\quad 记 $\mathbf{1}_\omega$ 为第 $\omega$ 分量为 1、其余为 0 的列向量（状态 $\omega$ 阿罗证券的支付）。若 $\mathrm{rank}(X)=\Omega$ 且 $N\ge\Omega$，则对每一 $\omega$ 存在组合 $\theta_\omega$ 使 $X\theta_\omega=\mathbf{1}_\omega$。当 $X$ 为 $\Omega\times\Omega$ 可逆方阵时，解唯一：
\[
    \theta_\omega=X^{-1}\mathbf{1}_\omega.
\]
此时称不确定性由证券\textbf{生成}（span）；市场与第 3 章完全阿罗市场等价。

\begin{tcolorbox}[knowledge]
\textbf{定理 4.1}\quad 无摩擦时，证券市场\textbf{完全}当且仅当线性无关证券只数等于状态数（可取 $N=\Omega$ 且 $\mathrm{rank}(X)=\Omega$，此时 $X$ 可逆）。这与定义 3.2、\textcolor{red}{\textbf{行满秩}}判据一致。
\end{tcolorbox}

\begin{tcolorbox}[example,title=PPT 补充：先复制阿罗证券，再复制任意支付]
两状态、两证券的支付矩阵为
\[
 X=\begin{bmatrix}2&3\\5&1\end{bmatrix},\qquad
 X^{-1}=\frac1{13}\begin{bmatrix}-1&3\\5&-2\end{bmatrix}.
\]
因此 $\theta_1=[-1,5]^{\mathsf T}/13$、$\theta_2=[3,-2]^{\mathsf T}/13$ 分别复制 $[1,0]^{\mathsf T}$、$[0,1]^{\mathsf T}$。要复制支付 $y=[5,2]^{\mathsf T}$，先把它写成五份第一状态阿罗证券加两份第二状态阿罗证券，再换成可交易证券：
\[
 \theta=X^{-1}y=5\theta_1+2\theta_2
 =\frac1{13}\begin{bmatrix}1\\21\end{bmatrix},\qquad
 X\theta=\frac1{13}\begin{bmatrix}2+63\\5+21\end{bmatrix}
 =\begin{bmatrix}5\\2\end{bmatrix}.
\]
复制成本为 $(S_1+21S_2)/13$。阿罗证券不必实际挂牌交易，能用已有证券复制即可。
\end{tcolorbox}
\end{tcolorbox}

\begin{tcolorbox}[knowledge,title=4.2 套利]
记交易证券的 0 期价格向量为 $S=\begin{bmatrix}S_1 \\ \vdots \\ S_N\end{bmatrix}$。组合 $\theta=\begin{bmatrix}\theta_1 \\ \vdots \\ \theta_N\end{bmatrix}$ 的 0 期成本为 $S^{\mathsf T}\theta$，1 期支付为 $X\theta$。从支付 $X$ 到价格 $S$ 的映射称为\textbf{资产定价关系}（资产定价模型）。若 $X\theta\ge 0$，称该组合在 1 期满足\textbf{有限责任}。

\begin{tcolorbox}[knowledge]
\textbf{定义 4.1}\quad 组合 $\theta$ 构成\textbf{套利机会}，若同时满足：
\begin{enumerate}
    \item $S^{\mathsf T}\theta\le 0$（0 期净支出非正）；
    \item $X\theta\ge 0$（1 期各状态支付非负）；
    \item 上述两式中\textbf{至少有一个严格不等式}成立（0 期严格收入或某状态严格正支付）。
\end{enumerate}
等价地，扩展支付向量 $\begin{bmatrix}-S^{\mathsf T}\theta \\ X\theta\end{bmatrix}$ 非负且至少一分量严格为正。
\end{tcolorbox}

\textbf{向量不等号约定：}$z\ge0$ 指各分量非负，$z>0$ 指 $z\ge0$ 且 $z\ne0$（至少一个分量为正），$z\gg0$ 才指每个分量都严格为正。因此套利不要求每个未来状态都获利，只要求没有亏损并在某时点、某状态有正所得。

常分为三类（记严格不等方向与教材一致）：
\begin{itemize}
    \item \textbf{第 1 类套利：}$S^{\mathsf T}\theta<0$ 且 $X\theta=0$——0 期白捡钱、未来无负债；
    \item \textbf{第 2 类套利：}$S^{\mathsf T}\theta=0$ 且 $X\theta>0$——\textbf{套利组合}（零成本、未来非负且非零支付）；
    \item \textbf{第 3 类套利：}$S^{\mathsf T}\theta<0$ 且 $X\theta>0$——0 期收入为正且未来非负、非零。
\end{itemize}

初始成本为零时，买入恰由卖空收入融资；初始成本为负时，卖空收入支付买入成本后还有余额。这就是本章“不需要额外初始资源”的含义。PPT 把前者称作自融资；动态模型中的自融资还要求各期调仓不追加外部资金，不能仅凭初始成本为零就判断。

\begin{tcolorbox}[example]
\textbf{例 4.1}\quad 三只证券的价格写在树的根上，支付写在各枝：
\begin{center}
\begin{tabular}{@{}c@{\quad}c@{\quad}c@{}}
\textbf{A} & \textbf{B} & \textbf{C} \\[0.35em]
\WJEventTree[left={$1$},stem]{1,1,1} &
\WJEventTree[left={$1$},stem]{0,2,2} &
\WJEventTree[left={$2$},stem]{2,0,0}
\end{tabular}
\end{center}
即 $S^{\mathsf T}=[1,\,1,\,2]$，$X=\begin{bmatrix}1&0&2\\1&2&0\\1&2&0\end{bmatrix}$。

组合 $\theta_1=[2;\,-1;\,-1]$：买入 $2$ 单位 $A$、卖出 $1$ 单位 $B$ 和 $1$ 单位 $C$。$S^{\mathsf T}\theta_1=-1$，$X\theta_1=0$，为\textbf{第 1 类套利}。支付 $[2;2;2]$ 可由 $2$ 单位 $A$ 以成本 $2$ 复制，也可由 $1$ 单位 $B$ 加 $1$ 单位 $C$ 以成本 $3$ 复制——同支付不同价。

另一种看法是先比较 $B+C$ 与 $2A$：二者都支付 $[2,2,2]^{\mathsf T}$，但成本分别为 $3$ 和 $2$，故应卖出昂贵组合、买入便宜组合。

组合 $\theta_2=[2;\,0;\,-1]$：买入 $2$ 单位 $A$、卖出 $1$ 单位 $C$。$S^{\mathsf T}\theta_2=0$，$X\theta_2=[0;\,2;\,2]>0$，为\textbf{第 2 类套利}。
\end{tcolorbox}
\end{tcolorbox}

\begin{tcolorbox}[knowledge,title=4.3 无套利原理]
套利机会只依赖公开的支付 $X$ 与价格 $S$，不依赖私人信息或对状态概率的主观信念。

\begin{tcolorbox}[knowledge]
\textbf{定理 4.2}\quad 市场均衡中\textbf{不存在}套利机会。
\end{tcolorbox}

\textbf{证明（按 PPT 的消费增量写法）：}反设均衡价格下存在套利 $\theta$。给任一参与者原有证券头寸再加上 $\theta$，同时把初始净收入用于 0 期消费，消费可由 $c_k$ 提高为
\[
 c'_k=c_k+\begin{bmatrix}-S^{\mathsf T}\theta\\X\theta\end{bmatrix}
 \ge c_k,\qquad c'_k\ne c_k.
\]
该交易的初始成本非正，不需要追加初始财富；未来也没有负支付。若至少一名参与者对这种非负、非零的消费增加严格偏好，则 $U_k(c'_k)>U_k(c_k)$，与均衡时其消费最优矛盾。\textbf{推出矛盾只需一次可行改进，无须先诉诸无限放大。}无限放大是在无交易限制下说明原价格无法维持的直觉。

\textbf{前提与边界：}市场无摩擦，可按给定价格买卖并实施所需空头；偏好须对上述消费增量具有严格单调性（本章不满足性假设用于此处）。仅有“某处还想多消费”的局部不满足性，并不足以保证对任意套利支付都严格偏好。无套利也不同于信息有效：这里检验的是价格与逐状态支付是否相容，不要求价格已经反映全部信息。均衡 $\Rightarrow$ 无套利是必要条件；无套利 $\not\Rightarrow$ 均衡，还需要偏好最优化与市场出清。\textbf{无摩擦是交易制度假设，无套利是给定支付与价格的约束，两者不等价}：例 4.1 即可在无摩擦条件下呈现错误价格与套利。

\begin{tcolorbox}[example,title=PPT 补充：跨市场的一价定律]
同一时刻、同一可交割黄金在两个市场分别报价 $1900$ 与 $1905$ 美元。假设可同时成交、自由卖空和交割，且没有交易与运输成本，买入便宜的黄金、卖出昂贵的黄金可立即收入 $5$ 美元，净实物头寸为零，是第 1 类套利。买入压力与卖出压力解释价差如何趋于消失，但不能据此推出最终价格必为两报价的中点。有摩擦时须扣除可锁定的全部成本；承担尚未对冲的汇率或交割风险的交易不能仅凭预期价差称为套利。这里数字只是教学假设。
\end{tcolorbox}

\begin{tcolorbox}[knowledge]
\textbf{定义 4.2}\quad 证券市场上不存在套利机会，称为满足\textbf{无套利原理}（no-arbitrage principle）。
\end{tcolorbox}
\end{tcolorbox}

\begin{tcolorbox}[knowledge,title=4.4 资产定价基本定理]
资产定价关系是从支付 $x$ 到价格的映射。记\textbf{定价算子}（估值算子）$V$：
\[
    S=V(x).
    \tag{4.1}
\]
严格地说，先在\textbf{可达支付空间} $\mathcal M=\{X\theta:\theta\in\mathrm R^N\}$ 上定义 $V(X\theta)=S^{\mathsf T}\theta$。若同一支付的不同复制组合成本不同，这个映射就没有唯一的值；一价定律正是保证它定义良好的条件。第 $n$ 只证券的价格即 $S_n=V(X_n)$。无套利把 $V$ 收成正的线性泛函。

\begin{tcolorbox}[knowledge]
\textbf{定理 4.3（一价定律）}
\[
    x=y\implies V(x)=V(y).
    \tag{4.2}
\]
\end{tcolorbox}

没有第 1 类套利就得到一价定律。推论：未来支付为 $0$ 的证券价格为 $0$，即 $V(0)=0$。证明：若 $x=y$ 却 $V(x)>V(y)$，卖出 $x$、买入 $y$，未来净支付为零而初始净收入为正，构成第 1 类套利。一价定律弱于无套利：没有第 1 类套利推不出没有第 2、3 类。

\begin{tcolorbox}[knowledge]
\textbf{定理 4.4}\quad 支付为正的证券价格为正：
\[
    x>0\implies V(x)>0.
    \tag{4.3}
\]

\textbf{定理 4.5}\quad $V$ 递增：
\[
    x_1\ge x_2\implies V(x_1)\ge V(x_2).
    \tag{4.4}
\]
\end{tcolorbox}

若 $x>0$ 而 $V(x)\le0$，买入它即为第 2 或第 3 类套利，故必须 $V(x)>0$。若 $x_1\ge x_2$ 却 $V(x_1)<V(x_2)$，买入 $x_1$、卖出 $x_2$，初始净收入为正且未来支付非负，也构成套利。当 $x_1-x_2>0$ 时还有严格结论 $V(x_1)>V(x_2)$。

\begin{tcolorbox}[knowledge]
\textbf{定理 4.6}\quad 无摩擦且无套利市场中，可达支付上的定价算子是递增的线性算子：对任意 $a,b\in\mathrm{R}$ 以及支付 $x$、$y$，
\[
    V(ax+by)=aV(x)+bV(y).
    \tag{4.5}
\]
从而 $V(0)=0$，并且存在向量 $\phi\in\mathrm{R}^{\Omega}$ 使 $V(x)=\phi^{\mathsf T}x$。
\end{tcolorbox}

\textbf{线性的证明：}分别复制 $x$、$y$，再持有 $a$ 倍与 $b$ 倍组合，得到支付 $ax+by$，成本为 $aV(x)+bV(y)$；一价定律要求它等于 $V(ax+by)$。允许负系数用到了卖空不受限的无摩擦假设。

这个向量表示首先只要求在 $\mathcal M$ 上成立；若市场不完全，不能据此为所有未交易支付直接指定唯一价格。形式与第 3 章的状态价格公式相同。这里只知道 $\phi$ 是一个向量；下一条才把它收成严格正、并且说明这正是无套利。

\begin{tcolorbox}[knowledge]
\textbf{定理 4.7（资产定价基本定理）}\quad 市场不存在套利机会，当且仅当存在 $\phi\gg 0$，使得每只可交易证券的价格都是
\[
    S=X^{\mathsf T}\phi.
    \tag{4.6}
\]
教材写成 $S=(\phi^{\mathsf T}X)^{\mathsf T}$，与上式相同。
\end{tcolorbox}

\textbf{充分性：}若 $S=X^{\mathsf T}\phi$、$\phi\gg0$，则任意组合满足 $S^{\mathsf T}\theta=\phi^{\mathsf T}X\theta$。未来支付为零时成本必为零；未来支付非负、非零时成本必严格为正。因此不可能同时有非正成本与套利支付，三类套利都被排除。必要性是 Stiemke 引理的特例：对 $m\times n$ 矩阵 $X$，存在 $\phi\gg 0$ 使 $S=X^{\mathsf T}\phi$，当且仅当集合 $\{\theta:[-S^{\mathsf T}\theta;\,X\theta]>0\}$ 为空。该集合正是套利机会的集合。证明放在教材附录。也可把扩展支付矩阵写成 $A=[-S^{\mathsf T};X]$：Stiemke 二择一定理说，“存在 $\theta$ 使 $A\theta>0$”与“存在 $y\gg0$ 使 $A^{\mathsf T}y=0$”恰有一个成立。无套利排除前者；把后者写成 $y=[y_0;y_1]$，则 $-y_0S+X^{\mathsf T}y_1=0$，取 $\phi=y_1/y_0\gg0$ 即得定价公式。

因此无套利意味着存在严格正的状态价格，可为\textbf{所有}可交易证券定价。这样的 $\phi$ 一般不唯一；定理 4.7 没有给出唯一性。给已有证券定价的线性解可以含负分量（练习 4.5），但那一支不能用来代表无套利：无套利要求的是存在一支 $\phi\gg 0$。

\begin{tcolorbox}[example]
\textbf{例 4.2}\quad 市场上先只有证券 $A$，价格为 $1$，三个状态下支付都是 $1$。没有套利。任一满足 $\phi_1+\phi_2+\phi_3=1$ 的 $\phi\gg 0$ 都能为它定价：
\[
    \Phi_A=\{\phi\gg 0:\phi_1+\phi_2+\phi_3=1\}.
\]
再加入证券 $B$，价格为 $1$，支付为 $[0;\,2;\,2]$。能为它定价的是
\[
    \Phi_B=\{\phi\gg 0:2\phi_2+2\phi_3=1\}.
\]
同时给 $A$、$B$ 定价的状态价格是交集
\[
    \Phi_A\cap\Phi_B=\{\phi\gg 0:\phi_1=\tfrac12,\;\phi_2+\phi_3=\tfrac12\},
\]
解不唯一，故 $A$、$B$ 的价格里没有套利。

再加入证券 $C$，价格为 $2$，支付为 $[2;\,0;\,0]$。条件变成 $\phi_1=1$，
\[
    \Phi_C=\{\phi\gg 0:\phi_1=1\}.
\]
三者的交集是空集。不存在能为所有证券定价的 $\phi\gg 0$，与例 4.1 中这三只证券存在套利一致。
\end{tcolorbox}

\begin{tcolorbox}[knowledge]
\textbf{定理 4.8}\quad 无套利的完全市场上，状态价格向量唯一。
\end{tcolorbox}

设 $\theta_\omega$ 复制状态 $\omega$ 的阿罗证券，$X\theta_\omega=\mathbf{1}_\omega$。则 $\phi_\omega=S^{\mathsf T}\theta_\omega$，由市场价格唯一给定；即使存在冗余证券而复制组合不唯一，一价定律仍保证所有复制组合成本相同。若有两支不同的定价向量 $\phi,\psi$，必有状态 $\omega$ 满足 $\phi_\omega\ne\psi_\omega$；但该阿罗证券已可复制，两支向量却给同一复制支付不同价格，与一价定律矛盾。方阵可逆时，直接有 $\phi=(X^{\mathsf T})^{-1}S$。完全意味着任意支付都可通过市场交易复制，因而其价格唯一；不完全市场的已有证券仍可有唯一市场报价，只是不能据此唯一确定所有未交易支付的价格。
\end{tcolorbox}

\begin{tcolorbox}[knowledge,title=4.5 风险中性定价和鞅]
设可交易证券中有一只无风险债券，1 期支付为 $\iota=[1,\ldots,1]^{\mathsf T}$，价格为 $S_1$。由基本定理，存在 $\phi\gg 0$ 使
\[
    S_1=\phi^{\mathsf T}\iota=\sum_{\omega\in\Omega}\phi_{\omega}.
\]
$1$ 单位债券投资的净收益率称为\textbf{无风险利率} $r_F$：
\[
    S_1(1+r_F)=1,\qquad r_F=\frac{1-S_1}{S_1}.
\]
因而
\[
    \frac{1}{1+r_F}=\sum_{\omega\in\Omega}\phi_{\omega}.
    \tag{4.7}
\]
债券价格可以写成对其确定支付 $1$ 的贴现。$1/(1+r_F)$ 是折现因子，$r_F$ 是折现率。贴现只是用状态价格定价的一种写法。利率也是货币的时间价值：今天 $1$ 单位资源换取未来确定资源时，市场给出的回报。

\begin{tcolorbox}[example,title=PPT 补充：从状态价格读出利率]
晴、雨两状态价格为 $\phi=[0.3,0.5]^{\mathsf T}$ 时，单位无风险债券价格为 $B=0.3+0.5=0.8$，所以 $r_F=1/0.8-1=25\%$。状态价格不是概率，其和不必为 $1$；规范化后 $Q=[3/8,5/8]$。总和 $0.8$ 决定时间贴现，各状态的相对权重决定不同风险支付如何估值，利率一个数并没有概括整个状态价格向量。
\end{tcolorbox}

对其余证券 $n=2,\ldots,N$，支付 $x_n=[x_{1n},\ldots,x_{\Omega n}]^{\mathsf T}$，
\[
    S_n=\phi^{\mathsf T}x_n=\sum_{\omega\in\Omega}\phi_{\omega}x_{\omega n}.
\]
定义\textbf{风险中性概率}（规范化的状态价格）
\[
    q_{\omega}=\frac{\phi_{\omega}}{\sum_{\omega'}\phi_{\omega'}}.
    \tag{4.8}
\]
$q_{\omega}>0$ 且 $\sum q_{\omega}=1$，故 $Q=\{q_{\omega}\}$ 是 $\Omega$ 上的概率测度。定价公式改写成
\[
    S_n=S_1\sum_{\omega}q_{\omega}x_{\omega n},
    \qquad
    S_n=\frac{\mathrm{E}^{Q}[\tilde{x}_n]}{1+r_F},\quad n=1,\ldots,N.
    \tag{4.9}
\]
(4.9) 就是 (4.6) 的另一种写法，称为\textbf{风险中性定价}（risk-neutral pricing），$Q$ 称为\textbf{风险中性测度}。它适用于所有证券，并不区分风险大小：在 $Q$ 下，每只证券的价格都只是期望支付按无风险利率折现，好像只取得无风险收益。

\textbf{为什么不需要假设投资者真的风险中性？}风险偏好、禀赋等会影响基础证券的均衡价格；给定这些价格后，可复制支付的价格由无套利约束决定，不必重新询问投资者偏好。$Q$ 是把已有价格中的风险调整吸收到权重后的定价工具。若市场不完全，已有证券价格通常只约束一组可行 $Q$，无法由此给所有衍生支付一个唯一价格。

\begin{tcolorbox}[example,title=PPT 补充：套餐支付的价格由复制成本决定]
股票 $A$ 价格为 $10$，两状态支付为 $[20,0]^{\mathsf T}$；债券 $B$ 价格为 $5$，支付为 $[5,5]^{\mathsf T}$。衍生证券 $C$ 支付为 $[25,5]^{\mathsf T}$，等于 $A+B$，所以价格必须为 $15$。若报价为 $20$，卖出 $C$ 并买入 $A+B$，初始收入 $5$，未来净支付为零。

由 $20\phi_1=10$、$5(\phi_1+\phi_2)=5$，得 $\phi_1=\phi_2=1/2$，无风险利率为零，故 $C=\mathrm E^Q[\tilde C]=(25+5)/2=15$。复制法与风险中性法给出同一价格；这里 $Q$ 两状态相等，不代表真实概率也相等。
\end{tcolorbox}

$Q$ 不必等于真实概率 $P$。它由 (4.8) 从状态价格定义，风险已经折进 $q_{\omega}$。因此 (4.9) 是定价关系，不是用 $P$ 作的概率预期。

\begin{tcolorbox}[example]
\textbf{例 4.3}\quad 1 期有两个等可能状态 $a$、$b$。
\begin{center}
\begin{tabular}{@{}c@{\quad}c@{}}
\textbf{证券 1} & \textbf{证券 2} \\[0.35em]
\WJEventTree[left={$1$},stem]{1,1} &
\WJEventTree[left={$\tfrac12$},stem]{2,0}
\end{tabular}
\end{center}
证券 1 是无风险债券，$r_F=0$。在真实概率下两只证券的期望支付都是 $1$，但证券 2 的价格只有证券 1 的一半：期望值相同，确定支付与风险支付的价格不必相同。

由价格得 $\phi_a+\phi_b=1$、$2\phi_a=1/2$，故 $\phi_a=1/4$、$\phi_b=3/4$，$Q=\{1/4,\,3/4\}$。风险中性定价给出
\[
    \frac{\mathrm{E}^{Q}[\tilde{x}_1]}{1+0}
    =\tfrac14\cdot 1+\tfrac34\cdot 1=1,
    \qquad
    \frac{\mathrm{E}^{Q}[\tilde{x}_2]}{1+0}
    =\tfrac14\cdot 2+\tfrac34\cdot 0=\tfrac12,
\]
正是两只证券的价格。求期望时用的是 $Q$ 不是 $P$。
\end{tcolorbox}

因为 $S_1=1/(1+r_F)$ 且债券的 1 期支付 $\tilde{x}_1=1$，(4.9) 也可写成
\[
    \frac{S_n}{S_1}=\mathrm{E}^{Q}\Bigl[\frac{\tilde{x}_n}{\tilde{x}_1}\Bigr].
\]
$S_n$、$S_1$ 以现时消费品计价；$S_n/S_1$ 则是以无风险债券计价的价格。用债券做计价单位时，0 期价格是 $S_n/S_1$，1 期价格是 $\tilde{x}_n$（债券在 1 期的价格恒为 $1$）。这改变的是计价单位，使时间贴现体现在分母中；它并没有消除风险，风险调整仍包含在 $Q$ 中。记 $\hat{S}_{n,t}=S_{n,t}/S_{1,t}$（$t=0,1$），则
\[
    \hat{S}_{n,0}=\mathrm{E}^{Q}[\hat{S}_{n,1}].
    \tag{4.10}
\]

\begin{tcolorbox}[knowledge]
\textbf{定义 4.3}\quad 设 $\mathcal F_t$ 表示截至 $t$ 期已知的信息。随机过程 $z_t$ 若适应该信息、满足 $\mathrm E^Q|z_t|<\infty$，且
\[
 z_t=\mathrm E^Q[z_{t+1}\mid\mathcal F_t],
\]
则称为测度 $Q$ 下的\textbf{鞅}。当前信息下，下一期的预期变动为零；这不表示每次实现值都不变，也不表示真实概率下的预期收益为零。
\end{tcolorbox}

(4.10) 说的是：以债券计价后，证券价格在风险中性测度 $Q$ 下是鞅。因此 $Q$ 也称为\textbf{等价鞅测度}。这里的“等价”是指 $Q$ 与真实测度 $P$ 有相同的零概率集：$p_{\omega}=0$ 当且仅当 $q_{\omega}=0$。$\phi\gg 0$ 给出 $q\gg 0$；若真实概率也严格为正，二者等价。

\textbf{等价性的套利含义须落实到可交易支付。}若事件 $E$ 在真实世界可能发生，即 $P(E)>0$，而候选 $Q'$ 令 $Q'(E)=0$，那么按它给事件保险 $\mathbf1_E$ 定价会得到零价格。若该支付确实可交易（或可复制），免费买入便是第 2 类套利。\textbf{单纯“低估某项风险”或买入后可能获利，并不足以构成套利}；看跌期权也不一定只在破产时支付。若 $\mathbf1_E$ 在不完全市场中不可交易、不可复制，仅发现某候选定价测度漏掉 $E$，不能直接声称已有证券中存在套利，仍须寻找可行交易组合。

本书只有两个时点，式 (4.10) 只取 $t=0$。规范化价格服从鞅这一关系可以直接搬到动态框架；经济含义在这个静态模型里已经清楚。
\end{tcolorbox}

\begin{tcolorbox}[knowledge,title=本章小结]
\begin{enumerate}
    \item 无摩擦时，线性无关的证券组合给出市场结构的等价描述。市场完全当且仅当独立支付的证券数等于状态数。
    \item 套利是 $S^{\mathsf T}\theta\le 0$、$X\theta\ge 0$ 且至少一处严格。它只依赖公开的支付和价格，不依赖私有信息或真实概率；无摩擦时人人都可以做。
    \item 均衡中不存在套利。前提是市场无摩擦，且至少有人严格偏好非负、非零的消费增加。无套利是均衡的必要条件。
    \item 资产定价基本定理：无套利当且仅当存在 $\phi\gg 0$ 使 $S=X^{\mathsf T}\phi$。完全市场中这个 $\phi$ 唯一。
    \item 风险中性定价 $S_n=\mathrm{E}^{Q}[\tilde{x}_n]/(1+r_F)$ 是同一公式的改写。$Q$ 由状态价格规范化得到，不必等于真实概率。
    \item 用无风险债券计价后，价格在 $Q$ 下是鞅，$Q$ 是等价鞅测度。
\end{enumerate}
\end{tcolorbox}
\fi

% !TeX root = ../../金融经济学_王江_讲义及习题.tex
% 04_套利和资产定价 课后习题（由章节文件拆分）

\begin{problem}
\problemno{4.1} 经济在 1 期有 4 个可能状态。在市场中有 5 只可交易证券，它们的支付矩阵 $X$ 如下：
\[
    X=\begin{bmatrix}
        1 & 1 & 0 & 0 & 0 \\
        1 & 2 & 1 & 0 & 0 \\
        1 & 3 & 2 & 1 & 0 \\
        1 & 4 & 3 & 2 & 1
    \end{bmatrix}.
\]
\begin{enumerate}[label=(\alph*)]
    \item 证明市场是完全的。
    \item 假设五只证券的 0 期价格分别为 $1$，$2.5$，$1.5$，$0.75$，$0.25$。证明状态价格向量存在且唯一，并给出它的值。
    \item 证明市场上存在冗余证券。
    \item 选择一组足以保证市场完全的复合证券。把每一只 Arrow-Debreu 证券都表示成这些证券的组合。称这些组合为状态或有组合。
    \item 用状态或有组合把任一由五只交易证券组成的组合表示出来。
    \item 证明状态或有组合可生成状态空间。
\end{enumerate}
\end{problem}
\begin{solution}
记 5 只交易证券的支付向量分别为
\[
    X_{1}=\begin{bmatrix}1 \\ 1 \\ 1 \\ 1\end{bmatrix},\quad X_{2}=\begin{bmatrix}1 \\ 2 \\ 3 \\ 4\end{bmatrix},\quad X_{3}=\begin{bmatrix}0 \\ 1 \\ 2 \\ 3\end{bmatrix},\quad X_{4}=\begin{bmatrix}0 \\ 0 \\ 1 \\ 2\end{bmatrix},\quad X_{5}=\begin{bmatrix}0 \\ 0 \\ 0 \\ 1\end{bmatrix}
\]
\begin{enumerate}[label=(\alph*)]
    \item 因为 $X_{1}$、$X_{2}$、$X_{3}$、$X_{4}$ 的支付是互相独立的四只证券，而状态数等于 4，因而证券市场是完全的；
    \item 假设状态价格向量为 $\phi=\begin{bmatrix}\phi_{1} \\ \phi_{2} \\ \phi_{3} \\ \phi_{4}\end{bmatrix}$，那么有：
    \[
        \begin{aligned}
            \phi_{1}+\phi_{2}+\phi_{3}+\phi_{4}&=1\\
            \phi_{1}+2\phi_{2}+3\phi_{3}+4\phi_{4}&=2.5\\
            \phi_{2}+2\phi_{3}+3\phi_{4}&=1.5\\
            \phi_{3}+2\phi_{4}&=0.75\\
            \phi_{4}&=0.25
        \end{aligned}
    \]
    求解上述方程组可以得到唯一的解为 $\phi=\begin{bmatrix}\tfrac14 \\ \tfrac14 \\ \tfrac14 \\ \tfrac14\end{bmatrix}$；
    \item 买入 1 单位的证券 2 卖出 1 单位的证券 3，得到的支付是 $\begin{bmatrix}1 \\ 1 \\ 1 \\ 1\end{bmatrix}$，恰好就是证券 1 的支付，因而市场上存在冗余证券；
    \item 证券 2、3、4、5 可以保证市场完全，而且我们有
    \[
        \begin{aligned}
            X_{2}-2X_{3}+X_{4}&=\begin{bmatrix}1 \\ 0 \\ 0 \\ 0\end{bmatrix}\\
            X_{3}-2X_{4}+X_{5}&=\begin{bmatrix}0 \\ 1 \\ 0 \\ 0\end{bmatrix}\\
            X_{4}-2X_{5}&=\begin{bmatrix}0 \\ 0 \\ 1 \\ 0\end{bmatrix}\\
            X_{5}&=\begin{bmatrix}0 \\ 0 \\ 0 \\ 1\end{bmatrix}
        \end{aligned}
    \]
    \item 假设组合中 5 只交易证券的头寸分别为 $\theta_{1}$、$\theta_{2}$、$\theta_{3}$、$\theta_{4}$、$\theta_{5}$，那么状态或有组合的头寸分别为 $\theta_{1}+\theta_{2}$，$\theta_{1}+2\theta_{2}+\theta_{3}$，$\theta_{1}+3\theta_{2}+2\theta_{3}+\theta_{4}$，$\theta_{1}+4\theta_{2}+3\theta_{3}+2\theta_{4}+\theta_{5}$；
    \item 因为状态或有组合的支付是 $\mathrm{R}^{4}$ 空间的一组基，因而状态或有组合可生成支付空间。
\end{enumerate}
\end{solution}

\begin{problem}
\problemno{4.2} 银行 1 的存贷利率为 $r_{1}$ 而银行 2 的为 $r_{2}$，且 $r_{1}>r_{2}$。证明存在套利机会。
\end{problem}
\begin{solution}
从银行 2 借钱并全部存入银行 1，参与者就可以获得第 2 类套利。
\end{solution}

\begin{problem}
\problemno{4.3} 银行 $i$ 的贷款利率为 $r_{ia}$（即顾客借钱的利率）而存款利率为 $r_{ib}$，$i=1,2$。
\begin{enumerate}[label=(\alph*)]
    \item 证明无套利要求 $r_{ia}\ge r_{ib}$。
    \item 两个银行的存贷利率应该满足什么样的条件才能使市场上不存在套利机会。
\end{enumerate}
\end{problem}
\begin{solution}
\begin{enumerate}[label=(\alph*)]
    \item 设 $r_{ia}<r_{ib}$。从银行 $i$ 借入 $1$ 并存入该银行：0 期收到贷款 $1$、支付存款 $1$，净现金流为 $0$；1 期收回 $1+r_{ib}$、偿还 $1+r_{ia}$，净现金流为
    \[
        (1+r_{ib})-(1+r_{ia})=r_{ib}-r_{ia}>0.
    \]
    这是第 2 类套利。因此无套利要求 $r_{ia}\ge r_{ib}$。
    \item 无套利当且仅当最高存款利率不超过最低贷款利率：
    \[
        \max(r_{1b},r_{2b})\le\min(r_{1a},r_{2a}),
    \]
    即 $r_{1a}\ge r_{1b}$、$r_{2a}\ge r_{2b}$、$r_{1a}\ge r_{2b}$、$r_{2a}\ge r_{1b}$。

    若 $r_{1a}<r_{2b}$，从银行 1 借入 $1$ 并存入银行 2，0 期净现金流为 $0$，1 期净现金流为 $r_{2b}-r_{1a}>0$。若 $r_{2a}<r_{1b}$，从银行 2 借入并存入银行 1，1 期净现金流为 $r_{1b}-r_{2a}>0$。本行的情形已由 (a) 给出。

    反过来，设上述四个不等式成立，记 $r_{b}^{*}=\max(r_{1b},r_{2b})$、$r_{a}^{*}=\min(r_{1a},r_{2a})$，则 $r_{b}^{*}\le r_{a}^{*}$。以 $\theta_{i}\ge 0$ 表示在银行 $i$ 的借款、$\delta_{i}\ge 0$ 表示存款。0 期净流入为 $\sum\theta_{i}-\sum\delta_{i}$，1 期净流入为
    \[
        \sum_{i}\delta_{i}(1+r_{ib})-\sum_{i}\theta_{i}(1+r_{ia})
        \le(1+r_{b}^{*})\sum\delta_{i}-(1+r_{a}^{*})\sum\theta_{i}.
    \]
    若 0 期净流入为 $0$，则 $\sum\theta_{i}=\sum\delta_{i}$，1 期净流入不超过 $(r_{b}^{*}-r_{a}^{*})\sum\delta_{i}\le 0$。若 0 期净流入为正，则 $\sum\theta_{i}>\sum\delta_{i}$，1 期净流入严格为负。因此不存在套利。
\end{enumerate}
\end{solution}

\begin{problem}
\problemno{4.4} 考虑练习 4.1 中定义的市场。假设 5 只交易证券的价格分别为 $1$，$2$，$3/2$，$1$，$1/2$。
\begin{enumerate}[label=(\alph*)]
    \item 在这样的价格和支付下存在套利机会吗？
    \item 你能构建第 1、2 以及 3 类套利机会吗？
\end{enumerate}
\end{problem}
\begin{solution}
\begin{enumerate}[label=(\alph*)]
    \item 买入 1 单位的证券 4，卖出 2 单位的证券 5，0 期支付为 0，而 1 期支付为 $\begin{bmatrix}0 \\ 0 \\ 1 \\ 0\end{bmatrix}>0$，因而存在套利机会。
    \item 上面构造了第 2 类套利机会。卖出 1 单位的证券 1 和 1 单位的证券 3，同时买入 1 单位的证券 2，那么 0 期支付为 $\tfrac12$，而 1 期支付为 0，这是第 1 类套利机会。同时进行上述两种交易就可以得到第 3 类套利机会。
\end{enumerate}
\end{solution}

\begin{problem}
\problemno{4.5} 在 1 期有三个可能状态：$\{\omega_{1},\omega_{2},\omega_{3}\}$。在经济中有三只交易证券，$a$，$b$，$c$，它们的支付向量如下：
\[
    X_{a}=\begin{bmatrix}1\\1\\0\end{bmatrix},\quad
    X_{b}=\begin{bmatrix}0\\1\\1\end{bmatrix},\quad
    X_{c}=\begin{bmatrix}1/2\\1\\1/2\end{bmatrix}.
\]
这些证券的 0 期均衡价格为：$S_{a}=S_{b}=S_{c}=1$。令 $\phi=\begin{bmatrix}\phi_{1} \\ \phi_{2} \\ \phi_{3}\end{bmatrix}$ 为状态价格向量。
\begin{enumerate}[label=(\alph*)]
    \item 证明 $\begin{bmatrix}2 \\ -1 \\ 2\end{bmatrix}$ 是一个可能的状态价格向量，因为它给出了三支现有证券的价格。
    \item “资产定价基本定理”指出状态价格必须是非负的。在一个两期经济里，对于证券市场这意味着什么？并给出其解释。
    \item 上面的例子是这个原理的反例吗？为什么？
    \item 证明 $\begin{bmatrix}1/2 \\ 1/2 \\ 1/2\end{bmatrix}$ 是另外一个可能的状态价格向量。在这种情况下，所有状态价格都是正的。
    \item 解释为什么存在多个可能的状态价格向量，它们对已有证券都给出相同的价格。
\end{enumerate}
\end{problem}
\begin{solution}
\begin{enumerate}[label=(\alph*)]
    \item 因为
    \[
        2\times 1+(-1)\times 1=1,\quad (-1)\times 1+2\times 1=1,\quad 2\times\tfrac12+(-1)\times 1+2\times\tfrac12=1
    \]
    因而 $\begin{bmatrix}2 \\ -1 \\ 2\end{bmatrix}$ 是一个可能的状态价格向量；
    \item 资产定价基本定理：没有套利，当且仅当存在严格正的状态价格 $\phi\gg 0$，使每只可交易证券的价格都是 $S=X^{\mathsf T}\phi$。因此“状态价格必须非负”在这里的含义是均衡价格里不能有套利：若每一支给现有证券定价的 $\phi$ 都有负分量，市场就存在套利。
    \item 不是。因而它只是可能的状态价格之一，并不一定是真实的状态价格；定理要求的是存在一支 $\phi\gg 0$，而不是每一个解都严格为正。下面 (d) 给出的正向量说明无套利仍然成立；
    \item 因为
    \[
        \tfrac12\times 1+\tfrac12\times 1=1,\quad \tfrac12\times 1+\tfrac12\times 1=1,\quad \tfrac12\times\tfrac12+\tfrac12\times 1+\tfrac12\times\tfrac12=1
    \]
    因而 $\begin{bmatrix}\tfrac12 \\ \tfrac12 \\ \tfrac12\end{bmatrix}$ 是一个可能的状态价格向量；
    \item 因为市场是不完全的，因而可能存在多个状态价格，它们都给出已有证券的价格。
\end{enumerate}
\end{solution}

\begin{problem}
\problemno{4.6} 远期合约（Forward）是在未来（即 1 期）以当期（即 0 期）确定的价格交易某一证券或商品的合约，而在当期没有现金流产生。在当期约定的未来交易价格称作远期价格。比方说，考虑一份在 1 期以 $F_{0}$ 的价格购买 1 盎司黄金的远期合约。合约的买方同意在 1 期（合约的到期日）从卖方买入 1 盎司黄金。在当期没有现金交换。令 $S_{t}$ 为 $t=0,1$ 期黄金的价格，而 $r_{F}$ 为利率。考虑在 1 期得到 1 盎司黄金的两种方法：
\begin{itemize}
    \item 现在购买一份远期合约。在 1 期，支付远期价格 $F_{0}$ 并获得黄金。
    \item 现在支付 $S_{0}$ 买入 1 盎司黄金，并且持有它直到 1 期。假设从现在起持有黄金直到 1 期没有另外的成本或收益。
\end{itemize}
在这两种情况下，我们在 1 期同样得到 1 盎司黄金，而那时它的价值为 $S_{1}$。
\begin{enumerate}[label=(\alph*)]
    \item 由于这两种策略给予我们相同的支付 $S_{1}$，无套利原理意味着它们的成本必须一样。用当期的价格 $S_{1}$ 和 $r_{F}$ 给出远期价格。
    \item 假设持有 1 盎司黄金不仅得到它的 1 期出售收益，在 1 期还得到额外收益 $y$。$y$ 将如何影响远期价格？
\end{enumerate}
\end{problem}
\begin{solution}
\begin{enumerate}[label=(\alph*)]
    \item 第一种方法的成本是 $\dfrac{F_{0}}{1+r_{F}}$，而第二种方法的成本是 $S_{0}$，根据无套利原理，$\dfrac{F_{0}}{1+r_{F}}=S_{0}$，因此，$F_{0}=(1+r_{F})S_{0}$；
    \item 由于第二种方法还得到额外收益，由无套利原理原理有 $F_{0}<(1+r_{F})S_{0}$。（如果收益是确定性的，那么应该有 $F_{0}+y=(1+r_{F})S_{0}$）
\end{enumerate}
\end{solution}

\begin{problem}
\problemno{4.7} 英镑现时对美元的价格（即英镑/美元汇率）为 $S_{0}$（即一英镑的美元价格为本 $S_{0}$），美元利率为 $r_{A}$ 而英镑利率为 $r_{B}$。考虑一份在 1 期购买 1 英镑的远期合约。远期价格（即远期汇率）应该是多少？提示：考虑在 1 期得到 1 英镑的两种策略。
\begin{itemize}
    \item 购买一份远期合约。在 1 期，支付远期价格 $F_{0}$ 购入 1 英镑。
    \item 现在以当前汇率 $S_{0}$ 买入 $x$ 英镑并以 $r_{B}$ 的利率进行投资，在 1 期得到 1 英镑。
\end{itemize}
得到的远期价格、现时汇率和利率之间的关系就叫作抛补利率平价理论（Covered Interest Rate Parity）。
\end{problem}
\begin{solution}
第一种策略的 0 期支付是 $\dfrac{F_{0}}{1+r_{A}}$ 美元，而第二种策略的 0 期支付是 $\dfrac{S_{0}}{1+r_{B}}$ 美元，根据无套利原理，$\dfrac{F_{0}}{1+r_{A}}=\dfrac{S_{0}}{1+r_{B}}$，即 $F_{0}=\dfrac{1+r_{A}}{1+r_{B}}S_{0}$。
\end{solution}

\begin{problem}
\problemno{4.8} 考虑练习 4.1 定义的市场。假设证券 1，2，4，5 的价格分别是 $0.8$，$2$，$0.6$，$0.2$。
\begin{enumerate}[label=(\alph*)]
    \item 计算状态价格和利率。
    \item 构建等价鞅测度。
    \item 证明 5 只证券的价格都具有鞅性质。
\end{enumerate}
\end{problem}
\begin{solution}
记练习 4.1 中五只证券的支付列向量为 $X_{1},\ldots,X_{5}$。状态价格 $\phi=\begin{bmatrix}\phi_{1} \\ \phi_{2} \\ \phi_{3} \\ \phi_{4}\end{bmatrix}$ 满足 $S_{n}=\phi^{\mathsf T}X_{n}$，即
\[
    S_{1}=\phi^{\mathsf T}X_{1},\quad S_{2}=\phi^{\mathsf T}X_{2},\quad S_{4}=\phi^{\mathsf T}X_{4},\quad S_{5}=\phi^{\mathsf T}X_{5}.
\]
\begin{enumerate}[label=(\alph*)]
    \item 代入 $S_{1}=0.8$，$S_{2}=2$，$S_{4}=0.6$，$S_{5}=0.2$ 得
    \[
        \begin{aligned}
            \phi_{1}+\phi_{2}+\phi_{3}+\phi_{4}&=0.8,\\
            \phi_{1}+2\phi_{2}+3\phi_{3}+4\phi_{4}&=2,\\
            \phi_{3}+2\phi_{4}&=0.6,\\
            \phi_{4}&=0.2.
        \end{aligned}
    \]
    由后两式 $\phi_{3}=0.2$；前两式相减得 $\phi_{2}=0.2$，$\phi_{1}=0.2$。故
    \[
        \phi=\begin{bmatrix}0.2 \\ 0.2 \\ 0.2 \\ 0.2\end{bmatrix},\qquad
        S_{3}=\phi^{\mathsf T}X_{3}=1.2.
    \]
    无风险利率 $r_{F}=\bigl(\sum_{i}\phi_{i}\bigr)^{-1}-1=0.25$。
    \item 等价鞅测度 $Q_{i}=\phi_{i}/\sum_{j}\phi_{j}$，即 $Q=\begin{bmatrix}0.25 \\ 0.25 \\ 0.25 \\ 0.25\end{bmatrix}$。
    \item 两期模型里的鞅性质是风险中性定价：用等价鞅测度 $Q$ 对 1 期支付取期望，再按无风险利率折现，应当回到 0 期价格，即 $S_{n}=\mathrm{E}^{Q}[X_{n}]/(1+r_{F})$。由 $Q_{i}=\phi_{i}/\sum_{j}\phi_{j}$ 与 $1+r_{F}=1/\sum_{j}\phi_{j}$，
    \[
        \frac{\mathrm{E}^{Q}[X_{n}]}{1+r_{F}}=\phi^{\mathsf T}X_{n}.
    \]
    证券 $1$、$2$、$4$、$5$ 的价格正是用来解出 $\phi$ 的，故 $\phi^{\mathsf T}X_{n}=S_{n}$。证券 $3$ 的价格取为 $S_{3}=\phi^{\mathsf T}X_{3}$。在 $Q_{i}=1/4$、$1+r_{F}=5/4$ 下，
    \[
        \begin{aligned}
            \frac{\mathrm{E}^{Q}[X_{1}]}{1+r_{F}}&=\frac{(1+1+1+1)/4}{5/4}=0.8=S_{1},\\
            \frac{\mathrm{E}^{Q}[X_{2}]}{1+r_{F}}&=\frac{(1+2+3+4)/4}{5/4}=2=S_{2},\\
            \frac{\mathrm{E}^{Q}[X_{3}]}{1+r_{F}}&=\frac{(0+1+2+3)/4}{5/4}=1.2=S_{3},\\
            \frac{\mathrm{E}^{Q}[X_{4}]}{1+r_{F}}&=\frac{(0+0+1+2)/4}{5/4}=0.6=S_{4},\\
            \frac{\mathrm{E}^{Q}[X_{5}]}{1+r_{F}}&=\frac{(0+0+0+1)/4}{5/4}=0.2=S_{5}.
        \end{aligned}
    \]
\end{enumerate}
\end{solution}



% !TeX root = ../../金融经济学_王江_讲义及习题.tex
% 知识点：references/金融经济学/金融经济学-王江_261013.pdf，印刷页 67–81
% 习题：references/金融经济学/史上最全最完整的金融经济学答案(王江版).pdf，页 20–25（印刷页 19–24）

\zhchapter{期权：一个套利定价的例子}

\ifdefined\WJHandoutKnowledge

第 4 章给出无套利与资产定价基本定理的一般形式。本章把它们用到\textbf{期权}上：先只靠无套利得到价格的上下界与平价关系，再在完全市场（尤其是二叉树）中得到精确价格；最后说明期权本身也可以把市场完全化。

\begin{tcolorbox}[knowledge,title=5.1 期权]
记标的证券（下文统称股票，也可以是任意证券）在 1 期的支付为 $\tilde{X}$，0 期价格为 $S$。

\begin{tcolorbox}[knowledge]
\textbf{定义 5.1}\quad \textbf{欧式看涨期权}（European call）赋予买方在到期日、以执行价格 $K$ 从期权出售者处买入 1 单位标的证券的权利。\textbf{欧式看跌期权}（European put）则是以执行价格 $K$ 卖出 1 单位标的证券的权利。

在本书的两期框架里只有一个到期日，就是 1 期。$K$ 称为\textbf{执行价格}（exercise price）。因为是权利而不是义务，只有支付为正时才会执行。记欧式看涨、看跌在 1 期的支付为 $\tilde{c}$、$\tilde{p}$，则
\[
    \tilde{c}=[\tilde{X}-K]_{+},\qquad \tilde{p}=[K-\tilde{X}]_{+},
    \tag{5.1}
\]
其中 $[x]_{+}=\max\{0,x\}$。
\end{tcolorbox}

\begin{tcolorbox}[knowledge]
\textbf{定义 5.2}\quad 执行价格为 $K$、到期日为 $t$ 的\textbf{美式看涨（看跌）期权}赋予买方在到期日及此前任一日期，以价格 $K$ 买入（卖出）1 单位标的证券的权利。
\end{tcolorbox}

$S-K$ 称为看涨期权的\textbf{内在价值}（intrinsic value），看跌期权的内在价值是 $K-S$。内在价值大于、等于、小于 $0$ 时，分别称为\textbf{实值}（in the money）、\textbf{平值}（at the money）、\textbf{虚值}（out of the money）。

期权是买方与出售者（\textbf{期权立权人}，option writer）之间的零和合约：买方得到的正支付恰为出售者的付出，净供给恒为 $0$。支付又完全由标的证券的价格决定。满足这两条的证券称为\textbf{衍生证券}（derivative security）。

记 $c(S,K)$、$p(S,K)$ 为欧式看涨、看跌的 0 期价格，$C(S,K)$、$P(S,K)$ 为相应的美式价格。定价分两步：无套利给出基本性质与上下界；完全市场给出精确公式。
\end{tcolorbox}

\begin{tcolorbox}[knowledge,title=5.2 期权价格的性质和界]
欧式期权的支付只在到期日发生，价格由支付决定。影响价格的因素主要有：标的证券的价格与支付、合同条款（执行价格；本书只有一个到期日）、以及决定货币时间价值的利率。

\begin{tcolorbox}[knowledge]
\textbf{定理 5.1}\quad $c(S,K)$ 与 $p(S,K)$ 非负。

支付非负，再由无套利即得。

\textbf{定理 5.2}\quad $c(S,K)$ 对 $K$ 非增，$p(S,K)$ 对 $K$ 非减。

对 $K_1>K_2$，$[\tilde{X}-K_1]_{+}\le[\tilde{X}-K_2]_{+}$，无套利（定价算子单调）给出 $c(S,K_1)\le c(S,K_2)$。看跌类似。

\textbf{定理 5.3}\quad $c(S,K)$ 与 $p(S,K)$ 是 $K$ 的凸函数。

$[\,\cdot\,]_{+}$ 对 $K$ 凸：对 $\alpha\in(0,1)$，
\[
    [\tilde{X}-(\alpha K_1+(1-\alpha)K_2)]_{+}
    \le\alpha[\tilde{X}-K_1]_{+}+(1-\alpha)[\tilde{X}-K_2]_{+}.
\]
定价算子线性且单调，故 $c(S,K)$ 对 $K$ 凸。$[K-\tilde{X}]_{+}$ 同样。

\textbf{定理 5.4}\quad 设 $\theta\ge 0$ 为 $N$ 只证券的组合，价格向量 $S=\begin{bmatrix}S_1\\\vdots\\S_N\end{bmatrix}\gg 0$，执行价格向量 $K=\begin{bmatrix}K_1\\\vdots\\K_N\end{bmatrix}\gg 0$。则
\[
    c(S^{\mathsf T}\theta,\,K^{\mathsf T}\theta)\le\sum_{i=1}^{N}\theta_i\,c(S_i,K_i),
    \qquad
    p(S^{\mathsf T}\theta,\,K^{\mathsf T}\theta)\le\sum_{i=1}^{N}\theta_i\,p(S_i,K_i).
\]
以组合为标的的期权，不比“各证券期权按同样权重组成的组合”更贵。证明来自支付的凸性：
\[
    [(\tilde{X}-K)^{\mathsf T}\theta]_{+}
    \le\sum_{i=1}^{N}\theta_i[\tilde{X}_i-K_i]_{+}.
\]

\textbf{定理 5.5}\quad $S\ge c(S,K)$。

$\tilde{X}\ge[\tilde{X}-K]_{+}$，无套利即得。

\textbf{定理 5.6}\quad 若存在收益率为 $r_F$ 的无风险证券，则
\[
    c(S,K)\ge\Bigl[S-\frac{K}{1+r_F}\Bigr]_{+}.
\]
买入 1 股、卖出面值 $K$（现值为 $K/(1+r_F)$）的无风险证券，现值为 $S-K/(1+r_F)$，1 期支付为 $\tilde{X}-K$。而 $[\tilde{X}-K]_{+}\ge\tilde{X}-K$，故 $c(S,K)\ge S-K/(1+r_F)$。再与 $c\ge 0$ 合并即得。
\end{tcolorbox}

欧式看涨的上下界因此是
\[
    \Bigl[S-\frac{K}{1+r_F}\Bigr]_{+}\le c(S,K)\le S.
    \tag{5.2}
\]

\begin{tcolorbox}[knowledge]
\textbf{定理 5.7（看涨—看跌平价）}\quad 若存在收益率为 $r_F$ 的无风险证券，则
\[
    c(S,K)+\frac{K}{1+r_F}=p(S,K)+S.
    \tag{5.3}
\]
\end{tcolorbox}

比较两个组合：（1）1 份执行价为 $K$ 的看涨，加上面值为 $K$ 的无风险证券（现值为 $K/(1+r_F)$）；（2）1 份执行价为 $K$ 的看跌，加上 1 股股票。
\[
\begin{array}{c|cc}
 & \tilde{X}\le K & \tilde{X}>K \\ \hline
\text{组合 1} & [\tilde{X}-K]_{+}+K=K & [\tilde{X}-K]_{+}+K=\tilde{X} \\
\text{组合 2} & [K-\tilde{X}]_{+}+\tilde{X}=K & [K-\tilde{X}]_{+}+\tilde{X}=\tilde{X}
\end{array}
\]
1 期支付相同，0 期成本必相同，即 (5.3)。看跌的界可由平价从看涨的界推出。
\end{tcolorbox}

\begin{tcolorbox}[knowledge,title=5.3 美式期权以及提前执行]
\textbf{提前执行}（early exercise）是权利不是义务。持有至到期时，美式期权与相应欧式期权支付相同；既然可以等到更优时再执行，美式价格不会低于欧式价格：
\[
    C(S,K)\ge c(S,K),\qquad P(S,K)\ge p(S,K).
    \tag{5.4}
\]
若提前执行以正概率发生，则不等式严格。是否提前执行，与标的在到期前是否支付\textbf{股利}有关。此前默认证券只在 1 期有支付；0 期的支付就是股利。

\begin{tcolorbox}[example]
\begin{question}
\textbf{例 5.1}\quad 0 期股价为 $100$ 的股票支付股利 $100$（如同清算：资产变现 $100$ 并全部作为股利付出）。除息后股价降为 $0$。以此股票为标的、执行价格为 $80$ 的欧式看涨与相应美式看涨，价格分别是多少？提前执行是否最优？
\end{question}
\begin{solution}
1 期股价为 $0$，欧式看涨支付为 $0$，故其价格为 $0$。美式看涨可以在股利发放前执行：支付 $80$ 买入股票，立刻得到股利 $100$，净得 $20$。提前执行严格优于持有至到期，美式看涨价值为 $20$。
\end{solution}
\end{tcolorbox}

因此，分析提前执行时必须把股利考虑进去。先看无股利，再看有股利。

\tcbline

\textbf{A.\ 无股利时的提前执行}

无股利时，美式看涨\textbf{不会}提前执行。立即执行只得到 $S-K$，而
\[
    S-K\le S-\frac{K}{1+r_F}
    \le\Bigl[S-\frac{K}{1+r_F}\Bigr]_{+}
    \le c(S,K)\le C(S,K).
\]
提前执行的所得不超过把同一期权当作欧式期权卖掉的价值；且当 $r_F>0$ 时，在 $S>K$ 的状态下第一个不等式严格，提前执行是次优的。因而
\[
    C(S,K)=c(S,K).
\]
代价来自两处：一是货币的时间价值——行权策略“到期无论股价如何都执行”得到 $\tilde{X}-K$，现值 $S-K/(1+r_F)$，已优于现在就付出 $K$；二是到期可以不执行，$[\tilde{X}-K]_{+}$ 优于 $\tilde{X}-K$。

无股利的美式看跌则可能提前执行。提前执行放弃的是“支付可以更高时不执行”的权利，得到的是现在就收到 $K$、而不是以后才收到。由平价，
\[
    P(S,K)
    =\max\Bigl[K-S,\;p(S,K)\Bigr]
    =\max\Bigl[K-S,\;\frac{K}{1+r_F}-S+c(S,K)\Bigr].
\]
当 $K-S>K/(1+r_F)-S+c(S,K)$，即
\[
    \frac{r_F}{1+r_F}K>c(S,K)
\]
时，提前执行最优。看跌深度实值（$K\gg S$，从而看涨几乎无价值）时，这个不等式成立。

\tcbline

\textbf{B.\ 有股利时的提前执行}

设股票在 $0$ 期还支付股利 $D$，记 $S$ 为\textbf{除息价}。美式看涨在 $0$ 期可以：（1）支付 $K$ 执行并拿到股利，得到 $S+D-K$；或（2）持有至 1 期。最优策略下
\[
    C(S,D,K)=\max\bigl[S+D-K,\;c(S,K)\bigr].
\]
看跌相应为
\[
    P(S,D,K)=\max\bigl[K-S-D,\;p(S,K)\bigr].
\]
由 (5.4)，$C(S,D,K)\ge c(S,K)$、$P(S,D,K)\ge p(S,K)$。因此对美式期权，\textcolor{red}{\textbf{股利促使持有者提前执行看涨、推迟执行看跌}}。

有股利时平价也要改。比较：（1）1 份执行价为 $K$ 的看涨，加上现值为 $K/(1+r_F)+D$ 的无风险证券；（2）1 份执行价为 $K$ 的看跌，加上 1 股股票（$0$ 期得到股利 $D$，并按无风险利率持有至 1 期）。
\[
\begin{array}{c|cc}
 & \tilde{X}\le K & \tilde{X}>K \\ \hline
\text{组合 1} & K+D(1+r_F) & \tilde{X}+D(1+r_F) \\
\text{组合 2} & K+D(1+r_F) & \tilde{X}+D(1+r_F)
\end{array}
\]
支付相同，故
\[
    c(S,K)+\frac{K}{1+r_F}+D=S+p(S,K).
\]
\end{tcolorbox}

\begin{tcolorbox}[knowledge,title=5.4 完全市场中的期权定价]
若证券市场完全，状态价格 $\phi\gg 0$ 唯一。记 $S_\omega$ 为标的在状态 $\omega$ 的支付，则
\[
    c(S,K)
    =\sum_{\omega}\phi_{\omega}[S_{\omega}-K]_{+}
    =\frac{1}{1+r_F}\,\mathrm{E}^{Q}[\tilde{S}-K]_{+},
\]
其中 $Q$ 为风险中性测度。要得到更具体的公式，需要对股价过程再加假设。

\tcbline

\textbf{二叉树。}股价现值为 $S$。下一期以概率 $p$ 变为 $uS$、以概率 $1-p$ 变为 $dS$，设 $u>d$：
\begin{center}
\WJEventTree[left={$S$},stem]{uS\text{，概率 }p, dS\text{，概率 }1-p}
\end{center}
另有无风险证券，1 期支付 $1$，价格 $B=1/(1+r_F)$，利率为 $r_F$。

无套利要求 $d<1+r_F<u$。反证：若 $d\ge 1+r_F$，买入 1 股（支出 $S$），同时卖出 $S(1+r_F)$ 份无风险证券（收入恰为 $S$），0 期净支出为 $0$。1 期上涨时支付 $[u-(1+r_F)]S>0$，下跌时 $[d-(1+r_F)]S\ge 0$，构成套利。故 $d<1+r_F$；同理 $u>1+r_F$。

两个状态、股票与债券的支付线性无关，市场完全，状态价格唯一。$\phi=\begin{bmatrix}\phi_u\\\phi_d\end{bmatrix}\gg 0$ 满足
\[
    S=\phi_u\,uS+\phi_d\,dS,\qquad B=\phi_u+\phi_d,
\]
解得
\[
    \phi_u=\frac{1}{1+r_F}\frac{1+r_F-d}{u-d},
    \qquad
    \phi_d=\frac{1}{1+r_F}\frac{u-(1+r_F)}{u-d}.
\]
欧式看涨价格为
\[
    c=\phi_u[uS-K]_{+}+\phi_d[dS-K]_{+}.
    \tag{5.5}
\]

\tcbline

\textbf{复制。}只有两个状态时，股票和债券使任何其他证券成为冗余证券，其价格等于复制组合的价格。到期支付
\[
    c_u=[uS-K]_{+},\qquad c_d=[dS-K]_{+}.
\]
持有 $\theta_S$ 股与 $\theta_B$ 份面值 $1$ 的债券，使两个状态下的支付分别等于 $c_u$、$c_d$：
\[
    \theta_S\,uS+\theta_B=c_u,\qquad \theta_S\,dS+\theta_B=c_d.
\]
解为
\[
    \theta_S=\frac{c_u-c_d}{(u-d)S},
    \qquad
    \theta_B=\frac{uc_d-dc_u}{u-d}.
\]
$\theta_S$ 称为\textbf{对冲比}（hedge ratio，期权的 $\delta$）。组合在 0 期的价值为 $v=\theta_S S+\theta_B/(1+r_F)$。无套利要求 $c=v$，代回即得 (5.5)。

\tcbline

\textbf{风险中性概率。}定义等价鞅测度
\[
    q=\frac{\phi_u}{\phi_u+\phi_d}=\frac{1+r_F-d}{u-d},
    \qquad
    1-q=\frac{\phi_d}{\phi_u+\phi_d}=\frac{u-(1+r_F)}{u-d}.
\]
则
\[
    c=\frac{1}{1+r_F}\bigl[qc_u+(1-q)c_d\bigr]
    =\frac{\mathrm{E}^{Q}[c_1]}{1+r_F}.
\]
记 $\hat{c}\equiv c/B$、$\hat{S}\equiv S/B$（用无风险债券价格把价格规范化）。在 $Q$ 下 $\hat{c}$ 与 $\hat{S}$ 是鞅：$\hat{c}_t=\mathrm{E}^{Q}[\hat{c}_{t+1}]$，$\hat{S}_t=\mathrm{E}^{Q}[\hat{S}_{t+1}]$。这就是\textbf{二叉树期权定价模型}（binomial model）：完全市场里，期权价格由股票和债券的价格经无套利唯一决定，不需要真实概率 $p$。
\end{tcolorbox}

\begin{tcolorbox}[knowledge,title=5.5 期权与市场完全化]
完全市场中期权是冗余证券，价格由原生证券复制。若期权不是冗余的，引入期权可以扩大可生成的支付，甚至把市场完全化。

\tcbline

\textbf{A.\ 蝴蝶头寸}\quad 同一标的、同一到期日、执行价格不同的欧式看涨可以拼出特定支付。\textbf{蝴蝶头寸}（butterfly）为：买入 1 份执行价 $K-\delta$ 的看涨，卖出 2 份执行价 $K$ 的看涨，买入 1 份执行价 $K+\delta$ 的看涨。1 期支付为
\[
    [\tilde{S}-(K-\delta)]_{+}-2[\tilde{S}-K]_{+}+[\tilde{S}-(K+\delta)]_{+}
    =
    \begin{cases}
        0, & \tilde{S}\le K-\delta,\\
        \tilde{S}-(K-\delta), & K-\delta<\tilde{S}\le K,\\
        -\tilde{S}+(K+\delta), & K<\tilde{S}<K+\delta,\\
        0, & \tilde{S}\ge K+\delta.
    \end{cases}
\]
只有 $\tilde{S}\in(K-\delta,\,K+\delta)$ 时支付非零，形状与状态或有证券类似。

\tcbline

\textbf{B.\ 用期权完全化市场}\quad 若一只证券（或组合）的支付 $\{X_\omega:\omega\in\Omega\}$ 满足 $\omega\neq\omega'$ 时 $X_\omega\neq X_{\omega'}$，称为具有\textbf{状态分离}（state separating）的支付，该证券称为\textbf{状态指数证券}（state-index security）。不妨按支付从小到大编号：$X_1<X_2<\cdots<X_\Omega$。

以它为标的、执行价格为 $X_\omega$ 的欧式看涨，支付为
\[
    [\tilde{X}-X_\omega]_{+}
    =\bigl[0,\ldots,0,\,X_{\omega+1}-X_\omega,\,\ldots,\,X_\Omega-X_\omega\bigr]^{\mathsf T}.
\]
下面 $\Omega$ 只证券的支付矩阵满秩，因而市场完全：1 份状态指数证券，加上执行价格分别为 $X_1,\ldots,X_{\Omega-1}$ 的 $\Omega-1$ 份欧式看涨，
\[
    X=
    \begin{bmatrix}
        X_1 & 0 & 0 & \cdots & 0 \\
        X_2 & X_2-X_1 & 0 & \cdots & 0 \\
        X_3 & X_3-X_1 & X_3-X_2 & \cdots & 0 \\
        \vdots & \vdots & \vdots & \ddots & \vdots \\
        X_\Omega & X_\Omega-X_1 & X_\Omega-X_2 & \cdots & X_\Omega-X_{\Omega-1}
    \end{bmatrix}.
\]

取一个方便计算的例子：状态 $\omega=1,\ldots,\Omega$ 下指数证券支付为 $\delta,2\delta,\ldots,\Omega\delta$。执行价格为 $0,\delta,2\delta,\ldots,(\Omega-1)\delta$ 的看涨（执行价 $0$ 就是标的本身）的支付矩阵为下三角满秩阵
\[
    X=
    \begin{bmatrix}
        \delta & 0 & \cdots & 0 \\
        2\delta & \delta & \cdots & 0 \\
        \vdots & \vdots & \ddots & \vdots \\
        \Omega\delta & (\Omega-1)\delta & \cdots & \delta
    \end{bmatrix}.
\]
执行价为 $n\delta$ 的看涨记为 $X_{\cdot,n}$。蝴蝶头寸“买入执行价 $0$、卖出 2 份执行价 $\delta$、买入执行价 $2\delta$”的支付为
\[
    X_p=(X_{\cdot,0}-X_{\cdot,1})-(X_{\cdot,1}-X_{\cdot,2})
    =[\delta,\delta,\ldots,\delta]^{\mathsf T}-[0,\delta,\ldots,\delta]^{\mathsf T}
    =[\delta,0,\ldots,0]^{\mathsf T},
\]
即 $\delta$ 份状态 $1$ 的或有证券。一般地，执行价 $(i-1)\delta$、$i\delta$、$(i+1)\delta$ 的蝴蝶头寸只在状态 $i$ 支付 $\delta$。因此这组期权可以复制任意状态或有要求权。

记 $\phi_i$ 为状态 $i$ 的状态价格，$K_i=i\delta$。由蝴蝶头寸的价格，
\[
    \phi_i
    =\frac{1}{\delta}\Bigl\{\bigl[c(K_{i-1})-c(K_i)\bigr]-\bigl[c(K_i)-c(K_{i+1})\bigr]\Bigr\}
    =\frac{c(K_{i-1})-2c(K_i)+c(K_{i+1})}{\delta}.
\]
支付为 $\{X_{\omega_i}\}$ 的资产价格为 $S=\sum_i\phi_i X_{\omega_i}$，也即
\[
    S=\sum_i\frac{c(K_{i+1})-2c(K_i)+c(K_{i-1})}{\delta^{2}}\,X_{\omega_i}\,\delta.
\]
状态与 $K_i$ 一一对应，支付可写成 $K$ 的函数 $X(K)$。指数证券连续分布且 $c$ 二阶可导时，取 $\delta=\mathrm{d}K$，二阶差商趋于 $\partial^{2}c/\partial K^{2}$，
\[
    S=\int_{0}^{\infty}\frac{\partial^{2}c(K)}{\partial K^{2}}X(K)\,\mathrm{d}K.
\]
这是 Breeden--Litzenberger 公式。定义风险中性密度
\[
    q(K)
    =\frac{\partial^{2}c(K)/\partial K^{2}}{\displaystyle\int_{0}^{\infty}\partial^{2}c(K)/\partial K^{2}\,\mathrm{d}K}.
\]
无风险债券价格为
\[
    B=\int_{0}^{\infty}\frac{\partial^{2}c(K)}{\partial K^{2}}\,\mathrm{d}K=\frac{1}{1+r_F},
\]
故规范化价格是鞅：
\[
    \frac{S}{B}=\int_{0}^{\infty}q(K)\,X(K)\,\mathrm{d}K.
\]
\end{tcolorbox}

\begin{tcolorbox}[knowledge,title=本章小结]
\begin{enumerate}
    \item 欧式看涨（看跌）是在到期日以执行价格 $K$ 买入（卖出）标的的权利，支付为 $[\tilde{X}-K]_{+}$（$[K-\tilde{X}]_{+}$）。美式期权允许到期日前执行。
    \item 无套利给出欧式看涨的界 $\bigl[S-K/(1+r_F)\bigr]_{+}\le c(S,K)\le S$。
    \item 看涨—看跌平价：$c(S,K)+K/(1+r_F)=p(S,K)+S$。有股利 $D$（$S$ 为除息价）时改为 $c+K/(1+r_F)+D=S+p$。
    \item 美式价格不低于欧式价格。无股利时美式看涨不会提前执行，$C=c$；美式看跌在深度实值时可能提前执行。股利促使提前执行看涨、推迟执行看跌。
    \item 市场完全时，期权由状态价格定价。二叉树中 $d<1+r_F<u$，股票与债券唯一决定状态价格、对冲比和风险中性概率 $q=(1+r_F-d)/(u-d)$，真实概率 $p$ 不进入价格。
    \item 期权不是冗余证券时，状态指数证券上不同执行价的看涨（蝴蝶头寸）可以复制阿罗证券，从而把市场完全化；连续情形下状态价格密度由看涨价格对执行价的二阶导数给出。
\end{enumerate}
\end{tcolorbox}
\fi

% !TeX root = ../../金融经济学_王江_讲义及习题.tex
% 05_期权 课后习题（由章节文件拆分）

\begin{problem}
\problemno{5.1} 1 期有两个可能状态 $a$ 和 $b$，发生概率分别为 $\pi$ 和 $1-\pi$。存在一只收益率为 $r_{F}$ 的无风险证券。还存在一只风险证券，它的现时价格为 $S$，在 1 期支付红利 $\delta S$。它在 1 期的除权价（即红利后的价格）在状态 $a$ 下是 $(1+r_{F}+u)S$、在状态 $b$ 下是 $(1+r_{F}-d)S$。
\begin{enumerate}[label=(\alph*)]
    \item 为使市场中不存在套利机会，$\pi$，$r_{F}$，$u$，$d$ 和 $\delta$ 应该满足什么条件？
    \item 1 期到期、执行价为 $K$ 的欧式买权的 0 期价格是多少？
    \item 美式买权的价格是多少？比较它与相应的欧式买权价格。什么时候提前执行是最优的？
    \item 将买权换成卖权，回答 (b) 和 (c)。
\end{enumerate}
\end{problem}
\begin{solution}
\begin{enumerate}[label=(\alph*)]
    \item $u+\delta>1+r_{F}>d+\delta$；
    \item 令状态价格向量为 $\phi=\begin{bmatrix}\phi_{a} \\ \phi_{b}\end{bmatrix}$，那么有
    \[
        \begin{aligned}
            \phi_{a}+\phi_{b}&=\frac{1}{1+r_{F}}\\
            (u+\delta)S\phi_{a}+(d+\delta)S\phi_{b}&=S
        \end{aligned}
    \]
    求解方程组可得状态价格为
    \[
        \begin{aligned}
            \phi_{a}&=\frac{1}{1+r_{F}}\frac{(1+r_{F})-(d+\delta)}{u-d}\\
            \phi_{b}&=\frac{1}{1+r_{F}}\frac{(u+\delta)-(1+r_{F})}{u-d}
        \end{aligned}
    \]
    因而期权的价格为
    \[
        c=\phi_{a}[uS-K]_{+}+\phi_{b}[dS-K]_{+}=
        \begin{cases}
            0 & uS<K\\
            \phi_{a}(uS-K) & dS\le K\le uS\\
            \phi_{a}(uS-K)+\phi_{b}(dS-K) & K<dS
        \end{cases}
    \]
    \item 对于美式期权来说，如果在 0 期执行，那么获得的收益是 $[S-K]_{+}$，如果在 1 期执行，那么收益和欧式期权一样，因而 $C=\max\{[S-K]_{+},c\}$。当且仅当 $[S-K]_{+}>c$ 时提前执行是最优的。当 $\delta$ 很大时，股价 $S$ 中含有股利的成份越多，因而越有可能提前执行。
    \item 对于欧式看跌期权，价格为
    \[
        p=\phi_{a}[K-uS]_{+}+\phi_{b}[K-dS]_{+}=
        \begin{cases}
            0 & K<dS\\
            \phi_{b}(K-dS) & dS\le K\le uS\\
            \phi_{a}(K-uS)+\phi_{b}(K-dS) & K>uS
        \end{cases}
    \]
    而对于美式看跌期权，$P=\max\{[K-S]_{+},p\}$。当 $\delta$ 越大时，股价 $S$ 中含有股利的成份越多，因而越有可能不提前执行。
\end{enumerate}
\end{solution}

\begin{problem}
\problemno{5.2} 在 1 期有三个可能状态：$\{\omega_{1},\omega_{2},\omega_{3}\}$，在经济中有三只交易证券，它们的支付向量如下：
\[
    X_{a}=\begin{bmatrix}1\\1\\0\end{bmatrix},\quad
    X_{b}=\begin{bmatrix}0\\1\\1\end{bmatrix},\quad
    X_{c}=\begin{bmatrix}1/2\\1\\1/2\end{bmatrix}.
\]
这些证券的 0 期均衡价格是 $S_{a}=S_{b}=S_{c}=1$。令 $\phi=\begin{bmatrix}\phi_{1} \\ \phi_{2} \\ \phi_{3}\end{bmatrix}$ 为状态价格向量。[这与前一章练习 4.5 定义的经济一样。]考虑一份标的资产证券为 $a$、执行价格为 $1/2$ 的欧式买权。它的支付是什么？你能确定它的价格吗？
\end{problem}
\begin{solution}
它的支付是 $\begin{bmatrix}1/2\\1/2\\0\end{bmatrix}$，因为它的支付恰为证券 $a$ 的一半，因而它的价格也是一半即 $\tfrac12$。
\end{solution}

\begin{problem}
\problemno{5.3} 蝴蝶头寸是看涨期权的组合：（1）买入 1 份执行价格为 $K-\delta$ 的看涨期权，（2）卖出 2 份执行价格为 $K$ 的看涨期权，（3）买入 1 份执行价格为 $K+\delta$ 的看涨期权。因此，蝴蝶头寸可由两个参数 $K$ 和 $\delta$ 刻划。比较两个蝴蝶头寸，$K$ 相同但 $\delta$ 不同。证明 $\delta$ 较高的蝴蝶头寸的价值较高。
\end{problem}
\begin{solution}
假设 $\delta_{1}>\delta_{0}$，比较两只期权的支付。当 1 期价格 $S$ 大于 $K+\delta$ 或者小于 $K-\delta$ 时，两者的支付都为 0；当 $S\in[K-\delta_{1},K-\delta_{0}]$ 时 $X_{1}=S-(K-\delta_{1})>X_{0}=0$；当 $S\in[K-\delta_{0},K]$ 时 $X_{1}=S-(K-\delta_{1})>S-(K-\delta_{0})=X_{0}$；当 $S\in[K,K+\delta_{0}]$ 时 $X_{1}=-S+(K+\delta_{1})>-S+(K+\delta_{0})=X_{0}$；当 $S\in[K+\delta_{0},K+\delta_{1}]$ 时 $X_{1}=-S+(K+\delta_{1})>0=X_{0}$；根据无套利原理，$\delta$ 较高的期权价格较高。
\end{solution}

\begin{problem}
\problemno{5.4} 对于参数为 $K$ 和 $\delta$ 的蝴蝶头寸，给出它的价格的一个上界和下界。
\end{problem}
\begin{solution}
一个上界是 $[S-(K-\delta)]_{+}-[S-(K+\delta)]_{+}$，因为期权的支付是非负的，因而 0 是它的一个下界。
\end{solution}

\begin{problem}
\problemno{5.5} 股票价格具有如下形式的二叉树结构：
\begin{center}
\begin{tikzpicture}[x=1.35em,y=1.15em,baseline=(current bounding box.center)]
    \node[anchor=east] at (-0.15,0) {$10$};
    \draw[thick] (0,0) -- (1.1,0.7);
    \draw[thick] (0,0) -- (1.1,-0.7);
    \draw[thick] (1.1,0.35) -- (1.1,1.05);
    \draw[thick] (1.1,0.7) -- (1.85,0.7);
    \draw[thick] (1.1,-0.7) -- (1.85,-0.7);
    \node[anchor=west] at (2.0,0.7) {$12$};
    \node[anchor=west] at (2.0,-0.7) {$8$};
\end{tikzpicture}
\end{center}
两个可能结果的发生概率相等，无风险利率是 $10\%$。
\begin{enumerate}[label=(\alph*)]
    \item 为一个以此股票为标的资产、执行价格为 9 的看跌期权定价。
    \item 假设你得到一个新的利好信息：股票价格上升到 12 的概率大于 $1/2$，而股票价格仍然是 10。你愿意支付的看跌期权的价格是上涨还是下跌？解释原因。
\end{enumerate}
\end{problem}
\begin{solution}
\begin{enumerate}[label=(\alph*)]
    \item 由第 1 题中公式可以得到状态价格为 $\phi=[\tfrac{15}{22},\tfrac{5}{22}]$，因而看跌期权的价格为 $\dfrac{5}{22}\times[9-8]=\dfrac{5}{22}$；
    \item 由于股价不变，状态价格也不会变化，因而看跌期权的价格是不变的。
\end{enumerate}
\end{solution}

\begin{problem}
\problemno{5.6} 二叉树模型可以拓展到多于两个时期的情形。考虑拓展到三个时期 0，1，2 的情况。股票价格服从如下的二叉树过程：
\begin{center}
\begin{tikzpicture}[x=1.6em,y=1.15em,baseline=(current bounding box.center)]
    \node[anchor=east] at (0,0) {$10$};
    \draw[thick] (0.15,0) -- (1.5,1.6);
    \draw[thick] (0.15,0) -- (1.5,-1.6);
    \node[anchor=east] at (1.5,1.6) {$12$};
    \node[anchor=east] at (1.5,-1.6) {$8$};
    \draw[thick] (1.65,1.6) -- (3.2,2.7);
    \draw[thick] (1.65,1.6) -- (3.2,0.7);
    \draw[thick] (1.65,-1.6) -- (3.2,-0.7);
    \draw[thick] (1.65,-1.6) -- (3.2,-2.7);
    \draw[thick] (3.2,2.2) -- (3.2,3.2);
    \draw[thick] (3.2,2.7) -- (4.0,2.7);
    \draw[thick] (3.2,0.7) -- (4.0,0.7);
    \draw[thick] (3.2,-0.7) -- (4.0,-0.7);
    \draw[thick] (3.2,-2.7) -- (4.0,-2.7);
    \node[anchor=west] at (4.15,2.7) {$14.4$};
    \node[anchor=west] at (4.15,0.7) {$9.6$};
    \node[anchor=west] at (4.15,-0.7) {$9.6$};
    \node[anchor=west] at (4.15,-2.7) {$6.4$};
\end{tikzpicture}
\end{center}
每一期间（即，从 0 期到 1 期或从 1 期到 2 期）的利率都是 $10\%$。考虑一份在 2 期到期、执行价格为 9 的看涨期权。
\begin{enumerate}[label=(\alph*)]
    \item 描述这份看涨期权在 2 期的支付。
    \item 求解看涨期权在 1 期、股票价格分别为 12 或 8 时的价格。
    \item 求解看涨期权在 0 期的价格。
    \item 如果是美式看涨期权，它在 0 期的价格是多少？
\end{enumerate}
\end{problem}
\begin{solution}
\begin{enumerate}[label=(\alph*)]
    \item 支付是
    \[
        \begin{bmatrix}5.4\\0.6\\0.6\\0\end{bmatrix}
    \]
    \item 可以求得每个节点的状态价格都和第 5 题中的一样。当股价为 12 时，期权价格为 $\dfrac{15}{22}\times 5.4+\dfrac{5}{22}\times 0.6=3.82$；当股价为 8 时，期权价格为 $\dfrac{15}{22}\times 0.6=0.41$；
    \item 看涨期权的 0 期价格为 $\dfrac{15}{22}\times 3.82+\dfrac{5}{22}\times 0.41=2.70$；
    \item 因为没有股利，美式看涨期权不会提前执行，价格也是 2.70。
\end{enumerate}
\end{solution}

\begin{problem}
\problemno{5.7*} 在练习 5.6 描述的经济中，考虑一份以股票为标的资产、在 2 期到期、执行价格为 14 的看跌期权。
\begin{enumerate}[label=(\alph*)]
    \item 如果它是欧式期权，求解它的 0 期价格。
    \item 如果它是美式期权，求解它的 0 期价格。
\end{enumerate}
\end{problem}
\begin{solution}
\begin{enumerate}[label=(\alph*)]
    \item 期权的 2 期支付是 $\begin{bmatrix}0\\4.4\\4.4\\7.6\end{bmatrix}$；在 1 期，当股价为 12 时，期权价格为 $\dfrac{5}{22}\times 4.4=1$；当股价为 8 时，期权价格为 $\dfrac{15}{22}\times 4.4+\dfrac{5}{22}\times 7.6=4.73$；因而期权的 0 期价格为 $\dfrac{15}{22}\times 1+\dfrac{5}{22}\times 4.73=1.76$。
    \item 期权的 2 期支付是 $\begin{bmatrix}0\\4.4\\4.4\\7.6\end{bmatrix}$；在 1 期如果不提前执行，当股价为 12 时，期权价格为 $\dfrac{5}{22}\times 4.4=1$；当股价为 8 时，期权价格为 $\dfrac{15}{22}\times 4.4+\dfrac{5}{22}\times 7.6=4.73$；而如果在 1 期提前执行，那么期权的支付分别为 2 和 6，因而参与者会提前执行，期权在 1 期的支付为 $\begin{bmatrix}2 \\ 6\end{bmatrix}$；如果在 0 期不提前执行，期权的价格为 $\dfrac{15}{22}\times 2+\dfrac{5}{22}\times 6=2.73$，而如果在 0 期提前执行，参与者的收益是 4，因此参与者在 0 期就执行期权，期权在 0 期的价格是 4。
\end{enumerate}
\end{solution}

\begin{problem}
\problemno{5.8*} 这里的经济与练习 5.6 描述的经济一样，只是这里的参数更一般性：
\begin{center}
\begin{tikzpicture}[x=1.6em,y=1.15em,baseline=(current bounding box.center)]
    \node[anchor=east] at (0,0) {$S$};
    \draw[thick] (0.15,0) -- (1.5,1.6);
    \draw[thick] (0.15,0) -- (1.5,-1.6);
    \node[anchor=east] at (1.5,1.6) {$uS$};
    \node[anchor=east] at (1.5,-1.6) {$dS$};
    \draw[thick] (1.65,1.6) -- (3.2,2.7);
    \draw[thick] (1.65,1.6) -- (3.2,0.7);
    \draw[thick] (1.65,-1.6) -- (3.2,-0.7);
    \draw[thick] (1.65,-1.6) -- (3.2,-2.7);
    \draw[thick] (3.2,2.2) -- (3.2,3.2);
    \draw[thick] (3.2,2.7) -- (4.0,2.7);
    \draw[thick] (3.2,0.7) -- (4.0,0.7);
    \draw[thick] (3.2,-0.7) -- (4.0,-0.7);
    \draw[thick] (3.2,-2.7) -- (4.0,-2.7);
    \node[anchor=west] at (4.15,2.7) {$u^{2}S$};
    \node[anchor=west] at (4.15,0.7) {$udS$};
    \node[anchor=west] at (4.15,-0.7) {$duS$};
    \node[anchor=west] at (4.15,-2.7) {$d^{2}S$};
\end{tikzpicture}
\end{center}
每一期间的利率都是 $r_{F}$ 且 $0<d<1+r_{F}<u$。
\begin{enumerate}[label=(\alph*)]
    \item 求解在每一时期考虑下一期的两种可能情形（即股价上升或下降）的风险中性概率。
    \item 用 $S$，$u$，$d$ 和 $r_{F}$ 表示欧式看跌期权在 0 期的价格。
    \item 证明在 2 期到期的美式看涨期权的价格比具有相同执行价格、但在 1 期到期的美式看涨期权的价格要高。
\end{enumerate}
\end{problem}
\begin{solution}
\begin{enumerate}[label=(\alph*)]
    \item 状态价格在每期都是一样的，
    \[
        \begin{aligned}
            \phi_{u}&=\frac{1}{1+r_{F}}\frac{(1+r_{F})-d}{u-d}\\
            \phi_{d}&=\frac{1}{1+r_{F}}\frac{u-(1+r_{F})}{u-d}
        \end{aligned}
    \]
    因此，每一个节点处的风险中性概率也是一样的，
    \[
        q_{u}=\frac{(1+r_{F})-d}{u-d},\quad q_{d}=\frac{u-(1+r_{F})}{u-d}
    \]
    \item 欧式看跌期权在 0 期的价格为
    \[
        \begin{aligned}
            &\phi_{u}\bigl(\phi_{u}[K-u^{2}S]_{+}+\phi_{d}[K-udS]_{+}\bigr)+\phi_{d}\bigl(\phi_{u}[K-udS]_{+}+\phi_{d}[K-d^{2}S]_{+}\bigr)\\
            &=\phi_{u}^{2}[K-u^{2}S]_{+}+2\phi_{u}\phi_{d}[K-udS]_{+}+\phi_{d}^{2}[K-d^{2}S]_{+}
        \end{aligned}
    \]
    \item 因为美式期权可以提前执行，在 2 期到期的期权既可以选在 1 期执行也可以选择在 2 期执行，而在 1 期到期的美式看涨期权只会选择在 1 期执行（虽然都可以在 0 期执行，但是在没有股利的情况下，美式看涨期权是不会提前执行的），因而在 2 期到期的期权价格更高。
\end{enumerate}
\end{solution}

\begin{problem}
\problemno{5.9*} 考虑练习 5.8 描述的经济。假设一只证券在 2 期的支付是股票价格从 0 期到 2 期所达到的最大值。（这是一个关于回望期权的例子，它的支付依赖于标的证券的价格路径）。
\begin{enumerate}[label=(\alph*)]
    \item 写出它在 2 期的支付。
    \item 导出它的 0 期价格。
    \item 证明用在 2 期支付一美元的无风险债券的价格正规化以后，这只证券的价格过程在等价鞅测度下是鞅。
\end{enumerate}
\end{problem}
\begin{solution}
这里只考虑 $ud<1$ 的情形，其它情形可以类似考虑。
\begin{enumerate}[label=(\alph*)]
    \item 期权在 2 期的支付是
    \[
        \begin{bmatrix}u^{2}S\\uS\\S\end{bmatrix}
    \]
    \item 0 期价格为 $\phi_{u}^{2}u^{2}S+\phi_{u}\phi_{d}(u+1)S+\phi_{d}^{2}S$；
    \item 在 1 期，当股价为 $uS$ 时，
    \[
        \mathrm{E}^{Q}[\tilde{X}_{2}]=\frac{1}{\phi_{u}+\phi_{d}}(\phi_{u}u^{2}S+\phi_{d}uS)
    \]
    而正规化后的价格为
    \[
        (\phi_{u}u^{2}S+\phi_{d}uS)\Big/\frac{1}{1+r_{F}}=(1+r_{F})(\phi_{u}u^{2}S+\phi_{d}uS)=\frac{1}{\phi_{u}+\phi_{d}}(\phi_{u}u^{2}S+\phi_{d}uS)
    \]
    可以验证其它各节点均如此，因而期权的价格过程为鞅。
\end{enumerate}
\end{solution}

\begin{problem}
\problemno{5.10} 经济在 1 期有 3 个状态。状态一指数证券在三个状态下的支付分别为 1，2，3：
\begin{center}
\WJEventTree[stem]{1,2,3}
\end{center}
并有两只以指数证券为标的资产、执行价格分别为 1 和 2 的看涨期权。用指数证券和上述期权复制如下支付：
\begin{center}
\WJEventTree[stem]{1,1,1}
\qquad
\WJEventTree[stem]{3,2,1}
\qquad
\WJEventTree[stem]{1,0,0}
\end{center}
\end{problem}
\begin{solution}
记状态指数证券的支付向量为 $X_{0}=\begin{bmatrix}1\\2\\3\end{bmatrix}$，执行价格为 1 的指数期权的支付向量为 $X_{1}=\begin{bmatrix}0\\1\\2\end{bmatrix}$，执行价格为 2 的指数期权的支付向量为 $X_{2}=\begin{bmatrix}0\\0\\1\end{bmatrix}$ 那么我们可得：
\[
    X_{0}-X_{1}=\begin{bmatrix}1\\1\\1\end{bmatrix},\quad
    3X_{0}-4X_{1}=\begin{bmatrix}3\\2\\1\end{bmatrix},\quad
    X_{0}+X_{2}-2X_{1}=\begin{bmatrix}1\\0\\0\end{bmatrix}.
\]
\end{solution}



% !TeX root = ../../金融经济学_王江_讲义及习题.tex
% 知识点：references/金融经济学/金融经济学-王江_261013.pdf，印刷页 87–95
% 课堂补充：references/金融经济学/金融经济学_第6章 期望效用函数.pptx（夏晖，2026 年 10 月），共 34 页；含图片公式核对。
% 习题：references/金融经济学/史上最全最完整的金融经济学答案(王江版).pdf，页 26–27（印刷页 25–26）

\lecture{期望效用函数}

\ifdefined\WJHandoutKnowledge

无套利给出的是证券价格之间的关系：由可交易证券的价格和支付可以定出状态价格的集合，也可以给冗余证券定价。它几乎不依赖经济的具体特征，因此也回答不了这些价格本身由什么决定。以二叉树为例，无套利能给出期权价格以及 $u>1+r_F>d$，但不能告诉股票过程和无风险利率从何而来，也不能把所有证券的价格都定出来。要追问价格水平以及证券在资源配置中的作用，必须回到第 2 章的均衡分析，而对偏好、风险分布和市场结构作出更具体的假设。本章从偏好开始。

第 2 章用不满足性、连续性和凸性等条件刻画偏好；一般效用表示的拟凹性与本章伯努利效用的凹性需要区分，$U(c)$ 的具体形式仍很宽。一般的 $U(c)$ 同时依赖两件事：未来状态的概率，以及各状态下对消费的偏好。本章用\textbf{独立性公理}把二者分开，得到期望效用；再加上状态独立、时间可加等假设，得到后面各章使用的形式；最后指出这些假设排除了哪些行为，以及文献里常见的拓展。

\begin{tcolorbox}[knowledge,title=6.1 期望效用函数]
先看 1 期没有不确定性、两期消费都确定的情形，效用记为 $u(c_0,c_1)$。若 1 期有两个等可能状态 $a$ 与 $b$，消费分别为 $c_{1a}$、$c_{1b}$。状态 $a$ 下消费路径 $(c_0,c_{1a})$ 的效用记为 $u_a(c_0,c_{1a})$，下标表示这条路径的效用可以依赖状态。一个直接的做法是用概率作权数，取各消费路径效用的平均值：
\[
    \tfrac12 u_a(c_0,c_{1a})+\tfrac12 u_b(c_0,c_{1b}).
\]
一般地，若上述分离成立，效用就是各消费路径效用的期望，
\[
    U(c)=\sum_{\omega\in\Omega}p_{\omega}\,u_{\omega}(c_0,c_{1\omega}).
    \tag{6.1}
\]
教材 (6.1) 把概率记为 $\pi_\omega$。这种 $U$ 称为\textbf{期望效用函数}（expected utility function）。概率以线性权重进入总效用；本书另假定这些概率外生给定、为所有人共同知晓，对消费本身的喜好则因人而异。期望效用也可配合不同人的主观概率使用，“共同概率”并不是期望效用形式本身强制的。它还蕴涵下面的公理。

\begin{tcolorbox}[knowledge]
\textbf{公理 4（独立性公理）}\quad 设消费计划 $c$ 与 $c'$ 在某一状态上的消费路径同为 $x$，且 $c\succeq c'$。把这个共同路径 $x$ 换成另一路径 $y$ 后，$c$ 与 $c'$ 的排序不变。
\end{tcolorbox}

含义：若 $c$ 与 $c'$ 只在状态 $\omega$ 上同为路径 $x$，且 $c\succeq c'$，则剩余状态下路径的综合比较支持 $c$，不要求它在每一个剩余状态都更好。把 $x$ 换成 $y$ 后排序不变，说明其他状态下的比较不受状态 $\omega$ 上发生什么的影响。各互斥状态下的效用彼此独立。

两状态、等概率时可以直接验证。设
\[
    c=\begin{bmatrix}c_0\\ c_{1a}\\ c_{1b}\end{bmatrix},\qquad
    c'=\begin{bmatrix}c_0\\ c'_{1a}\\ c_{1b}\end{bmatrix},
\]
二者在状态 $b$ 上的路径相同。若 $c\succeq c'$，则
\[
    U(c)=\tfrac12\bigl[u_a(c_0,c_{1a})+u_b(c_0,c_{1b})\bigr]
    \ge\tfrac12\bigl[u_a(c_0,c'_{1a})+u_b(c_0,c_{1b})\bigr]=U(c').
\]
把状态 $b$ 的共同路径换成 $(c_0,c'_{1b})$，得到 $c''$、$c'''$。$u_b$ 在两边的贡献仍相同，故 $c''\succeq c'''$。

各消费路径与互斥的状态相对应：\textbf{一个状态发生意味着其余状态不发生，因此不同路径之间可以分开比较}。这是独立性公理说得通的一个理由。

% PPT 第 6–8 页：替换共同后果的含义与数字验证。
\textbf{独立性的对象是偏好排序。}这里将两个计划在某状态的\textbf{共同后果同时替换}，其余不同部分保持不变；它不表示“只改变一个计划的消费也不影响排序”，也不要求证券收益等随机变量在统计上独立。互斥性说明为何可以考虑这种分离，但独立性本身仍是一项可检验、也可能被违反的偏好假设。

\textbf{课堂数字验证（PPT 第 8 页）：}固定相同的当前消费，未来两状态等概率，$u(x)=\sqrt{x}$。原来的未来消费为 $c=(100,16)$、$c'=(64,16)$；把共同的雨天消费 16 同时换成 36，得到 $c''=(100,36)$、$c'''=(64,36)$。
\[
\begin{array}{c|cc|c}
 \text{计划}&\text{晴天消费}&\text{雨天消费}&\text{未来期望效用}\\\hline
 c&100&16&\tfrac12(10+4)=7\\
 c'&64&16&\tfrac12(8+4)=6\\
 c''&100&36&\tfrac12(10+6)=8\\
 c'''&64&36&\tfrac12(8+6)=7
\end{array}
\]
两个比较的效用差均为 $\tfrac12(10-8)=1$。共同部分改变了总效用水平，却没有改变两计划之间的排序。


\begin{tcolorbox}[example]
\begin{question}
\textbf{例 6.1}\quad 明天或晴或雨。计划 $c$：晴则在海滩玩 $4$ 小时，雨则看 $4$ 小时电视。计划 $c'$：晴则海滩 $2$ 小时再看电视 $2$ 小时，雨则仍看 $4$ 小时电视。这里出现两种消费品（海滩与电视），但比较的仍是整条消费计划。参与者更喜欢海滩，故晴天时更喜欢 $c$ 而不是 $c'$，即 $c\succeq c'$。

另给两个计划：$\tilde{c}$ 为晴则海滩 $4$ 小时、雨则海滩 $4$ 小时；$\tilde{c}'$ 为晴则海滩 $2$ 小时再电视 $2$ 小时、雨则海滩 $4$ 小时。独立性公理要求他选择 $\tilde{c}$ 而不是 $\tilde{c}'$。只要雨天的安排相同，雨天具体做什么不影响晴天的偏好。
\end{question}
\begin{solution}
两个计划在雨天的消费路径相同，排序已由晴天的路径决定。把雨天的共同路径从“看电视”换成“去海滩”，晴天的比较原样保留。
\end{solution}
\end{tcolorbox}

\begin{tcolorbox}[example]
\begin{question}
\textbf{例 6.2}\quad 仍在例 6.1 的晴天里比较两个消费组合，而不是两条跨状态的消费计划。组合 $c_1$：海滩 $4$ 小时并喝 $1$ 升水；组合 $c_2$：电视 $4$ 小时并喝 $1$ 升水。若他选择 $c_1$，把那 $1$ 升水去掉之后，是否仍选择海滩 $4$ 小时？
\end{question}
\begin{solution}
不一定。酷暑里他可能因此改选待在室内。海滩和水是互补品，总效用取决于两种消费品的搭配，而不只是单件消费品。独立性公理讲的是\textbf{互斥状态}之间的独立，不是同一状态里不同商品之间的独立。
\end{solution}
\end{tcolorbox}

独立性是期望效用形式的必要条件，也是充分条件。

\[
\text{独立性} \Leftrightarrow \text{期望效用}
\]

\begin{tcolorbox}[knowledge]
\textbf{定理 6.1（Debreu）}\quad 在独立性公理下，定理 2.1 给出的效用函数具有 (6.1) 的期望效用形式。
\end{tcolorbox}

解释：不确定状态下的效用，是每一可能状态下消费路径效用的加权平均，权重为该状态的概率。 教材在既定消费空间和前述偏好条件下陈述此表示结果；后文用到的彩票混合独立性还需以彩票为选择对象，不能把仅有连续性或互斥性本身当作期望效用表示的充分条件。
\end{tcolorbox}

\begin{tcolorbox}[knowledge,title=6.2 附加假设]
(6.1) 仍相当一般。下面几条简化是为了在最优消费、组合选择和均衡价格中得到更具体的结论。除特别声明外，后文都采用这些形式。

\tcbline

\textbf{A.\ 状态独立}

(6.1) 允许同一消费路径在不同状态下给出不同效用。这是否合理，取决于状态和消费怎样描述。用“晴天喝 $1$ 升水”和“多云喝 $1$ 升水”来写，效用除了水量还会依赖天气；若改用“阳光的消费量”刻画晴和阴，扣掉阳光之后，两个状态可以变成一样的，同时喝 $1$ 升水加享受 $1$ 小时阳光的效用就不必再依赖“晴或阴”这个标签。对市场总体而言状态仍有区别（可消费的阳光总量不同，会影响价格）；若只关心个人消费和效用，就可以认为没有区别。于是状态独立的期望效用是自然的：
\[
    u_{\omega}(c_0,c_{1\omega})=u(c_0,c_{1\omega}),
    \qquad
    U(c)=\sum_{\omega\in\Omega}p_{\omega}\,u(c_0,c_{1\omega}).
    \tag{6.2}
\]

\textbf{课堂解释（PPT 第 13 页）：}把完整消费篮子写为 $(x,s)$，$x$ 是冰淇淋数量，$s$ 是阳光数量。晴天可消费 $(1,1)$，阴天为 $(1,0)$；用同一个 $v(x,s)$，可以有 $v(1,1)=10$、$v(1,0)=5$。效用差异来自\textbf{消费篮子不同}，无需另给“晴天效用函数”和“阴天效用函数”。这不表示阳光可自由交易，也不意味着任何状态依赖都能在既定商品定义下消除；它说明状态独立性依赖模型怎样完整描述消费后果。

若不加限制，状态依赖的效用非常灵活，几乎任何行为都可以被解释成“他只是不喜欢在某个状态下消费”。这种写法方便，但假设与结论之间几乎没有距离，本书对此持审慎态度，后文默认 (6.2)。

\tcbline

\textbf{B.\ 时间可加性}

(6.2) 中，1 期的效用仍可以依赖 0 期消费，反之亦然，即交叉导数 $\partial^{2}u/(\partial c_0\partial c_{1\omega})$ 可以非零。实物消费里这很常见：今天吃了很多牛肉，明天对牛肉的胃口下降；今天多抽烟，明天更想抽。但本书配置的是财富和资源，而不是特定实物。为简单起见，假定
\[
    \frac{\partial^{2}u(c_0,c_1)}{\partial c_0\,\partial c_1}=0,
\]
从而
\[
    u(c_0,c_1)=u_0(c_0)+u_1(c_1).
\]
\textbf{为什么交叉偏导为零会得到可加形式？}在矩形消费定义域内，若 $u$ 二阶连续可微且 $\partial_{01}u=0$，则 $\partial_0u(c_0,c_1)$ 不随 $c_1$ 变化，可记为 $f(c_0)$。对 $c_0$ 积分，
\[
 u(c_0,c_1)=\int f(c_0)\,dc_0+g(c_1)=u_0(c_0)+u_1(c_1).
\]
反过来，可加形式的交叉偏导显然为零。它排除的是一期消费对另一期\textbf{边际效用}的影响，并不表示两期消费在预算中互不约束。

一条消费路径的效用等于各期效用之和，称为\textbf{时间可加}（time additive）或\textbf{时间可分}（time separable）。期望效用化为
\[
    U(c)=u_0(c_0)+\sum_{\omega\in\Omega}p_{\omega}\,u_1(c_{1\omega}).
    \tag{6.3}
\]

有时再设两期的效用函数只差一个正的倍数，
\[
    u_0(c_0)+u_1(c_1)=u(c_0)+\rho\,u(c_1),\qquad\rho>0.
\]
$\rho$ 称为\textbf{时间偏好系数}（time preference coefficient），也称为时间折扣系数。$\rho=1$ 时，参与者不关心消费发生在哪一期；$\rho<1$ 时，同样的消费越早越好，推迟一期使等量消费的效用降为原来的 $\rho$ 倍。此时
\[
    U(c)=u(c_0)+\rho\sum_{\omega\in\Omega}p_{\omega}\,u(c_{1\omega}).
    \tag{6.4}
\]
三个因素完全分开：状态概率 $P=\{p_\omega\}$、时间性 $\rho$、消费本身的效用 $u(\cdot)$。称呼上常把 $u(\cdot)$ 也叫效用函数，但须与总效用 $U$ 分开。

\textbf{与资产定价的联系（PPT 第 5、15 页）：}可加形式给出 $\partial_0U=u'(c_0)$、$\partial_\omega U=\rho p_\omega u'(c_{1\omega})$。内点最优时，状态价格满足
\[
 \phi_\omega=\rho p_\omega\frac{u'(c_{1\omega})}{u'(c_0)},\qquad
 m_\omega=\frac{\phi_\omega}{p_\omega}
 =\rho\frac{u'(c_{1\omega})}{u'(c_0)}.
\]
$m$ 是随机折现因子，证券定价写为 $S_n=\mathrm E[m\widetilde X_n]$。概率、时间偏好与消费边际效用的作用由此分别显现。若效用时间不可加，边际效用会依赖整条消费路径，推导仍可进行，但不再自动具有上述单期边际效用比的简洁形式；是否存在解析解取决于具体模型。


$U$ 是序数的：正单调变换不改变消费计划的排序，数值大小本身不重要。$u$ 不是这样。对 $u$ 作正单调变换，可能改变消费计划的排序。

\begin{tcolorbox}[example]
\begin{question}
\textbf{例 6.3}\quad 两个状态等可能。消费计划
\begin{center}
\begin{tabular}{@{}c@{\quad}c@{}}
\textbf{计划 A} & \textbf{计划 B} \\[0.35em]
\WJEventTree[left={$1$},stem]{1,1} &
\WJEventTree[left={$1$},stem]{3,0}
\end{tabular}
\end{center}
偏好由 $u(c_0)+\frac12\bigl[u(c_{11})+u(c_{12})\bigr]$ 表示。先取 $u(c)=c$，再取 $v(c)=\sqrt{u(c)}=\sqrt{c}$。两种效用函数下，A 与 B 的排序是否相同？
\end{question}
\begin{solution}
$u(c)=c$ 时，
\[
    U(A)=1+\tfrac12(1+1)=2,
    \qquad
    U(B)=1+\tfrac12(3+0)=\tfrac52.
\]
$B$ 优于 $A$。换成 $v(c)=\sqrt{c}$ 后，
\[
    U(A)=1+\tfrac12(\sqrt{1}+\sqrt{1})=2,
    \qquad
    U(B)=1+\tfrac12(\sqrt{3}+\sqrt{0})=1+\tfrac{\sqrt{3}}{2}.
\]
$1+\sqrt{3}/2<2$，故 $A$ 优于 $B$。排序反过来了。
\end{solution}
\end{tcolorbox}

\textbf{课堂对照（PPT 第 18 页）：}令 $C$ 确定获得 1.5，$D$ 等概率获得 1 或 2。在 $u(x)=x$ 下，两者效用都为 1.5；换成 $v(x)=\sqrt{x}$ 后，
\[
 V(C)=\sqrt{1.5}\approx1.224745,
 \qquad V(D)=\tfrac12(1+\sqrt2)\approx1.207107.
\]
$C$ 严格优于 $D$。虽然平方根保持每个确定结果的排序，却改变了对彩票的排序。对\textbf{总效用}变换成 $\sqrt{U}$ 与先对\textbf{单期效用}变换后再取期望不同：一般有
$\sqrt{\mathrm E[u(\widetilde W)]}\ne\mathrm E[\sqrt{u(\widetilde W)}]$。

因此 $u$ 具有\textbf{基数性}（cardinal）：不同消费水平上效用的相对差值会进入期望，从而影响排序。在期望效用形式下，偏好由 $u$ 唯一描述；$u$ 的正线性变换（$au+b$，$a>0$）与 $u$ 等价，一般的正单调变换则不行。这与定理 2.1 里序数的 $U$ 不同。 对 (6.4) 使用同一个变换 $v=au+b$，可直接验证
\[
 V(c)=aU(c)+b(1+\rho),\qquad a>0,
\]
所以所有计划的排序不变。“基数性”不表示效用具有可直接观察的绝对单位，也不保证不同人的效用数值可以直接比较。
\end{tcolorbox}

\begin{tcolorbox}[knowledge,title=6.3 期望效用函数的拓展]
(6.4) 加上单调和凹，对偏好的刻画简洁，也因此排除了不少行为。当由它推出的结论与观察不一致时，自然会放宽假设。\textbf{行为金融学}（behavioral finance）讨论这些实证偏离及其含义，超出本书范围。这里只说明它们怎样改写效用函数。例子包括（远不止这些）：

\begin{itemize}
    \item \textbf{习惯}（habit formation）：效用依赖消费相对于过去水平的增量，是时间不可分的一种形式。例如
    \[
        u(c_0,c_1)=u(c_0)+\rho\,u(c_1-hc_0),\qquad h\ge 0.
    \]
    1 期效用取决于相对于前期消费的增量。若采用要求正消费的对数形式，须有 $c_1-hc_0>0$，且接近该边界时效用趋于 $-\infty$。只有边际效用 $u'(0)=\infty$ 本身不能推出零消费效用是 $-\infty$，例如平方根效用在零处仍有限。
    \item \textbf{攀比}（catching up with the Joneses）：效用不仅依赖自己的消费，还依赖别人的消费，
    \[
        u(c_0,c_1)=u(c_0-kC_0)+\rho\,u(c_1-kC_1),\qquad 0\le k<1,
    \]
    其中 $C_0$、$C_1$ 为两期的人均总消费。实际效用取决于个人消费与人均消费之差。给定状态下总禀赋是给定的，这也可以看成状态依赖的一种特殊形式。
    \item \textbf{状态依赖}（state dependence）：例如效用直接依赖经济总体的消费水平。由第 3 章，边际效用高的状态会被配置更多资源。若“经济情况好”时边际效用更高，参与者会把更多资源放到好状态、给差状态留下更少，与通常把财富从好状态转移到坏状态的保险直觉相反。这类偏好可以解释为过于乐观，高估了好状态发生的可能，也是经济与行为金融讨论的对象。
    \item \textbf{一阶风险厌恶}（first order risk aversion）：效用在参考点处不可微。例如
    \[
        u(c)=
        \begin{cases}
            a_{+}(c-\bar{c}), & c\ge\bar{c},\\
            a_{-}(c-\bar{c}), & c<\bar{c},
        \end{cases}
        \qquad a_{-}>a_{+}>0.
    \]
    $u$ 在 $c=\bar{c}$ 处不可微，跨越参考点时损失与收益的边际价值不同。这可以刻画一种参考点不对称，但本例两侧均线性，并不是包含损失域凸性、概率加权等要素的完整前景理论。其一阶风险厌恶含义在下一章讨论。
    \item \textbf{不确定性厌恶}（uncertainty aversion）：可引出超出标准期望效用的模型。要区分已知概率下对风险的态度与概率本身不明确时的模糊厌恶；违反独立性公理不能自动等同于模糊厌恶。下节按 PPT 补充已知概率的彩票反例。
\end{itemize}

\textbf{习惯与攀比的区别（PPT 第 19 页）：}习惯的参照来自自己的过去，攀比的参照来自他人的消费。在习惯模型 $u(c_0)+\rho u(c_1-hc_0)$ 中，
\[
 \partial_0U=u'(c_0)-\rho h u'(c_1-hc_0),\qquad
 \partial_1U=\rho u'(c_1-hc_0).
\]
今天多消费不仅提供当期效用，还提高明天的习惯基准，因此边际收益要扣除后一项。攀比模型在个人把 $C_t$ 视为给定时，边际效用为 $u'(c_t-kC_t)$，基准来自社会环境。两类偏好可影响组合与价格，但并不单凭函数形式就必然推出“追涨杀跌”或“羊群行为”。

这些形式都还落在第 2 章的基本框架里，只是比 (6.3)、(6.4) 更灵活。本书此后只用 (6.3) 或 (6.4)。基本框架一旦清楚，再把效用换成上面的某一类拓展是直接的。
\end{tcolorbox}

\begin{tcolorbox}[knowledge,title=6.4 不确定性下的决策：圣彼得堡悖论]
% PPT 第 21–24 页，包含图片中的奖金表、级数与 CE。
掷一枚公平硬币，直到第一次正面出现。若首次正面在第 $i$ 次出现，奖金为 $X_i=2^{i-1}$，概率为 $2^{-i}$：
\[
\begin{array}{c|ccccc}
 \text{路径}&H&TH&TTH&TTTH&\cdots\\\hline
 \text{概率}&1/2&1/4&1/8&1/16&\cdots\\
 \text{奖金}&1&2&4&8&\cdots
\end{array}
\]
若只按期望奖金评价，
\[
 \mathrm E[X]=\sum_{i=1}^{\infty}2^{-i}2^{i-1}
 =\tfrac12\sum_{i=1}^{\infty}1=\infty.
\]
但无限的期望奖金并不意味着人们愿意支付无限入场费。课堂用凹效用说明，关心效用与只关心期望金额会给出不同评价。

\textbf{A.\ 对数效用的计算}\quad 暂把奖金本身视为最终财富，$u(x)=\ln x$。则
\[
 \mathrm E[u(X)]=\ln2\sum_{i=1}^{\infty}2^{-i}(i-1).
\]
从几何级数开始，在 $|x|<1$ 时逐项求导：
\[
 \sum_{i=0}^{\infty}x^i=\frac1{1-x},\qquad
 \sum_{i=1}^{\infty}ix^{i-1}=\frac1{(1-x)^2}.
\]
令 $x=1/2$ 并乘以 $1/2$，得到
\[
 \sum_{i=1}^{\infty}i2^{-i}=2,\qquad
 \sum_{i=1}^{\infty}2^{-i}=1.
\]
所以 $\mathrm E[\ln X]=\ln2(2-1)=\ln2$，确定性等值为
\[
 CE_{\ln}=\exp(\mathrm E[\ln X])=2.
\]

\textbf{B.\ 平方根效用的计算}\quad 若 $u(x)=\sqrt{x}$，
\[
\begin{aligned}
 \mathrm E[\sqrt X]
 &=\sum_{i=1}^{\infty}2^{-i}2^{(i-1)/2}
 =\frac12\sum_{j=0}^{\infty}(1/\sqrt2)^j\\
 &=\frac{1/2}{1-1/\sqrt2}=1+\frac1{\sqrt2}.
\end{aligned}
\]
故
\[
 CE_{\sqrt{\cdot}}=\left(1+\frac1{\sqrt2}\right)^2
 =\frac32+\sqrt2\approx2.914214.
\]
两种凹效用都使该彩票的期望效用和 CE 有限，但 CE 数值不同。凹性并不保证\textbf{任意}无界彩票的期望效用都有限；结论须针对具体分布与效用检验。

\textbf{CE 与入场费须区分：}上述计算没有纳入初始财富和支付费用。若初始财富为 $w$，入场费为 $f$，无差异条件应为
\[
 \sum_{i=1}^{\infty}2^{-i}u(w-f+2^{i-1})=u(w),
\]
且每个结果须在效用定义域内。因此不能把 2 或 2.914214 直接当作所有财富者的最高入场费。有限财富、主办者偿付能力与奖金上限也会影响实际决策。
\end{tcolorbox}

\begin{tcolorbox}[knowledge,title=6.4（续）阿莱斯悖论与独立性]
% PPT 第 25–27 页；p3/p4 沿用 PPT 第 26 页的编号。
奖金单位为\textbf{万美元}，四种彩票为：
\[
\begin{array}{c|ccc}
 &0&100&500\\\hline
 P_1&0&1&0\\
 P_2&0.01&0.89&0.10\\
 P_3&0.90&0&0.10\\
 P_4&0.89&0.11&0
\end{array}
\]
课堂讨论的一种选择模式是 $P_1\succ P_2$，同时 $P_3\succ P_4$。第一对偏好确定得到 100；第二对则偏好 10\% 得到 500，而非 11\% 得到 100。

\textbf{为什么不可能由同一个期望效用表示？}若 $P_1\succ P_2$，就有
\[
 u(100)>0.01u(0)+0.89u(100)+0.10u(500),
\]
即
\[
 0.11u(100)>0.01u(0)+0.10u(500).\tag{6.A}
\]
而 $P_3\succ P_4$ 要求
\[
 0.90u(0)+0.10u(500)>0.89u(0)+0.11u(100),
\]
即 $0.01u(0)+0.10u(500)>0.11u(100)$，与 (6.A) 矛盾。这与 $u$ 是否凹无关；仅改变风险厌恶程度不能消除这个矛盾。

\textbf{共同后果是怎样被替换的？}令 $Q$ 为“以 $1/11$ 概率得到 0、以 $10/11$ 概率得到 500”的彩票，$\delta_x$ 表示确定得到 $x$。则
\[
\begin{aligned}
 P_1&=0.11\delta_{100}+0.89\delta_{100},&
 P_2&=0.11Q+0.89\delta_{100},\\
 P_4&=0.11\delta_{100}+0.89\delta_0,&
 P_3&=0.11Q+0.89\delta_0.
\end{aligned}
\]
彩票混合独立性要求，对任意 $0<\alpha<1$，
\[
 L\succeq M\ \Longleftrightarrow\
 \alpha L+(1-\alpha)R\succeq\alpha M+(1-\alpha)R.
\]
把双方共有的 89\%“得到 100”同时换为“得到 0”，原来的排序应保持，故 $P_1\succ P_2$ 应推出 $P_4\succ P_3$。观察到的相反选择体现\textbf{共同后果效应或确定性效应}。这是已知概率的彩票，不是因为人们不知道概率才产生的悖论。

PPT 列出几种回应：坚持独立性作为规范要求；研究接近 0、1 的概率是否有特殊作用；引入后悔；或放宽独立性、构建非期望效用模型。悖论说明这组选择与特定公理不能同时成立，并不单凭一个例子就证明所有人、所有情形都遵循另一种模型。
\end{tcolorbox}

\begin{tcolorbox}[knowledge,title=6.4（续）秩依效用与前景理论]
\textbf{A.\ 秩依效用（PPT 第 28 页）}\quad 秩依效用保留最终结果的效用，但按结果排序后的\textbf{累计概率}构造决策权重。设 $x_1\le\cdots\le x_n$、对应概率为 $p_i$，取递增函数 $g:[0,1]\to[0,1]$，满足 $g(0)=0,g(1)=1$，可写成
\[
 V=\sum_{i=1}^n\pi_i u(x_i),\qquad
 \pi_i=g\!\left(\sum_{j=i}^np_j\right)
       -g\!\left(\sum_{j=i+1}^np_j\right).
\]
$\pi_i\ge0$ 且 $\sum\pi_i=1$；$g(p)=p$ 时退回期望效用。权重依赖结果的\textbf{秩}，不是把每个 $p_i$ 分别变换后直接相加。某些加权函数可刻画高估小概率等现象，但这些性质不是所有秩依效用模型都自动具备的；概率为零的事件仍有零权重。

\tcbline
\textbf{B.\ 前景理论（PPT 第 29–31 页）}\quad 前景理论的关键是\textbf{参考点、收益与损失的不对称及概率决策权重}。价值函数作用于相对参考点 $r$ 的变化 $x-r$，而不只作用于最终财富。常见示意形式为
\[
 v(y)=\begin{cases}y^\alpha,&y\ge0,\\-\lambda(-y)^\alpha,&y<0,\end{cases}
 \qquad0<\alpha<1,\quad\lambda>1.
\]
收益域凹、损失域凸；$\lambda>1$ 使同额损失的价值变化绝对值更大。损失厌恶与风险厌恶不是同一个概念：前者比较收益和损失，后者比较同一域内确定与不确定结果。完整模型还需指定概率权重；仅给这个价值函数并不能推出所有选择。

\textbf{课堂例：同样最终财富，不同参考点}\quad 明确两项选择是\textbf{最终财富}：$A$ 确定为 50 元，$B$ 等概率为 100 元或 0 元。设 $\alpha=1/2,\lambda=2$，暂用原概率权重，只考察参考点的影响。
\[
\begin{array}{c|cc|c}
 \text{参考点}&V(A)&V(B)&\text{排序}\\\hline
 r=0&\sqrt{50}\approx7.071&\tfrac12\sqrt{100}+\tfrac12\sqrt0=5&A\succ B\\
 r=100&-2\sqrt{50}\approx-14.142&\tfrac12v(0)+\tfrac12v(-100)=-10&B\succ A
\end{array}
\]
相同最终财富在第一种框架下是收益，在第二种框架下是损失，因而可能发生偏好反转。若题目改为“在现有财富之外\textbf{新增}50 元”，则必须先加回现有财富再减参考点，不能把新增收益和最终财富混用。
\end{tcolorbox}

\begin{tcolorbox}[example]
\textbf{课堂例：丢票与丢钱（PPT 第 32–33 页）}\quad 两张演出票总价为 $K$，仍能买到同样的票且没有流动性约束。比较：已经买好的票丢失后是否重买；与准备买票前丢失了 $K$ 元现金后是否仍买。两种情形的经济损失均为 $K$，再次购票都需另付 $K$。

如果只考虑最终财富、同样的演出体验和同样的后续机会，两种情形的可选结果相同，标准期望效用模型应给出相同的购票决定。但心理账户可能使参与者给出不同回答：丢票后的重买被计入“看演出”账户，感觉这场演出的成本变成 $2K$；丢现金则记入一般现金损失，演出账户仍只记 $K$，因而更容易接受购买。

这说明\textbf{损失怎样归类和描述}会影响选择。心理账户与参考点、框架效应相关，但不等于前景理论的全部定义；若两种情形实际存在不同的找回概率、票价或预算约束，也应先控制这些差异。
\end{tcolorbox}

\begin{tcolorbox}[knowledge,title=本章小结]
\begin{enumerate}
    \item 独立性公理下，效用可以写成各消费路径效用的期望 $U(c)=\sum_{\omega}p_{\omega}u_{\omega}(c_0,c_{1\omega})$。独立性是这一形式的充分必要条件（定理 6.1）。它针对的是互斥状态，不是同一状态中不同商品。
    \item 状态独立把 $u_{\omega}$ 收成同一个 $u$；时间可加再把它拆成 $u_0(c_0)+\sum p_{\omega}u_1(c_{1\omega})$。若两期只差时间偏好系数 $\rho$，则 $U(c)=u(c_0)+\rho\sum p_{\omega}u(c_{1\omega})$。总效用 $U$ 是序数的，单期效用 $u$ 是基数的，只在正线性变换下保持排序。
    \item 习惯以自己的过去消费为参照，攀比以他人消费为参照；参考点不对称、非线性概率权重等也可拓展标准模型。后文的组合选择与均衡定价主要使用 (6.3) 或 (6.4)。
    \item 圣彼得堡悖论区分期望金额与期望效用；CE 与包含初始财富和费用的最高入场费须分别计算。
    \item 阿莱斯例中的两对选择若反向排序，就与共同后果独立性矛盾；秩依效用、前景理论和心理账户分别强调概率权重、参考点及框架等作用，不能混为同一个假设。
\end{enumerate}
\end{tcolorbox}
\fi

% !TeX root = ../../金融经济学_王江_讲义及习题.tex
% 06_期望效用函数 课后习题（由章节文件拆分）

\begin{problem}
\problemno{6.1} 证明如下形式的期望效用函数：
\[
    U(c)=u_{0}(c_{0})+\sum_{\omega\in\Omega}\pi_{\omega}u_{1}(c_{1\omega})
\]
满足独立性公理。
\end{problem}
\begin{solution}
设两个计划 $c,c'$ 在状态 $\omega'$ 的消费路径同为 $(c_0,x)$，并且 $c\succeq c'$。由效用表示，
\[
 U(c)-U(c')=
 \sum_{\omega\ne\omega'}\pi_\omega
 \bigl[u_1(c_{1\omega})-u_1(c'_{1\omega})\bigr]\ge0.
\]
把共同路径换成 $(c_0,y)$，记新计划为 $d,d'$。两边共同的 $u_0(c_0)$ 和 $\pi_{\omega'}u_1(y)$ 仍抵消，故
\[
 U(d)-U(d')=U(c)-U(c')\ge0.
\]
因此 $d\succeq d'$，排序保持不变，独立性公理成立。
\end{solution}

\begin{problem}
\problemno{6.2*} 对于由练习 6.1 给出的效用函数表示的偏好，证明连续性意味着 $u_{0}(\cdot)$ 和 $u_{1}(\cdot)$ 都是连续的。
\end{problem}
\begin{solution}
由喜多厌少，$u_0,u_1$ 严格递增，且 $\sum_\omega\pi_\omega=1$。

考虑未来各状态消费均为 $y$ 的计划，其效用为
\[
 F(x,y)=u_0(x)+u_1(y).
\]
原偏好限制在这些计划上仍连续。下面证明：若 $F(x,y)=f(x)+g(y)$ 表示连续偏好，且 $f,g$ 严格递增，则 $f$ 连续。

单调函数在内点具有左右极限，且只有至多可数个不连续点。取 $g$ 的一个内点连续点 $y$。
\begin{enumerate}[label=(\roman*)]
\item 若 $f$ 在 $x$ 有左跳跃，记 $\delta=f(x)-f(x-)>0$。由 $g$ 在 $y$ 连续且严格递增，可取 $h>0$ 使
\[
 0<g(y)-g(y-h)<\delta.
\]
取 $x_n\uparrow x$，则
\[
 F(x_n,y)\le f(x-)+g(y)<f(x)+g(y-h)=F(x,y-h).
\]
由下等值集闭性，令 $n\to\infty$ 得 $(x,y)\preceq(x,y-h)$，与 $g$ 严格递增矛盾。
\item 若 $f$ 在 $x$ 有右跳跃，记 $\delta=f(x+)-f(x)>0$。取 $h>0$ 使
\[
 0<g(y+h)-g(y)<\delta.
\]
对 $x_n\downarrow x$，有 $F(x_n,y)>F(x,y+h)$。上等值集闭性推出 $(x,y)\succeq(x,y+h)$，同样矛盾。
\end{enumerate}
所以左右跳跃均不存在，$f$ 连续；在定义域端点只需相应的单侧论证。令 $(f,g)=(u_0,u_1)$ 得 $u_0$ 连续；交换二者，得 $u_1$ 连续。

\end{solution}

\begin{problem}
\problemno{6.3**} 如果 $a\succeq b$ 意味着对于 $\forall\,\alpha\in[0,1]$ 都有 $\alpha a+(1-\alpha)b\succeq b$，那么称偏好关系 $\succ$ 是弱凸的（weakly convex）。这与第 2 章中的强凸性有所不同。假设偏好关系是连续的并且有练习 6.1 给出的期望效用函数表示（因而 $u_{0}(\cdot)$ 和 $u_{1}(\cdot)$ 都是连续函数）。证明此时强凸性包含弱凸性。
\end{problem}
\begin{solution}
按本题的强凸性条件，若 $a\succ b$，则对任意 $0<\alpha<1$，有
\[
 \alpha a+(1-\alpha)b\succ b.
\]
要证弱凸性，只须补上 $a\sim b$ 的情形；$\alpha=0,1$ 以及 $a\succ b$ 的情形都已成立。

设 $a\sim b$，$0<\alpha<1$，记 $m=\alpha a+(1-\alpha)b$。反设 $m\prec b$，则
\[
 U(m)<U(a)=U(b).
\]
由于 $U$ 连续，沿着从 $a$ 到 $m$ 的线段，由介值定理可取
\[
 d=(1-t)a+tm,\quad 0<t<1,\qquad
 U(m)<U(d)<U(b).
\]
于是 $b\succ d$。另一方面，$m$ 位于 $d$ 与 $b$ 之间。具体地，令
\[
 \beta=\frac{\alpha}{1-t+t\alpha}\in(0,1),
 \qquad m=\beta d+(1-\beta)b.
\]
由 $b\succ d$ 和强凸性，应有 $m\succ d$，与 $U(m)<U(d)$ 矛盾。因此 $m\succeq b$，弱凸性成立。

\end{solution}

\begin{problem}
\problemno{6.4**} 假设定义在随机支付上的效用函数具有练习 6.1 中的形式，其中 $\Omega$ 是任意的状态空间而 $P$ 是相应的概率测度，并且偏好满足弱凸性。
\begin{enumerate}[label=(\alph*)]
    \item 证明 $u_{0}(\cdot)$ 和 $u_{1}(\cdot)$ 都是拟凹的。
    \item 它们是否必为凹的？若否，给出反例。
\end{enumerate}
\end{problem}
\begin{solution}
\begin{enumerate}[label=(\alph*)]
    \item $u$ 拟凹是指：若 $u(x)\ge u(y)$，则对任意 $\lambda\in[0,1]$ 有 $u(\lambda x+(1-\lambda)y)\ge u(y)$。也就是较高的那个函数值不会被凸组合拉到更低的值之下。只证 $u_{1}$，$u_{0}$ 同样。任取 $x,y$ 与 $\lambda\in[0,1]$，并取定 $c_{0}$ 以及 $1$ 期除某一状态 $\omega$ 外的消费。令 $a$ 在状态 $\omega$ 的消费为 $x$、$b$ 在状态 $\omega$ 的消费为 $y$，其余消费与 $a$ 相同。若 $u_{1}(x)\ge u_{1}(y)$，则 $U(a)\ge U(b)$，弱凸性给出 $U(\lambda a+(1-\lambda)b)\ge U(b)$，即 $u_{1}(\lambda x+(1-\lambda)y)\ge u_{1}(y)$。故 $u_{1}$ 拟凹。
    \item 不必凹。取 $u_{1}\equiv 0$、$u_{0}(c)=c^{3}$（$c\ge 0$），概率任意。此时 $U$ 只依赖于 $c_{0}$ 且对 $c_{0}$ 严格递增，故 $U(a)\ge U(b)$ 意味着 $a$ 的 $0$ 期消费不低于 $b$，其凸组合亦然，弱凸性成立。但 $u_{0}(c)=c^{3}$ 在 $[0,\infty)$ 上是凸的，不是凹的。
\end{enumerate}
\end{solution}



% !TeX root = ../../金融经济学_王江_讲义及习题.tex
% 知识点：references/金融经济学/金融经济学-王江_261013.pdf，印刷页 99–107
% 课堂补充：references/金融经济学/金融经济学_第7章 风险厌恶.pptx（夏晖，2026 年 10 月），共 32 页；公式图片与正文逐页核对。
% 习题：references/金融经济学/史上最全最完整的金融经济学答案(王江版).pdf，页 28–31（印刷页 27–30）

\lecture{风险厌恶}

\ifdefined\WJHandoutKnowledge

第 2 章要求偏好凸。在第 6 章的期望效用 (6.4) 下，这个要求落在单期效用 $u$ 上：边际效用递减，并且面对公平赌博时宁可要期望值。本章把这件事说成风险厌恶，再用 Arrow--Pratt 系数度量它，并比较不同的 $u$。除特别声明外，设 $u$ 可微；不满足性即 $u'>0$。

\begin{tcolorbox}[knowledge,title=7.1 边际效用递减]
\begin{tcolorbox}[knowledge]
\textbf{定义 7.1}\quad 若对任意 $x$、$y$ 与 $\alpha\in[0,1]$，
\[
    u\bigl(\alpha x+(1-\alpha)y\bigr)\ge \alpha u(x)+(1-\alpha)u(y),
\]
则称 $u$ 为\textbf{凹}函数。当 $x\ne y$ 且 $0<\alpha<1$ 时不等式严格，则称为严格凹。
\end{tcolorbox}

凹性就是 Jensen 不等式的方向：函数在弦的上方。二阶可微时它等价于边际效用递减。

% PPT 第 3–4 页：凹性与拟凹性的区别。
\textbf{凹函数与拟凹函数不能混同。}一般效用表示 $U$ 的\textbf{拟凹性}要求
\[
 U(\alpha x+(1-\alpha)y)\ge\min\{U(x),U(y)\},
\]
它表示两种消费方案的混合不差于较差的那个，对应偏好的凸性。\textbf{凹性}则要求混合的效用不低于效用的加权平均，条件更强。凹函数必拟凹，反之不成立；例如 $U(c)=c^2$ 在 $c>0$ 上递增且拟凹，却是凸函数。

这里须分清\textbf{一般消费效用表示 $U$}与\textbf{期望效用中的伯努利效用 $u$}。一般序数效用可作任意严格递增变换，函数的曲率未必保持；期望效用的伯努利效用只在正仿射变换下保持全部彩票排序。定理 7.1 是在期望效用结构下得到的结论，不能仅凭任意一个一维效用表示的拟凹性就推出风险厌恶。


\begin{tcolorbox}[knowledge]
\textbf{定理 7.1}\quad 若凸的连续偏好能写成 (6.4) 的期望效用，则相应的 $u(\cdot)$ 是凹的。

\textbf{定理 7.2}\quad 若凹函数 $u$ 二阶可微，则 $u''\le 0$。
\end{tcolorbox}

\textbf{定理 7.1 的二阶可微情形：}这是从总效用的\textbf{拟凹性}推出伯努利效用凹性的关键，不能直接把两个概念等同。取一个正概率状态 $\omega$，记 $k=\rho p_\omega>0$。在当前消费及该状态消费都为内点 $w$、其余消费固定的计划上，考虑
\[
 c^+(t):(c_0,c_{1\omega})=(w+kt,w-t),\qquad
 c^-(t):(c_0,c_{1\omega})=(w-kt,w+t).
\]
两计划的中点是原计划 $c$。沿这一方向的边际效用恰好抵消：
\[
 ku'(w)-ku'(w)=0.
\]
若反设 $u''(w)>0$，泰勒展开给出
\[
 U(c^\pm(t))-U(c)=\tfrac12(k^2+k)u''(w)t^2+o(t^2)>0
\]
（$t$ 充分小）。两端都严格优于中点，与凸偏好要求
$U(c)\ge\min\{U(c^+(t)),U(c^-(t))\}$ 矛盾。故每个内点都须有 $u''(w)\le0$，从而 $u$ 凹。此论证明确使用了 (6.4)、正概率和二阶可微条件，且没有把拟凹性直接替换成凹性。

\textbf{定理 7.2 的证明：}在 $z$ 左右各取距离 $\delta>0$ 的一点，凹性要求中点的效用不低于两端效用的平均值：
\[
 u(z)\ge\tfrac12[u(z-\delta)+u(z+\delta)].
\]
两边乘 2、移项，再除以 $\delta^2>0$，直接得到
\[
 \frac{u(z+\delta)-2u(z)+u(z-\delta)}{\delta^2}\le0.
\]
左边正是二阶中心差商。令 $\delta\to0$，它趋于 $u''(z)$，所以 $u''(z)\le0$。这也说明凹性的几何图像与边际效用递减是同一曲率性质的两种表述。

因此在 (6.4) 中，$u'$ 是消费的边际效用，不满足性要求 $u'>0$；偏好凸则 $u''\le 0$，多出来的一单位消费带来的效用越来越少。一块钱对穷人比对富人更重要。
\end{tcolorbox}

\begin{tcolorbox}[knowledge,title=7.2 风险厌恶的定义]
凹性还有一层含义：在期望值相同的确定支付和不确定支付之间，参与者选择确定的那个。

\begin{tcolorbox}[knowledge]
\textbf{定义 7.2}\quad 不确定支付 $\tilde{g}$ 若满足 $\mathrm{E}[\tilde{g}]=0$，称为\textbf{公平赌博}。

\textbf{定义 7.3}\quad 效用为 $u$ 的参与者是（严格）\textbf{风险厌恶}的，若对一切定义域内的公平赌博，
\[
    \mathrm{E}\bigl[u(w+\tilde{g})\bigr]\le(<)\,\mathrm{E}\bigl[u(w)\bigr]=u(w).
\]
严格风险厌恶的严格不等式只针对非退化赌博；$\tilde g\equiv0$ 时两边相等。
\end{tcolorbox}

\begin{tcolorbox}[knowledge]
\textbf{定理 7.3}\quad 参与者（严格）风险厌恶，当且仅当 $u$ 是（严格）凹函数。
\end{tcolorbox}

\textbf{证明：}

\textbf{充分性：}凹函数满足 Jensen 不等式。对随机变量 $\tilde{x}$，
\[
    \mathrm{E}\bigl[u(\tilde{x})\bigr]\le u\bigl(\mathrm{E}[\tilde{x}]\bigr);
\]
$u$ 严格凹且 $\tilde{x}$ 不恒为常数时，不等式严格。取 $\tilde{x}=w+\tilde{g}$。公平赌博给出 $\mathrm{E}[\tilde{g}]=0$，故 $\mathrm{E}[\tilde{x}]=w+\mathrm{E}[\tilde{g}]=w$，从而
\[
    \mathrm{E}\bigl[u(w+\tilde{g})\bigr]\le u(w).
\]
$w$ 是确定的，$\mathrm{E}[u(w)]=u(w)$，这就是定义 7.3。$\tilde{g}$ 不恒为零时，严格凹对应严格风险厌恶。

\textbf{必要性（按 PPT 第 8–9 页的构造思路）：}现在只知道参与者拒绝公平赌博，要据此证明 $u$ 凹。任取 $w_2>w_1$、$p\in(0,1)$，目标是得到
\[
 pu(w_1)+(1-p)u(w_2)\le u(pw_1+(1-p)w_2).
\]

\textbf{第一步：选择初始财富。}把两个最终结果的期望值设为初始财富：
\[
 w=pw_1+(1-p)w_2.
\]
这样只需设计一个赌博，把财富从 $w$ 分别变成 $w_1$、$w_2$，它就会自动公平。

\textbf{第二步：用“最终财富减初始财富”确定净支付。}记 $\Delta=w_2-w_1>0$，则
\[
\begin{aligned}
 g_1&=w_1-w=-(1-p)\Delta<0,\qquad\text{概率 }p,\\
 g_2&=w_2-w=p\Delta>0,\qquad\text{概率 }1-p.
\end{aligned}
\]
验证公平性：
\[
 \mathrm E[\tilde g]=p[-(1-p)\Delta]+(1-p)p\Delta=0.
\]
注意坏结果的概率是 $p$，并不要求输赢概率相同；净支付的大小正好补偿概率差异。

\textbf{第三步：应用风险厌恶定义。}因为 $w+g_1=w_1$、$w+g_2=w_2$，
\[
\begin{aligned}
 pu(w_1)+(1-p)u(w_2)
 &=\mathrm E[u(w+\tilde g)]\\
 &\le u(w)=u(pw_1+(1-p)w_2).
\end{aligned}
\]
这恰好是要证明的凹性不等式。若对所有非退化公平赌博严格拒绝，则上式严格，得到严格凹性。

所以，偏好能写成期望效用时，\textbf{偏好凸等价于风险厌恶}。

% PPT 第 5–7 页：三种风险态度及净支付口径。
\textbf{三种风险态度：}对非退化公平赌博，风险厌恶者偏好确定财富 $w$；风险喜好者偏好赌博；风险中性者无差异。相应地，$u$ 分别为凹、凸、仿射函数。二阶可微时，弱风险厌恶对应 $u''\le0$，弱风险喜好对应 $u''\ge0$，风险中性对应 $u''=0$。处处 $u''<0$ 是严格凹的充分条件，但严格凹不要求二阶导数在每一点都严格为负。
\end{tcolorbox}

\begin{tcolorbox}[example]
\textbf{课堂例：入局费与公平赌博（PPT 第 5 页）}\quad 掷公平硬币，正面获得 200 元，反面获得 0 元，入局费为 $c$。判断公平性须计算\textbf{扣除费用后的净支付}：
\[
 \tilde g=\begin{cases}200-c,&\text{概率 }1/2,\\-c,&\text{概率 }1/2,\end{cases}
 \qquad \mathrm E[\tilde g]=100-c.
\]
因此 $c=100$ 时是公平赌博；$c=50$ 或 $c=0$ 时净期望收益分别为 50 元和 100 元，均不是公平赌博。严格风险厌恶者拒绝前者的非退化公平赌博，却\textbf{可能接受有正期望收益的风险}；“厌恶风险”不等于“拒绝所有不确定投资”。
\end{tcolorbox}

\newpage
\begin{tcolorbox}[knowledge,title=7.3 风险厌恶的度量]
风险厌恶的程度既依赖赌博本身，也依赖 $u$ 和财富。下面的度量针对\textbf{小}风险：取值范围（最大值减最小值）很小。大风险不一定适用。

\tcbline

\textbf{A.\ 绝对风险厌恶}

\begin{tcolorbox}[knowledge]
\textbf{定义 7.4}\quad 参与公平赌博 $\tilde{g}$ 所要求的\textbf{风险溢价} $\pi$ 由
\[
    \mathrm{E}\bigl[u(w+\tilde{g})\bigr]=u(w-\pi)
    \tag{7.1}
\]
定义。它是 \textbf{为了消除风险而愿意放弃的财富}。赌博 $\tilde{g}$ 的确定性等值是 $-\pi$；随机财富 $w+\tilde{g}$ 的确定性等值是 $w-\pi$。

另一种写法是为承担风险所要求的最小补偿 $\hat{\pi}$：
\[
    \mathrm{E}\bigl[u(w+\tilde{g}+\hat{\pi})\bigr]=u(w).
\]
一般 $\hat{\pi}\neq\pi$。在光滑效用的小风险二阶近似下，两者具有相同的主导项。

\textbf{定义 7.5}\quad 取值范围很小的赌博称为小风险。
\end{tcolorbox}

% PPT 第 11–17 页：CE 的含义及泰勒展开。
\textbf{确定性等值不等于期望财富。}对一般随机财富 $\widetilde W$，若 $u$ 严格递增，则
\[
 u(CE)=\mathrm E[u(\widetilde W)],\qquad
 CE=u^{-1}(\mathrm E[u(\widetilde W)]),\qquad
 \pi=\mathrm E[\widetilde W]-CE.
\]
先对\textbf{效用}取期望，再用反函数换回财富单位；不能把 $u(\mathrm E[\widetilde W])$ 当成期望效用。公平赌博时 $\mathrm E[\widetilde W]=w$，所以 $CE=w-\pi$。对非退化风险，严格风险厌恶、风险中性、严格风险喜好分别给出 $CE<w$、$CE=w$、$CE>w$，即 $\pi>0$、$\pi=0$、$\pi<0$。

\textbf{图解（PPT 第 12 页）：}横轴是财富，纵轴是效用，图中两个结果等概率。先连接两个可能结果的效用点；弦上对应期望财富的高度就是期望效用。从该高度水平移到效用曲线，再向下投影，得到 $CE$。凹曲线位于弦上方，故 $CE$ 在期望财富左边，水平距离就是风险溢价。不同效用的比较须看 CE 或 Arrow--Pratt 系数，不能只凭图像的纵向高低判断。

\begin{center}
\begin{tikzpicture}[x=1.35cm,y=1.35cm,>=stealth,font=\small]
 \draw[->] (0,0)--(5.7,0) node[right]{财富};
 \draw[->] (0,0)--(0,2.65) node[above]{效用};
 \draw[thick,blue,domain=.25:5.3,samples=70] plot (\x,{sqrt(\x)});
 \node[blue] at (5.25,2.52) {$u$};
 \draw (1,1)--(5,2.2361);
 \draw[dashed] (1,0)--(1,1) (5,0)--(5,2.2361);
 \draw[dashed] (0,1.618)--(3,1.618) (2.618,0)--(2.618,1.618);
 \draw[dashed] (3,0)--(3,1.7321);
 \fill (1,1) circle (1.4pt) (5,2.2361) circle (1.4pt)
       (3,1.618) circle (1.4pt) (2.618,1.618) circle (1.4pt);
 \node[left] at (0,1.618) {$\mathrm E[u(\widetilde W)]$};
 \node[below] at (1,0) {$w-\delta$};
 \node[below left] at (2.618,0) {$CE$};
 \node[below right] at (3,0) {$w$};
 \node[below] at (5,0) {$w+\delta$};
 \draw[<->] (2.618,-.55)--(3,-.55) node[midway,below] {$\pi$};
\end{tikzpicture}
\end{center}

\textbf{泰勒展开的每一步：}设 $\tilde g=\varepsilon\tilde z$，$\mathrm E[\tilde z]=0$、$\tilde z$ 有界，令 $\varepsilon\to0$。假设 $u$ 在 $w$ 附近二阶连续可微且 $u'(w)>0$，则
\[
\begin{aligned}
 u(w+\varepsilon\tilde z)
 &=u(w)+u'(w)\varepsilon\tilde z
   +\tfrac12u''(w)\varepsilon^2\tilde z^2+o(\varepsilon^2),\\
 \mathrm E[u(w+\varepsilon\tilde z)]
 &=u(w)+\underbrace{u'(w)\varepsilon\mathrm E[\tilde z]}_{=0}
   +\tfrac12u''(w)\varepsilon^2\mathrm E[\tilde z^2]+o(\varepsilon^2).
\end{aligned}
\]
\textbf{左边为什么只剩方差？}公平性使一阶项为零，而且 $\mathrm E[\tilde z^2]=\operatorname{Var}(\tilde z)$。因此赌博造成的效用变化为
\[
 \mathrm E[u(w+\varepsilon\tilde z)]-u(w)
 =\tfrac12u''(w)\varepsilon^2\operatorname{Var}(\tilde z)+o(\varepsilon^2).
\]

\textbf{右边为什么只保留 $\pi$ 的一次项？}记风险溢价为 $\pi_\varepsilon$。由定义和中值定理，在 $w-\pi_\varepsilon$ 与 $w$ 之间存在 $\xi_\varepsilon$，使
\[
\begin{aligned}
 u'(\xi_\varepsilon)\pi_\varepsilon
 &=u(w)-u(w-\pi_\varepsilon)\\
 &=u(w)-\mathrm E[u(w+\varepsilon\tilde z)]
 =O(\varepsilon^2).
\end{aligned}
\]
风险缩小时 $CE\to w$，故 $\xi_\varepsilon\to w$，且 $u'(\xi_\varepsilon)\to u'(w)>0$。于是 $\pi_\varepsilon=O(\varepsilon^2)$、$\pi_\varepsilon^2=O(\varepsilon^4)$。右边的二次项比保留的二阶风险项更小，所以
\[
 u(w-\pi_\varepsilon)=u(w)-u'(w)\pi_\varepsilon+o(\varepsilon^2).
\]

\textbf{令左右两边相等，再换算为财富单位：}
\[
 \underbrace{-u'(w)\pi_\varepsilon}_{\text{放弃确定财富的效用变化}}
 =\underbrace{\tfrac12u''(w)\varepsilon^2\operatorname{Var}(\tilde z)}_{\text{承担公平风险的效用变化}}
 +o(\varepsilon^2).
\]
除以 $-u'(w)$，得到
\[
 \pi_\varepsilon=-\frac{u''(w)}{2u'(w)}\varepsilon^2\operatorname{Var}(\tilde z)
 +o(\varepsilon^2).
\]
记回 $\tilde g=\varepsilon\tilde z$，就有
\[
 \pi\approx\tfrac12 A(w)\operatorname{Var}(\tilde g),\qquad
 A(w)=-\frac{u''(w)}{u'(w)}.\tag{7.2}
\]
这里除以边际效用，是把\textbf{效用损失换成财富损失}；负号使风险厌恶时的溢价为正。这与后文折点例中“效用损失除以左侧斜率”的换算逻辑相同。

\textbf{中心矩展开（PPT 第 15 页）：}若一般随机财富的均值为 $m$，可写为 $\widetilde W=m+\varepsilon\tilde z$，其中 $\mathrm E[\tilde z]=0$、$\tilde z$ 有界。若 $u$ 在 $m$ 附近四阶连续可微，则
\[
\begin{aligned}
 \mathrm E[u(\widetilde W)]
 &=u(m)+\frac{u''(m)}2\varepsilon^2\mathrm E[\tilde z^2]\\
 &\quad+\frac{u'''(m)}6\varepsilon^3\mathrm E[\tilde z^3]
 +\frac{u^{(4)}(m)}{24}\varepsilon^4\mathrm E[\tilde z^4]+o(\varepsilon^4).
\end{aligned}
\]
$\mathrm E[\tilde z^j]$ 是这里的 $j$ 阶中心矩：一阶为零，二阶是方差，三阶反映不对称，四阶反映更高阶的分布形态。二阶近似只保留方差，不表示所有风险都能只用方差刻画。PPT 的无穷泰勒级数还需解析性、收敛性及交换期望与求和的条件，有限阶小风险展开不依赖这些更强要求。

方差度量风险的大小，溢价与它成正比。$u'>0$、$u''\le 0$ 保证 $A(w)\ge 0$。$A(w)$ 是 Arrow--Pratt \textbf{绝对风险厌恶}系数：它等于小风险下风险溢价与方差之比的两倍（$A(w)\approx2\pi/\operatorname{Var}(\tilde g)$），既依赖 $u$，也依赖财富 $w$。其倒数是\textbf{风险容忍}系数
\[
    T(w)=\frac{1}{A(w)}=-\frac{u'(w)}{u''(w)}.
    \tag{7.3}
\]

\tcbline

\textbf{B.\ 相对风险厌恶}

$A(w)$ 衡量的是给定绝对大小的赌博，比如抛硬币赢或输 $100$。财富不同时，同样 $100$ 元的相对分量不同。若盈亏与总财富成比例，写成 $w\tilde{g}$，
\[
    \mathrm{E}\bigl[u\bigl(w(1+\tilde{g})\bigr)\bigr]=u\bigl(w(1-\pi_R)\bigr).
\]
\textbf{相对风险溢价的推导（PPT 第 19–20 页）：}令 $\tilde g=\varepsilon\tilde z$，$\mathrm E[\tilde z]=0$。风险金额为 $w\varepsilon\tilde z$，确定财富减少量为 $w\pi_R$，故
\[
\begin{aligned}
 \mathrm E[u(w+w\varepsilon\tilde z)]
 &=u(w)+\tfrac12u''(w)w^2\varepsilon^2\mathrm E[\tilde z^2]+o(\varepsilon^2),\\
 u(w-w\pi_R)&=u(w)-u'(w)w\pi_R+o(\varepsilon^2).
\end{aligned}
\]
第二式利用 $\pi_R=O(\varepsilon^2)$。令两边相等，再除以 $wu'(w)>0$，得到
\[
 \pi_R=-\frac12\frac{wu''(w)}{u'(w)}\varepsilon^2\mathrm E[\tilde z^2]
 +o(\varepsilon^2).
\]
因此
\[
    \pi_R\approx\tfrac12 R(w)\,\mathrm{Var}[\tilde{g}],
    \qquad
    R(w)\equiv-\frac{wu''(w)}{u'(w)}=wA(w).
    \tag{7.5}
\]
对同一个按比例的赌博，货币风险溢价与比例风险溢价满足\textbf{精确关系} $\pi_{\mathrm{money}}=w\pi_R$。其二阶近似为
\[
 \pi_{\mathrm{money}}\approx\tfrac12 A(w)w^2\operatorname{Var}(\tilde g),\qquad
 \pi_R\approx\tfrac12 R(w)\operatorname{Var}(\tilde g).
\]
$R(w)$ 是\textbf{相对风险厌恶}。相对溢价同时取决于 $R(w)$ 和风险相对财富的大小。方差仍然是这里对小风险本身的度量。

\textbf{单位与标准化：}$A(w)$ 的单位是财富单位的倒数，$R(w)$ 无量纲。若 $\bar u=\alpha+\beta u$、$\beta>0$，则 $-\bar u''/\bar u'=-u''/u'$，故度量不受效用刻度影响。相反，单看 $u''$ 的大小会受效用单位影响，不能据此比较两个人。
\end{tcolorbox}

\begin{tcolorbox}[example]
\textbf{课堂例：对数效用下的大风险（PPT 第 13–14 页）}\quad 初始财富 100\,000 元，等概率变为 150\,000 元或 50\,000 元，$u(w)=\ln w$。净赌博为 $\pm50\,000$，期望为零。
\[
\begin{aligned}
 \text{盈利的效用增加}&=\ln(150000/100000)=\ln1.5\approx0.405465,\\
 \text{亏损的效用减少}&=\ln(100000/50000)=\ln2\approx0.693147.
\end{aligned}
\]
同额亏损的效用损失更大，故
\[
 \mathrm E[u(\widetilde W)]-u(100000)
 =\tfrac12\ln1.5-\tfrac12\ln2
 =\tfrac12\ln0.75\approx-0.143841<0.
\]
确定性等值与风险溢价应使用未提前舍入的数值计算：
\[
 CE=\exp\bigl(\tfrac12\ln150000+\tfrac12\ln50000\bigr)
 =\sqrt{150000\times50000}\approx86602.54\text{ 元},
\]
\[
 \pi=100000-CE\approx13397.46\text{ 元}.
\]
本例若套用二阶近似，只得到 $\pi\approx\tfrac12(1/100000)50000^2=12500$ 元；风险达到初始财富的 50\%，近似与精确值存在明显差异。

\tcbline
\textbf{课堂例：对数效用下的小风险（PPT 第 18 页）}\quad $w=20\,000$ 元，等概率赢或输 10 元。由于 $A(w)=1/w$、$\operatorname{Var}(\tilde g)=100$，
\[
 \pi\approx\frac{100}{2\times20000}=0.0025\text{ 元}.
\]
精确计算为
\[
\begin{aligned}
 CE&=\sqrt{20010\times19990},\\
 \pi&=20000-\sqrt{20000^2-10^2}\\
 &=\frac{100}{20000+\sqrt{20000^2-100}}
 \approx0.002500000156\text{ 元}.
\end{aligned}
\]
小风险下近似十分接近精确值。风险溢价是\textbf{财富金额}，这里不是 $0.25\%$，而是约四分之一分钱。
\end{tcolorbox}

\begin{tcolorbox}[example]
\textbf{课堂例：相同财富比例的风险（PPT 第 21 页）}\quad 两人财富分别为 10 万元、1000 万元，均具有 $\gamma=2$ 的 CRRA 效用 $u(w)=-1/w$，面临等概率赢或输财富的 10\% 的赌博。令 $\tilde g=\pm0.1$，则
\[
 R(w)=2,\qquad \operatorname{Var}(\tilde g)=0.01,\qquad
 \pi_R\approx\tfrac12\times2\times0.01=0.01.
\]
本例还能精确求解：
\[
 \mathrm E[u(w(1+\tilde g))]
 =-\frac{1}{2w}\left(\frac1{1.1}+\frac1{0.9}\right)
 =-\frac{1}{0.99w}=u(0.99w).
\]
所以 $CE=0.99w$，$\pi_R=1\%$ 在本例恰好也是精确值。两人消除风险愿付的金额分别为 1000 元、10 万元，但占财富的比例相同。

\textbf{与 CARA 对照：}对 $u(w)=-e^{-aw}$ 和固定金额公平赌博 $\tilde g$，
\[
 -e^{-a(w-\pi)}=-e^{-aw}\mathrm E[e^{-a\tilde g}]
 \quad\Longrightarrow\quad
 \pi=\frac1a\ln\mathrm E[e^{-a\tilde g}],
\]
故同一 $a$ 的两人面对\textbf{同一个固定金额风险}时，风险溢价金额相同，不随初始财富改变。比如都面对 $\pm1$ 万元的赌博，两人的\textbf{赌注}分别占财富的 10\% 和 0.1\%；这两个百分比不是尚未计算的风险溢价比例。
\end{tcolorbox}

\begin{tcolorbox}[knowledge,title=7.4 风险厌恶的几个例子]
下面几个 $u$ 在后面的市场模型里反复出现。第 2 章只记了它们的形状；这里补上 $A$ 和 $R$。若 $A'(w)>0$（$<0$），称为绝对风险厌恶递增（递减），记为 IARA（DARA）。若 $R'(w)>0$（$<0$），称为相对风险厌恶递增（递减），记为 IRRA（DRRA）。

\begin{enumerate}
    \item \textbf{线性（风险中性）}\quad $u(w)=w$。则 $A(w)=R(w)=0$，风险容忍为无穷。期望支付相同的结果被看成一样，不管风险。

    \item \textbf{负指数（CARA）}\quad $u(w)=-e^{-aw}$，$a>0$。
    \[
        A(w)=a,\qquad R(w)=aw.
    \]
    绝对风险厌恶是常数，故称 CARA（constant absolute risk aversion）。相对风险厌恶随财富上升（IRRA）。

    \item \textbf{平方}\quad $u(w)=w-\frac12 aw^{2}$，$a>0$，定义在 $w\in[0,1/a)$ 上，以保证 $u'(w)=1-aw>0$。
    \[
        A(w)=\frac{a}{1-aw},\qquad R(w)=\frac{aw}{1-aw}.
    \]
    $A(w)$ 随 $w$ 上升，是 IARA：财富越高，风险容忍 $T(w)=(1-aw)/a$ 越低。\textcolor{red}{\textbf{教材该段写成“绝对风险厌恶随财富递减”，与这个公式相反。}}

    \item \textbf{幂（CRRA）}\quad $u(w)=\dfrac{w^{1-\gamma}}{1-\gamma}$，$\gamma>0$、$\gamma\neq 1$，$w>0$。
    \[
        A(w)=\frac{\gamma}{w},\qquad R(w)=\gamma.
    \]
    绝对风险厌恶随财富下降（DARA），相对风险厌恶是常数，故称 CRRA（constant relative risk aversion）。风险容忍 $T(w)=w/\gamma$ 与财富成正比。

    \item \textbf{对数}\quad $u(w)=\log w$。将幂效用作不影响偏好的常数平移后，
    \[
      \lim_{\gamma\to1}\frac{w^{1-\gamma}-1}{1-\gamma}=\ln w.
    \]
    因此对数效用是正仿射标准化后的幂效用极限，仍属 CRRA：
    \[
        A(w)=\frac{1}{w},\qquad R(w)=1.
    \]

    \item \textbf{双曲线绝对风险厌恶（HARA）}\quad 直接由容忍系数规定。$\gamma\ge -1$ 时
    \[
        A(w)=\frac{1}{d+w/\gamma},
        \qquad
        T(w)=d+\frac{w}{\gamma}.
    \]
    风险容忍对财富是线性的，前面几类都是它的特例：
    \begin{itemize}
        \item 风险中性：$d=\infty$；
        \item 平方：$\gamma=-1$ 且 $d=1/a$；
        \item 负指数：$\gamma\to\infty$ 且 $d=1/a$；
        \item 幂：$d=0$，$\gamma>0$ 且 $\gamma\neq 1$；
        \item 对数：$d=0$ 且 $\gamma\to 1$。
    \end{itemize}
\end{enumerate}

% PPT 第 23–24 页：参数解释及统一推导。
\textbf{HARA 参数怎样影响风险态度？}须在 $T(w)=d+w/\gamma>0$ 的定义域中讨论，且公式取 $\gamma\ne0$。$d$ 是容忍度直线的截距，$1/\gamma$ 是斜率；$d$ 可解释为独立于当前财富的基础容忍度，但不是无需模型就能识别的实际保障金额。只有零财富属于定义域时，$T(0)=d$ 才能直接解释为零财富处的容忍度。
\[
 T'(w)=\frac1\gamma,\qquad
 A'(w)=-\frac{1/\gamma}{(d+w/\gamma)^2},\qquad
 R(w)=\frac{w}{d+w/\gamma},\qquad
 R'(w)=\frac{d}{(d+w/\gamma)^2}.
\]
因此 $\gamma>0$ 时容忍度随财富上升、绝对风险厌恶下降；$\gamma=-1$ 的二次效用则相反。$d=0$ 给出常相对风险厌恶；$d>0$ 时相对风险厌恶随财富上升，尽管绝对风险厌恶可以同时下降。\textbf{DARA 与 IRRA 可以同时成立。}

从容忍度还可反推效用。因为
\[
 \frac{d}{dw}\ln u'(w)=\frac{u''(w)}{u'(w)}=-\frac1{d+w/\gamma},
\]
积分得到 $u'(w)=C(d+w/\gamma)^{-\gamma}$，$C>0$。再积分，$\gamma\ne1$ 时
\[
 u(w)=K+\frac{C\gamma}{1-\gamma}(d+w/\gamma)^{1-\gamma};
\]
$\gamma=1$ 时为 $u(w)=K+C\ln(d+w)$。这说明不同经典效用来自对风险容忍度结构的不同假设；固定 $d>0$、令 $\gamma\to\infty$ 时，$T(w)\to d$，相应效用在适当正仿射标准化后趋向 CARA 形式。
\end{tcolorbox}

\begin{tcolorbox}[knowledge,title=7.5 风险厌恶的比较]
$A(w)$ 可以用来比较两个人，或同一个人在不同财富上的态度。设 $u_1$、$u_2$ 递增且二阶可微，$A_1$、$A_2$ 为相应的绝对风险厌恶。

\begin{tcolorbox}[knowledge]
\textbf{定理 7.4（Pratt）}\quad 下列四条等价：
\begin{enumerate}
    \item $A_1(w)\ge A_2(w)$ 对一切 $w$ 成立；
    \item $u_1\bigl(u_2^{-1}(z)\bigr)$ 是 $z$ 的凹函数；
    \item 存在 $f$，满足 $f'>0$、$f''\le 0$，使 $u_1(w)=f\bigl(u_2(w)\bigr)$；
    \item 对一切财富和一切公平赌博，风险溢价 $\pi_1\ge\pi_2$。
\end{enumerate}
若 $A_1(w)>A_2(w)$ 对一切 $w$ 成立，则称参与者 1 \textbf{总体上}比参与者 2 更厌恶风险。
\end{tcolorbox}

% PPT 第 26–27 页：补齐链式求导与反向证明。
\textbf{第 1、2、3 条的等价：}由于 $u_2'>0$，反函数存在。令 $z=u_2(w)$、$f(z)=u_1(u_2^{-1}(z))$，链式法则给出
\[
 f'(z)=\frac{u_1'(w)}{u_2'(w)}>0,
\]
\[
\begin{aligned}
 f''(z)
 &=\frac{u_1''(w)u_2'(w)-u_1'(w)u_2''(w)}{[u_2'(w)]^3}\\
 &=\frac{u_1'(w)}{[u_2'(w)]^2}
 \left(\frac{u_1''(w)}{u_1'(w)}-\frac{u_2''(w)}{u_2'(w)}\right)\\
 &=-\frac{u_1'(w)}{[u_2'(w)]^2}[A_1(w)-A_2(w)].
\end{aligned}
\]
前面的比例因子为正，所以 $A_1\ge A_2$ 当且仅当 $f''\le0$。且 $f(u_2(w))=u_1(w)$，这就把“更厌恶风险”与“对另一个效用再作递增凹变换”连起来。

\textbf{第 3 条推出第 4 条：}对公平赌博 $\tilde g$，
\begin{align*}
    u_1(w-\pi_1)
    &=\mathrm{E}\bigl[u_1(w+\tilde{g})\bigr]
    =\mathrm{E}\bigl[f\bigl(u_2(w+\tilde{g})\bigr)\bigr]\\
    &\le f\bigl(\mathrm{E}[u_2(w+\tilde{g})]\bigr)
    =f\bigl(u_2(w-\pi_2)\bigr)
    =u_1(w-\pi_2).
\end{align*}

$u_1$ 递增，故 $w-\pi_1\le w-\pi_2$，即 $\pi_1\ge\pi_2$。凹变换把已经风险厌恶的 $u_2$ 变得更厌恶风险；Arrow--Pratt 系数更大，为消除同一风险愿意放弃的财富更多。

\textbf{两个人的 CE 比较（PPT 第 12 页）：}对同一个随机财富 $\widetilde W$，记 $CE_i=u_i^{-1}(\mathrm E[u_i(\widetilde W)])$，$\pi_i=\mathrm E[\widetilde W]-CE_i$。Pratt 定理给出
\[
 A_1\ge A_2\ \text{处处成立}
 \quad\Longrightarrow\quad CE_1\le CE_2,\qquad\pi_1\ge\pi_2.
\]
图中越靠左的 CE 对应越大的风险溢价。\textbf{同一个赌局}的 CE 排序只是这一个比较；要推出“总体上更厌恶风险”，必须对所有相关财富和赌博都成立，不能只看一张图或一次选择。

\textbf{数字验证：}两人都面对等概率为 1 或 9 的同一随机财富，期望财富为 5；第一人取 $u_1(w)=\ln w$，第二人取 $u_2(w)=\sqrt w$。则
\[
 A_1(w)=\frac1w>\frac1{2w}=A_2(w),\qquad w>0.
\]
按各自效用计算而不是直接比较效用数值：
\[
\begin{array}{c|c|c|c}
 &\text{期望效用}&CE&\pi=5-CE\\\hline
 u_1=\ln w&\tfrac12\ln1+\tfrac12\ln9=\ln3&3&2\\
 u_2=\sqrt w&\tfrac12\sqrt1+\tfrac12\sqrt9=2&4&1
\end{array}
\]
第一人的 CE 更靠左、为消除同一风险愿意放弃更多财富，与 Pratt 定理一致；两人的期望效用数字本身不具有跨人比较意义。

\textbf{第 4 条推出第 1 条：}固定任意财富 $w$，取等概率的 $\pm\varepsilon$ 公平赌博，则
\[
 \pi_i(w,\varepsilon)=\tfrac12A_i(w)\varepsilon^2+o(\varepsilon^2).
\]
若对所有赌博都有 $\pi_1\ge\pi_2$，两边除以 $\varepsilon^2/2>0$，再令 $\varepsilon\to0$，得到 $A_1(w)\ge A_2(w)$。$w$ 任意，故总体比较成立。这里用的是\textbf{小风险极限}，而非把二阶近似当成对任意大风险的精确公式。
\end{tcolorbox}

\begin{tcolorbox}[knowledge,title=7.6 一阶风险厌恶]
上面的 $A(w)$、$R(w)$ 都假设 $u$ 可微，并且只刻画小风险的二阶效应。若 $u$ 在参考点不可微，溢价的阶数会变。第 6 章的折点效用是
\[
    u(w)=
    \begin{cases}
        a_{+}(w-\bar{w}), & w\ge\bar{w},\\
        a_{-}(w-\bar{w}), & w<\bar{w},
    \end{cases}
    \qquad a_{-}>a_{+}>0.
\]

% PPT 第 29–30 页：门槛含义及一阶效应的范围。
$\bar w$ 可代表债务偿还门槛、最低保障目标或作为参照的初始财富；同额损失的效用变化率 $a_-$ 大于同额收益的变化率 $a_+$，反映跨越目标时的不对称。\textbf{在任一线性分支内部，$u''=0$，对不跨越折点的赌博是风险中性的。}一阶风险厌恶来自折点处边际效用的跳跃，不能将左侧斜率较大解释成左侧曲率更大。

光滑效用下，对 $\pm\delta$ 风险，$\pi=O(\delta^2)$，即与\textbf{方差}同阶；并非与“方差的平方”同阶。折点处则有 $\pi=O(\delta)$。两种溢价都随 $\delta\to0$ 趋于零，但前者 $\pi/\delta\to0$，后者该比值保持正值。若初始财富与折点有固定距离，足够小且始终处于同一线性分支的风险不会产生这一折点效应。
\end{tcolorbox}

\begin{tcolorbox}[example]
\begin{question}
【可能要考】\textbf{折点处的公平赌博}\quad 财富恰为 $\bar{w}$。公平的 Bernoulli 赌博以 $1/2$ 的概率得到 $\delta$、以 $1/2$ 的概率得到 $-\delta$。按定义 7.4，风险溢价是多少？它与可微情形的阶数有何不同？
\end{question}
\begin{solution}
% 课堂板书：用户提供的折点风险溢价照片，按板书等式顺序整理。
\textbf{按课堂板书，先写定义，再代入两种结果，最后代入分段效用。}定义 7.4 为
\[
 \mathrm E[u(w+\tilde g)]=u(w-\pi).
\]
本题 $w=\bar w$，且 $\tilde g$ 以相同概率取 $\delta$、$-\delta$（$\delta>0$），故
\[
 \frac12u(\bar w+\delta)+\frac12u(\bar w-\delta)
 =u(\bar w-\pi).
\]
赢时在折点右侧，用 $a_+$；输时在折点左侧，用 $a_-$。左边的期望效用为
$\tfrac12a_+\delta+\tfrac12a_-(-\delta)<0=u(\bar w)$，所以确定性等值低于 $\bar w$、$\pi>0$。因此右边也在折点左侧，须用 $a_-$。于是板书的关键一步是
\[
 \underbrace{\frac12a_+\delta+\frac12a_-(-\delta)}_{\text{赌博的期望效用}}
 =\underbrace{a_-(-\pi)}_{\text{确定财富的效用}}.
\]
保留输时的 $(-\delta)$ 与放弃财富的 $(-\pi)$，便能直接看清符号。整理得
\[
 \frac\delta2(a_+-a_-)=-a_-\pi
 \quad\Longrightarrow\quad
 \boxed{\pi=\frac{\delta}{2a_-}(a_--a_+).}
\]
\textbf{为什么除以 $a_-$？}左边先算出的是效用损失
$\tfrac12(a_--a_+)\delta$；确定财富每减少一元，左侧分支的效用减少 $a_-$，所以除以 $a_-$ 才换算成财富金额。这正是板书从第三行到最后一行的含义。

\textbf{再比较风险溢价的阶数。}
赌博的方差为
\[
    \mathrm{Var}[\tilde{g}]
    =\mathrm{E}[\tilde{g}^{2}]-(\mathrm{E}[\tilde{g}])^{2}
    =\tfrac12\delta^{2}+\tfrac12(-\delta)^{2}
    =\delta^{2}.
\]

对于二阶连续可微且 $u'(w)>0$ 的效用函数，小风险下
\[
    \pi\approx\tfrac12 A(w)\delta^{2},
    \qquad A(w)=-\frac{u''(w)}{u'(w)},
\]
是风险规模 $\delta$ 的二阶量。本题折点处的溢价则恰好与 $\delta$ 成正比，是一阶量，称为\textbf{一阶风险厌恶}。例如，风险规模减半时，本题的溢价也减半；光滑效用下非零的二阶近似则减为原来的四分之一。

\textcolor{red}{\textbf{教材印作 $\pi=\tfrac12(a_{-}-a_{+})\delta$，按定义 7.4 还应除以左侧斜率 $a_{-}$。}}与 $\delta$ 而非 $\delta^{2}$ 成正比这一点不受影响。
\end{solution}
\end{tcolorbox}


\begin{tcolorbox}[knowledge,title=7.7 风险厌恶度量的应用]
% PPT 第 30–31 页：参与门槛与配置公式，补上所需条件。
\textbf{A.\ 一阶风险厌恶与参与门槛}\quad 在折点财富 $\bar w$ 处，考虑等概率带来 $\varepsilon(\mu+\delta)$ 或 $\varepsilon(\mu-\delta)$ 的小额投资，$\varepsilon>0$、$0<\mu<\delta$。这项投资的期望盈利是 $\varepsilon\mu>0$，但期望效用变化为
\[
 \tfrac{\varepsilon}{2}\{a_+(\mu+\delta)+a_-(\mu-\delta)\}
 =\tfrac{\varepsilon}{2}\{(a_++a_-)\mu-(a_--a_+)\delta\}.
\]
故参与条件为
\[
 \frac{\mu}{\delta}\ge\frac{a_--a_+}{a_-+a_+}.
\]
缩小投资规模 $\varepsilon$ 不改变这个门槛：即使期望收益为正，只要补偿相对损失幅度不足，参与者仍会拒绝。光滑效用且 $u'(w)>0$ 时，同一投资的效用变化为
\[
 u'(w)\varepsilon\mu+O(\varepsilon^2),
\]
固定 $\mu>0$、充分小的 $\varepsilon$ 下为正。这解释了折点为何可能形成不参与区间；是否出现持仓跳跃，还取决于完整投资机会与约束，不能仅凭一个折点断言。

\textbf{保险解释（PPT 第 30 页）：}若随机损失为 $L\ge0$，全额保险保费为 $P$，无保险财富为 $w-L$、投保后为 $w-P$。最高愿付保费 $P_{\max}$ 满足
\[
 u(w-P_{\max})=\mathrm E[u(w-L)].
\]
令 $m=w-\mathrm E[L]$、$\tilde g=\mathrm E[L]-L$，就有 $\mathrm E[\tilde g]=0$、$w-L=m+\tilde g$。按风险溢价定义，
\[
 \mathrm E[u(w-L)]=u(m-\pi)
 \quad\Longrightarrow\quad P_{\max}=\mathrm E[L]+\pi.
\]
最高愿付保费等于\textbf{精算公平保费加消除风险的溢价}。光滑偏好对小风险的额外支付通常是二阶量；若风险跨越折点，该额外支付可为一阶量。两种模型都可能愿付高于期望损失的保费，是否投保还要比较实际报价，不能说只有一阶风险厌恶者才愿支付附加保费。

\tcbline
\textbf{B.\ 风险厌恶与投资金额、比例}\quad 令 $q=\tilde r-r_F$ 为超额收益率，$b=w(1+r_F)$ 为全投无风险证券时的期末财富，$a$ 为风险投资金额。对\textbf{小的风险敞口}在 $b$ 处展开，
\[
 \mathrm E[u(b+aq)]\approx u(b)+u'(b)a\mathrm E[q]
                  +\tfrac12u''(b)a^2\mathrm E[q^2].
\]
对 $a$ 求导，得到局部近似
\[
 a^*\approx\frac{\mathrm E[q]}{A(b)\mathrm E[q^2]}.
\]
若风险溢价较小，$\mathrm E[q^2]=\operatorname{Var}(q)+(\mathrm E[q])^2
\approx\sigma^2$，且忽略短期无风险增长、令 $b\approx w$，则
\[
 a^*\approx\frac{\mathrm E[\tilde r]-r_F}{A(w)\sigma^2},\qquad
 \alpha^*=\frac{a^*}{w}\approx
 \frac{\mathrm E[\tilde r]-r_F}{R(w)\sigma^2}.
\]
因此同样投资机会下，更高风险厌恶降低风险敞口；CARA 对应风险投资金额不随财富变化，CRRA 对应风险投资比例不随财富变化，其精确条件在第 8、9 章讨论。

\textbf{与 Merton 公式的关系：}PPT 的 $\alpha^*=(\mu-r_F)/(\gamma\sigma^2)$，在单只风险资产服从几何布朗运动、参数恒定、CRRA 偏好、无额外交易限制且无外生劳动收入等经典连续时间条件下是精确的最优比例。这里的 $\mu$ 是单位时间收益漂移，$\sigma$ 是扩散波动率；不能将它直接当成任意单期收益分布下的精确结论。给定同样机会，$\gamma=5$ 时的最优比例是 $\gamma=2$ 时的 $2/5$，但不能脱离模型把某个 $\gamma$ 数值固定归类为“高度”或“中度”风险厌恶。
\end{tcolorbox}

\begin{tcolorbox}[knowledge,title=本章小结]
\begin{enumerate}
    \item 凸偏好写成期望效用后，$u$ 是凹的；二阶可微时 $u''\le 0$，即边际效用递减。
    \item 风险厌恶：公平赌博（$\mathrm{E}[\tilde{g}]=0$）不如它的期望值。这与 $u$ 凹等价。
    \item 小风险的溢价 $\pi\approx\tfrac12 A(w)\,\mathrm{Var}[\tilde{g}]$，$A(w)=-u''(w)/u'(w)$ 为绝对风险厌恶，其倒数为风险容忍。与财富成比例的风险用 $R(w)=wA(w)$。
    \item 线性效用风险中性；负指数是 CARA；幂效用和对数是 CRRA；平方效用是 IARA，且只在 $w<1/a$ 上边际效用为正。HARA 把这些情形收成风险容忍对财富线性。
    \item Pratt 定理：绝对风险厌恶处处更高、效用是另一个效用的凹变换、以及对一切公平赌博要求更高溢价，三者等价。
    \item 确定性等值由 $u(CE)=\mathrm E[u(\widetilde W)]$ 求得，风险溢价是期望财富与 CE 的差；应先计算期望效用，再反解财富。
    \item 本章分段线性效用在折点处对称小风险的溢价与风险规模成正比，称为一阶风险厌恶；在同一线性分支内则风险中性。风险厌恶度量应用于资产配置时，须同时保留收益分布与交易约束等模型条件。
\end{enumerate}
\end{tcolorbox}
\fi

% !TeX root = ../../金融经济学_王江_讲义及习题.tex
% 07_风险厌恶 课后习题（由章节文件拆分）

\begin{problem}
\problemno{7.1**} 令 $R(w)\equiv -wu''(w)/u'(w)$ 而
\[
    a(w)=\operatorname{ArgMax}E\bigl[u(w+a\tilde{r})\bigr]
\]
为相对于效用函数 $u(\cdot)$ 的最优投资选择。假设 $w>0$、$u'(\cdot)>0$、$u''(\cdot)\le 0$ 且 $E[\tilde{r}]>0$。证明以下命题：
\begin{enumerate}[label=(\alph*)]
    \item 如果 $\lambda>1$，给定效用函数 $u(\cdot)$ 且 $R'(w)\ge 0$。定义
    \[
        u_{1}(w)\equiv u(\lambda w),\qquad u_{2}\equiv u(w).
    \]
    那么，$R_{1}(w)\ge R_{2}(w)$。
    \item 给定以上结论，我们有 $a_{1}(w)\le a_{2}(w)$，其中 $a_{1}(w)$ 和 $a_{2}(w)$ 分别是相对于 $u_{1}$ 和 $u_{2}$ 的最优投资选择。
    \item 定义 $\widehat{a}(w)\equiv a(w)/w$。那么，对于 $w_{1}\le w_{2}$ 有 $\widehat{a}(w_{1})\ge\widehat{a}(w_{2})$。（提示：令 $\lambda=w_{2}/w_{1}$。）
\end{enumerate}
最后一个结论说的是，对于具有单调递增的相对风险厌恶的投资者，他在风险资产上的投资额占他财富的比例随着财富的增加而减少。
\end{problem}
\begin{solution}
设最优投资额有限且唯一，可行投资的期末财富为正，相关期望存在。以下先讨论内点解，边界解由凹目标函数的单侧导数条件得到。
\begin{enumerate}[label=(\alph*)]
\item 由链式法则，
\[
 u_1'(z)=\lambda u'(\lambda z),\qquad
 u_1''(z)=\lambda^2u''(\lambda z).
\]
因此
\[
 R_1(z)=-\frac{z u_1''(z)}{u_1'(z)}
 =-\frac{\lambda z u''(\lambda z)}{u'(\lambda z)}
 =R(\lambda z)\ge R(z)=R_2(z),
\]
其中用到 $z>0$、$\lambda>1$ 和 $R'\ge0$。

\item 记绝对风险厌恶系数为 $A_i(z)=-u_i''(z)/u_i'(z)$。由 (a)，
\[
 A_1(z)=\frac{R_1(z)}{z}\ge\frac{R_2(z)}{z}=A_2(z).
\]
令 $h(z)=u_1'(z)/u_2'(z)>0$，则
\[
 \frac{h'(z)}{h(z)}=-A_1(z)+A_2(z)\le0,
\]
所以 $h$ 非增。对给定的初始财富 $w$，记
\[
 F_i(a)=E[u_i(w+a\tilde r)],\qquad
 F_i'(a)=E[u_i'(w+a\tilde r)\tilde r].
\]
由于 $F_i'(0)=u_i'(w)E[\tilde r]>0$，有限的内点最优额为正。令 $a_2=a_2(w)>0$，其一阶条件为 $F_2'(a_2)=0$。当 $\tilde r\ge0$ 时，$h(w+a_2\tilde r)\le h(w)$；当 $\tilde r<0$ 时，不等式反向，但乘以 $\tilde r$ 后又反向。因此逐状态有
\[
 h(w+a_2\tilde r)\tilde r\le h(w)\tilde r.
\]
乘以正数 $u_2'(w+a_2\tilde r)$ 并取期望，得到
\[
 \begin{aligned}
 F_1'(a_2)
 &=E[h(w+a_2\tilde r)u_2'(w+a_2\tilde r)\tilde r]\\
 &\le h(w)F_2'(a_2)=0.
 \end{aligned}
\]
$F_1$ 为凹函数，故其唯一最优点不在 $a_2$ 右侧，即 $a_1(w)\le a_2(w)$。

\item 若 $w_1=w_2$，结论显然。若 $0<w_1<w_2$，令 $\lambda=w_2/w_1>1$。在初始财富 $w_1$ 下，
\[
 E[u_1(w_1+a\tilde r)]
 =E[u(w_2+\lambda a\tilde r)].
\]
用 $b=\lambda a$ 换元，并在可行集随财富同比例缩放的条件下，有
\[
 a_1(w_1)=\frac{a(w_2)}{\lambda},\qquad a_2(w_1)=a(w_1).
\]
由 (b) 得
\[
 \frac{a(w_2)}{\lambda}\le a(w_1)
 \quad\Longrightarrow\quad
 \boxed{\frac{a(w_2)}{w_2}\le\frac{a(w_1)}{w_1}}.
\]
因此相对风险厌恶随财富非减时，最优风险投资比例随财富非增。
\end{enumerate}
\end{solution}

\begin{problem}
\problemno{7.2} 参与者的初始财富为 $w$ 且他的绝对风险厌恶系数为常数 $a$。现在他必须承担风险 $\tilde{g}$，$\tilde{g}$ 是一公平赌博（即，他的财富变成 $w+\tilde{g}$）。当 $\tilde{g}$ 具有如下分布时计算确定性等价：
\begin{enumerate}[label=(\alph*)]
    \item 取值为 $-b$ 和 $b$ 的二项分布；
    \item 区间 $[-c,c]$ 上的均匀分布；
    \item 均值为 0、标准差为 $\sigma$ 的正态分布。
    \item 讨论在以上情形下，确定性等价如何依赖于初始财富 $w$。从中你能得到什么结论？并解释你的结论。
\end{enumerate}
\end{problem}
\begin{solution}
假设效用函数为 $U(w)=-e^{-aw}$，
\begin{enumerate}[label=(\alph*)]
    \item 由定义有
    \[
        \begin{aligned}
            U(\pi_{1}+w)
            &=-e^{-a(w+\pi_{1})}=E\bigl[U(w+\tilde{g})\bigr]=\tfrac12 U(w+b)+\tfrac12 U(w-b)\\
            &=-\tfrac12 e^{-a(w+b)}-\tfrac12 e^{-a(w-b)}\\
            \Rightarrow\pi_{1}&=-\frac{1}{a}\log\bigl(\tfrac12 e^{-ab}+\tfrac12 e^{ab}\bigr)
        \end{aligned}
    \]

    \item 由定义有
    \[
        \begin{aligned}
            U(w+\pi_{2})
            &=-e^{-a(w+\pi_{2})}=E\bigl[U(w+\tilde{g})\bigr]=\frac{1}{2c}\int_{-c}^{c}-e^{-a(w+x)}\,dx\\
            &=-\frac{1}{2ac}e^{-aw}(e^{ac}-e^{-ac})\\
            \Rightarrow\pi_{2}&=-\frac{1}{a}\log\Bigl(\frac{e^{ac}-e^{-ac}}{2ac}\Bigr)
        \end{aligned}
    \]

    \item 由定义有
    \[
        \begin{aligned}
            U(w+\pi_{3})
            &=-e^{-a(w+\pi_{3})}=E\bigl[U(w+\tilde{g})\bigr]\\
            &=\int_{-\infty}^{\infty}-e^{-a(w+x)}\frac{1}{\sqrt{2\pi\sigma}}\exp\Bigl(-\frac{x^{2}}{2\sigma^{2}}\Bigr)dx\\
            &=-e^{-aw+\frac{a^{2}\sigma^{2}}{2}}\\
            \Rightarrow\pi_{3}&=-\frac{a\sigma^{2}}{2}
        \end{aligned}
    \]

    \item 在以上各种情形下，确定性等价都不依赖于初始财富 $w$，也就是说，参与者对待公平赌博的态度与财富无关。
\end{enumerate}
\end{solution}

\begin{problem}
\problemno{7.3} 参与者的初始财富为 $w$ 且他的绝对风险厌恶系数为 1。现在他必须承担与他的财富成比例的风险 $w\tilde{g}$，$w\tilde{g}$ 是一公平赌博。当 $\tilde{g}$ 具有如下分布时，计算确定性等价占初始财富的比例——即相对确定性等价：
\begin{enumerate}[label=(\alph*)]
    \item 取值为 $-b$ 和 $b$ 的二项分布，其中 $0<b<1$；
    \item 区间 $[-c,c]$ 上的均匀分布，其中 $0<c<1$；
    \item 讨论在以上情形下，相对确定性等价如何依赖于初始财富 $w$。从中你能得到什么结论？并解释你的结论。
\end{enumerate}
\end{problem}
\begin{solution}
假设效用函数为 $U(w)=\log(w)$，
\begin{enumerate}[label=(\alph*)]
    \item 由定义有
    \[
        \begin{aligned}
            U(w\pi_{1}+w)&=\log(w+w\pi_{1})=E\bigl[U(w+w\tilde{g})\bigr]\\
            &=\tfrac12 U(w+b)+\tfrac12 U(w-b)=\log w+\tfrac12\log(1-b^{2})\\
            \Rightarrow\pi_{1}&=\sqrt{1-b^{2}}-1
        \end{aligned}
    \]

    \item 由定义有
    \[
        \begin{aligned}
            U(w+w\pi_{2})
            &=\log(w+w\pi_{2})=E\bigl[U(w+w\tilde{g})\bigr]=\frac{1}{2c}\int_{-c}^{c}\log(w+wx)\,dx\\
            &=\log w+\frac{1}{2c}\int_{-c}^{c}\log(1+x)\,dx\\
            \Rightarrow\pi_{2}&=\exp\Bigl(\frac{1}{2c}\int_{-c}^{c}\log(1+x)\,dx\Bigr)-1
        \end{aligned}
    \]

    \item 在以上各种情形下，确定性等价都不依赖于初始财富 $w$，也就是说，参与者对待与初始财富成比例的公平赌博的态度与初始财富无关。
\end{enumerate}
\end{solution}

\begin{problem}
\problemno{7.4} 参与者的初始财富为 $w$ 且相对风险厌恶系数为 2。现在他要承担风险 $\tilde{g}$，$\tilde{g}$ 是有两个可能取值 $-b$ 和 $b$ 的公平赌博，$0<b<w$。我们能用两种方式定义风险溢价：
\begin{itemize}
    \item $E[u(w+\tilde{g})]=u(w-\pi)$；
    \item $E[u(w+\tilde{g}+\pi)]=u(w)$。
\end{itemize}
用这两种定义计算风险溢价。它们相等吗？解释所得到的结论。
\end{problem}
\begin{solution}
相对风险厌恶系数为 2，即
\[
 -\frac{x u''(x)}{u'(x)}=2.
\]

积分得 $u'(x)=Kx^{-2}$（$K>0$），故可取
\[
 u(x)=-\frac1x,\qquad x>0.
\]

公平二项赌博的两个结果各以 $1/2$ 的概率发生。

\textbf{第一种定义。}由 $E[u(w+\tilde g)]=u(w-\pi_1)$，
\[
 -\frac12\left(\frac1{w+b}+\frac1{w-b}\right)
 =-\frac{w}{w^2-b^2}=-\frac1{w-\pi_1}.
\]

所以 $w-\pi_1=(w^2-b^2)/w>0$，从而
\[
  \pi_1=\frac{b^2}{w}.
\]

\textbf{第二种定义。}由 $E[u(w+\tilde g+\pi_2)]=u(w)$，
\[
 -\frac{w+\pi_2}{(w+\pi_2)^2-b^2}=-\frac1w.
\]

令 $z=w+\pi_2$，得到 $z^2-wz-b^2=0$。取满足正财富约束的正根，得
\[
 z=\frac{w+\sqrt{w^2+4b^2}}2>w>b,
 \qquad
 \pi_2=\frac{\sqrt{w^2+4b^2}-w}{2}.
\]

比较可得
\[
 \pi_2=\frac{2b^2}{\sqrt{w^2+4b^2}+w}
 <\frac{b^2}{w}=\pi_1.
\]

两者不相等。$\pi_1$ 是在财富 $w$ 下为避免风险愿意支付的金额；$\pi_2$ 是为接受风险所需的事前补偿。由于 $A(x)=2/x$ 随财富增加而下降，加入补偿后承担同一风险的绝对风险厌恶程度降低，所以所需补偿小于避免风险的支付意愿。等价地，第二个方程给出 $\pi_2=b^2/(w+\pi_2)<b^2/w$。小风险下两者的一阶非零近似均为 $b^2/w$，但精确值不同。

\end{solution}

\begin{problem}
\problemno{7.5*} 参与者的初始财富为 $w$ 且是严格风险厌恶的。$\tilde{g}_{1}$ 和 $\tilde{g}_{2}$ 是两个独立的公平赌博，具有相同的、定义于 $\{-b,b\}$ 上的二项分布，其中 $0<b<w/2$。
\begin{enumerate}[label=(\alph*)]
    \item 假设他必须承担风险 $\tilde{g}_{1}$，即他的财富变成 $w+\tilde{g}_{1}$。在这种情况下，他的期望效用记为 $V_{1}=E[u(w+\tilde{g}_{1})]$。证明风险使得他的情况恶化，即 $V_{1}<u(w)$。
    \item 现在假设他必须承担的是分散化的风险 $\tilde{g}=\frac12(\tilde{g}_{1}+\tilde{g}_{2})$。记这种情形下的期望效用为 $V_{2}=E[u(w+\tilde{g})]$。$V_{2}$ 是否总是高于 $V_{1}$？
\end{enumerate}
\end{problem}
\begin{solution}
\begin{enumerate}[label=(\alph*)]
\item 严格风险厌恶意味着伯努利效用 $u$ 严格凹。由于 $b>0$，$w+b$ 与 $w-b$ 不同，严格 Jensen 不等式给出
\[
 V_1=\tfrac12u(w+b)+\tfrac12u(w-b)
 <u\left(\tfrac12(w+b)+\tfrac12(w-b)\right)=u(w).
\]

\item 两个赌博独立且各以 $1/2$ 的概率取 $-b,b$，故平均风险的分布为
\[
 \Pr(\tilde g=b)=\tfrac14,\quad
 \Pr(\tilde g=0)=\tfrac12,\quad
 \Pr(\tilde g=-b)=\tfrac14.
\]
于是
\[
 \begin{aligned}
 V_2&=\tfrac14u(w+b)+\tfrac12u(w)+\tfrac14u(w-b)\\
 &=\tfrac12V_1+\tfrac12u(w),
 \end{aligned}
\]
从而
\[
  V_2-V_1=\tfrac12\bigl[u(w)-V_1\bigr]>0.
\]
所以在题设条件下，总有 $V_2>V_1$。分散化改善了处境，但平均后的风险仍非退化，故仍有 $V_2<u(w)$；完整排序为 $V_1<V_2<u(w)$。
\end{enumerate}
\end{solution}



% !TeX root = ../../金融经济学_王江_讲义及习题.tex
% 知识点：references/金融经济学/金融经济学-王江_261013.pdf，印刷页 111–120（小结续至 121）
% 习题：references/金融经济学/史上最全最完整的金融经济学答案(王江版).pdf，页 32–36（印刷页 31–35）

\zhchapter{组合选择（一）}

\ifdefined\WJHandoutKnowledge

第 6、7 章给出了期望效用与风险厌恶的描述。本章在给定证券价格和支付的条件下，讨论参与者怎样选择消费与投资组合：先说明最优解为什么存在，再推导最优条件，最后分析风险溢价、风险厌恶与财富对投资的影响。具体效用函数下的显式解留到第 9 章。

\begin{tcolorbox}[knowledge,title=8.1 解的存在性及其特征]
\textbf{A.\ 消费与投资的联合选择}\quad 仍令 $X$ 为 $\Omega\times N$ 支付矩阵，$S$ 为价格列向量，$\theta$ 为证券持有\textbf{数量}。参与者的初始禀赋为 $(e_0,e_1)$，状态概率为 $\pi_\omega>0$，优化问题为
\[
\begin{aligned}
 \max_\theta\quad &u_0(c_0)+\sum_{\omega=1}^{\Omega}\pi_\omega u_1(c_{1\omega})\\
 \text{s.t.}\quad &c_0=e_0-S^{\mathsf T}\theta,\qquad c_1=e_1+X\theta,\\
 &c_0\ge0,\qquad c_1\ge0.
\end{aligned}\tag{8.1}
\]
投资占用今天的消费，却增加未来各状态下的消费。这里价格是参与者接受的外生条件，尚未讨论市场出清。

\begin{tcolorbox}[knowledge]
\textbf{定理 8.1}\quad 在有限状态、无摩擦且允许自由交易的模型中，若可行消费集合非空、效用连续并对每个时期及状态的消费严格递增，则消费与组合选择问题有最优解，当且仅当市场无套利。
\end{tcolorbox}

\textbf{证明（必要性）：}若 $\theta^*$ 最优，而 $\eta$ 是套利组合，则
\[
 d=\begin{bmatrix}-S^{\mathsf T}\eta\\X\eta\end{bmatrix}\ge0,\qquad d\ne0.
\]
选择 $\theta^*+\eta$ 后，消费由 $c^*$ 变为 $c^*+d$。没有任何分量减少，至少一项增加，故效用严格提高，与最优性矛盾。

\textbf{证明（充分性）：}无套利保证存在 $\phi\gg0$ 使 $S=X^{\mathsf T}\phi$。令 $\widehat\phi=(1,\phi^{\mathsf T})^{\mathsf T}$，任意可行消费满足
\[
 c_0+\phi^{\mathsf T}c_1=e_0+\phi^{\mathsf T}e_1=:W.
\]
由于消费非负且每个状态价格严格为正，
\[
 0\le c_0\le W,\qquad 0\le c_{1\omega}\le W/\phi_\omega.
\]
所以可行消费集合有界。它又是非负正交象限与仿射子空间
$e+\operatorname{Im}[-S^{\mathsf T};X]$ 的交集，因而闭。连续效用在这个非空有界闭集合上取得最大值，且该消费可由某个组合实现。

\textbf{说明：}有界的是可行\textbf{消费}，证券头寸未必有界；冗余证券允许不同头寸实现同一消费。严格凹效用保证最优消费唯一，组合数量还需支付映射无冗余才唯一。对数效用在零处无定义时，须在其定义域内处理边界；可利用趋近零时效用趋于 $-\infty$ 排除最优消费落在该边界。

\tcbline
\textbf{B.\ Euler 方程}\quad 把预算约束代入目标函数。若最优消费位于内点且效用可微，对第 $n$ 只证券的数量求导得
\[
 -u_0'(c_0)S_n+\sum_\omega\pi_\omega u_1'(c_{1\omega})x_{\omega,n}=0.
\]
所以
\[
 \underbrace{u_0'(c_0)S_n}_{\text{今天少消费的边际效用损失}}
 =\underbrace{\mathrm E[u_1'(\tilde c_1)\widetilde X_n]}_{\text{未来支付带来的期望边际效用}},
 \quad n=1,\ldots,N.\tag{8.2}
\]
这就是\textbf{Euler 方程}。当 $S_n\ne0$ 时，也可写为
\[
 \mathrm E\left[\frac{u_1'(\tilde c_1)}{u_0'(c_0)}
                 \frac{\widetilde X_n}{S_n}\right]=1.\tag{8.3}
\]
含义是：在最优处，将最后一元用于今天消费，或用于投资后在未来消费，两者的边际效用相等。令 $m_\omega=u_1'(c_{1\omega})/u_0'(c_0)$，则 $S_n=\mathrm E[m\widetilde X_n]$；$\pi_\omega m_\omega$ 给出由该参与者最优消费支持的一组状态价格。

\textbf{凹性为什么使一阶条件充分？}对任意头寸变动 $h$，目标函数的二阶方向导数为
\[
 u_0''(c_0)(S^{\mathsf T}h)^2+
 \sum_\omega\pi_\omega u_1''(c_{1\omega})(Xh)_\omega^2\le0.
\]
因此整个目标函数对 $\theta$ 凹，内点一阶条件给出全局最优。教材逐只证券的二阶条件为
\[
 S_n^2u_0''(c_0)+\mathrm E[u_1''(\tilde c_1)\widetilde X_n^2]\le0.\tag{8.4}
\]
多只证券时应同时看所有方向的凹性；只有各个对角项非正一般不足以判断最优。若消费非负、禁止卖空或借贷约束绑定，则应使用相应约束下的最优条件，不能直接套用内点等式。
\end{tcolorbox}

\begin{tcolorbox}[example]
\textbf{例 8.1（两状态下的对数效用）}\quad 两状态等概率，唯一证券价格为 $S>0$，支付为 $[x;1]$，其中 $x\ge1$。初始禀赋为 $e_0=w_0>0$、$e_1=0$，效用为
\[
 \log c_0+\frac{\rho}{2}(\log c_{1a}+\log c_{1b}),\qquad\rho>0.
\]
持有 $\theta$ 单位证券时，$c_0=w_0-S\theta$、$c_{1a}=\theta x$、$c_{1b}=\theta$，所以
\[
 \max_{0<\theta<w_0/S}\left\{\log(w_0-S\theta)+\rho\log\theta+
 \frac{\rho}{2}\log x\right\}.
\]
对 $\theta$ 求导，Euler 方程为
\[
 \frac{S}{w_0-S\theta}=\frac{\rho}{\theta}
 \ \Longrightarrow\ S\theta=\rho(w_0-S\theta)
 \ \Longrightarrow\ \theta^*=\frac{\rho}{1+\rho}\frac{w_0}{S}.
\]
二阶导数 $-S^2/(w_0-S\theta)^2-\rho/\theta^2<0$，且两端效用趋于 $-\infty$，故这是唯一最优解。相应消费为
\[
 c_0^*=\frac{w_0}{1+\rho},\qquad
 c_{1a}^*=\frac{\rho w_0x}{(1+\rho)S},\qquad
 c_{1b}^*=\frac{\rho w_0}{(1+\rho)S}.
\]
\textbf{经济含义：}投资支出 $S\theta^*=\rho w_0/(1+\rho)$；越重视未来（$\rho$ 越大），储蓄比例越高。财富增加使消费和投资同比增加；给定价格，提高 $x$ 只改变目标函数中的常数项，故不改变本例的最优数量。这是本例对数偏好与零未来禀赋的性质。

\textbf{核对边际效用：}
\[
 u_0'(c_0^*)=\frac{1+\rho}{w_0}
 =\frac{\rho}{2}\left(\frac{x}{Sc_{1a}^*}+\frac{1}{Sc_{1b}^*}\right).
\]
\end{tcolorbox}

\begin{tcolorbox}[knowledge,title=8.2 投资组合的选择]
\textbf{A.\ 先选组合，再选储蓄}\quad 令
\[
 w=e_0-c_0=S^{\mathsf T}\theta.\tag{8.5}
\]
这里 $w$ 是扣除当前消费后投入证券的\textbf{金融财富}。给定 $w$，最优组合能达到的期望效用称为\textbf{间接效用函数}：
\[
 v_1(w)=\max_{\theta:S^{\mathsf T}\theta=w}
       \mathrm E[u_1(\tilde e_1+\widetilde X\theta)].\tag{8.6}
\]
其中 $\widetilde X$ 是随机状态实现下的支付行向量，非负消费约束仍须满足。完整问题等价于
\[
 \max_w\{u_0(e_0-w)+v_1(w)\}.\tag{8.7}
\]
内层解决“同样的钱怎样投”，外层解决“今天消费多少、留下多少投资”。后文设 $e_1=0$，略去时间下标，重点研究
\[
 v(w)=\max_{\theta:S^{\mathsf T}\theta=w}\mathrm E[u(\widetilde X\theta)].\tag{8.8}
\]
不可交易的未来收入会改变组合选择，不能直接并入可自由投资的 $w$。

\tcbline
\textbf{B.\ 从证券数量改写成投资金额}\quad 设第 $N$ 只证券无风险，净利率为 $r_F$，$1+r_F>0$。将证券取向和计量单位适当调整后，设 $S_n>0$，定义
\[
 a_n=S_n\theta_n,\qquad
 \tilde x_n=\frac{\widetilde X_n}{S_n},\tag{8.9}
\]
\[
 \tilde r_n=\frac{\widetilde X_n-S_n}{S_n}=\tilde x_n-1.\tag{8.10}
\]
$a_n$ 是\textbf{金额}，$\theta_n$ 是\textbf{数量}，$a_n/w$ 才是投资\textbf{权重}；$\tilde x_n$ 是总收益率，$\tilde r_n$ 是净收益率。投资金额满足 $\sum_n a_n=w$，从而期末财富为
\[
\begin{aligned}
 \tilde w&=\sum_{n=1}^N a_n(1+\tilde r_n)\\
 &=w(1+r_F)+\sum_{n=1}^{N-1}a_n(\tilde r_n-r_F).
\end{aligned}\tag{8.11}
\]
$\tilde r_n-r_F$ 是\textbf{超额收益率}，其期望 $\mathrm E[\tilde r_n]-r_F$ 是本章的\textbf{风险溢价}（收益率之差），与第 7 章用确定性等价定义的货币风险溢价须区分。

记 $a=(a_1,\ldots,a_{N-1})^{\mathsf T}$，$\iota$ 为全 1 列向量，组合问题变成
\[
 \max_a\ \mathrm E\bigl[u(w(1+r_F)+a^{\mathsf T}(\tilde r-r_F\iota))\bigr].\tag{8.12}
\]
无风险投资为 $a_N=w-\sum_{n<N}a_n$；负数表示借款。无约束内点解满足
\[
 \mathrm E[u'(\tilde w)(\tilde r_n-r_F)]=0,
 \qquad n=1,\ldots,N-1.\tag{8.13}
\]
这表示在风险证券与无风险证券之间再转移一元，期望边际效用的净变化为零。
\end{tcolorbox}

\begin{tcolorbox}[knowledge,title=8.3 最优组合的性质：单只风险证券]
以下设 $u'>0$、$u''<0$，风险证券收益非退化，允许卖空与借贷，并假定有限的内点最优解存在、期望与求导可交换。令
\[
 z=\tilde r-r_F,\qquad b=w(1+r_F),\qquad
 F(a)=\mathrm E[u(b+az)].
\]
则 $F'(a)=\mathrm E[u'(b+az)z]$，$F''(a)=\mathrm E[u''(b+az)z^2]<0$。

\begin{tcolorbox}[knowledge]
\textbf{定理 8.2}\quad 最优风险投资额 $a^*$ 与风险溢价同号：
\[
 \begin{array}{c|ccc}
 \mathrm E[\tilde r]-r_F&=0&>0&<0\\\hline
 a^*&=0&>0&<0
 \end{array}
\]
\end{tcolorbox}
\textbf{证明：}在 $a=0$ 时，$F'(0)=u'(b)\mathrm E[z]$，其符号与 $\mathrm E[z]$ 相同。严格凹性使 $F'$ 严格递减，故 $F'(a^*)=0$ 的根分别在零点、零点右侧或左侧。正风险溢价使参与者愿意从零开始买入，但持仓增大时风险成本最终抵消收益。

\tcbline
\textbf{A.\ 绝对风险厌恶决定投资金额怎样随财富变化}\quad 以下设 $\mathrm E[z]>0$，所以 $a^*(w)>0$。记
$A(t)=-u''(t)/u'(t)$。
\begin{tcolorbox}[knowledge]
\textbf{定理 8.3}\quad 在相关财富区间内：CARA 使最优风险投资\textbf{金额}不随财富变化；DARA 使其随财富增加；IARA 使其随财富减少。
\end{tcolorbox}
\textbf{证明：}把一阶条件 $\mathrm E[u'(b+a^*z)z]=0$ 对 $w$ 求导，记 $\tilde w^*=b+a^*z$，得
\[
 (1+r_F)\mathrm E[u''(\tilde w^*)z]
 +\frac{da^*}{dw}\mathrm E[u''(\tilde w^*)z^2]=0,
\]
即
\[
 \frac{da^*}{dw}=-(1+r_F)
 \frac{\mathrm E[u''(\tilde w^*)z]}{\mathrm E[u''(\tilde w^*)z^2]}.
\]
分母负，故导数与分子 $\mathrm E[u''(\tilde w^*)z]$ 同号。利用一阶条件减掉常数项，
\[
 \mathrm E[u''(\tilde w^*)z]
 =-\mathrm E[(A(\tilde w^*)-A(b))u'(\tilde w^*)z].
\]
DARA 下，$z>0$ 时 $\tilde w^*>b$、$A(\tilde w^*)<A(b)$；$z<0$ 时则反向。两种情况下括号与 $z$ 的乘积均为负，前面再取负号，得到 $da^*/dw>0$。CARA 时差为零；IARA 时符号相反。

这里比较的是 $A$ 在可能期末财富\textbf{整个区间}上的变化，不能仅凭初始财富处一个 $A'(w)$ 的符号推断任意风险下的结论。

\tcbline
\textbf{B.\ 相对风险厌恶决定投资比例怎样随财富变化}\quad 设 $w>0$、期末财富为正，令
\[
 \alpha^*(w)=\frac{a^*(w)}{w},\qquad
 e(w)=\frac{w}{a^*(w)}\frac{da^*}{dw}.
\]
$e$ 是风险投资金额对财富的\textbf{弹性}，且
\[
 \frac{d\alpha^*}{dw}=\frac{a^*}{w^2}(e-1).
\]
\begin{tcolorbox}[knowledge]
\textbf{定理 8.4}\quad 记 $R(t)=tA(t)$。在相关财富区间内，CRRA 对应 $e=1$、风险投资比例不变；DRRA 对应 $e>1$、比例随财富增加；IRRA 对应 $e<1$、比例随财富减少。
\end{tcolorbox}
\textbf{证明：}记 $D=\mathrm E[u''(\tilde w^*)z^2]<0$。由前式及 $\tilde w^*=b+a^*z$，
\[
\begin{aligned}
 e-1&=-\frac{b\,\mathrm E[u''(\tilde w^*)z]+a^*D}{a^*D}\\
 &=-\frac{\mathrm E[u''(\tilde w^*)\tilde w^*z]}{a^*D}
 =\frac{\mathrm E[R(\tilde w^*)u'(\tilde w^*)z]}{a^*D}.
\end{aligned}
\]
再用一阶条件，将分子写成
$\mathrm E[(R(\tilde w^*)-R(b))u'(\tilde w^*)z]$。
DRRA 下该分子为负，分母也为负，故 $e-1>0$；CRRA 时为零，IRRA 时相反。\textbf{金额增加不等于比例增加：}例如 CRRA 下金额随财富同比增加，而比例不变。
\end{tcolorbox}

\begin{tcolorbox}[knowledge,title=8.3（续）最优组合的性质：多只风险证券]
\begin{tcolorbox}[knowledge]
\textbf{定理 8.5}\quad 所有风险证券的风险溢价都为零，当且仅当全投无风险证券的 $a=0$ 是最优组合之一。
\end{tcolorbox}
\textbf{证明：}若 $a=0$ 最优，内点一阶条件给出
$u'(b)\mathrm E[\tilde r_n-r_F]=0$，故每只证券风险溢价均为零。反之，若均为零，则任意组合都有 $\mathrm E[\tilde w]=b$。由 Jensen 不等式，
\[
 \mathrm E[u(\tilde w)]\le u(\mathrm E[\tilde w])=u(b),
\]
而 $a=0$ 恰好达到 $u(b)$，故最优。若存在能完全抵消的风险头寸，最优金额向量未必唯一；严格凹性排除的是带有非零消费风险的同效用方案。

\begin{tcolorbox}[knowledge]
\textbf{定理 8.6}\quad 最优组合的期望收益率不低于无风险收益率。若至少一只证券风险溢价非零，且允许在零头寸附近自由调整、最优解存在，则最优组合的期望收益率严格更高。
\end{tcolorbox}
\textbf{证明：}全投无风险证券是可行的，所以
\[
 u(\mathrm E[\tilde w^*])\ge\mathrm E[u(\tilde w^*)]\ge u(b).
\]
由 $u'>0$，$\mathrm E[\tilde w^*]\ge b$，即
\[
 (a^*)^{\mathsf T}(\mathrm E[\tilde r]-r_F\iota)\ge0.
\]
若某一风险溢价非零，在 $a=0$ 处沿该证券风险溢价的同号方向小幅交易，效用的一阶变化严格为正，所以最优效用严格大于 $u(b)$，上述期望财富不等式也严格成立。$w>0$ 时除以 $w$，得到期望组合净收益率大于 $r_F$。

\textbf{为什么不能逐只套用定理 8.2？}多资产下还要看证券对\textbf{整个组合风险}的影响。即使每只证券的风险溢价都为正，最优组合也可能卖空其中某只证券，以对冲其余证券的风险；只能保证最优组合整体的期望超额收益非负。
\end{tcolorbox}

\begin{tcolorbox}[knowledge,title=8.4 本章结论]
\begin{enumerate}
 \item 无套利使可行消费集合有界；连续效用在非空闭集合上达到最优。冗余会影响头寸唯一性。
 \item 内点 Euler 方程把今天投资的边际成本与未来消费的期望边际效用相等；凹性使一阶条件充分。
 \item 消费与投资可分两步：给定储蓄求最优组合，再用间接效用决定储蓄。
 \item 单风险资产下，最优头寸与风险溢价同号；绝对风险厌恶控制投资\textbf{金额}的财富效应，相对风险厌恶控制投资\textbf{比例}的财富效应。
 \item 多风险资产下，零风险溢价使全投无风险证券最优；存在非零风险溢价时，最优组合整体有正的期望超额收益，但不能据此推断每只证券的持仓符号。
\end{enumerate}
\end{tcolorbox}

\fi

% !TeX root = ../../金融经济学_王江_讲义及习题.tex
% 08_组合选择一 课后习题（由章节文件拆分）

\begin{problem}
\problemno{8.1} 假设有 $N$ 项资产，收益率为 $\tilde{r}_{i}$，$i=1,\cdots,N$，它们是独立同分布的。对于一个具有不满足性和凹性效用函数的投资者，计算他的最优投资组合。
\end{problem}
\begin{solution}
设组合权重 $\theta=(\theta_{1},\ldots,\theta_{N})$，$\sum_{i}\theta_{i}=1$，期望效用目标先写为
\[
 \max_\theta E\!\left[u\!\left(
 \underbrace{w\sum_i\theta_i\tilde r_i}_{\tilde w\text{：期末财富}}
 \right)\right].
\]
其一阶条件为：对任意 $j,k$，
\[
    E\!\left[u'\!\left(w\sum_i\theta_i\tilde r_i\right)(\tilde r_j-\tilde r_k)\right]=0.
\]
由于 $\{\tilde{r}_{i}\}$ 独立同分布，目标函数对 $\theta$ 对称；严格凹效用下最优解唯一，故必有 $\theta_{i}^{*}=1/N$。即最优组合是所有风险资产的等权重组合。
\end{solution}

\begin{problem}
\problemno{8.2*} 假设你必须将 $w$ 投资于两只证券：一只无风险债券和一只风险证券。无风险债券的收益率确定为 $r_{F}$。风险证券的收益率为 $\tilde{r}$，它的均值为 $\bar{r}$ 而方差为 $\sigma^{2}$。假设你的目标就是最大化你的二次效用函数的期望值：
\[
    E\Bigl[\tilde{w}-\tfrac12 a\tilde{w}^{2}\Bigr].
\]
\begin{enumerate}[label=(\alph*)]
    \item 求解最优组合。
    \item 讨论最优组合如何依赖于股票和债券的期望收益率 $r_{F}$ 和 $\bar{r}$，股票收益率的波动率 $\sigma$，以及偏好系数 $a$。为这些依赖关系提供经济解释。
\end{enumerate}
\end{problem}
\begin{solution}
\begin{enumerate}[label=(\alph*)]
    \item 假设投资于股票的资产比例为 $x$，那么有
    \[
        \tilde{w}=w\bigl(1+r_{F}+x(\tilde{r}-r_{F})\bigr)
    \]
    参与者的优化问题是
    \[
        \max_x E\!\left[
 \underbrace{w\bigl(1+r_F+x(\tilde r-r_F)\bigr)}_{\tilde w}
 -\frac a2\bigl(w(1+r_F+x(\tilde r-r_F))\bigr)^2\right].
    \]
    先计算期末财富的一阶、二阶矩：
    \[
    \begin{aligned}
        E[\tilde w]&=w\bigl(1+r_F+x(\bar r-r_F)\bigr),\\
        E[\tilde w^2]&=w^2\bigl[(1+r_F+x(\bar r-r_F))^2+x^2\sigma^2\bigr].
    \end{aligned}
    \]
    代入目标函数后对 $x$ 求导，一阶条件为
    \[
        w(\bar r-r_F)-aw^2\bigl[(1+r_F+x(\bar r-r_F))(\bar r-r_F)+x\sigma^2\bigr]=0.
    \]
    整理得
    \[
        (\bar r-r_F)\bigl[1-aw(1+r_F)\bigr]
        =aw\bigl((\bar r-r_F)^2+\sigma^2\bigr)x.
    \]
    解得最优权重
    \[
        x^{*}=\frac{(\bar{r}-r_{F})\bigl(1-aw(1+r_{F})\bigr)}{aw\bigl(\sigma^{2}+(\bar{r}-r_{F})^{2}\bigr)}
    \]

    \item $x^{*}$ 随着 $\bar{r}$ 的增加而增加，随着 $r_{F}$、$\sigma^{2}$、$a$ 的增加而减小。股票的期望收益率增加使得股票的需求增加，而无风险收益率增加使得债券的需求增加；$\sigma^{2}$ 增大，股票的风险变大，对股票的需求变小；最后，$a$ 增大表明参与者更加厌恶风险，因而对股票的需求减小。
\end{enumerate}
\end{solution}

\begin{problem}
\problemno{8.3*} 考虑与练习 8.2 一样的问题。不过现在假设效用函数具有常数绝对风险厌恶系数：
\[
    E\bigl[-e^{-a\tilde{w}}\bigr].
\]
另外，假设股票收益率 $\tilde{r}$ 是正态分布的。
\begin{enumerate}[label=(\alph*)]
    \item 求解最优投资组合。
    \item 讨论最优组合如何依赖于股票和债券的期望收益率 $r_{F}$ 和 $\bar{r}$，股票收益率的波动率 $\sigma$，以及偏好系数 $a$。为这些依赖性提供经济解释。[如果需要的话，你可以添加其它条件。]
\end{enumerate}
\end{problem}
\begin{solution}
\begin{enumerate}[label=(\alph*)]
    \item 假设投资于股票的资产比例为 $x$，那么有
    \[
        \tilde{w}=w\bigl(1+r_{F}+x(\tilde{r}-r_{F})\bigr)
    \]
    参与者的优化问题是
    \[
        \max_x E\!\left[-\exp\!\left(-a
 \underbrace{w(1+r_F+x(\tilde r-r_F))}_{\tilde w}\right)\right].
    \]
    因 $\tilde{r}$ 正态，$\tilde{w}$ 亦正态。对正态随机变量有
    \[
        E\bigl[-e^{-a\tilde{w}}\bigr]=-\exp\Bigl(-aE[\tilde{w}]+\tfrac12 a^{2}\operatorname{Var}(\tilde{w})\Bigr)
    \]
    其中 $E[\tilde{w}]=w\bigl(1+r_{F}+x(\bar{r}-r_{F})\bigr)$，$\operatorname{Var}(\tilde{w})=w^{2}x^{2}\sigma^{2}$。最大化上式等价于最大化确定性等价，完整写为
    \[
        \max_x\left\{w\bigl(1+r_F+x(\bar r-r_F)\bigr)
        -\frac a2 w^2x^2\sigma^2\right\}.
    \]
    对 $x$ 的一阶条件为
    \[
        w(\bar{r}-r_{F})-a\,w^{2}x\sigma^{2}=0
    \]
    因而最优权重（CARA 下风险资产的金额 $wx^{*}$ 与财富无关）为
    \[
        x^{*}=\frac{\bar{r}-r_{F}}{aw\sigma^{2}}
    \]
    \item $x^{*}$ 随着 $\bar{r}$ 的增加而增加，随着 $r_{F}$、$\sigma^{2}$、$a$ 以及财富 $w$ 的增加而减小。超额收益上升提高股票需求；$\sigma^{2}$ 或绝对风险厌恶 $a$ 增大则降低股票需求。与二次效用不同：CARA 下风险资产金额 $wx^{*}=(\bar{r}-r_{F})/(a\sigma^{2})$ 不随 $w$ 变化，故权重 $x^{*}$ 随财富下降。
\end{enumerate}
\end{solution}

\begin{problem}
\problemno{8.4**} 赋予投资者初始金融财富 $w$ 以及 $\lambda$ 份不可交易收入 $\tilde{y}$，其中 $\lambda\ge 0$ 且 $E[\tilde{y}]=0$。他可以交易一只无风险证券和一只风险证券。无风险证券的收益率为 0 而风险证券的收益率 $\tilde{r}$ 的均值为 $\bar{r}$。他具有不满足性且是风险厌恶的。他的绝对风险厌恶系数随着他的财富增加而减少。令 $a$ 为风险证券上的投资额。他的期望效用函数为：
\[
    \max E\bigl[u(w+a\tilde{r}+\lambda\tilde{y})\bigr].
\]
\begin{enumerate}[label=(\alph*)]
    \item 当 $\lambda=0$ 时，风险资产的最优头寸 $a_{0}^{*}$ 如何随着 $w$ 变化而变化？
    \item 假设 $\tilde{y}$ 和 $\tilde{r}$ 是独立的。当 $\lambda>0$ 时，记风险资产的最优头寸为 $a_{\lambda}^{*}$。$a_{\lambda}^{*}$ 是大于还是小于 $a_{0}^{*}$？
    \item 假设 $\tilde{y}=c(\tilde{r}-\bar{r})+\tilde{\varepsilon}$ 以及 $E[\tilde{\varepsilon}\mid\tilde{r}]=0$。因此，如果 $c\neq 0$，$\tilde{y}$ 和 $\tilde{r}$ 是相关的。当 $\lambda$ 极小时，他持有的股票头寸如何随着 $w$ 变化而变化？解释所得到的结论。
\end{enumerate}
\end{problem}
\begin{solution}
\begin{enumerate}[label=(\alph*)]
    \item 当 $\lambda=0$ 时，优化的一阶条件是
    \[
        E\bigl[u'\!\left(\underbrace{w+a_0^*\tilde r}_{\tilde w}\right)\tilde{r}\bigr]=0
    \]
    两边对 $w$ 求导可得
    \[
        \frac{da_{0}^{*}}{dw}=\frac{E[u''(\tilde{w})\tilde{r}]}{-E[u''(\tilde{w})\tilde{r}^{2}]}
    \]
    这里绝对风险厌恶系数沿用正文的记号，写为
    \[
        \underbrace{-\frac{u''(t)}{u'(t)}}_{A(t)\text{：绝对风险厌恶系数}}.
    \]
    又
    \[
        \begin{aligned}
            E[u''(\tilde{w})\tilde{r}]
            &=-E[u'(\tilde{w})A(\tilde{w})\tilde{r}]\\
            &\ge -E[u'(\tilde{w})A(w)\tilde{r}]\\
            &=-A(w)E[u'(\tilde{w})\tilde{r}]\\
            &=0
        \end{aligned}
    \]
    故 $\dfrac{da_{0}^{*}}{dw}\ge 0$，即最优头寸随着 $w$ 的增加而增加；

    \item 参与者优化的一阶条件是
    \[
        E\bigl[u'\!\left(\underbrace{w+a_\lambda^*\tilde r+\lambda\tilde y}_{\tilde w}\right)\tilde{r}\bigr]=0
    \]
    两边对 $\lambda$ 求导可得
    \[
        \frac{da_{\lambda}^{*}}{d\lambda}=\frac{E[u''(\tilde{w})\tilde{r}\tilde{y}]}{-E[u''(\tilde{w})\tilde{r}^{2}]}
    \]
    将下面的条件期望记为 $K$：
    \[
        \underbrace{E[-A(\tilde w)\tilde y\mid\tilde r]}_{K(w+a_\lambda^*\tilde r)}.
    \]
    可以证明 $K(\cdot)>0$，$K'(\cdot)>0$，因此
    \[
        \begin{aligned}
 E[u''(\tilde w)\tilde r\tilde y]
 &=E[-A(\tilde w)u'(\tilde w)\tilde r\tilde y]\\
 &=E[K(w+a_\lambda^*\tilde r)u'(\tilde w)\tilde r]\\
 &\ge K(w)E[u'(\tilde w)\tilde r]=0.
 \end{aligned}
    \]
    最后一个等式是因为一阶条件，因此，$\dfrac{da_{\lambda}^{*}}{d\lambda}\ge 0$，故 $a_{\lambda}^{*}\ge a_{0}^{*}$；

    \item 由 (a) 有
    \[
        \frac{da_{0}^{*}}{dw}=\frac{E[u''(\tilde{w})\tilde{r}]}{-E[u''(\tilde{w})\tilde{r}^{2}]}=\frac{-E[u'(\tilde{w})A(\tilde{w})\tilde{r}]}{-E[u''(\tilde{w})\tilde{r}^{2}]}
    \]
    而
    \[
        \begin{aligned}
 E[u'(\tilde w)A(\tilde w)\tilde r]
 &=E[A(w+a\tilde r)u'(\tilde w)\tilde r]\\
 &\quad+\lambda E[A'(w+a\tilde r)u'(\tilde w)\tilde r\tilde y]
 +o(\lambda^2).
 \end{aligned}
    \]
    由 (a) 知道 $E\bigl[A(w+a\tilde{r})u'(\tilde{w})\tilde{r}\bigr]\le 0$，另外 $\lambda$ 又极小，因此
    \[
        E[u'(\tilde{w})A(\tilde{w})\tilde{r}]\le 0,\qquad \frac{da^{*}}{dw}\ge 0
    \]
    即最优头寸随着初始财富 $w$ 的增加而增加。这是因为 $\lambda$ 极小，因而不可交易收入不会改变 (a) 中的结论。
\end{enumerate}
\end{solution}

\begin{problem}
\problemno{8.5**} 假设你是对冲基金中交易政府债的交易员。你有 $w$ 的现金额度。
\begin{itemize}
    \item 在 0 期，当美联储拍卖新债券时，你可以提交报价而买入政府债。在 1 期，你可以在二级市场上以价格 $\tilde{v}$ 再把它们卖掉，其中 $\tilde{v}=\bar{v}+\tilde{\varepsilon}$，$\bar{v}>0$ 且 $E[\tilde{\varepsilon}]=0$。假设无风险收益率为零。
    \item 你购买债券的报价是 $p$。你购买到的数量 $x(p)$ 是 $p$ 的递增函数且 $x(0)=0$。选择 $p$ 以最大化你的期望效用函数：
    \[
        V(p)=E[u(\tilde{w})],\qquad \tilde{w}=w+x(p)(\bar{v}+\tilde{\varepsilon}-p).
    \]
    假设你是不满足的且严格风险厌恶的。为简单起见，假设 $V(p)$ 是严格凹的。
\end{itemize}
\begin{enumerate}[label=(\alph*)]
    \item 写出关于最优报价 $p^{*}$ 的一阶条件。
    \item 证明 $0<p^{*}<\bar{v}$。
    \item 对于两个绝对风险厌恶系数分别为 $A_{1}(w)$ 和 $A_{2}(w)$ 的交易员，如果 $\forall\,w$，$A_{1}(w)\ge A_{2}(w)$，那么 $p_{1}^{*}\le p_{2}^{*}$，其中 $p_{i}^{*}$ 是交易员 $i$ 的最优报价，$i=1,2$。
    \item 假设 $u(\cdot)$ 是绝对风险厌恶递减的。证明初始财富 $w$ 的增加使得 $p^{*}$ 也增加。
\end{enumerate}
\end{problem}
\begin{solution}
\begin{enumerate}[label=(\alph*)]
    \item 一阶条件是
    \[
        \begin{aligned}
 V'(p^*)&=E\!\left[
 u'\!\left(\underbrace{w+x(p^*)(\bar v+\tilde\varepsilon-p^*)}_{\tilde w}\right)\right.\\
 &\hspace{3em}\left.\cdot\bigl(x'(p^*)(\bar v+\tilde\varepsilon-p^*)-x(p^*)\bigr)\right]=0.
 \end{aligned}
    \]

    \item
    \[
        \begin{aligned}
            V'(0)&=E\bigl[u'(w)\bigl(x'(0)(\bar{v}+\tilde{\varepsilon})\bigr)\bigr]=u'(w)x'(0)\bar{v}>0\\
            V(\bar{v})&=E\bigl[u(w+x(\bar{v})\tilde{\varepsilon})\bigr]<u(w)=V(0)
        \end{aligned}
    \]
    又因为 $V(p)$ 是凹的，故 $0<p^{*}<\bar{v}$；

    \item 由定理 (7.4) 可知，存在 $f'>0$，$f''<0$ 使得 $u_{1}=f(u_{2})$。参与者的优化问题的一阶条件为：
    \[
        \begin{aligned}
 E\!\left[
 u_2'\!\left(\underbrace{w+x(p_2)(\bar v+\tilde\varepsilon-p_2)}_{\tilde w_2}\right)\right.&\\
 \left.\cdot\bigl(x'(p_2)(\bar v+\tilde\varepsilon-p_2)-x(p_2)\bigr)\right]&=0.
 \end{aligned}
    \]
    考虑
    \[
        V_{1}'(p_{2})=E\bigl[f'(u_{2}(\tilde{w}_{2}))u_{2}'(\tilde{w}_{2})\bigl(x'(p_{2})(\bar{v}+\tilde{\varepsilon}-p_{2})-x(p_{2})\bigr)\bigr]
    \]
    将报价的边际财富变化标记为
    \[
        \underbrace{x'(p_2)(\bar v+\tilde\varepsilon-p_2)-x(p_2)}_{\tilde r}.
    \]
    那么 $\tilde{w}_{2}$ 是 $\tilde{r}$ 的增函数，又 $f'(u_{2}(\tilde{w}_{2}))>0$ 且是 $\tilde{w}_{2}$ 的减函数，因此，$V_{1}'(p_{2})\le 0$，另外 $V(p)$ 是凹函数，因此 $p_{1}^{*}\le p_{2}^{*}$。

    \item 假设 $w_{1}>w_{0}$，当 $w=w_{0}$ 时优化问题的一阶条件为
    \[
        \begin{aligned}
 E\!\left[u'\!\left(\underbrace{w_0+x(p_0)(\bar v+\tilde\varepsilon-p_0)}_{\tilde w_0}\right)\right.&\\
 \left.\cdot\bigl(x'(p_0)(\bar v+\tilde\varepsilon-p_0)-x(p_0)\bigr)\right]&=0.
 \end{aligned}
    \]
    以下仍将绝对风险厌恶系数写为
    \[
        \underbrace{-\frac{u''(t)}{u'(t)}}_{A(t)}.
    \]
    等式两边对 $w_{0}$ 求导可得
    \[
        \begin{aligned}
            &E\bigl[u''(\tilde{w}_{0})\bigl(x'(p_{0})(\bar{v}+\tilde{\varepsilon}-p_{0})-x(p_{0})\bigr)\bigr]\\
            &=E\bigl[-A(\tilde{w}_{0})u'(\tilde{w}_{0})\bigl(x'(p_{0})(\bar{v}+\tilde{\varepsilon}-p_{0})-x(p_{0})\bigr)\bigr]\\
            &\ge 0
        \end{aligned}
    \]
    因此，
    \[
        E\bigl[u'\bigl(w_{1}+x(p_{0})(\bar{v}+\tilde{\varepsilon}-p_{0})\bigr)\bigl(x'(p_{0})(\bar{v}+\tilde{\varepsilon}-p_{0})-x(p_{0})\bigr)\bigr]\ge 0
    \]
    而 $V(p)$ 又是凹的，因此 $p_{1}\ge p_{0}$，即初始财富增加报价也会增加。
\end{enumerate}
\end{solution}



% !TeX root = ../../金融经济学_王江_讲义及习题.tex
% 源 PDF：references/金融经济学/史上最全最完整的金融经济学答案(王江版).pdf，页 38–42（印刷页 37–41）

\zhchapter{组合选择（二）}

\ifdefined\WJHandoutKnowledge
\fi

% !TeX root = ../../金融经济学_王江_讲义及习题.tex
% 09_组合选择二 课后习题（由章节文件拆分）

\begin{problem}
\problemno{9.1} 假设有四个等概率的可能状态。状态空间是 $\Omega=\{\omega_{1},\cdots,\omega_{4}\}$。考虑风险资产 $A$ 和 $B$，它们的收益率如下：
\begin{center}
\begin{tabular}{c|cccc}
    & $\omega_{1}$ & $\omega_{2}$ & $\omega_{3}$ & $\omega_{4}$ \\
    \hline
    $\tilde{r}_{A}$ & 0.5 & 0.5 & 0.7 & 0.7 \\
    $\tilde{r}_{B}$ & 0.9 & 0.8 & 0.4 & 0.3
\end{tabular}
\end{center}
在随机占优的意义上，你能为它们排序吗？
\end{problem}
\begin{solution}
\[
    F_{A}(x)=\begin{cases}
        0 & 0\le x<0.5\\
        0.5 & 0.5\le x<0.7\\
        1 & 0.7\le x
    \end{cases}
\]
\[
    F_{B}(x)=\begin{cases}
        0 & 0\le x<0.3\\
        0.25 & 0.3\le x<0.4\\
        0.5 & 0.4\le x<0.8\\
        0.75 & 0.8\le x<0.9\\
        1 & 0.9\le x
    \end{cases}
\]
\[
    S(y)=\int_{0}^{y}\bigl(F_{A}(x)-F_{B}(x)\bigr)\,dx=\begin{cases}
        0 & 0\le y<0.3\\
        -0.25(y-0.3) & 0.3\le y<0.4\\
        -0.025-0.5(y-0.4) & 0.4\le y<0.5\\
        -0.075 & 0.5\le y<0.7\\
        -0.075+0.5(y-0.7) & 0.7\le y<0.8\\
        -0.025+0.25(y-0.8) & 0.8\le y<0.9\\
        0 & 0.9\le y
    \end{cases}
\]
可以验证 $S(y)\le 0$，$y\in[0,1]$；且 $E[x_{A}]=E[x_{B}]$，因此 $A\succeq_{\mathrm{SSD}}B$。
\end{solution}

\begin{problem}
\problemno{9.2*} 两只风险资产的收益率分别为 $\tilde{r}_{A}$ 和 $\tilde{r}_{B}$。考虑一个投资者，他的初始财富为 $w=1$、效用函数为 $u(\cdot)$，其中 $u'(\cdot)>0$ 且 $u''(\cdot)<0$。他最大化期望效用 $E[u(\tilde{w})]$。令他的最优组合为 $[a,1-a]$。
\begin{enumerate}[label=(\alph*)]
    \item 假设 $\tilde{r}_{A}=\tilde{x}$，$\tilde{r}_{B}=\tilde{x}+\tilde{e}$，其中 $\tilde{e}$ 和 $\tilde{x}$ 互相独立且 $E[\tilde{e}]=0$。
    \begin{itemize}
        \item $\tilde{r}_{A}$ 是否二阶随机占优于 $\tilde{r}_{B}$？
        \item 对于投资者的最优投资组合 $a$，你能得到什么样的结论？
    \end{itemize}
    \item 假设 $\tilde{r}_{A}=\tilde{x}$，$\tilde{r}_{B}=\tilde{y}+\tilde{e}$，$\tilde{x}$、$\tilde{y}$、$\tilde{e}$ 互相独立，$\tilde{x}\sim_{d}\tilde{y}$ 且 $E[\tilde{e}]=0$。
    \begin{itemize}
        \item $\tilde{r}_{A}$ 是否二阶随机占优于 $\tilde{r}_{B}$？
        \item 证明投资者在资产 $A$ 上的投资额要大于他资产 $B$ 上的投资额，即 $a>1/2$。
    \end{itemize}
    \item 解释上面两种情形下结论之间的差异。
\end{enumerate}
\end{problem}
\begin{solution}
\begin{enumerate}[label=(\alph*)]
    \item 根据定义可以判断出
    \[
        A\succeq_{\mathrm{SSD}}B
    \]
    \[
        \begin{aligned}
            E[u(\tilde{w})]
            &=E\bigl[u\bigl(\tilde{x}+(1-a)\tilde{e}\bigr)\bigr]\\
            &=E\bigl[E\bigl[u\bigl(\tilde{x}+(1-a)\tilde{e}\bigr)\bigm|\tilde{x}\bigr]\bigr]\\
            &\le E\bigl[u\bigl(E\bigl[\tilde{x}+(1-a)\tilde{e}\bigm|\tilde{x}\bigr]\bigr)\bigr]\\
            &=E\bigl[u(\tilde{x})\bigr]\\
            &=E\bigl[u(\tilde{r}_{A})\bigr]
        \end{aligned}
    \]
    当 $a=1$ 时等式成立，因此 $a=1$。

    \item 因为 $\tilde{x}$，$\tilde{y}$ 是独立同分布的，因而 $\tilde{r}_{B}\sim_{d}\tilde{r}_{A}+\tilde{e}$，且 $E[\tilde{e}\mid\tilde{r}_{A}]=0$，故 $A\succeq_{\mathrm{SSD}}B$。令 $V(a)=E\bigl[u\bigl(a\tilde{x}+(1-a)\tilde{y}+(1-a)\tilde{e}\bigr)\bigr]$，那么
    \[
        \begin{aligned}
            V'(a)&=E\bigl[u'\bigl(a\tilde{x}+(1-a)\tilde{y}+(1-a)\tilde{e}\bigr)(\tilde{x}-\tilde{y}-\tilde{e})\bigr]\\
            V''(a)&=E\bigl[u''\bigl(a\tilde{x}+(1-a)\tilde{y}+(1-a)\tilde{e}\bigr)(\tilde{x}-\tilde{y}-\tilde{e})^{2}\bigr]\\
            V'\bigl(\tfrac12\bigr)&=E\bigl[u'\bigl(\tfrac12\tilde{x}+\tfrac12\tilde{y}+\tfrac12\tilde{e}\bigr)(\tilde{x}-\tilde{y}-\tilde{e})\bigr]\\
            &=E\bigl[E\bigl[u'\bigl(\tfrac12\tilde{x}+\tfrac12\tilde{y}+\tfrac12\tilde{e}\bigr)(\tilde{x}-\tilde{y}-\tilde{e})\bigm|\tilde{x}+\tilde{y}\bigr]\bigr]\\
            &=E\bigl[u'\bigl(\tfrac12\tilde{x}+\tfrac12\tilde{y}+\tfrac12\tilde{e}\bigr)E[\tilde{x}-\tilde{y}\mid\tilde{x}+\tilde{y}]-E\bigl[u'\bigl(\tfrac12\tilde{x}+\tfrac12\tilde{y}+\tfrac12\tilde{e}\bigr)\tilde{e}\bigm|\tilde{x}+\tilde{y}\bigr]\bigr]\\
            &=-E\bigl[u'\bigl(\tfrac12\tilde{x}+\tfrac12\tilde{y}+\tfrac12\tilde{e}\bigr)\tilde{e}\bigr]\\
            &=-\operatorname{Cov}\bigl(u'\bigl(\tfrac12\tilde{x}+\tfrac12\tilde{y}+\tfrac12\tilde{e}\bigr),\tilde{e}\bigr)>0
        \end{aligned}
    \]
    故，$a^{*}>\frac12$，同样的方法可以验证 $a^{*}<1$。

    \item 因为在第一种情形下，$\tilde{r}_{B}$ 只是比 $\tilde{r}_{A}$ 多一个噪音项 $\tilde{e}$，引入 $\tilde{r}_{B}$ 的效果是增加了收益中的噪音，因此 $a=1$；而在第二种情形下，因为 $\tilde{x}$，$\tilde{y}$ 是独立同分布的，引入 $\tilde{r}_{B}$ 可以分散风险，但由于 $\tilde{r}_{B}$ 比 $\tilde{r}_{A}$ 多一个噪音项，因此 $1>a^{*}>\frac12$。
\end{enumerate}
\end{solution}

\begin{problem}
\problemno{9.3} 证券市场最初由 $N$ 只证券构成。证明：当把另外 $M$ 只证券添加到市场中去时，如果投资者的最优投资组合中包含新证券，那么他的福利（即他的期望效用）会增加（也就是说，更多选择总是好的）。
\end{problem}
\begin{solution}
因为在添加了新证券以后，投资者总是可以选择原来的最优组合，只有当他的效用增加时他才会选择新的证券去构成新的组合。
\end{solution}

\begin{problem}
\problemno{9.4} 假设有 $N$ 只证券，它们的供给都是无限的、收益率为 $\tilde{r}_{i}$，$i=1,\cdots,N$。经济中的参与者都是不满足的、风险中性的且具有正的禀赋。市场不允许卖空。在什么条件下 $K$ 基金分离成立（$K<N$）？
\end{problem}
\begin{solution}
$K$ 基金分离是指：存在 $K$ 只共同基金，使每个投资者的最优组合都只是这 $K$ 只基金的某个持有量，而不必单独持有全部 $N$ 只证券。

风险中性且不能卖空时，投资者只愿意持有期望收益率最高的那些证券。因此当且仅当期望收益率最大的证券数等于 $K$ 时，$K$ 基金分离成立。
\end{solution}

\begin{problem}
\problemno{9.5} 假设在市场上有 $N$ 只证券，它们的收益率是独立同分布的。证明单基金分离成立。
\end{problem}
\begin{solution}
单基金分离是指所有投资者的最优风险组合都是同一只基金，彼此只差持有数量。由练习 8.1：收益率独立同分布时，严格凹效用下最优组合必为等权重。所有投资者持有这一同一组合，故单基金分离成立。
\end{solution}

\begin{problem}
\problemno{9.6*} 证券市场由一只无风险债券和 $N$ 只风险证券组成。无风险证券的收益率为 $r_{F}$ 而风险证券的收益率向量为 $\tilde{r}=\begin{bmatrix}\tilde{r}_{1} \\ \ldots \\ \tilde{r}_{N}\end{bmatrix}$。另外，$E[\tilde{r}]=\bar{r}=\begin{bmatrix}\bar{r}_{1} \\ \ldots \\ \bar{r}_{N}\end{bmatrix}$ 且 $E[(\tilde{r}-\bar{r})(\tilde{r}-\bar{r})^{\top}]=\Sigma$（$\Sigma$ 是风险证券收益率的协方差阵）。考虑一个财富为 $w$ 的参与者的投资问题。假设他具有二次效用函数 $u(\tilde{w})=E[\tilde{w}]-\tfrac12 a E[\tilde{w}^{2}]$。
\begin{enumerate}[label=(\alph*)]
    \item 求解他的最优投资组合。
    \item 证明对于具有二次效用函数的参与者，两基金货币分离成立。
\end{enumerate}
\end{problem}
\begin{solution}
\begin{enumerate}[label=(\alph*)]
    \item 令风险资产权重为 $x$，并把期末财富写成与练习 8.2 相同的形式
    \[
        \tilde{w}=w\bigl(1+r_{F}+x^{\top}(\tilde{r}-r_{F}\iota)\bigr)
    \]
    记 $\mu=\bar{r}-r_{F}\iota$。则 $E[\tilde{w}]=w(1+r_{F}+x^{\top}\mu)$，且
    \[
        E[\tilde{w}^{2}]=w^{2}\bigl[(1+r_{F}+x^{\top}\mu)^{2}+x^{\top}\Sigma x\bigr]
    \]
    对 $x$ 求导，一阶条件为
    \[
        w\mu-aw^{2}\bigl[(1+r_{F}+x^{\top}\mu)\mu+\Sigma x\bigr]=0
    \]
    即 $\bigl[1-aw(1+r_{F})\bigr]\mu=aw(\mu\mu^{\top}+\Sigma)x$，解得最优权重
    \[
        x^{*}=\frac{1-aw(1+r_{F})}{aw}\bigl[(\bar{r}-r_{F}\iota)(\bar{r}-r_{F}\iota)^{\top}+\Sigma\bigr]^{-1}(\bar{r}-r_{F}\iota)
    \]
    其中，$\iota$ 是元素全为 1 的 $N$ 维列向量。风险组合中各资产的相对权重由 $(\mu\mu^{\top}+\Sigma)^{-1}\mu$ 决定，与财富、风险厌恶的水平只差一个共同标量。

    \item 因为投资者持有的风险组合的权重之比恒为常数，因而对于所有的投资者来说，风险组合是相同的，只是投资于风险组合的总金额不一样而已；另外，投资者还持有无风险资产，因此两基金分离成立。
\end{enumerate}
\end{solution}

\begin{problem}
\problemno{9.7*} 考虑练习 9.6 中描述的市场。另外，我们假设风险收益率是联合正态分布的且参与者的绝对风险厌恶系数恒为 $a$。
\begin{enumerate}[label=(\alph*)]
    \item 求解他的最优投资组合。
    \item 证明对于具有常数绝对风险厌恶的参与者，两基金货币分离成立。
\end{enumerate}
\end{problem}
\begin{solution}
\begin{enumerate}[label=(\alph*)]
    \item 仍令风险资产权重为 $x$，
    \[
        \tilde{w}=w\bigl(1+r_{F}+x^{\top}(\tilde{r}-r_{F}\iota)\bigr)
    \]
    $\tilde{r}$ 联合正态则 $\tilde{w}$ 正态。与练习 8.3 相同，
    \[
        E\bigl[-e^{-a\tilde{w}}\bigr]=-\exp\Bigl(-aE[\tilde{w}]+\tfrac12 a^{2}\operatorname{Var}(\tilde{w})\Bigr)
    \]
    其中 $E[\tilde{w}]=w(1+r_{F}+x^{\top}\mu)$，$\operatorname{Var}(\tilde{w})=w^{2}x^{\top}\Sigma x$，$\mu=\bar{r}-r_{F}\iota$。对 $x$ 的一阶条件为
    \[
        w\mu-aw^{2}\Sigma x=0
    \]
    因而最优权重为
    \[
        x^{*}=\frac{1}{aw}\Sigma^{-1}(\bar{r}-r_{F}\iota)
    \]
    其中，$\iota$ 是元素全为 1 的 $N$ 维列向量。相对权重由 $\Sigma^{-1}\mu$ 决定，与 $a$、$w$ 无关。

    \item 因为投资者持有的风险组合的权重之比恒为常数，因而对于所有的投资者来说，风险组合是相同的，只是投资于风险组合的总金额不一样而已；另外，投资者还持有无风险资产，因此两基金分离成立。
\end{enumerate}
\end{solution}

\begin{problem}
\problemno{9.8*} 继续练习 9.7 中的问题。
\begin{enumerate}[label=(\alph*)]
    \item 计算在最优投资组合下参与者的期望效用。
    \item 假设现在有 $M$ 支新风险证券加入到证券市场中。（现有 $N$ 只证券的收益率不能生成新证券的收益率。）所有风险证券的收益率都是联合正态分布的。证明，如果参与者的最优投资组合中含有新的证券，无论 $M$ 只新证券的收益率分布如何，参与者的福利总是增加。
    \item 当引入新证券后，计算他的福利增加的确定性等价值。
\end{enumerate}
\end{problem}
\begin{solution}
\begin{enumerate}[label=(\alph*)]
    \item 由第七题我们知道最优组合的风险资产权重为 $x^{*}=\dfrac{1}{aw}\Sigma^{-1}(\bar{r}-r_{F}\iota)$，因而期望效用为
    \[
        \begin{aligned}
            E[u(\tilde{w})]
            &=-\exp\Bigl(-aE[\tilde{w}]+\tfrac12 a^{2}\operatorname{Var}(\tilde{w})\Bigr)\\
            &=-\exp\Bigl(-aw(1+r_{F})-aw\,{x^{*}}^{\top}(\bar{r}-r_{F}\iota)+\tfrac12 a^{2}w^{2}{x^{*}}^{\top}\Sigma x^{*}\Bigr)\\
            &=-\exp\Bigl(-aw(1+r_{F})-\tfrac12(\bar{r}-r_{F}\iota)^{\top}\Sigma^{-1}(\bar{r}-r_{F}\iota)\Bigr)
        \end{aligned}
    \]

    \item 由第三题可以得到此结论；

    \item 假设原来的风险资产的收益率和协方差矩阵分布为 $\bar{r}_{\mathrm{old}}$，$\Sigma_{\mathrm{old}}$，加入新证券的风险资产的收益率和协方差矩阵分布为 $\bar{r}_{\mathrm{new}}$，$\Sigma_{\mathrm{new}}$，那么福利增加的确定性等价值为
    \[
        \tfrac12(\bar{r}_{\mathrm{new}}-r_{F}\iota)^{\top}\Sigma_{\mathrm{new}}^{-1}(\bar{r}_{\mathrm{new}}-r_{F}\iota)-\tfrac12(\bar{r}_{\mathrm{old}}-r_{F}\iota)^{\top}\Sigma_{\mathrm{old}}^{-1}(\bar{r}_{\mathrm{old}}-r_{F}\iota)
    \]
\end{enumerate}
\end{solution}

\begin{problem}
\problemno{9.9*} 在练习 9.7 中，假设只有一只新证券（即 $M=1$），它的收益率是 $\tilde{r}_{N+1}$。令 $\bar{r}_{N+1}$ 为它的预期收益率、$\sigma_{N+1}$ 为收益率波动率而 $\sigma_{\cdot,N+1}$ 为它与 $N$ 只原有证券之间的协方差向量。另外，证券 $N+1$ 上的最优组合权重记作 $z_{N+1}$。
\begin{enumerate}[label=(\alph*)]
    \item 假设 $\bar{r}_{N+1}-r_{F}>(<)\,0$。我们能得到结论 $z_{N+1}>(<)\,0$ 吗？给出解释。
    \item 假设 $\bar{r}_{N+1}<\bar{r}_{n}$ 且 $\sigma_{N+1}>\sigma_{n}$，$\forall\,n=1,\ldots,N$。也就是说，其它所有证券在收益率和方差上都优于证券 $N+1$。这是否意味着 $z_{N+1}=0$ 即最优组合不应含有证券 $N+1$？给出解释。
\end{enumerate}
\end{problem}
\begin{solution}
\begin{enumerate}[label=(\alph*)]
    \item 由第七题的结论我们知道 $z_{N+1}=(\Sigma^{-1})_{N+1}(\bar{r}-r_{F}\iota)$。其中 $(\Sigma^{-1})_{N+1}$ 表示 $\Sigma^{-1}$ 的第 $N+1$ 行。在构造组合的时候，一方面我们需要考虑资产对组合收益的影响，另一方面我们还得考虑资产对组合风险的影响，如果 $\bar{r}_{N+1}-r_{F}>(<)\,0$，那么买入（卖出）证券 $N+1$ 会增加组合的收益，但同时可能会提高组合的风险，因而总体效应是不确定的，即不能得出 $z_{N+1}>(<)\,0$ 的结论。

    \item 单个资产的权重不但与自身的收益率、标准差有关，还与其它的资产的收益率以及它与其它资产的相关性有关，如果资产 $N+1$ 与其它资产是负（正）相关的，那么 $z_{N+1}$ 有可能大于（小于）0。
\end{enumerate}
\end{solution}



% !TeX root = ../../金融经济学_王江_讲义及习题.tex
% 知识点：references/金融经济学/金融经济学-王江_261013.pdf，印刷页 143–163（PDF 页 158–178）
% 习题：references/金融经济学/史上最全最完整的金融经济学答案(王江版).pdf，页 44–53（印刷页 43–52）

\lecture{完全市场中的资源配置与资产价格}

\ifdefined\WJHandoutKnowledge

第 3 章的 Arrow--Debreu 经济直接交易状态或有消费。本章从证券市场出发：当证券支付能够生成每个状态的支付时，两种市场给出相同的预算集合和均衡配置。以下依次讨论完全市场均衡、风险分担、代表性参与者及消费定价。

\begin{tcolorbox}[knowledge,title=10.1 完全市场中的均衡]
设 1 期有 $\Omega$ 个正概率状态，交易证券的支付矩阵为 $X$，价格向量为 $S$。若 $\operatorname{rank}(X)=\Omega$，任意或有支付 $y$ 都能由组合 $\theta$ 复制：$X\theta=y$。无套利时存在唯一严格正的状态价格向量 $\phi$，满足
\[
 S=X^{\mathsf T}\phi,\qquad V(y)=\phi^{\mathsf T}y.
\]
若 $X$ 是方阵，$\phi=X^{-\mathsf T}S$；若证券多于状态，满行秩仍保证状态价格唯一。状态 $\omega$ 的阿罗证券支付为 $\mathbf e_\omega$，价格为 $\phi_\omega$。

参与者 $k$ 的两期禀赋为 $(e_{k,0},e_{k,1\omega})$，消费为 $(c_{k,0},c_{k,1\omega})$。在无交易摩擦、可自由复制或有支付时，预算集合写成
\[
 c_{k,0}+\sum_{\omega=1}^{\Omega}\phi_\omega c_{k,1\omega}
 \le w_k:=e_{k,0}+\sum_{\omega=1}^{\Omega}\phi_\omega e_{k,1\omega}.
\]
它与 Arrow--Debreu 市场的预算集合相同，个人禀赋的\textbf{市值}而非具体状态分布决定预算。均衡要求各人最优化且每期每状态总消费等于总禀赋。这给出\textbf{定理 10.1}：相同财富分布下，完全证券市场与 Arrow--Debreu 市场的均衡配置和状态价格相同。财富分布条件不可省略。

若 $U_k=u_{k,0}(c_{k,0})+\sum_\omega p_\omega u_{k,1}(c_{k,1\omega})$，$p_\omega>0$，效用严格递增、严格凹，内点最优条件为
\[
 u'_{k,0}(c_{k,0})=\lambda_k,\qquad
 p_\omega u'_{k,1}(c_{k,1\omega})=\lambda_k\phi_\omega.
 \tag{10.2}
\]
因此\textbf{定理 10.2}给出跨状态边际替代率：
\[
 \frac{p_\omega u'_{k,1}(c_{k,1\omega})}
 {p_{\omega'}u'_{k,1}(c_{k,1\omega'})}
 =\frac{\phi_\omega}{\phi_{\omega'}}.
 \tag{10.3}
\]
给定价格和个人财富，严格凹性使消费选择唯一；再由预算和市场出清确定均衡。\textbf{定理 10.3}陈述本章模型的均衡存在且唯一。一般经济的总均衡唯一性不能仅由个人效用严格凹推出，还须检验教材设定的市场、禀赋与内点条件。

\begin{tcolorbox}[example]
\textbf{例 10.1}\quad 两状态等概率；债券支付 $(1,1)$，股票支付 $(2,1)$，价格分别为 $B,S$。状态价格满足
\[
 B=\phi_a+\phi_b,\quad S=2\phi_a+\phi_b,\qquad
 \phi_a=S-B,\quad\phi_b=2B-S.
\]
无套利要求 $B<S<2B$。参与者 1 当期禀赋为 1，参与者 2 持有一单位股票；前者每期效用为 $\log c$，后者为 $2\sqrt c$。逐人求最优消费，再令三处总消费等于 $(1,2,1)$，得到
\[
 \phi_a=\frac{1+\sqrt{17}}{16},\quad \phi_b=\frac12,\qquad
 S=\frac{5+\sqrt{17}}8,\quad B=\frac{9+\sqrt{17}}{16}.
\]
可直接核对债券、股票定价式与两个状态价格的正性。
\end{tcolorbox}
\end{tcolorbox}

\begin{tcolorbox}[knowledge,title=10.2 最优配置与风险分担]
\textbf{定理 10.4}\quad 在凸偏好与内点条件下，帕累托最优配置可由一组正权重 $\mu_k$ 的规划问题表示：
\[
 \max_{\{c_k\}}\sum_k\mu_k\left[u_{k,0}(c_{k,0})+
 \sum_\omega p_\omega u_{k,1}(c_{k,1\omega})\right],\quad
 \text{s.t. }\sum_k c_{k,t\omega}=C_{t\omega}.
 \tag{10.7}
\]
将资源约束乘子记为 $\phi'_0,\phi'_\omega$，一阶条件是
\[
 \mu_k u'_{k,0}(c_{k,0})=\phi'_0,\qquad
 \mu_k p_\omega u'_{k,1}(c_{k,1\omega})=\phi'_\omega.
 \tag{10.8}
\]
严格凹性使内点一阶条件充分。令市场状态价格 $\phi_\omega=\phi'_\omega/\phi'_0$，并调整个人初始财富以支持此配置，就得到\textbf{定理 10.5（第二福利经济学定理）}：适当财富再分配后，帕累托最优配置可由完全市场竞争均衡实现。财富再分配是逆命题的一部分。

给定权重 $\mu$ 和 1 期状态总消费 $C_{1\omega}$，最优\textbf{分享规则}为
\[
 c_{k,1\omega}=f_{k,1}(C_{1\omega};\mu),\qquad
 \sum_k f_{k,1}(C;\mu)=C.
 \tag{10.10}
\]
0 期类似。\textbf{定理 10.6}：若 $C_{1\omega}>C_{1\omega'}$，则每个 $k$ 都有 $c_{k,1\omega}>c_{k,1\omega'}$。反证时，若某人消费不增，共同的加权边际效用条件与严格凹性意味着所有人消费均不增，与总资源增加矛盾。完全市场可以消除与总资源无关的个人禀赋风险，但不能消除总资源风险。

\textbf{定理 10.7（线性分享规则）}\quad 共同指数的 HARA 边际效用
\[
 u'_{k,1}(c)=\rho_k(d_k+c/\gamma)^{-\gamma},
 \qquad\rho_k>0,\quad\gamma\ne0,
\]
在消费处于定义域内时，足以产生
\[
 f_{k,1}(C;\mu)=b_k(C+\gamma\delta_k),\quad
 b_k=\frac{(\mu_k\rho_k)^{1/\gamma}}
 {\sum_j(\mu_j\rho_j)^{1/\gamma}},\quad
 \delta_k=\sum_jd_j-\frac{d_k}{b_k}.
 \tag{10.11}
\]
由规划者条件解出各人消费即可验证；$\sum_k b_k=1$、$\sum_k b_k\delta_k=0$ 保证资源出清。教材把必要性留在附录，本处只把共同指数作为充分条件；$\gamma=0$ 的极限情形须另行处理。
\end{tcolorbox}

\begin{tcolorbox}[knowledge,title=10.3 代表性参与者]
将规划问题按各期各状态分别求解，定义总消费的\textbf{间接效用}：
\[
 u_0(C_0)=\max_{\sum_kc_{k,0}=C_0}\sum_k\mu_ku_{k,0}(c_{k,0}),\quad
 u_1(C)=\max_{\sum_kc_{k,1}=C}\sum_k\mu_ku_{k,1}(c_{k,1}).
 \tag{10.12}
\]
由包络定理，$u'_0(C_0)=\phi'_0$、$p_\omega u'_1(C_{1\omega})=\phi'_\omega$；\textbf{定理 10.8}指出 $u_0,u_1$ 严格递增且严格凹。于是代表性参与者的期望效用为
\[
 U(C)=u_0(C_0)+\sum_\omega p_\omega u_1(C_{1\omega}),
 \quad \phi_\omega=\frac{p_\omega u'_1(C_{1\omega})}{u'_0(C_0)}.
 \tag{10.15}
\]
这里规划权重由均衡中的财富分布决定，因此代表性效用函数一般随财富分布变化。

\begin{tcolorbox}[example]
\textbf{例 10.2}\quad 对例 10.1，取参与者 1 权重为 1、参与者 2 权重为 $\mu$，聚合效用是
\[
 u(C)=\max_{c_1+c_2=C}\{\log c_1+2\mu\sqrt{c_2}\}.
\]
一阶条件 $1/c_1=\mu/\sqrt{c_2}$ 给出
\[
 c_1(C)=\frac{\sqrt{1+4\mu^2C}-1}{2\mu^2},\qquad
 u'(C)=\frac1{c_1(C)}.
\]
将总资源 $(C_0,C_a,C_b)=(1,2,1)$ 代入 (10.15)，再用个人预算确定 $\mu$，即可求回例 10.1 的状态价格。
\end{tcolorbox}

\textbf{定理 10.9（简单加总）}\quad 若个人效用同属平移幂效用族
\[
 u_k(c)=\frac{(d_k+c)^{1-\gamma}}{1-\gamma},\qquad\gamma>0,
\]
则聚合效用经正仿射标准化可写成同一形式
\[
 u(C)=\frac{(d+C)^{1-\gamma}}{1-\gamma},\qquad d=\sum_kd_k.
\]
$\gamma=1$ 取对数极限。\textbf{聚合形式}不依赖财富份额，但各人的消费分配仍随规划权重变化；须满足效用定义域与内点条件。
\end{tcolorbox}

\begin{tcolorbox}[knowledge,title=10.5 基于消费的资本资产定价模型]
代表性参与者的一阶条件给出随机支付 $\widetilde X_n$ 的 0 期价格：
\[
 S_n=\sum_\omega p_\omega
 \frac{u'_1(C_{1\omega})}{u'_0(C_0)}X_{n\omega}
 =\mathrm E[\widetilde m\widetilde X_n],\qquad
 \widetilde m:=\frac{u'_1(\widetilde C_1)}{u'_0(C_0)}.
 \tag{10.16}
\]
$\widetilde m$ 是随机贴现因子；状态价格为 $\phi_\omega=p_\omega m_\omega$。无风险支付 1 的债券价格为
\[
 B=\mathrm E[\widetilde m]=\frac1{1+r_F}.
 \tag{10.17}
\]
用协方差分解价格，得到
\[
 S_n=\frac{\mathrm E[\widetilde X_n]}{1+r_F}
 +\operatorname{Cov}(\widetilde m,\widetilde X_n).
 \tag{10.18}
\]
在总消费低、边际效用高的状态提供更多支付的证券有保险作用，价格相对更高。定价相关的是支付与\textbf{总消费边际效用}的协方差，而非支付方差本身。

若 $\widetilde r_n=\widetilde X_n/S_n-1$，由 $\mathrm E[\widetilde m(1+\widetilde r_n)]=1$ 得精确风险溢价关系
\[
 \bar r_n-r_F=-(1+r_F)\operatorname{Cov}(\widetilde m,\widetilde r_n),
 \qquad\bar r_n=\mathrm E[\widetilde r_n].
 \tag{10.21}
\]
对应的\textbf{现值公式}是
\[
 S_n=\frac{\mathrm E[\widetilde X_n]}{1+\bar r_n}.
 \tag{10.22}
\]
资产自己的期望收益率随风险而异，不能将所有风险资产一律按 $r_F$ 折现。

\textbf{消费增长与利率。}为作局部近似，设 $u_0=u$、$u_1=\rho u$，$\widetilde C_1=C_0(1+\widetilde g)$，$\widetilde g$ 较小。令 $R=-C_0u''(C_0)/u'(C_0)$，一阶泰勒展开给出
\[
 \widetilde m\approx\rho(1-R\widetilde g),\quad
 r_F\approx\left(\frac1\rho-1\right)+\frac{R}{\rho}\bar g,
 \qquad\bar g=\mathrm E[\widetilde g].
 \tag{10.24}
\]
同一局部近似还给出
\[
 \bar r_n-r_F\approx(1+r_F)\rho R
 \operatorname{Cov}(\widetilde g,\widetilde r_n).
\]
只有近似条件成立时，才能用消费增长协方差替换精确的边际效用协方差。

\textbf{风险测度与 CAPM。}完全市场中总消费索取权可由状态证券复制，但其收益率不必与消费增长率完全相同。若进一步用市场组合 $m$ 的收益率代表总消费风险，则可得到形式上类似 CAPM 的关系
\[
 \bar r_n-r_F\approx\beta_{n,m}(\bar r_m-r_F),\qquad
 \beta_{n,m}=\frac{\operatorname{Cov}(\widetilde r_n,\widetilde r_m)}
 {\operatorname{Var}(\widetilde r_m)}.
 \tag{10.25}
\]
精确定价核始终是边际效用；以可观测市场组合为风险代理还需要额外假设。第 13 章再讨论 CAPM 自身的条件。

\begin{tcolorbox}[example]
\textbf{例 10.3}\quad 若 $\bar g=0.02$、$\rho=0.98$、$R=2$，则 (10.24) 给出
\[
 r_F\approx(1/0.98-1)+(2/0.98)\times0.02\approx6.12\%.
\]
这是小消费增长风险下的近似计算；风险较大时应回到 $B=\mathrm E[\widetilde m]$ 的精确式。
\end{tcolorbox}
\end{tcolorbox}

\begin{tcolorbox}[summary,title=本章小结]
完全市场使或有支付可复制，并把证券市场与 Arrow--Debreu 经济接通。最优分享规则把个人消费写成总消费的函数；规划权重决定代表性参与者的效用。状态价格是概率乘以边际效用之比，证券价格等于随机贴现因子与支付乘积的期望。风险溢价由收益与边际效用的协方差决定；消费增长和市场组合的公式各有近似或附加条件。
\end{tcolorbox}

\fi

% !TeX root = ../../金融经济学_王江_讲义及习题.tex
% 10_完全市场 课后习题（由章节文件拆分）

\begin{problem}
\problemno{10.1*} 考虑一个两期（$t=0,\,1$）、两个状态（$\omega=a,\,b$），存在 $K$ 个参与者的经济。状态 $a$ 发生的概率为 $\pi_{a}$ 而状态 $b$ 发生的概率为 $\pi_{b}$。所有参与者都是不满足的且严格风险厌恶的。他们期望效用有如下形式：
\[
    u_{k}(c_{k,0})+\pi_{a}u_{k}(c_{k,a})+\pi_{b}u_{k}(c_{k,b}),\quad k=1,\dots,K.
\]
参与者 $k$ 的禀赋是 $\begin{bmatrix}e_{k,0} \\ e_{k,a} \\ e_{k,b}\end{bmatrix}$。整个经济的总禀赋是 $\begin{bmatrix}C_{0} \\ C_{a} \\ C_{b}\end{bmatrix}$，其中 $C_{s}=\sum_{k}e_{k,s}$，$s=0,a,b$。假设 $C_{a}>C_{b}$。假设市场上交易的 Arrow--Debreu 状态或有要求权的集合是完全的。令 $\phi_{a}$ 和 $\phi_{b}$ 分别是状态 $a$ 和状态 $b$ 的状态价格。
\begin{enumerate}[label=(\alph*)]
    \item 定义这个经济的均衡。简要讨论均衡的存在性和唯一性。
    \item 令代表性参与者的（间接）期望效用函数为
    \[
        u(C_{0})+\pi_{a}u(C_{a})+\pi_{b}u(C_{b})
    \]
    且 $u'''(\cdot)<0$。当以下变量变化而其它变量保持不变时，状态价格（$\phi_{a}$ 和 $\phi_{b}$）、它们的比例（$\phi_{a}/\phi_{b}$）和无风险利率如何变化（上升或下降）？给出证明。
    \begin{itemize}
        \item $C_{0}$ 上升；
        \item $C_{a}$ 上升；
        \item $C_{b}$ 上升；
        \item $C_{a}$ 和 $C_{b}$ 同比例上升；
        \item $C_{a}-C_{b}$ 上升（禀赋的风险上升）。
    \end{itemize}
\end{enumerate}
\end{problem}
\begin{solution}
\begin{enumerate}[label=(\alph*)]
    \item 这个经济的均衡是存在状态价格 $\begin{bmatrix}\phi_{a} \\ \phi_{b}\end{bmatrix}$，使得
    \begin{enumerate}[label=（\arabic*）]
        \item 给定状态价格和自己的禀赋，每个参与者都最优化地选择自己的消费；
        \item 市场出清：$\sum c_{k,0}=C_{0}$，$\sum c_{k,a}=C_{a}$，$\sum c_{k,b}=C_{b}$。由于市场是完全的，而参与者的效用函数是严格凹的，因而市场均衡是存在且唯一的。
    \end{enumerate}
    \item 我们知道
    \[
        \begin{aligned}
            \phi_{a}&=\pi_{a}u'(C_{a})/u'(C_{0}),\quad
            \phi_{b}=\pi_{b}u'(C_{b})/u'(C_{0}),\quad
            \frac{\phi_{a}}{\phi_{b}}=\frac{\pi_{a}u'(C_{a})}{\pi_{b}u'(C_{b})}\\
            r_{F}&=\frac{1}{\phi_{a}+\phi_{b}}-1=\frac{u'(C_{0})}{\pi_{a}u'(C_{a})+\pi_{b}u'(C_{b})}-1\\
            u'&>0,\quad u''<0,\quad u'''<0
        \end{aligned}
    \]
    \begin{itemize}
        \item 因为 $u''<0$，故 $C_{0}$ 上升使得 $u'(C_{0})$ 减小；$\phi_{a}$，$\phi_{b}$ 上升，$\phi_{a}/\phi_{b}$ 不变，$r_{F}$ 下降；
        \item 因为 $C_{a}$ 上升使得 $u'(C_{a})$ 下降，因而 $\phi_{a}$，$\phi_{a}/\phi_{b}$ 下降，$\phi_{b}$ 不变，$r_{F}$ 上升；
        \item 类似地，因为 $C_{b}$ 上升使得 $u'(C_{b})$ 下降，因而 $\phi_{b}$ 下降，$\phi_{a}/\phi_{b}$ 上升，$\phi_{a}$ 不变，$r_{F}$ 上升；
        \item $\phi_{a}$，$\phi_{b}$ 下降，$r_{F}$ 上升。现在考虑 $\phi_{a}/\phi_{b}$：
        \[
            \frac{\phi_{a}'/\phi_{b}'}{\phi_{a}/\phi_{b}}=\frac{u'(kC_{a})/u'(kC_{b})}{u'(C_{a})/u'(C_{b})}=\frac{u'(kC_{a})/u'(C_{a})}{u'(kC_{b})/u'(C_{b})}
        \]
        令 $f(x)=\dfrac{u'(kx)}{u'(x)}$，那么 $f'(x)=\dfrac{u'(kx)}{xu'(x)}(-R(kx)+R(x))$，其中 $R(x)$ 是相对风险厌恶系数，可以证明当 $u'''<0$ 时，$R(x)$ 是递增的，因而 $f'(x)<0$。因此，如果 $C_{a}>C_{b}$，那么 $\phi_{a}/\phi_{b}$ 会下降，反之则会上升。
        \item $\phi_{a}$，$\phi_{a}/\phi_{b}$ 下降，$\phi_{b}$ 上升，现在考虑 $r_{F}$。不妨假设 $C_{a}>C_{b}$，现在 $C_{a}$ 上升，$C_{b}$ 下降使得 $\phi_{a}C_{a}+\phi_{b}C_{b}=\tilde{C}$ 保持不变。考虑 $\phi_{a}u'(C_{a})+\phi_{b}u'(C_{b})=\phi_{a}u'(C_{a})+\phi_{b}u'(\tilde{C}-\pi_{a}C_{a})$，对 $C_{a}$ 求导可得 $\pi_{a}[u''(C_{a})-u''(C_{b})]$，因为 $u'''<0$，故 $r_{F}$ 上升。
    \end{itemize}
\end{enumerate}
\end{solution}

\begin{problem}
\problemno{10.2**} 考虑一个含有参与者 $k=1,\dots,K$ 的 Arrow--Debreu 经济。参与者 $k$ 具有形式为 $u_{k,0}(c_{0})+\mathrm{E}[u_{k,1}(c_{1})]$ 的效用函数，对于所有 $k$，$u_{0,k}$ 和 $u_{k,1}$ 是严格递增和严格凹的。$f_{k}^{0}$ 和 $f_{k}$ 是在给定权重集合 $\lambda=\{\lambda_{1},\dots,\lambda_{K}\}$ 下对于参与者 $k$ 的 Pareto 最优分配规则。定义 $u_{0}$ 和 $u_{1}$ 如下
\[
    u_{0}'(x)=\lambda_{k}u_{k,0}(f_{k}^{0}(x)),\qquad u_{1}'(x)=\lambda_{k}u_{k,1}(f_{k}(x)).
\]
\begin{enumerate}[label=(\alph*)]
    \item 证明 $u_{0}$ 和 $u_{1}$ 是严格递增和严格凹的。
    \item 证明 $u_{0}(c_{0})+\mathrm{E}[u_{1}(c_{1})]$ 是代表性参与者的效用函数。
    \item 证明如果所有参与者都是绝对风险厌恶递减的，那么以上定义的代表性参与者的效用函数也是绝对风险厌恶递减的。
\end{enumerate}
\end{problem}
\begin{solution}
\begin{enumerate}[label=(\alph*)]
    \item $u_{1}'(x)=\lambda_{k}u_{k,1}'(f_{k}(x))>0$，$u_{1}''(x)=\lambda_{k}u_{k,1}''(f_{k}(x))f_{k}'(x)<0$，这是因为 $u_{k}''<0$ 而 $f_{k}'>0$（定理 10.6），因而 $u_{1}$ 是严格递增且凹的，同样可以证明 $u_{0}$ 是严格递增且凹的；
    \item 由 (10.14) 式我们知道，代表性参与者的效用函数满足
    \[
        u_{0}'(x)=\lambda_{k}u_{k,0}(f_{k}^{0}(x)),\qquad u_{1}'(x)=\lambda_{k}u_{k,1}(f_{k}(x)),\quad\forall k
    \]
    因此（除相差一个常数为）$u_{0}(c_{0})+\mathrm{E}[u_{1}(c_{1})]$ 就是代表性参与者的效用函数。
    \item 由于 $u_{0}$、$u_{1}$ 的证明方法完全一样，因而我们忽略下标 $0$、$1$。如果参与者是绝对风险厌恶递减的，那么有 $[u_{k}'']^{2}-u_{k}'u_{k}'''<0$。对于代表性参与者有
    \[
        A(x)=-\frac{u_{k}''(f_{k}(x))f_{k}'(x)}{u_{k}'(f_{k}(x))}
    \]
    那么
    \[
        A'(x)=\frac{[u_{k}'']^{2}[f_{k}']^{2}-u_{k}'u_{k}''f_{k}''-u_{k}'u_{k}'''[f_{k}']^{2}}{[u_{k}']^{2}}
    \]
    因此，
    \[
        f_{k}''=\frac{A'(x)[u_{k}']^{2}+u_{k}'u_{k}'''[f_{k}']^{2}-[u_{k}'']^{2}[f_{k}']^{2}}{-u_{k}'u_{k}''}
    \]
    对于分配规则，我们有
    \[
        \sum f_{k}(x)=x,\quad \sum f_{k}'(x)=1,\quad \sum f_{k}''(x)=0
    \]
    因此我们有
    \[
        A'(x)=\sum\frac{[f_{k}']^{2}\bigl([u_{k}'']^{2}-u_{k}'u_{k}'''\bigr)}{-u_{k}'u_{k}''}\Big/\sum\frac{-u_{k}'}{u_{k}''}<0
    \]
    因而代表性参与者的效用函数也是绝对厌恶递减的。
\end{enumerate}
\end{solution}

\begin{problem}
\problemno{10.3*} 在 1 期，经济有两个可能状态 $a$ 和 $b$。它们的发生的概率分别为 $\pi$ 和 $1-\pi$。市场有两只交易证券，证券 1 和 2，支付如下：
\begin{center}
\begin{tabular}{@{}c@{\hspace{2.5em}}c@{}}
证券 1：
\WJEventTree{1,1}
&
证券 2：
\WJEventTree{uC,dC}
\end{tabular}
\end{center}
其中 $C>0$ 且 $0<d<u$。因此，证券 1 是无风险债券而证券 2 是风险证券，我们可以称它为股票。经济中有参与者 $k=1,\dots,K$。参与者 $k$ 的禀赋是 $s_{k}C$ 单位的当期消费、0 单位的证券 1 和 $s_{k}$ 单位的证券 2，其中 $s_{k}>0$ 且 $\sum_{k}s_{k}=1$。所有参与者具有如下的期望效用函数：
\[
    U_{k}(c_{k})=\log c_{k,0}+\rho\bigl[\pi\log c_{k,1a}+(1-\pi)\log c_{k,1b}\bigr]
\]
其中 $\rho>0$。
\begin{enumerate}[label=(\alph*)]
    \item 证明证券市场是完全的。
    \item 求解每一参与者的最优消费/投资选择。
    \item 求解均衡证券价格。
    \item 将价格过程与期权定价中的二叉树模型相比较，你能得到什么结论？
\end{enumerate}
\end{problem}
\begin{solution}
\begin{enumerate}[label=(\alph*)]
    \item 市场上有两只支付独立的证券，而状态数目恰为 2，因而市场是完全的。
    \item 参与者 $k$ 面临的问题是
    \[
        \begin{aligned}
            &\max_{\{c_{k,0},c_{k,1a},c_{k,1b}\}}\ \log c_{k,0}+\rho\bigl[\pi\log c_{k,1a}+(1-\pi)\log c_{k,1b}\bigr]\\
            &\quad\text{s.t.}\quad c_{k,0}+\phi_{a}c_{k,1a}+\phi_{b}c_{k,1b}=S_{k}C(1+\phi_{a}u+\phi_{b}d)
        \end{aligned}
    \]
    记财富 $w_{k}=S_{k}C(1+\phi_{a}u+\phi_{b}d)$。拉格朗日函数为
    \[
        \mathcal{L}=\log c_{k,0}+\rho\bigl[\pi\log c_{k,1a}+(1-\pi)\log c_{k,1b}\bigr]+\lambda\bigl(w_{k}-c_{k,0}-\phi_{a}c_{k,1a}-\phi_{b}c_{k,1b}\bigr)
    \]
    一阶条件（Euler 方程）为
    \[
        \frac{1}{c_{k,0}}=\lambda,\qquad
        \frac{\rho\pi}{c_{k,1a}}=\lambda\phi_{a},\qquad
        \frac{\rho(1-\pi)}{c_{k,1b}}=\lambda\phi_{b}
    \]
    因而 $c_{k,1a}^{*}=\rho\pi/(\lambda\phi_{a})$，$c_{k,1b}^{*}=\rho(1-\pi)/(\lambda\phi_{b})$。代入预算
    \[
        \frac{1}{\lambda}+\frac{\rho\pi}{\lambda}+\frac{\rho(1-\pi)}{\lambda}=w_{k}\implies c_{k,0}^{*}=\frac{w_{k}}{1+\rho}
    \]
    得到最优消费
    \[
        \begin{aligned}
            c_{k,0}^{*}&=\frac{S_{k}C(1+\phi_{a}u+\phi_{b}d)}{1+\rho}\\
            c_{k,1a}^{*}&=\frac{\rho\pi S_{k}C(1+\phi_{a}u+\phi_{b}d)}{\phi_{a}(1+\rho)}\\
            c_{k,1b}^{*}&=\frac{\rho(1-\pi)S_{k}C(1+\phi_{a}u+\phi_{b}d)}{\phi_{b}(1+\rho)}
        \end{aligned}
    \]
    当市场达到均衡时：
    \[
        \begin{aligned}
            \sum c_{k,0}^{*}&=\sum S_{k}C=C\\
            \sum c_{k,1a}^{*}&=\sum S_{k}uC=uC\\
            \sum c_{k,1b}^{*}&=\sum S_{k}dC=dC
        \end{aligned}
    \]
    因此求得
    \[
        \begin{aligned}
            \phi_{a}&=\frac{\rho\pi}{u},\qquad\phi_{b}=\frac{\rho(1-\pi)}{d}\\
            c_{k,0}^{*}&=S_{k}C,\qquad c_{k,1a}^{*}=S_{k}uC,\qquad c_{k,1b}^{*}=S_{k}dC
        \end{aligned}
    \]
    \item 债券价格为 $\dfrac{\rho\pi}{u}+\dfrac{\rho(1-\pi)}{d}$，股票价格为 $\rho C$。
    \item 在二叉树模型中，债券、股票的价格都是给定的，因而状态价格只由债券、股票的支付给定，得到的是部分均衡的结果；在这个模型中，价格是通过市场均衡求得的，因而状态价格还与参与者的效用函数、时间偏好系数、状态发生的概率有关，得到的是一般均衡的结果。
\end{enumerate}
\end{solution}

\begin{problem}
\problemno{10.4**} 继续考虑练习 10.3 定义的经济，
\begin{enumerate}[label=(\alph*)]
    \item 构建代表性参与者。
    \item 计算在代表性参与者（即中央计划者）的效用函数中每个参与者的权重。它依赖于什么？为什么？
    \item 证明代表性参与者的偏好与禀赋在参与者之间的分布无关。
    \item 用基于消费的 CAPM 为两只证券定价。
    \item 计算无风险收益和股票的风险溢价。
    \item 利率和股票的风险溢价依赖于什么？解释它们对经济原生变量：当前消费水平、消费增长率的期望值 $\mathrm{E}[\tilde{C}_{1}/C_{0}]$、消费增长率的波动率 $\mathrm{Var}[\tilde{C}_{1}/C_{0}]$ 以及时间折现系数 $\rho$ 的依赖性？
\end{enumerate}
\end{problem}
\begin{solution}
\begin{enumerate}[label=(\alph*)]
    \item 由第二题我们知道，代表性参与者的效用函数满足
    \[
        u_{0}'(C)=\mu_{k}u_{k,0}'(c_{k})=\frac{\mu_{k}}{c_{k}}=\frac{\mu_{k}}{S_{k}C},\quad\forall k
    \]
    因此 $\mu_{k}=S_{k}$，$u_{0}'(C)=\dfrac{1}{C}$，$u_{0}(C)=\log C$，同样地，$u_{1}(C)=\log C$；因而代表性参与者的效用函数为
    \[
        \log C_{0}+\rho\bigl[\pi\log C_{1a}+(1-\pi)\log C_{1b}\bigr]
    \]
    \item 参与者 $k$ 的权重是 $S_{k}$，即他的禀赋占总禀赋的比例。如果参与者拥有的资源愈多，在均衡条件下，他的消费越多，权重也越大；
    \item 代表性参与者的偏好与 $S_{k}$ 无关，即与禀赋在参与者之间的分布无关；
    \item 由 (10.16) 式我们有：
    \[
        \begin{aligned}
            \text{债券价格为：}&\quad
            \pi\frac{\rho C}{uC}+(1-\pi)\frac{\rho C}{dC}=\rho\Bigl[\frac{\pi}{u}+\frac{1-\pi}{d}\Bigr]\\
            \text{股票价格为：}&\quad
            \pi\frac{\rho C}{uC}uC+(1-\pi)\frac{\rho C}{dC}dC=\rho C
        \end{aligned}
    \]
    \item 无风险收益率为
    \[
        r_{F}=1\big/\Bigl(\rho\Bigl[\frac{\pi}{u}+\frac{1-\pi}{d}\Bigr]\Bigr)-1
    \]
    风险溢价为
    \[
        p=\Bigl(\frac{\pi}{u}+\frac{1-\pi}{d}\Bigr)\big/\rho-1\big/\Bigl(\rho\Bigl[\frac{\pi}{u}+\frac{1-\pi}{d}\Bigr]\Bigr)
    \]
    \item 如果 $u$，$d$ 与 1 的距离很小，我们可以把风险溢价写成 $p=\dfrac{\mathrm{E}[\tilde{C}_{1}]/C_{0}}{\rho}-r_{F}$，其中 $r_{F}=1\big/\bigl\{\rho\bigl[\mathrm{Var}[\tilde{C}_{1}/C_{0}]-2\mathrm{E}[\tilde{C}_{1}/C_{0}]+3\bigr]\bigr\}$。无风险利率依赖于时间偏好系数 $\rho$、消费增长率的波动率 $\mathrm{Var}[\tilde{C}_{1}/C_{0}]$、以及消费增长率的期望值 $\mathrm{E}[\tilde{C}_{1}/C_{0}]$。当 $\rho$ 减小时，参与者偏好于当前消费，而为了吸引投资者储蓄以平衡当前消费和未来消费，利率 $r_{F}$ 应该上升；当未来消费的风险上升时，为规避未来消费的风险，参与者对债券的需求上升，其价格上升、利率下降；当 $\mathrm{E}[\tilde{C}_{1}/C_{0}]$ 上升时，人们倾向于在时间上平滑消费水平，即增加当前消费，为吸引投资，利率上升。至于风险溢价，当 $\rho$ 上升时，参与者倾向于未来消费，因而债券、股票的需求量都会上升，利率和股票收益率都减小，而且对两者的影响相对一样，因而风险溢价变小（风险溢价一般都是正的）。当 $\mathrm{E}[\tilde{C}_{1}/C_{0}]$ 变大时，参与者倾向于当前消费，因而对债券和股票的需求量都减小，两者的价格都下降、收益率都上升，对风险溢价的影响是不确定的，但当 $\mathrm{E}[\tilde{C}_{1}/C_{0}]$ 较小时，利率上升更快，风险溢价会变小。当 $\mathrm{Var}[\tilde{C}_{1}/C_{0}]$ 变大时，因为股票的风险和禀赋的风险完全一致，为了减少未来消费的波动，人们对债券的需求增加而对股票的需求减少，因而债券收益率下降而股票收益上升，风险溢价上升。
\end{enumerate}
\end{solution}

\begin{problem}
\problemno{10.5} 继续前面练习中的经济。
\begin{enumerate}[label=(\alph*)]
    \item 求解以证券 2 为标的资产、执行价格为 $K$ 的看涨期权的价格。
    \item 证明基于消费的 CAPM 也适用于期权。
\end{enumerate}
\end{problem}
\begin{solution}
\begin{enumerate}[label=(\alph*)]
    \item 不失一般性，假设 $dC\le K\le uC$，期权的支付为 $\begin{bmatrix}uC-K \\ 0\end{bmatrix}$，用复制的方法可以求得期权的价格为 $\dfrac{\rho\pi}{u}(uC-K)$；
    \item 而用 CCAPM 定价，期权的价格为
    \[
        \pi\frac{\rho C}{uC}[uC-K]_{+}+(1-\pi)\frac{\rho C}{dC}[dC-K]_{+}=\frac{\rho\pi}{u}(uC-K)
    \]
    因而 CCAPM 成立。
\end{enumerate}
\end{solution}

\begin{problem}
\problemno{10.6**} 在 1 期，经济有两个可能状态 $a$ 和 $b$。它们的发生概率分别为 $\pi$ 和 $1-\pi$。市场有两只交易证券，证券 1 和 2，支付如下：
\begin{center}
\begin{tabular}{@{}c@{\hspace{2.5em}}c@{}}
证券 1：
\WJEventTree{1,1}
&
证券 2：
\WJEventTree{C_{1a},C_{1b}}
\end{tabular}
\end{center}
其中 $C>0$ 且 $0<d<u$。因此，证券 1 是无风险债券而证券 2 是风险证券，我们可以称它为股票。经济中有参与者 $k=1,\dots,K$。参与者 $k$ 的禀赋是 $s_{k}C_{0}$ 单位的当期消费、0 单位的证券 1 和 $s_{k}$ 单位的证券 2，其中 $s_{k}>0$ 且 $\sum_{k}s_{k}=1$。所有参与者具有如下的期望效用函数：
\[
    U_{k}(c_{k})=-e^{-\alpha c_{k,0}}-\rho\bigl[\pi e^{-\alpha c_{k,1a}}+(1-\pi)e^{-\alpha c_{k,1b}}\bigr]
\]
其中 $\rho>0$。
\begin{enumerate}[label=(\alph*)]
    \item 证明证券市场是完全的。
    \item 求解每一参与者的最优消费/投资选择。
    \item 求解均衡证券价格。
\end{enumerate}
\end{problem}
\begin{solution}
\begin{enumerate}[label=(\alph*)]
    \item 有 2 个可能状态，而支付独立的状态数目恰为 2，因此证券市场是完全的；
    \item 假设状态价格为 $\begin{bmatrix}\phi_{a} \\ \phi_{b}\end{bmatrix}$；考虑参与者 $k$，他的优化问题是
    \[
        \begin{aligned}
            &\max_{c_{k,0},\,c_{k,1a},\,c_{k,1b}}\ -e^{-\alpha c_{k,0}}-\rho\bigl[\pi e^{-\alpha c_{k,1a}}+(1-\pi)e^{-\alpha c_{k,1b}}\bigr]\\
            &\quad\text{s.t.}\quad c_{k,0}+\phi_{a}c_{k,1a}+\phi_{b}c_{k,1b}=s_{k}C_{0}+\phi_{a}s_{k}C_{1a}+\phi_{b}s_{k}C_{1b}
        \end{aligned}
    \]
    令 $w_{k}=s_{k}W$、$W=C_{0}+\phi_{a}C_{1a}+\phi_{b}C_{1b}$ 为经济中的总财富。拉格朗日函数为
    \[
        \mathcal{L}=-e^{-\alpha c_{k,0}}-\rho\bigl[\pi e^{-\alpha c_{k,1a}}+(1-\pi)e^{-\alpha c_{k,1b}}\bigr]+\lambda\bigl(w_{k}-c_{k,0}-\phi_{a}c_{k,1a}-\phi_{b}c_{k,1b}\bigr)
    \]
    一阶条件
    \[
        \alpha e^{-\alpha c_{k,0}}=\lambda,\qquad
        \rho\pi\alpha e^{-\alpha c_{k,1a}}=\lambda\phi_{a},\qquad
        \rho(1-\pi)\alpha e^{-\alpha c_{k,1b}}=\lambda\phi_{b}
    \]
    给出最优消费 $c_{k,0}^{*}=\dfrac{1}{\alpha}\log(\alpha/\lambda)$，$c_{k,1a}^{*}=\dfrac{1}{\alpha}\log(\rho\pi\alpha/(\lambda\phi_{a}))$，$c_{k,1b}^{*}=\dfrac{1}{\alpha}\log(\rho(1-\pi)\alpha/(\lambda\phi_{b}))$。代入预算并记 $\beta=\log\alpha+\phi_{a}\log\dfrac{\rho\pi\alpha}{\phi_{a}}+\phi_{b}\log\dfrac{\rho(1-\pi)\alpha}{\phi_{b}}$，得
    \[
        \begin{aligned}
            c_{k,0}^{*}&=\frac{1}{\alpha}\log\alpha-\frac{\beta-\alpha s_{k}W}{\alpha(1+\phi_{a}+\phi_{b})}\\
            c_{k,1a}^{*}&=\frac{1}{\alpha}\log\frac{\rho\pi\alpha}{\phi_{a}}-\frac{\beta-\alpha s_{k}W}{\alpha(1+\phi_{a}+\phi_{b})}\\
            c_{k,1b}^{*}&=\frac{1}{\alpha}\log\frac{\rho(1-\pi)\alpha}{\phi_{b}}-\frac{\beta-\alpha s_{k}W}{\alpha(1+\phi_{a}+\phi_{b})}
        \end{aligned}
    \]
    \item 市场出清条件是
    \[
        \begin{aligned}
            \sum c_{k,0}&=C_{0}\\
            \sum c_{k,1a}&=C_{1a}\\
            \sum c_{k,1b}&=C_{1b}
        \end{aligned}
    \]
    由此可以求得均衡状态价格为
    \[
        \begin{aligned}
            \phi_{a}&=\rho\pi e^{-\frac{\alpha}{K}(C_{1a}-C_{0})}\\
            \phi_{b}&=\rho(1-\pi)e^{-\frac{\alpha}{K}(C_{1b}-C_{0})}
        \end{aligned}
    \]
    因而债券的均衡价格为 $\phi_{a}+\phi_{b}$，股票的均衡价格为 $\phi_{a}C_{1a}+\phi_{b}C_{1b}$；在均衡条件下，参与者 $k$ 的最优消费选择为
    \[
        \begin{aligned}
            c_{k,0}^{*}&=\frac{1}{K}C_{0}+\Bigl(s_{k}-\frac{1}{K}\Bigr)\frac{W}{1+\phi_{a}+\phi_{b}}\\
            c_{k,1a}^{*}&=\frac{1}{K}C_{1a}+\Bigl(s_{k}-\frac{1}{K}\Bigr)\frac{W}{1+\phi_{a}+\phi_{b}}\\
            c_{k,1b}^{*}&=\frac{1}{K}C_{1b}+\Bigl(s_{k}-\frac{1}{K}\Bigr)\frac{W}{1+\phi_{a}+\phi_{b}}
        \end{aligned}
    \]
    在 0 期他持有的无风险债券、股票的头寸分别为
    \[
        \begin{aligned}
            \theta_{B}&=\Bigl(s_{k}-\frac{1}{K}\Bigr)\frac{W}{1+\phi_{a}+\phi_{b}}\\
            \theta_{S}&=\frac{1}{K}-s_{k}.
        \end{aligned}
    \]
\end{enumerate}
\end{solution}

\begin{problem}
\problemno{10.7**} 继续考虑练习 10.6 定义的经济
\begin{enumerate}[label=(\alph*)]
    \item 构建代表性参与者。
    \item 计算在代表性参与者（即中央计划者）的效用函数中每个参与者的权重。它依赖于什么？为什么？
    \item 证明代表性参与者的偏好与禀赋在参与者之间的分布无关。
    \item 用基于消费的 CAPM 为两只证券定价。
    \item 计算无风险收益和股票的风险溢价，现在风险溢价的定义为 $\mathrm{E}[\tilde{S}_{1}]-(1+r_{F})S_{0}$。即一股股票的绝对超额收益。
    \item 利率和股票的风险溢价依赖于什么？解释它们对经济原生变量：当前消费水平、消费增长的期望值 $\mathrm{E}[\tilde{C}_{1}/C_{0}]$、消费增长的波动率 $\mathrm{Var}[\tilde{C}_{1}/C_{0}]$ 以及时间折现系数 $\rho$ 的依赖性？
    \item 求解以证券 2 为标的资产、执行价格为 $K$ 的看涨期权的价格。
\end{enumerate}
\end{problem}
\begin{solution}
\begin{enumerate}[label=(\alph*)]
    \item 因为每个参与者的 0、1 期效用函数只差时间偏好系数 $\rho$，因而代表性参与者的 0、1 期效用函数只差时间偏好系数 $\rho$；假设中央计划者问题中参与者 $k$ 的权重为 $\mu_{k}>0$，代表性参与者的 0 期效用函数为 $u(C_{0})$，那么
    \[
        \begin{aligned}
            u'(C_{0})
            &=\mu_{k}u_{k}'(c_{k,0})\\
            &=\mu_{k}\alpha e^{-\alpha\bigl(\frac{1}{K}C_{0}+\bigl(s_{k}-\frac{1}{K}\bigr)\frac{W}{1+\phi_{a}+\phi_{b}}\bigr)}\\
            &=e^{-\alpha\frac{C_{0}}{K}}\Bigl(\alpha\mu_{k}e^{-\alpha\bigl(s_{k}-\frac{1}{K}\bigr)\frac{W}{1+\phi_{a}+\phi_{b}}}\Bigr),\quad\forall k
        \end{aligned}
    \]
    因此，$\alpha\mu_{k}e^{-\alpha\bigl(s_{k}-\frac{1}{K}\bigr)\frac{W}{1+\phi_{a}+\phi_{b}}}$ 是与 $k$ 无关的常数，即 $\mu_{k}$ 与 $e^{\frac{\alpha}{1+\phi_{a}+\phi_{b}}s_{k}}$ 成正比，而且代表性参与者的效用函数具有负指数形式 $u(C_{0})=-e^{-\alpha\frac{C_{0}}{K}}$，因此，代表性参与者的效用函数为
    \[
        -e^{-\alpha\frac{C_{0}}{K}}-\rho\Bigl[\pi e^{-\alpha\frac{C_{1a}}{K}}+(1-\pi)e^{-\alpha\frac{C_{1b}}{K}}\Bigr]
    \]
    \item 权重 $\mu_{k}$ 与 $e^{\frac{\alpha}{1+\phi_{a}+\phi_{b}}s_{k}}$ 成正比，依赖于参与者的禀赋参数 $s_{k}$，即他的禀赋占总禀赋的比例。如果参与者拥有的资源愈多，在均衡条件下，他的消费越多，权重也越大；
    \item 代表性参与者的偏好与 $s_{k}$ 无关，即与禀赋在参与者之间的分布无关；
    \item
    \[
        \begin{aligned}
            \text{债券价格为：}&\quad
            \rho\pi\frac{\alpha e^{-\alpha\frac{C_{1a}}{K}}}{\alpha e^{-\alpha\frac{C_{0}}{K}}}+\rho(1-\pi)\frac{\alpha e^{-\alpha\frac{C_{1b}}{K}}}{\alpha e^{-\alpha\frac{C_{0}}{K}}}\\
            &\qquad=\rho\Bigl[\pi e^{-\alpha\frac{C_{1a}-C_{0}}{K}}+(1-\pi)e^{-\alpha\frac{C_{1b}-C_{0}}{K}}\Bigr]\\
            \text{股票价格为：}&\quad
            \rho\pi\frac{\alpha e^{-\alpha\frac{C_{1a}}{K}}}{\alpha e^{-\alpha\frac{C_{0}}{K}}}C_{1a}+\rho(1-\pi)\frac{\alpha e^{-\alpha\frac{C_{1b}}{K}}}{\alpha e^{-\alpha\frac{C_{0}}{K}}}C_{1b}\\
            &\qquad=\rho\Bigl[\pi e^{-\alpha\frac{C_{1a}-C_{0}}{K}}C_{1a}+(1-\pi)e^{-\alpha\frac{C_{1b}-C_{0}}{K}}C_{1b}\Bigr]
        \end{aligned}
    \]
    \item 无风险收益率为
    \[
        r_{F}=\frac{1}{\phi_{a}+\phi_{b}}-1=\frac{1}{\rho\Bigl[\pi e^{-\alpha\frac{C_{1a}-C_{0}}{K}}+(1-\pi)e^{-\alpha\frac{C_{1b}-C_{0}}{K}}\Bigr]}-1
    \]
    风险溢价为
    \[
        \begin{aligned}
            p&=\mathrm{E}[\tilde{S}_{1}]-(1+r_{F})S_{0}\\
            &=\pi C_{1a}+(1-\pi)C_{1b}-\frac{\phi_{a}C_{1a}+\phi_{b}C_{1b}}{\phi_{a}+\phi_{b}}\\
            &=\Bigl(\pi-\frac{\phi_{a}}{\phi_{a}+\phi_{b}}\Bigr)(C_{1a}-C_{1b})\\
            &=\frac{\pi(1-\pi)\bigl(e^{\alpha\frac{C_{1a}-C_{1b}}{K}}-1\bigr)}{\pi+(1-\pi)e^{\alpha\frac{C_{1a}-C_{1b}}{K}}}(C_{1a}-C_{1b})
        \end{aligned}
    \]
    \item 如果 $\tilde{g}=\tilde{C}_{1}/C_{0}-1$ 很小，而人均消费 $C_{0}/K$ 是很大，那么无风险利率可以写成
    \[
        r_{F}\doteq\frac{1}{\rho\Bigl[1-\alpha\frac{C_{0}}{K}\mathrm{E}[\tilde{g}]+\alpha^{2}\frac{C_{0}^{2}}{K^{2}}\bigl(\mathrm{Var}[\tilde{g}]+\mathrm{E}^{2}[\tilde{g}]\bigr)\Bigr]}
    \]
    可以看出，无风险利率依赖于时间偏好系数 $\rho$、消费增长率的波动率 $\mathrm{Var}[\tilde{C}_{1}/C_{0}]$、以及消费增长率的期望值 $\mathrm{E}[\tilde{C}_{1}/C_{0}]$。当 $\rho$ 减小时，参与者偏好于当前消费，而为了吸引投资者储蓄以平衡当前消费和未来消费，利率 $r_{F}$ 应该上升；当未来消费的风险上升时，为规避未来消费的风险，参与者对债券的需求上升，其价格上升、利率下降；当 $\mathrm{E}[\tilde{C}_{1}/C_{0}]$ 上升时，人们倾向于在时间上平滑消费水平，即增加当前消费，为吸引投资，利率上升。至于风险溢价，它只与 $C_{1a}-C_{1b}$ 有关，且是 $C_{1a}-C_{1b}$ 的增函数；某种意义上我们可以把 $C_{1a}-C_{1b}$ 作为未来禀赋风险的一种度量，也就是说，当未来禀赋的风险增加时，风险溢价增加，这是因为股票的风险和禀赋的风险完全一致，为了减少未来消费的波动，人们对债券的需求增加而对股票的需求减少，因而债券收益率下降而股票收益上升，风险溢价上升。
    \item 可以用 CCAPM 为期权定价
    \[
        \begin{aligned}
            c&=\rho\pi\frac{\alpha e^{-\alpha\frac{C_{1a}}{K}}}{\alpha e^{-\alpha\frac{C_{0}}{K}}}[C_{1a}-K]_{+}+\rho(1-\pi)\frac{\alpha e^{-\alpha\frac{C_{1b}}{K}}}{\alpha e^{-\alpha\frac{C_{0}}{K}}}[C_{1b}-K]_{+}\\
            &=\rho\Bigl[\pi e^{-\alpha\frac{C_{1a}-C_{0}}{K}}[C_{1a}-K]_{+}+(1-\pi)e^{-\alpha\frac{C_{1b}-C_{0}}{K}}[C_{1b}-K]_{+}\Bigr]
        \end{aligned}
    \]
\end{enumerate}
\end{solution}



% !TeX root = ../../金融经济学_王江_讲义及习题.tex
% 知识点：references/金融经济学/金融经济学-王江_261013.pdf，印刷页 169–175（PDF 页 184–190）
% 习题：references/金融经济学/史上最全最完整的金融经济学答案(王江版).pdf，页 54–61（印刷页 53–60）

\lecture{不完全市场中的资源配置和资产定价}

\ifdefined\WJHandoutKnowledge

完全市场中证券支付可以生成所有状态的支付；本章研究证券支付空间较小时，消费选择、资源配置和证券价格如何变化。关键是区分\textbf{可获得的配置}与所有物理可行的配置。

\begin{tcolorbox}[knowledge,title=11.1 消费选择与证券价格]
设支付矩阵 $X$ 有 $N$ 个独立列，状态数为 $\Omega>N$。交易可生成的支付空间为
\[
 M=\{X\theta:\theta\in\mathbb R^N\}\subsetneq\mathbb R^\Omega.
\]
参与者 $k$ 的 1 期消费仅能属于 $e_{k,1}+M$。\textbf{例 11.1}：两状态只有一只支付 $(2,1)$ 的证券，持仓 $\theta$ 后消费为 $(e_{k,a}+2\theta,e_{k,b}+\theta)$，可选点只在二维平面的一条直线上。

对持仓的内点优化给出 Euler 方程
\[
 S_nu'_{k,0}(c_{k,0})
 =\mathrm E[u'_{k,1}(\widetilde c_{k,1})\widetilde X_n],
 \qquad n=1,\ldots,N.
 \tag{11.1}
\]
记个人随机贴现因子 $m_k=u'_{k,1}(\widetilde c_{k,1})/u'_{k,0}(c_{k,0})$，等价地
\[
 S_n=\mathrm E[m_k\widetilde X_n],\qquad\forall k,n.
 \tag{11.2}
\]
所有参与者对\textbf{已交易证券}的边际估值一致，但他们的 $m_k$ 可因未能交易的风险而不同。将 $m_k$ 乘以状态概率，可得到给现有证券定价的正状态价格向量；由于 $N<\Omega$，这种向量通常不唯一。仅凭个人 Euler 方程不能像完全市场那样从总消费唯一确定所有或有支付的价格。
\end{tcolorbox}

\begin{tcolorbox}[knowledge,title=11.2 约束下的最优配置]
\textbf{定义 11.1}\quad 给定禀赋，若每人的消费变化都可由市场证券支付复制，即 $c_{k,1}-e_{k,1}\in M$，且 0 期预算可满足，则称配置\textbf{可获得}。市场不完全时，满足总资源约束的物理可行配置未必可获得。

\textbf{定义 11.2}\quad 若不存在另一可获得配置，在保持总资源不变时使所有人弱变好、至少一人严格变好，则配置\textbf{约束最优}。它只在现有证券结构允许的交易集合内比较。

\textbf{定理 11.1}\quad 在本章的无套利、局部不满足与均衡条件下，证券市场均衡是约束最优的。证明与第一福利定理类似：若存在可获得的帕累托改进，每个改进者所需的证券组合按均衡价格至少不便宜；对所有人加总，便与证券市场出清矛盾。约束最优\emph{不等于}无约束帕累托最优。
\end{tcolorbox}

\begin{tcolorbox}[knowledge,title=11.3 实质完全市场]
\textbf{定义 11.3}\quad 对\emph{给定}禀赋与偏好，若所有帕累托最优配置都可通过现有证券交易实现，市场称为\textbf{实质完全}。它是相对于该经济的概念，比 $\operatorname{rank}(X)=\Omega$ 的结构性完全更弱。

\textbf{定理 11.2}\quad 在本章均衡与福利定理的条件下，实质完全市场的均衡配置是帕累托最优的：若有更优的物理可行配置，实质完全性会使它也可获得，与约束最优矛盾。

\textbf{定理 11.3}\quad 若交易以市场组合为标的、执行价格齐备的期权，并满足无套利、可交易与支付可分辨的条件，期权的分段支付可复制市场组合的各个不同支付值。若市场组合只有 $N_d$ 个不同的状态支付值，所需独立支付的维数是 $N_d$，它可以小于 $\Omega$。当帕累托最优风险分担只依赖总禀赋，即只依赖市场组合支付时，复制这些支付值即可使市场实质完全。第 5 章以单个\emph{状态}为标的补齐期权支付的结构性完全化，与这里以\emph{市场组合支付}为标的的论证不同。

\textbf{定理 11.4}\quad 若有无风险债券，且所有参与者 1 期效用的边际效用同属共同指数的 HARA 形式
\[
 u'_{k,1}(c)=\rho_k(d_k+c/\gamma)^{-\gamma},
 \qquad \rho_k>0,\quad\gamma>-1,\quad\gamma\ne0,
\]
在分享规则和消费满足定义域、内点条件下，市场实质完全。由第 10 章的线性分享规则，每个人的最优消费是 $b_kC_1$ 加一个常数；总禀赋 $C_1$ 由市场组合支付生成，常数由债券复制。因此所有帕累托最优分享都可由证券实现。这里还需要市场组合与债券可交易，不能仅凭偏好形式推出结论。
\end{tcolorbox}

\begin{tcolorbox}[summary,title=本章小结]
不完全市场限制了个人的可交易支付空间。已交易证券的 Euler 方程仍成立，但不同参与者可以有不同的随机贴现因子，也可能存在多个能给现有证券定价的状态价格向量。均衡一般只保证约束最优；当现有证券足以实现该经济所有帕累托最优风险分担时，市场虽未结构性完全，均衡仍可达到帕累托最优。
\end{tcolorbox}

\fi

% !TeX root = ../../金融经济学_王江_讲义及习题.tex
% 11_不完全市场 课后习题（由章节文件拆分）

\begin{problem}
\problemno{11.1*} 经济有两期 $t=0,1$。有两个参与者，1 和 2，他们都有定义于消费路径上的对数效用函数 $c=(c_{0},c_{1})$：$u(c)=\log c_{0}+\log c_{1}$。在 $t=1$ 期有三个等概率的可能状态，$\omega_{1}$，$\omega_{2}$ 以及 $\omega_{3}$。在 $t=0$ 期的总禀赋是 1 而 $t=1$ 期的总禀赋是 $\begin{bmatrix}1 \\ 1 \\ 2\end{bmatrix}$。有两只交易证券，一只是无风险债券，在 $t=1$ 期的所有状态下都支付 1 单位的消费品。另外一只是股票，它的 $t=1$ 期支付和总禀赋完全一样，即它的支付向量是 $\begin{bmatrix}1 \\ 1 \\ 2\end{bmatrix}$。
\begin{enumerate}[label=(\alph*)]
    \item 假设参与者 1 和 2 的禀赋分别为 $\begin{bmatrix}0.4\\0.4\\0.4\\0.8\end{bmatrix}$ 和 $\begin{bmatrix}0.6\\0.6\\0.6\\1.2\end{bmatrix}$。求解均衡利率和均衡股价。均衡配置是 Pareto 最优的吗？解释原因。
    \item 假设参与者 $a$ 和 $b$ 的禀赋分别为 $\begin{bmatrix}0.5\\1\\0\\1\end{bmatrix}$ 和 $\begin{bmatrix}0.5\\0\\1\\1\end{bmatrix}$。导出求解均衡的一阶条件和市场出清条件。在不解出均衡的情况下，讨论均衡配置是否是 Pareto 最优的。解释原因。
\end{enumerate}
\end{problem}
\begin{solution}
\begin{enumerate}[label=(\alph*)]
    \item 假设债券股票的价格分别为 $B$，$S$。对于参与者 1 来说，假设他持有的债券和股票头寸分别为 $\theta_{11}$，$\theta_{12}$，那么他的优化问题是
    \[
        \begin{aligned}
            &\max_{\theta_{11},\,\theta_{12}}\ \log(c_{0})+\tfrac13\bigl[\log(c_{1,1})+\log(c_{1,2})+\log(c_{1,3})\bigr]\\
            &\quad\text{s.t.}\quad c_{0}=0.4-\theta_{11}B-\theta_{12}S\\
            &\qquad c_{1,1}=0.4+\theta_{11}+\theta_{12}\\
            &\qquad c_{1,2}=0.4+\theta_{11}+\theta_{12}\\
            &\qquad c_{1,3}=0.8+\theta_{11}+2\theta_{12}
        \end{aligned}
    \]
    可以求得
    \[
        \begin{aligned}
            \theta_{11}&=-\frac{1}{15}(5S^{2}+5S-6BS-6B)/(S^{2}-3BS+2B^{2})\\
            \theta_{12}&=-\frac{1}{15}(3S^{2}-3S-14BS+12B^{2}+4B)/(S^{2}-3BS+2B^{2})
        \end{aligned}
    \]
    同样地可以求得参与者 2 对债券、股票的需求量为
    \[
        \begin{aligned}
            \theta_{21}&=-\frac{1}{10}(5S^{2}+5S-6BS-6B)/(S^{2}-3BS+2B^{2})\\
            \theta_{22}&=-\frac{1}{10}(3S^{2}-3S-14BS+12B^{2}+4B)/(S^{2}-3BS+2B^{2})
        \end{aligned}
    \]
    市场出清条件为
    \[
        \begin{aligned}
            \theta_{11}+\theta_{21}&=0\\
            \theta_{12}+\theta_{22}&=0
        \end{aligned}
    \]
    因而可以求得债券、股票的价格分别为
    \[
        B=\frac{5}{6},\quad S=1
    \]
    均衡利率为 $r_{F}=1/B-1=20\%$；可以验证 $\dfrac{u_{1,1}'(\tilde{c}_{1,1})}{u_{1,0}'(\tilde{c}_{1,0})}=\dfrac{u_{2,1}'(\tilde{c}_{2,1})}{u_{2,0}'(\tilde{c}_{2,0})}$，即边际效用之比在各状态下均相等，因此均衡配置是 Pareto 最优的。
    \item 对于参与者 1 来说，最优化的一阶条件是
    \[
        \begin{aligned}
            \frac{-B}{0.5-\theta_{11}B-\theta_{12}S}+\frac13\Bigl[\frac{1}{1+\theta_{11}+\theta_{12}}+\frac{1}{\theta_{11}+\theta_{12}}+\frac{1}{1+\theta_{11}+2\theta_{12}}\Bigr]&=0\\
            \frac{-S}{0.5-\theta_{11}B-\theta_{12}S}+\frac13\Bigl[\frac{1}{1+\theta_{11}+\theta_{12}}+\frac{1}{\theta_{11}+\theta_{12}}+\frac{2}{1+\theta_{11}+2\theta_{12}}\Bigr]&=0
        \end{aligned}
    \]
    对于参与者 2 来说，最优化的一阶条件是
    \[
        \begin{aligned}
            \frac{-B}{0.5-\theta_{21}B-\theta_{22}S}+\frac13\Bigl[\frac{1}{\theta_{21}+\theta_{22}}+\frac{1}{1+\theta_{21}+\theta_{22}}+\frac{1}{1+\theta_{21}+2\theta_{22}}\Bigr]&=0\\
            \frac{-S}{0.5-\theta_{21}B-\theta_{22}S}+\frac13\Bigl[\frac{1}{\theta_{21}+\theta_{22}}+\frac{1}{1+\theta_{21}+\theta_{22}}+\frac{2}{1+\theta_{21}+2\theta_{22}}\Bigr]&=0
        \end{aligned}
    \]
    市场出清条件为
    \[
        \begin{aligned}
            \theta_{11}+\theta_{21}&=0\\
            \theta_{12}+\theta_{22}&=0
        \end{aligned}
    \]
    配置是 Pareto 最优的一个条件是边际效用之比在各状态之间相等，考虑
    \[
        \begin{aligned}
            \frac{\theta_{11}+\theta_{12}}{1+\theta_{11}+\theta_{12}}\Big/\frac{1+\theta_{21}+\theta_{22}}{\theta_{21}+\theta_{22}}
            &=\frac{\theta_{11}+\theta_{12}}{1+\theta_{11}+\theta_{12}}\Big/\frac{1-\theta_{11}-\theta_{12}}{-\theta_{11}-\theta_{12}}\\
            &=1\big/\bigl(1-\tfrac{1}{(\theta_{11}+\theta_{12})^{2}}\bigr)\ne 1
        \end{aligned}
    \]
    因而，均衡配置不是 Pareto 最优的。
\end{enumerate}
\end{solution}

\begin{problem}
\problemno{11.2*} 经济中有 $K$ 个具有 CARA 偏好的参与者。经济中只有两个企业，它们的资产在 1 期的支付具有联合正态分布。相关的参数如下：
\begin{center}
\begin{tabular}{lccc}
    \toprule
     & 均值 & 标准差 & 相关系数 \\
    \midrule
    企业 1 & 10 & 8 & \multirow{2}{*}{0.6} \\
    企业 2 & 12 & 11 & \\
    \bottomrule
\end{tabular}
\end{center}
利率为 $8\%$。推导两个企业 0 期市场价值之间的关系。
\end{problem}
\begin{solution}
考虑参与者 $k$，假设他的效用函数为 $-e^{-a_{k}\tilde{w}}$，那么他的优化问题是
\[
    \begin{aligned}
        &\max_{\theta_{k,1},\,\theta_{k,2}}\ \mathrm{E}\bigl[-e^{-a_{k}\tilde{w}_{k}}\bigr]\\
        &\quad\text{s.t.}\quad \tilde{w}_{k}=(w_{k}-P_{1}\theta_{k,1}-P_{2}\theta_{k,2})(1+r_{F})+\tilde{r}_{1}\theta_{k,1}+\tilde{r}_{2}\theta_{k,2}
    \end{aligned}
\]
其中 $P_{1}$，$P_{2}$ 分别是企业 1、2 的价格。支付联合正态则 $\tilde{w}_{k}$ 正态，与练习 9.7 相同：最大化 $E[\tilde{w}_{k}]-\dfrac{a_{k}}{2}\operatorname{Var}(\tilde{w}_{k})$。记超额支付 $\mu_{i}=\bar{r}_{i}-P_{i}(1+r_{F})$，$\mu=\begin{bmatrix}\mu_{1}\\\mu_{2}\end{bmatrix}$，协方差阵 $\Sigma=\begin{bmatrix}\sigma_{1}^{2}&\rho\sigma_{1}\sigma_{2}\\\rho\sigma_{1}\sigma_{2}&\sigma_{2}^{2}\end{bmatrix}$。一阶条件 $\mu=a_{k}\Sigma\theta_{k}$ 给出 $\theta_{k}=\dfrac{1}{a_{k}}\Sigma^{-1}\mu$。展开
\[
    \Sigma^{-1}=\frac{1}{1-\rho^{2}}\begin{bmatrix}1/\sigma_{1}^{2}&-\rho/(\sigma_{1}\sigma_{2})\\-\rho/(\sigma_{1}\sigma_{2})&1/\sigma_{2}^{2}\end{bmatrix}
\]
即
\[
    \begin{aligned}
        \theta_{k,1}&=\frac{1}{a_{k}(1-\rho^{2})}\Bigl[\frac{\bar{r}_{1}-P_{1}(1+r_{F})}{\sigma_{1}^{2}}-\frac{\rho}{\sigma_{1}\sigma_{2}}\bigl(\bar{r}_{2}-P_{2}(1+r_{F})\bigr)\Bigr]\\
        \theta_{k,2}&=\frac{1}{a_{k}(1-\rho^{2})}\Bigl[\frac{\bar{r}_{2}-P_{2}(1+r_{F})}{\sigma_{2}^{2}}-\frac{\rho}{\sigma_{1}\sigma_{2}}\bigl(\bar{r}_{1}-P_{1}(1+r_{F})\bigr)\Bigr]
    \end{aligned}
\]
在均衡条件下，企业资产的总需求都是 1，因而我们有
\[
    \frac{\bar{r}_{1}-P_{1}(1+r_{F})}{\sigma_{1}^{2}}-\frac{\rho}{\sigma_{1}\sigma_{2}}\bigl(\bar{r}_{2}-P_{2}(1+r_{F})\bigr)
    =\frac{\bar{r}_{2}-P_{2}(1+r_{F})}{\sigma_{2}^{2}}-\frac{\rho}{\sigma_{1}\sigma_{2}}\bigl(\bar{r}_{1}-P_{1}(1+r_{F})\bigr)
\]
即
\[
    \frac{\bar{r}_{1}-P_{1}(1+r_{F})}{\sigma_{1}^{2}+\rho\sigma_{1}\sigma_{2}}=\frac{\bar{r}_{2}-P_{2}(1+r_{F})}{\sigma_{2}^{2}+\rho\sigma_{1}\sigma_{2}}
\]
代入数值即得 $10-1.08P_{1}=0.77(12-1.08P_{2})$。
\end{solution}

\begin{problem}
\problemno{11.3**} 考虑一个可以以 0 成本生产煤的企业。
\begin{itemize}
    \item 煤的未来价格 $\tilde{p}$ 在当前看来是不确定的。令 $\mathrm{E}[\tilde{p}]=\bar{p}$。
    \item 销售额取决未来价格：$S(\tilde{p})$。
    \item 有一个煤的远期市场，远期价格为 $f$。企业可以用远期合约来对冲未来价格中的风险。比如，如果企业售出 $h$ 份远期合约，远期头寸的未来支付为 $h(f-\tilde{p})$。如果 $\tilde{p}$ 的实现值很低，企业从销售中得到的利润很低但是从远期头寸得到利润则很高。
    \item 公司经理是风险厌恶的，他的期望效用为
    \[
        \mathrm{E}[u(\tilde{\pi})],\qquad \tilde{\pi}=\tilde{p}S(\tilde{p})+h(f-\tilde{p})
    \]
    （$u'>0$，$u''<0$）。
    \item 市场上还存在投机者，他们是作为企业所售远期合约的买方而出现的。他们也具有风险厌恶度的偏好：$v(\cdot)$（$v'>0$，$v''<0$）。他们的效用函数为：
    \[
        \mathrm{E}\bigl[v\bigl(h(\tilde{p}-f)\bigr)\bigr]
    \]
    其中 $h$ 是他们买（卖）的远期合约数。
\end{itemize}
\begin{enumerate}[label=(\alph*)]
    \item 证明当且仅当 $f<\mathrm{E}[\tilde{p}]$（$f>\mathrm{E}[\tilde{p}]$）时投机者才会买入（卖出）远期合约。
    \item 假设 $f=\mathrm{E}[\tilde{p}]$。在什么时候企业会使用远期合约来对冲风险？它是买入还是卖出远期合约？[提示：考虑销售的弹性，$-(d\ln S)/(d\ln P)$。]
    \item 假设 $u(x)=v(x)=-e^{-ax}$，$S(p)=s_{0}+k/p$ 且 $\tilde{p}$ 具有如下分布：
    \[
        \text{概率密度函数}=e^{-(p-d)},\quad p\in[d,+\infty).
    \]
    假设共有 $N$ 个投机者。在远期市场上，企业和投机者都表现得像是在竞争市场上一样。均衡远期价格是多少？远期溢价/折扣：$f-\mathrm{E}[\tilde{p}]$ 是多少？解释所得到的结论。
\end{enumerate}
\end{problem}
\begin{solution}
\begin{enumerate}[label=(\alph*)]
    \item 投机者的最优化问题的一阶条件是
    \[
        \mathrm{E}\bigl[v'\bigl(h(\tilde{p}-f)\bigr)(\tilde{p}-f)\bigr]=0
    \]
    因而，
    \[
        \mathrm{Cov}\bigl(v'\bigl(h(\tilde{p}-f)\bigr),(\tilde{p}-f)\bigr)=-\mathrm{E}\bigl[v'\bigl(h(\tilde{p}-f)\bigr)\bigr]\mathrm{E}[\tilde{p}-f]
    \]
    而 $v'>0$，$v''<0$，因此
    \[
        h>0\,(<0)\Leftrightarrow \mathrm{E}[\tilde{p}]>(<)\,f
    \]
    即当且仅当 $f<\mathrm{E}[\tilde{p}]$（$f>\mathrm{E}[\tilde{p}]$）时投机者才会买入（卖出）远期合约。
    \item 厂商最优化问题的一阶条件是
    \[
        \mathrm{E}\bigl[u'\bigl(\tilde{p}S(\tilde{p})+h(\tilde{p}-f)\bigr)(\tilde{p}-f)\bigr]=0
    \]
    又由假设 $\mathrm{E}[\tilde{p}]=f$ 有
    \[
        \mathrm{Cov}\bigl(u'\bigl(\tilde{p}S(\tilde{p})+h(\tilde{p}-f)\bigr),(\tilde{p}-f)\bigr)=0
    \]
    令 $e(p)=-\dfrac{d\log S}{d\log p}$，那么 $d\bigl(pS(p)+h(f-p)\bigr)/dp=S(p)\bigl(1-e(p)\bigr)-h$；因此，如果 $e(p)<(>)1$，那么 $h>(<)0$。
    \item 给定 $f$，$h$，厂商的期望效用是
    \[
        \begin{aligned}
            \mathrm{E}[u(\tilde{\pi})]
            &=\mathrm{E}\bigl[u\bigl(\tilde{p}S(\tilde{p})+h(f-\tilde{p})\bigr)\bigr]\\
            &=\frac{e^{-a[S_{0}d+k+h(f-d)]}}{a(S_{0}-h)+1}
        \end{aligned}
    \]
    求导可得厂商的远期合约售出量为：$h_{1}^{*}=S_{0}+\dfrac{1}{a}-\dfrac{1}{a(f-d)}$。投机者的期望效用是
    \[
        \mathrm{E}\bigl[v\bigl(h(\tilde{p}-f)\bigr)\bigr]=\frac{e^{ah(f-d)}}{ah+1}
    \]
    求导可得投机者的远期合约需求量为：$h_{2}^{*}=\dfrac{1}{a(f-d)}-\dfrac{1}{a}$。市场出清时
    \[
        S_{0}+\frac{1}{a}-\frac{1}{a(f-d)}=N\Bigl(\frac{1}{a(f-d)}-\frac{1}{a}\Bigr)
    \]
    求得均衡远期价格为 $f^{*}=d+1/\bigl(1-\dfrac{aS_{0}}{N+1}\bigr)$，期权溢价为 $f-\mathrm{E}[\tilde{p}]=1/\bigl(1-\dfrac{aS_{0}}{N+1}\bigr)-1$。当风险厌恶系数增大时，溢价变大（或折扣变小），这是很显然的。当投机者增多时，溢价（折扣）的绝对值变小，即使得远期合约的价格向期望值靠拢，这是投机市场竞争加剧时：如果 $f<\mathrm{E}[\tilde{p}]$，此时投机者买入远期，厂商可以抬高价格；如果 $f>\mathrm{E}[\tilde{p}]$，此时投机者卖出远期，厂商可以压低价格。
\end{enumerate}
\end{solution}

\begin{problem}
\problemno{11.4*} 证券市场只含有一只风险证券（股票）。一单位股票的支付为 $\tilde{x}$，它服从均值为 $\bar{x}$、波动率为 $\sigma$ 的正态分布。市场上有参与者 1 和 2。参与者 1 的禀赋是 $x_{0}$ 单位的 0 期消费而参与者 2 的禀赋是一单位股票。两个参与者都具有如下形式的期望效用：
\[
    -e^{-\alpha c_{0}}-\rho\mathrm{E}\bigl[-e^{-\alpha\tilde{c}_{1}}\bigr].
\]
\begin{enumerate}[label=(\alph*)]
    \item 证明证券市场是不完全的。
    \item 写出均衡条件。
    \item 均衡配置是 Pareto 最优的吗？证明结论。
\end{enumerate}
\end{problem}
\begin{solution}
\begin{enumerate}[label=(\alph*)]
    \item 因为有无穷多个状态（支付是正态分布的），而只有一只证券，因此证券市场是不完全的；
    \item 均衡条件是存在股票价格 $S$ 使得
    \[
        \begin{aligned}
            \theta_{1}^{*}&=\operatorname{Argmax}_{\theta}\bigl[-e^{-\alpha(x_{0}-\theta S)}+\rho\mathrm{E}\bigl[-e^{-\alpha\theta\tilde{x}}\bigr]\bigr]\\
            \theta_{2}^{*}&=\operatorname{Argmax}_{\theta}\bigl[-e^{-\alpha(S-\theta S)}+\rho\mathrm{E}\bigl[-e^{-\alpha\theta\tilde{x}}\bigr]\bigr]\\
            \theta_{1}^{*}+\theta_{2}^{*}&=1
        \end{aligned}
    \]
    \item 参与者 1、2 的最优化问题的一阶条件分别是
    \[
        \begin{aligned}
            e^{-\alpha(x_{0}-\theta_{1}^{*}S)}S&=\rho e^{-\alpha\theta_{1}^{*}\bar{x}+\frac12\alpha^{2}(\theta_{1}^{*})^{2}\sigma^{2}}(\bar{x}-\alpha\theta_{1}^{*}\sigma^{2})\\
            e^{-\alpha(S-\theta_{2}^{*}S)}S&=\rho e^{-\alpha\theta_{2}^{*}\bar{x}+\frac12\alpha^{2}(\theta_{2}^{*})^{2}\sigma^{2}}(\bar{x}-\alpha\theta_{2}^{*}\sigma^{2})
        \end{aligned}
    \]
    如果均衡配置是 Pareto 最优的，那么两个参与者的边际效用之比应该相等，即
    \[
        \frac{-\rho e^{-\alpha\theta_{1}^{*}\tilde{x}}(-a)}{-e^{-\alpha(x_{0}-\theta_{1}^{*}S)}(-a)}=\frac{-\rho e^{-\alpha\theta_{2}^{*}\tilde{x}}(-a)}{-e^{-\alpha(S-\theta_{2}^{*}S)}(-a)}
    \]
    又因为 $\theta_{1}^{*}+\theta_{2}^{*}=1$，因而 $\theta_{1}^{*}=\theta_{2}^{*}=\tfrac12$，$S=x_{0}$。代入一阶条件有
    \[
        e^{-\frac12\alpha\tilde{x}}=\rho e^{-\frac12\alpha\tilde{x}+\frac18\alpha^{2}\sigma^{2}}(\tilde{x}-\tfrac12\alpha\sigma^{2})
    \]
    上述等式不一定成立，因而均衡配置不一定是 Pareto 最优的（实际上，只有当上述等式成立时才是 Pareto 最优的）。
\end{enumerate}
\end{solution}

\begin{problem}
\problemno{11.5**} 考虑练习 11.4 中的经济。现在引入一只无风险证券（债券），它在 1 期支付 1 单位消费品。参与者在债券上的禀赋都为 0。
\begin{enumerate}[label=(\alph*)]
    \item 证券市场完全吗？解释原因。
    \item 写出这个经济的均衡条件。
    \item 求解均衡。
    \item 比较存在无风险证券时的均衡配置和在练习 11.4 中不存在无风险证券时得到的均衡配置。它们一样吗？为什么一样或不一样？
    \item 均衡配置是 Pareto 最优的吗？证明并解释所得到的结论。
\end{enumerate}
\end{problem}
\begin{solution}
\begin{enumerate}[label=(\alph*)]
    \item 因为状态数是无穷的，但只有两只证券，因而市场不完全；
    \item 这个经济的均衡条件是存在债券价格 $B$ 和股票价格 $S$，使得
    \[
        \begin{aligned}
            \begin{bmatrix}\theta_{1,1}^{*} \\ \theta_{1,2}^{*}\end{bmatrix}&=\operatorname{Argmax}_{\{\theta_{1,1},\theta_{1,2}\}}\bigl[-e^{-\alpha(x_{0}-\theta_{1,1}B-\theta_{1,2}S)}+\rho\mathrm{E}\bigl[-e^{-\alpha(\theta_{1,1}+\theta_{1,2}\tilde{x})}\bigr]\bigr]\\
            \begin{bmatrix}\theta_{2,1}^{*} \\ \theta_{2,2}^{*}\end{bmatrix}&=\operatorname{Argmax}_{\{\theta_{2,1},\theta_{2,2}\}}\bigl[-e^{-\alpha(S-\theta_{2,1}B-\theta_{2,2}S)}+\rho\mathrm{E}\bigl[-e^{-\alpha(\theta_{2,1}+\theta_{2,2}\tilde{x})}\bigr]\bigr]\\
            \theta_{1,1}^{*}+\theta_{2,1}^{*}=0;&\quad\theta_{1,2}^{*}+\theta_{2,2}^{*}=1
        \end{aligned}
    \]
    \item 对于参与者 1 来说，最优化的一阶条件是
    \begin{align}
            -e^{-\alpha(x_{0}-\theta_{1,1}^{*}B-\theta_{1,2}^{*}S)}(\alpha B)
            &=\rho e^{-\alpha(\theta_{1,1}^{*}+\theta_{1,2}^{*}\tilde{x})+\frac12\alpha^{2}(\theta_{1,2}^{*})^{2}\sigma^{2}}(-\alpha)
            \tag{11.1}\\
            -e^{-\alpha(x_{0}-\theta_{1,1}^{*}B-\theta_{1,2}^{*}S)}(\alpha S)
            &=\rho e^{-\alpha(\theta_{1,1}^{*}+\theta_{1,2}^{*}\tilde{x})+\frac12\alpha^{2}(\theta_{1,2}^{*})^{2}\sigma^{2}}(-\alpha\bar{x}+\alpha^{2}\theta_{1,2}^{*}\sigma^{2})
            \tag{11.2}
    \end{align}
    参与者 2 的最优化的一阶条件是
    \begin{align}
            -e^{-\alpha(S-\theta_{2,1}^{*}B-\theta_{2,2}^{*}S)}(\alpha B)
            &=\rho e^{-\alpha(\theta_{2,1}^{*}+\theta_{2,2}^{*}\tilde{x})+\frac12\alpha^{2}(\theta_{2,2}^{*})^{2}\sigma^{2}}(-\alpha)
            \tag{11.3}\\
            -e^{-\alpha(S-\theta_{2,1}^{*}B-\theta_{2,2}^{*}S)}(\alpha S)
            &=\rho e^{-\alpha(\theta_{2,1}^{*}+\theta_{2,2}^{*}\tilde{x})+\frac12\alpha^{2}(\theta_{2,2}^{*})^{2}\sigma^{2}}(-\alpha\bar{x}+\alpha^{2}\theta_{2,2}^{*}\sigma^{2})
            \tag{11.4}
    \end{align}
    (11.2) 式/(11.1) 式，(11.4) 式/(11.3) 式可得
    \[
        \begin{aligned}
            \frac{S}{B}&=\bar{x}-\alpha\theta_{1,2}^{*}\sigma^{2}\\
            \frac{S}{B}&=\bar{x}-\alpha\theta_{2,2}^{*}\sigma^{2}
        \end{aligned}
    \]
    因此 $\theta_{1,2}^{*}=\theta_{2,2}^{*}=\tfrac12$。\\
    (11.3) 式/(11.1) 式可得
    \[
        e^{\alpha(x_{0}-S)+\alpha B(\theta_{2,1}^{*}-\theta_{1,1}^{*})}=e^{\alpha(\theta_{1,1}^{*}-\theta_{2,1}^{*})}
    \]
    因此，$\theta_{1,1}^{*}=-\theta_{2,1}^{*}=\dfrac{x_{0}}{2(1+B)}\widehat{S}$。利用 $\dfrac{\widehat{S}}{B}=\bar{x}-\tfrac12\alpha\sigma^{2}$，$\theta_{1,1}^{*}=-\theta_{2,1}^{*}=\dfrac{x_{0}-\widehat{S}}{2(1+B)}$ 以及 (11.1) 式就可求得债券价格 $B$ 以及股票价格 $S$。
    \item 因为存在了无风险债券，投资者的投资机会增加，使得投资者的选择不一样了；另外，这也改变了股票的投资价值，从而改变了参与者 2 的禀赋的相对值，这也使得均衡配置不一样。最明显的是在第 4 题中两个参与者的股票头寸不相等，而在本题当中则相等；
    \item 考虑两人的边际效用之间的比值
    \[
        \begin{aligned}
            &\frac{-\rho(-a)e^{-\alpha(\theta_{1,1}^{*}+\theta_{1,2}^{*}\tilde{x})}}{-(-a)e^{-\alpha(x_{0}-\theta_{1,1}^{*}B-\theta_{1,2}^{*}S)}}
            \Big/
            \frac{-\rho(-a)e^{-\alpha(\theta_{2,1}^{*}+\theta_{2,2}^{*}\tilde{x})}}{-(-a)e^{-\alpha(S-\theta_{2,1}^{*}B-\theta_{2,2}^{*}S)}}\\
            &\qquad=e^{-\alpha\bigl(x_{0}-S-(1+B)(\theta_{1,1}^{*}-\theta_{2,1}^{*})\bigr)}=1
        \end{aligned}
    \]
    因而均衡配置是 Pareto 最优的。
\end{enumerate}
\end{solution}

\begin{problem}
\problemno{11.6**} 在练习 11.5 描述的经济中，我们引入一只新的证券，一份以股票为标的资产、执行价格为 $\bar{x}$ 的期权。
\begin{enumerate}[label=(\alph*)]
    \item 现在，由股票、债券和期权所构成的证券市场是完全的吗？解释原因。
    \item 假设均衡配置和练习 11.5 中的一样。用参与者 1 的 Euler 方程为期权定价。同样地，用参与者 2 的 Euler 方程为期权定价。
    \item 使用前面得到的结论，分析练习 11.5 中得到的均衡配置是否会因为引入期权而改变。
    \item 计算期权的均衡价格。
\end{enumerate}
\end{problem}
\begin{solution}
\begin{enumerate}[label=(\alph*)]
    \item 由于状态数目大于证券数目，因而市场是不完全的；
    \item 由参与者 1 的 Euler 方程，期权价格为
    \[
        c_{1}=\mathrm{E}\Bigl[\frac{u_{1,1}'(\tilde{c}_{1,1})}{u_{1,0}'(c_{1,0})}[\tilde{x}-\bar{x}]\Bigr]
        =\mathrm{E}\Bigl[\frac{\rho e^{-\alpha(\theta_{1,1}^{*}+\theta_{1,2}^{*}\tilde{x})}}{e^{-\alpha(x_{0}-\theta_{1,1}^{*}B-\theta_{1,2}^{*}S)}}[\tilde{x}-\bar{x}]\Bigr]
    \]
    由参与者 2 的 Euler 方程，期权价格为
    \[
        c_{2}=\mathrm{E}\Bigl[\frac{u_{2,1}'(\tilde{c}_{2,1})}{u_{2,0}'(c_{2,0})}[\tilde{x}-\bar{x}]\Bigr]
        =\mathrm{E}\Bigl[\frac{\rho e^{-\alpha(\theta_{2,1}^{*}+\theta_{2,2}^{*}\tilde{x})}}{e^{-\alpha(S-\theta_{2,1}^{*}B-\theta_{2,2}^{*}S)}}[\tilde{x}-\bar{x}]\Bigr]
    \]
    \item 因为第五题中的均衡是 Pareto 最优的，因而边际效用之比相等，$c_{1}=c_{2}$，即参与者对期权的定价是一致的，因而他们不会交易期权，引入期权不会改变第 5 题中的均衡配置。
    \item 由前面的分析可得 $c=c_{1}=c_{2}$。
\end{enumerate}
\end{solution}



% !TeX root = ../../金融经济学_王江_讲义及习题.tex
% 知识点：references/金融经济学/金融经济学-王江_261013.pdf，印刷页 179–192（PDF 页 194–207）
% 习题：references/金融经济学/史上最全最完整的金融经济学答案(王江版).pdf，页 62–67（印刷页 61–66）

\lecture{在均值-方差偏好下的投资组合选择}

\ifdefined\WJHandoutKnowledge

上一章在一般偏好和不完全证券结构下，均衡选择依赖完整的支付分布。本章增加条件，使参与者只按收益分布的\textbf{均值}与\textbf{方差}选择组合；这把组合问题化为协方差矩阵上的二次优化。

\begin{tcolorbox}[knowledge,title=12.1 均值--方差偏好]
记终值财富为 $\widetilde w$，均值 $\bar w=\mathrm E[\widetilde w]$，标准差 $\sigma_w=\sqrt{\operatorname{Var}(\widetilde w)}$。若参与者的排序可由
\[
 \mathrm E[u(\widetilde w)]=v(\bar w,\sigma_w)
 \tag{12.1}
\]
完全表示，则称具有\textbf{均值--方差偏好}。这需要对效用函数或可选收益分布施加条件，而非一般期望效用的自动性质。

\textbf{定理 12.1}\quad 若 $u(w)=w-\tfrac12 a w^2$，则
\[
 \mathrm E[u(\widetilde w)]
 =\bar w-\tfrac12a(\bar w^2+\sigma_w^2).
\]
因此二次效用具有均值--方差表示。但 $u'(w)=1-aw$ 只在 $w<1/a$ 时为正；若财富越过该区间，二次效用会违反不满足性，使用时须限制财富域。

\textbf{定理 12.2}\quad 若可交易证券的联合支付为正态分布，则每个线性组合仍为正态分布，其分布完全由均值和方差决定。因此，只要期望效用存在，参与者对这些组合的排序具有均值--方差表示；不要求效用本身二次型。

\textbf{定理 12.3}\quad 对联合正态支付及递增、严格凹的效用，$v$ 对均值严格递增、对标准差严格递减；$v$ 对均值凹、对标准差凹，且具有相应的联合凹性。证明把财富写成 $\bar w+\sigma_w Z$，$Z\sim N(0,1)$：
\[
 v_1=\mathrm E[u'(\bar w+\sigma_wZ)]>0,
 \quad v_2=\mathrm E[u'(\bar w+\sigma_wZ)Z]<0.
\]
后一个符号可由 $Z$ 对称性与 $u'$ 递减得到。二阶导数的凹性由 $u''<0$ 及 Cauchy--Schwarz 不等式得到。正态分布有负财富尾部；若效用只定义在正财富域，必须另作定义域处理，不能直接套用。
\end{tcolorbox}

\begin{tcolorbox}[knowledge,title=12.2 无风险证券缺席时的前沿]
设 $N$ 只风险证券的净收益率向量为 $\widetilde r$，均值 $\bar r$，正定协方差矩阵 $\Sigma$。权重 $z$ 满足 $\mathbf1^{\mathsf T}z=1$，组合均值和方差为
\[
 \bar r_p=z^{\mathsf T}\bar r,\qquad
 \sigma_p^2=z^{\mathsf T}\Sigma z.
\]
在给定目标收益 $\bar r_p$ 下求最小方差，得到\textbf{均值--方差前沿}（MVF）组合：
\[
 \min_z\ \tfrac12 z^{\mathsf T}\Sigma z,
 \quad\text{s.t. }z^{\mathsf T}\bar r=\bar r_p,
 \quad z^{\mathsf T}\mathbf1=1.
 \tag{12.2}
\]
令 $a=\bar r^{\mathsf T}\Sigma^{-1}\bar r$、$b=\bar r^{\mathsf T}\Sigma^{-1}\mathbf1$、$c=\mathbf1^{\mathsf T}\Sigma^{-1}\mathbf1$、$d=ac-b^2>0$，一阶条件解为
\[
 z_p=\frac1d\Sigma^{-1}\bigl[(c\bar r_p-b)\bar r+(a-b\bar r_p)\mathbf1\bigr].
 \tag{12.3}
\]
$d>0$ 要求各风险证券的预期收益不全相同；若 $d=0$，目标收益约束退化，应单独处理。正定的 $\Sigma$ 保证解唯一。

把解写成 $z_p=z_0+\bar r_p(z_1-z_0)$，可见\textbf{定理 12.5}：任意 MVF 组合可由 $z_0,z_1$ 两个 MVF 组合形成；\textbf{推论 12.1--12.2}：任意两个不同的 MVF 组合足以生成整条前沿，它们的线性组合仍在前沿。这就是无风险证券缺席时的\textbf{两基金分离}（定理 12.6--12.7）。组合权重允许卖空；若另有权重非负约束，前沿会分段，公式须改写。

前沿在 $(\sigma,\bar r)$ 平面是一条双曲线的边界。顶点是\textbf{最小方差组合}（MVP），其均值为 $b/c$，权重为 $\Sigma^{-1}\mathbf1/c$。前沿的上半支在均值--方差偏好下有效；下半支被同方差、更高均值的组合支配。

\textbf{定理 12.8}\quad 对任意组合 $p$，$\operatorname{Cov}(\widetilde r_p,\widetilde r_{\rm MVP})=\sigma_{\rm MVP}^2$。证明：MVP 的一阶条件给出 $\Sigma z_{\rm MVP}=\mathbf1/c$，故 $z_p^{\mathsf T}\Sigma z_{\rm MVP}=1/c$。\textbf{定理 12.9}进一步指出，对非 MVP 的前沿组合 $p$，存在一个均值--方差前沿上的零协方差组合 $zcp$，用作后面的定价基准。

\textbf{定理 12.10}\quad 若 $p$ 是非 MVP 的前沿组合，则任意组合 $q$ 满足
\[
 \bar r_q-\bar r_{zcp}
 =\beta_{q,p}(\bar r_p-\bar r_{zcp}),\qquad
 \beta_{q,p}=\frac{\operatorname{Cov}(\widetilde r_q,\widetilde r_p)}{\sigma_p^2}.
 \tag{12.6}
\]
证明将 $\Sigma z_p=\lambda_r\bar r+\lambda_1\mathbf1$ 两边左乘任意权重 $z_q^{\mathsf T}$，再分别取 $q=p,zcp$ 消去乘子。这里是\textbf{零贝塔定价式}，其基准收益不一定等于无风险利率。
\end{tcolorbox}

\begin{tcolorbox}[knowledge,title=12.4 存在无风险证券时]
令无风险净利率为 $r_F$，风险证券的预期超额收益向量 $\eta=\bar r-r_F\mathbf1$。风险证券权重 $z$ 不必合计为 1，余下资金投向无风险证券。若 $\eta^{\mathsf T}\Sigma^{-1}\eta>0$，给定组合预期超额收益 $\bar\eta_p$，最小方差解是
\[
 z_p=\frac{\bar\eta_p}{\eta^{\mathsf T}\Sigma^{-1}\eta}\Sigma^{-1}\eta.
 \tag{12.7}
\]
将 $\bar\eta_p$ 取为使风险证券权重和为 1 的值，得到\textbf{切点组合}
\[
 z_T=\frac{\Sigma^{-1}\eta}{\mathbf1^{\mathsf T}\Sigma^{-1}\eta},
\]
前提是分母非零。所有此类前沿组合是无风险证券与 $T$ 的线性组合（\textbf{定理 12.11}）；这就是一基金加无风险证券的分离。若 $r_F$ 位于风险证券前沿的适当位置，切点 $T$ 在有效前沿上（\textbf{定理 12.12}）。

无风险证券与 $T$ 在 $(\sigma,\bar r)$ 平面形成\textbf{资本市场线}。风险组合 $q$ 的 Sharpe 比为 $(\bar r_q-r_F)/\sigma_q$；在允许自由借贷和卖空的上述条件下，切点组合具有最高 Sharpe 比。

\textbf{定理 12.13}\quad 若 $p$ 是有无风险证券时的均值--方差有效组合，任意组合 $q$ 满足
\[
 \bar r_q-r_F=\beta_{q,p}(\bar r_p-r_F),\qquad
 \beta_{q,p}=\frac{\operatorname{Cov}(\widetilde r_q,\widetilde r_p)}{\sigma_p^2}.
 \tag{12.9}
\]
特别地，取 $p=T$ 得切点组合的贝塔定价式（12.10）。这里的 $p$ 是\emph{有效组合}；下一章才在均衡条件下识别它与\emph{市场组合}的关系。
\end{tcolorbox}

\begin{tcolorbox}[summary,title=本章小结]
二次效用或联合正态支付等条件可使组合排序只依赖均值与方差。无风险证券缺席时，前沿由两个风险基金生成，并有零贝塔定价关系；存在无风险证券时，前沿由切点风险组合与无风险证券生成，风险溢价可由相对于有效组合的贝塔表示。上述分析以协方差矩阵正定、相关组合可交易及适当的定义域为条件。
\end{tcolorbox}

\fi

% !TeX root = ../../金融经济学_王江_讲义及习题.tex
% 12_均值方差偏好 课后习题（由章节文件拆分）

\begin{problem}
\problemno{12.1} 令交易证券的价格向量为 $p$、总收益率向量为 $\tilde{x}=\begin{bmatrix}\tilde{x}_{1} \\ \dots \\ \tilde{x}_{N}\end{bmatrix}$。假设 $\tilde{x}=\iota+\tilde{r}$ 服从均值为 $\iota+\bar{r}$、协方差阵为 $\Sigma$ 的正态分布。对于初始财富为 $w$ 的参与者，他的 1 期财富是 $\tilde{w}=wz^{\top}\tilde{x}$，其中 $z$ 是他的权重向量（$\iota^{\top}z=1$）。假设参与者的绝对风险厌恶系数恒为 $a$。
\begin{enumerate}[label=(\alph*)]
    \item 计算他的间接效用函数
    \[
        v=\mathrm{E}[u(\tilde{w})]
    \]
    其中 $u$ 是他的直接效用函数。
    \item 证明间接效用函数具有均值--方差表示，即它可以表示成 $\bar{w}=\mathrm{E}[\tilde{w}]$ 和 $\sigma_{w}^{2}=\mathrm{Var}[\tilde{w}]$ 的函数，即 $v(\bar{w},\sigma_{w})$。
    \item 证明均值方差效用函数 $v(\bar{w},\sigma_{w})$ 满足如下条件：$\partial_{1}v>0$，$\partial_{2}v<0$ 且 $\partial_{11}^{2}v<0$，$\partial_{22}^{2}v<0$ 且 $v$ 是凹的。
\end{enumerate}
\end{problem}
\begin{solution}
不妨假设效用函数为 $u(w)=-e^{-aw}$，那么 $v=\mathrm{E}[u(\tilde{w})]=\mathrm{E}[-e^{-a\tilde{w}}]$。因为 $\tilde{w}$ 是正态分布的，因而
\[
    v=-e^{-a\mathrm{E}\tilde{w}+\frac12 a^{2}\mathrm{Var}\tilde{w}}=-e^{-a\bar{w}+\frac12 a^{2}\sigma_{w}^{2}}=v(\bar{w},\sigma_{w})
\]
即间接效用函数具有均值方差表示，对 $v$ 求导可得
\[
    \begin{aligned}
        \partial_{1}v&=e^{-a\bar{w}+\frac12 a^{2}\sigma_{w}^{2}}a>0\\
        \partial_{2}v&=-e^{-a\bar{w}+\frac12 a^{2}\sigma_{w}^{2}}a^{2}\sigma_{w}<0\\
        \partial_{11}v&=e^{-a\bar{w}+\frac12 a^{2}\sigma_{w}^{2}}(-a^{2})<0\\
        \partial_{22}v&=-e^{-a\bar{w}+\frac12 a^{2}\sigma_{w}^{2}}(a^{4}\sigma_{w}^{2}+a^{2})<0\\
        \partial_{12}v&=e^{-a\bar{w}+\frac12 a^{2}\sigma_{w}^{2}}a^{3}\sigma_{w},\qquad
        \partial_{11}v\,\partial_{22}v-(\partial_{12}v)^{2}=v^{2}(\bar{w},\sigma_{w})a^{2}>0
    \end{aligned}
\]
因而，$v$ 还是凹的。
\end{solution}

\begin{problem}
\problemno{12.2} 考虑证券 1 和 2，它们的收益率分别是 $\tilde{r}_{1}$ 和 $\tilde{r}_{2}$。它们的波动率分别是 $\sigma_{1}$ 和 $\sigma_{2}$，并且它们的相关系数是 $\rho$。
\begin{enumerate}[label=(\alph*)]
    \item 假设 $\rho\ne 1$。用波动率度量风险，证明分散化即用证券形成组合可以降低风险。（也就是说，总是存在波动率小于两只单个证券的波动率的组合。）
    \item 找出收益率具有最小波动率的组合。
    \item 证明当 $\rho=1$ 时，分散化不能降低风险。
    \item 假设 $\sigma_{1}=\sigma_{2}$。证明无论 $\rho$ 的值是多少，等权重组合的风险是最小的。
\end{enumerate}
\end{problem}
\begin{solution}
\begin{enumerate}[label=(\alph*)]
    \item 设组合权重为 $z=\begin{bmatrix}z_{1} \\ 1-z_{1}\end{bmatrix}$，那么 $\tilde{r}_{p}=z_{1}\tilde{r}_{1}+(1-z_{1})\tilde{r}_{2}$，
    \[
        \sigma_{p}^{2}=z_{1}^{2}\sigma_{1}^{2}+2z_{1}(1-z_{1})\rho\sigma_{1}\sigma_{2}+(1-z_{1})^{2}\sigma_{2}^{2}
    \]
    对 $z_{1}$ 求导，$\partial\sigma_{p}^{2}/\partial z_{1}=0$ 给出
    \[
        z_{1}\sigma_{1}^{2}+(1-2z_{1})\rho\sigma_{1}\sigma_{2}-(1-z_{1})\sigma_{2}^{2}=0
    \]
    因而全局最小方差权重为
    \[
        z_{1}^{*}=\frac{\sigma_{2}^{2}-\rho\sigma_{1}\sigma_{2}}{\sigma_{1}^{2}-2\rho\sigma_{1}\sigma_{2}+\sigma_{2}^{2}}
    \]
    代回并配方，
    \[
        \sigma_{p}^{2}=\sigma_{2}^{2}-\frac{(\sigma_{2}^{2}-\rho\sigma_{1}\sigma_{2})^{2}}{\sigma_{1}^{2}-2\rho\sigma_{1}\sigma_{2}+\sigma_{2}^{2}}=\sigma_{1}^{2}-\frac{(\sigma_{1}^{2}-\rho\sigma_{1}\sigma_{2})^{2}}{\sigma_{1}^{2}-2\rho\sigma_{1}\sigma_{2}+\sigma_{2}^{2}}
    \]
    当 $\rho\ne 1$ 时分子为正，故 $\sigma_{p}<\min(\sigma_{1},\sigma_{2})$，分散化可以降低风险；
    \item 即上一问中的 $z^{*}$；
    \item 当 $\rho=1$ 时，$\sigma_{p}^{2}=(\sigma_{1}z_{1}+\sigma_{2}(1-z_{1}))^{2}\in[\min(\sigma_{1}^{2},\sigma_{2}^{2}),\,\max(\sigma_{1}^{2},\sigma_{2}^{2})]$，$z_{1}\in[0,1]$，即分散化不能降低风险。
    \item 当 $\sigma_{1}=\sigma_{2}$ 时，$z_{1}=(\sigma_{2}^{2}-\rho\sigma_{1}\sigma_{2})/(\sigma_{1}^{2}-2\rho\sigma_{1}\sigma_{2}+\sigma_{2}^{2})=\tfrac12$，即无论 $\rho$ 多少，等权重组合的风险最小。
\end{enumerate}
\end{solution}

\begin{problem}
\problemno{12.3*} 假设市场上有 $N$ 只证券，它们收益率波动率均为 $\sigma$，而协方差为 $c=\mathrm{Cov}[\tilde{r}_{i},\tilde{r}_{j}]$，$\forall\, i,j=1,\dots,N$ 且 $i\ne j$。
\begin{enumerate}[label=(\alph*)]
    \item 证明由 $N$ 只证券组成的等权组合的收益率具有最小波动率。
    \item 证明当 $N\to\infty$ 时等权重组合的方差趋近于 $c$。这和单个证券的波动率 $\sigma$ 无关。
    \item 解释上面得到的结论。
\end{enumerate}
\end{problem}
\begin{solution}
\begin{enumerate}[label=(\alph*)]
    \item 假设组合权重为 $z=\begin{bmatrix}z_{1} \\ z_{2} \\ \dots \\ z_{N}\end{bmatrix}$，那么组合的方差为
    \[
        \sigma_{p}^{2}=\mathrm{Var}(z_{1}\tilde{r}_{1}+z_{2}\tilde{r}_{2}+\dots+z_{N}\tilde{r}_{N})=\sigma^{2}\sum_{i}z_{i}^{2}+c\sum_{i\ne j}z_{i}z_{j}
    \]
    利用 $\sum_{i}z_{i}=1$ 改写为 $\sigma_{p}^{2}=(\sigma^{2}-c)\sum_{i}z_{i}^{2}+c$。拉格朗日函数为
    \[
        \mathcal{L}=(\sigma^{2}-c)\sum_{i}z_{i}^{2}+\lambda\Bigl(1-\sum_{i}z_{i}\Bigr)
    \]
    一阶条件 $2(\sigma^{2}-c)z_{i}=\lambda$（$\sigma^{2}-c>0$）迫使 $z_{i}$ 全相等，因此
    \[
        z_{i}=\frac{1}{N},\quad\forall i
    \]
    \item 把 $z_{i}=\dfrac{1}{N}$ 代入组合方差可得 $\sigma_{p}^{2}=\dfrac{1}{N}\sigma^{2}+\bigl(1-\dfrac{1}{N}\bigr)c\to c$，$N\to\infty$；
    \item 当 $N\to\infty$ 时，证券间不相关的风险都被分散掉了，只剩下了其中相关的部分，因而最终的组合风险只与协方差项有关。
\end{enumerate}
\end{solution}

\begin{problem}
\problemno{12.4} 考虑风险证券 1 和 2，它们的收益率为 $\tilde{r}_{1}$ 和 $\tilde{r}_{2}$。它们的期望收益率和协方差阵分别为 $\bar{r}_{1}$，$\bar{r}_{2}$，以及 $\Sigma=\begin{bmatrix}\sigma_{1}^{2} & \sigma_{12} \\ \sigma_{12} & \sigma_{2}^{2}\end{bmatrix}$。
\begin{enumerate}[label=(\alph*)]
    \item 推导由这两只证券构成的均值--方差有效前沿。
    \item 证明这两只证券都在前沿上。
    \item 考虑如下三种特别形式的前沿：$\rho_{12}=-1,\,0,\,1$，其中 $\rho_{12}$ 是两只证券收益率之间的相关系数。在 $\bar{r}$--$\sigma$ 平面画出这三种情形下的组合前沿。
\end{enumerate}
\end{problem}
\begin{solution}
\begin{enumerate}[label=(\alph*)]
    \item 为简便起见，假设 $\bar{r}_{1}\ne\bar{r}_{2}$；给定组合的期望收益率 $\bar{r}_{p}$ 有 $z\bar{r}_{1}+(1-z)\bar{r}_{2}=\bar{r}_{p}$，因而 $z=\dfrac{\bar{r}_{p}-\bar{r}_{2}}{\bar{r}_{1}-\bar{r}_{2}}$，此时，组合的波动率为
    \[
        \begin{aligned}
            \sigma_{p}^{2}&=\frac{1}{(\bar{r}_{1}-\bar{r}_{2})^{2}}\bigl[(\sigma_{1}^{2}+\sigma_{2}^{2}-2\sigma_{12})\bar{r}_{p}^{2}-(2\sigma_{1}^{2}\bar{r}_{2}+2\sigma_{2}^{2}\bar{r}_{1}-2(\bar{r}_{1}+\bar{r}_{2})\sigma_{12})\bar{r}_{p}\\
            &\qquad+(\sigma_{1}^{2}\bar{r}_{2}^{2}+\sigma_{2}^{2}\bar{r}_{1}^{2}-2\bar{r}_{1}\bar{r}_{2}\sigma_{12})\bigr]
        \end{aligned}
    \]
    $(\bar{r}_{p},\sigma_{p})$ 即为均值--方差前沿组合；
    \item 当 $\bar{r}_{p}=\bar{r}_{1}$ 时，前沿组合即为证券 1；当 $\bar{r}_{p}=\bar{r}_{2}$ 时，前沿组合即为证券 2；
    \item 当 $\rho_{12}=-1$ 时
    \[
        \sigma_{p}=\left|\frac{(\sigma_{1}+\sigma_{2})\bar{r}_{p}-(\sigma_{1}\bar{r}_{2}+\sigma_{2}\bar{r}_{1})}{\bar{r}_{1}-\bar{r}_{2}}\right|
    \]
    当 $\rho_{12}=1$ 时
    \[
        \sigma_{p}=\left|\frac{(\sigma_{1}-\sigma_{2})\bar{r}_{p}-(\sigma_{1}\bar{r}_{2}+\sigma_{2}\bar{r}_{1})}{\bar{r}_{1}-\bar{r}_{2}}\right|
    \]
    在这两种情形下，前沿组合由两条起点一直的射线构成。当 $\rho_{12}=0$ 时，
    \[
        \sigma_{p}^{2}=\frac{1}{(\bar{r}_{1}-\bar{r}_{2})^{2}}\bigl[(\sigma_{1}^{2}+\sigma_{2}^{2}-2\sigma_{12})\bar{r}_{p}^{2}-(2\sigma_{1}^{2}\bar{r}_{2}+2\sigma_{2}^{2}\bar{r}_{1})\bar{r}_{p}+(\sigma_{1}^{2}\bar{r}_{2}^{2}+\sigma_{2}^{2}\bar{r}_{1}^{2})\bigr]
    \]
    此时的前沿组合是 $\bar{r}$--$\sigma$ 平面中的一条抛物线。
\end{enumerate}
\end{solution}

\begin{problem}
\problemno{12.5} 令市场上风险证券的收益率向量为 $\tilde{r}$，它的均值为 $\bar{r}$，它的协方差阵为 $\Sigma$，$\Sigma$ 是正定的，而 $\iota$ 为同阶的元素全为 1 的向量。另外，令 $a=\iota^{\top}\Sigma^{-1}\bar{r}$，$b=\bar{r}^{\top}\Sigma^{-1}\bar{r}$，$c=\iota^{\top}\Sigma^{-1}\iota$ 且 $d=bc-a^{2}$。证明 $d=0$ 等价于 $\bar{r}=(b/a)\iota$，即所有风险证券具有相同的期望收益 $(b/a)$。
\end{problem}
\begin{solution}
当 $d=0$ 时，$bd=(a\bar{r}-b\iota)^{\top}\Sigma^{-1}(a\bar{r}-b\iota)=0$，因为 $\Sigma$，$\Sigma^{-1}$ 都是正定的，因而 $a\bar{r}-b\iota=0$，$\bar{r}=\dfrac{b}{a}\iota$，即所有风险证券具有相同的期望收益 $\dfrac{b}{a}$。
\end{solution}

\begin{problem}
\problemno{12.6} 证明最小方差组合（MVP）的期望收益率为 $a/c$，这里的 $a$ 和 $c$ 由练习 12.5 给定。
\end{problem}
\begin{solution}
我们知道最小方差组合的期望收益率为 $r_{\mathrm{MVP}}=-\dfrac{(z_{1}-z_{0})^{\top}\Sigma^{-1}z_{0}}{(z_{1}-z_{0})^{\top}\Sigma^{-1}(z_{1}-z_{0})}$，$z_{0}$，$z_{1}$ 由正文 2.3 节定义，把 $z_{0}$，$z_{1}$ 代入得
\[
    \begin{aligned}
        r_{\mathrm{MVP}}
        &=-\frac{d^{-2}(c\Sigma^{-1}\bar{r}-a\Sigma^{-1}\iota)^{\top}\Sigma(b\Sigma^{-1}\iota-a\Sigma^{-1}\bar{r})}{d^{-2}(c\Sigma^{-1}\bar{r}-a\Sigma^{-1}\iota)^{\top}\Sigma(c\Sigma^{-1}\bar{r}-a\Sigma^{-1}\iota)}\\
        &=-\frac{a(a^{2}-bc)}{c(bc-a^{2})}=\frac{a}{c}
    \end{aligned}
\]
\end{solution}

\begin{problem}
\problemno{12.7} 考虑 $\bar{r}=(b/a)\iota$ 的特殊情形。导出均值--方差前沿。在 $\bar{r}$--$\sigma$ 平面上画出这个组合前沿。
\end{problem}
\begin{solution}
当期望收益率都相等时，所有组合的期望收益率都是常数，因而均值--方差前沿组合退化成为一点即最小方差组合，此时的方差为 $\sigma_{\mathrm{MVP}}^{2}=\dfrac{1}{\iota^{\top}\Sigma^{-1}\iota}$，对应的组合权重为 $z_{\mathrm{MVP}}=\dfrac{\Sigma^{-1}\iota}{\iota^{\top}\Sigma^{-1}\iota}$。
\end{solution}

\begin{problem}
\problemno{12.8**} 令证券收益率向量为 $\tilde{r}$，它的均值为 $\bar{r}$、协方差阵为 $\Sigma$。在 $\bar{r}$--$\sigma$ 平面上，任一组合 $p$ 由一点 $(\bar{r}_{p},\sigma_{p})$ 给定。令所有可能组合的集合为 $A=\{(\bar{r}_{p},\sigma_{p}),\forall\, p\}$。
\begin{enumerate}[label=(\alph*)]
    \item 证明 $A$ 是闭集。
    \item $A$ 一定是凸集吗？若否，给出反例。
\end{enumerate}
\end{problem}
\begin{solution}
\begin{enumerate}[label=(\alph*)]
    \item 假设有一序列 $\bigl\{(r_{p}^{(n)},\sigma_{p}^{(n)})\bigr\}_{1}^{\infty}$，其极限为 $(r_{p},\sigma_{p})$。那么存在 $\theta^{(n)}\in S^{N-1}=\{\theta\in\mathrm{R}^{N}:\theta^{\top}\iota=1\}$ 使得
    \[
        \theta^{(n)\top}\bar{r}=r_{p}^{(n)},\qquad\theta^{(n)\top}\Sigma\theta^{(n)}=\sigma_{p}^{(n)2}
    \]
    显然，$\theta^{(n)}$ 是 $S^{N-1}$ 中的有界序列，存在收敛子列 $\{\theta^{(k_{n})}\}_{1}^{\infty}$，记其极限为 $\theta\in S^{N-1}$，那么
    \[
        \begin{aligned}
            \theta^{\top}\bar{r}&=\lim_{n}\theta^{(k_{n})\top}\bar{r}=\lim_{n}r_{p}^{(k_{n})}=r_{p}\\
            \theta^{\top}\Sigma\theta&=\lim_{n}\theta^{(k_{n})\top}\Sigma\theta^{(k_{n})}=\lim_{n}\sigma_{p}^{(k_{n})2}=\sigma_{p}^{2}
        \end{aligned}
    \]
    因而 $(r_{p},\sigma_{p})\in A$，$A$ 是闭集；
    \item 不一定。当 $N=2$ 且两只证券的期望收益不同、收益率不完全相关时，卖空下的组合在 $\bar{r}$--$\sigma$ 平面上只扫出均值--方差前沿，即一条双曲线。双曲线不是凸集：取前沿上期望收益不同的两点，连接它们的线段会落到前沿左侧，对应的 $(\bar{r},\sigma)$ 不可由任何组合达到。
\end{enumerate}
\end{solution}

\begin{problem}
\problemno{12.9**} 令 $p$ 为 MVF 组合，$zcp$ 是与 $p$ 的协方差为零的 MVF。
\begin{enumerate}[label=(\alph*)]
    \item 证明对任意证券 $i$，我们有
    \[
        \tilde{r}_{i}=\alpha+\beta_{i,zcp}\tilde{r}_{zcp}+\beta_{i,p}\tilde{r}_{p}+\tilde{u}_{i}
    \]
    其中 $\mathrm{Cov}[\tilde{r}_{zcp},\tilde{u}_{i}]=0$，$\mathrm{Cov}[\tilde{r}_{p},\tilde{u}_{i}]=0$，$\mathrm{E}[\tilde{u}_{i}]=0$，$\beta_{i,zcp}=\mathrm{Cov}[\tilde{r}_{i},\tilde{r}_{zcp}]/\mathrm{Var}[\tilde{r}_{zcp}]$ 且 $\beta_{i,p}=\mathrm{Cov}[\tilde{r}_{i},\tilde{r}_{p}]/\mathrm{Var}[\tilde{r}_{p}]$。
    \item 证明 $\beta_{i,zcp}=1-\beta_{i,p}$。
    \item 证明 $\alpha=0$ 及 $\bar{r}_{i}=(1-\beta_{i,p})\bar{r}_{zcp}+\beta_{i,p}\bar{r}_{p}$。
\end{enumerate}
\end{problem}
\begin{solution}
\begin{enumerate}[label=(\alph*)]
    \item 令 $\tilde{r}_{i}=\alpha+\beta_{i,1}\tilde{r}_{zcp}+\beta_{i,2}\tilde{r}_{p}+\tilde{u}_{i}$，由
    \[
        \mathrm{Cov}(\tilde{u}_{i},\tilde{r}_{p})=\mathrm{Cov}(\tilde{r}_{i}-\alpha-\beta_{i,1}\tilde{r}_{zcp}-\beta_{i,2}\tilde{r}_{p},\tilde{r}_{p})=0
    \]
    可得 $\beta_{i,2}=\beta_{i,p}$；同样地有 $\beta_{i,1}=\beta_{i,zcp}$，由 $\mathrm{E}\tilde{u}_{i}=0$ 可以得到 $\alpha=\bar{r}_{i}-\beta_{i,p}\bar{r}_{p}-\beta_{i,zcp}\bar{r}_{zcp}$；
    \item 由文中 (12.6) 式有
    \[
        \begin{aligned}
            \bar{r}_{i}-\bar{r}_{zcp}&=\beta_{i,p}(\bar{r}_{p}-\bar{r}_{zcp})\\
            \bar{r}_{i}-\bar{r}_{p}&=\beta_{i,zcp}(\bar{r}_{zcp}-\bar{r}_{p})
        \end{aligned}
    \]
    两式相减可以得到 $\beta_{i,p}+\beta_{i,zcp}=1$；
    \item 另外，由 (12.6) 式有
    \[
        \bar{r}_{i}=\bar{r}_{zcp}+\beta_{i,p}(\bar{r}_{p}-\bar{r}_{zcp})=\beta_{i,p}\bar{r}_{p}+(1-\beta_{i,p})\bar{r}_{zcp}
    \]
    因此 $\alpha=0$。
\end{enumerate}
\end{solution}

\begin{problem}
\problemno{12.10} 证明 MVE 组合的凸组合（即权重非负的组合）仍然是 MVE 组合。
\end{problem}
\begin{solution}
首先，我们知道 MVE 组合的组合是 MVF 组合；另外，由于是凸组合，因而所得 MVF 组合在均值--方差前沿的上半部，即为 MVE 组合。
\end{solution}

\begin{problem}
\problemno{12.11} 当存在一只收益率为 $r_{F}$ 的无风险债券时，记 $\bar{\eta}=\bar{r}-r_{F}\iota$ 是风险证券的超额收益向量。证明 $\bar{\eta}^{\top}\Sigma^{-1}\bar{\eta}=0$ 等价于 $\bar{r}=r_{F}\iota$。
\end{problem}
\begin{solution}
因为 $\Sigma^{-1}$ 是正定的，
\[
    \bar{\eta}^{\top}\Sigma^{-1}\bar{\eta}=0\Leftrightarrow\bar{\eta}=\bar{r}-r_{F}\iota=0\Leftrightarrow\bar{r}=r_{F}\iota
\]
\end{solution}

\begin{problem}
\problemno{12.12} 当存在无风险证券时，切点组合的期望超额收益为正等价于 $\bar{r}_{\mathrm{mvp}}-r_{F}>0$，即最小方差组合的期望收益为正。
\end{problem}
\begin{solution}
当切点组合的期望收益率为正时，
\[
    \begin{aligned}
        \iota^{\top}\Sigma^{-1}\bar{\eta}>0
        &\Leftrightarrow\iota^{\top}\Sigma^{-1}(\bar{r}-r_{F}\iota)=\iota^{\top}\Sigma^{-1}\bar{r}-r_{F}\iota^{\top}\Sigma^{-1}\iota=\iota^{\top}\Sigma^{-1}\iota\Bigl(\frac{\iota^{\top}\Sigma^{-1}\bar{r}}{\iota^{\top}\Sigma^{-1}\iota}-r_{F}\Bigr)\\
        &=\iota^{\top}\Sigma^{-1}\iota\Bigl(\frac{a}{c}-r_{F}\Bigr)=\iota^{\top}\Sigma^{-1}\iota(\bar{r}_{\mathrm{mvp}}-r_{F})>0\\
        &\Leftrightarrow\bar{r}_{\mathrm{mvp}}-r_{F}>0.
    \end{aligned}
\]
\end{solution}

\begin{problem}
\problemno{12.13*} 当存在无风险证券时，证明切点组合的 Sharpe 比最大。
\end{problem}
\begin{solution}
考虑任意组合 Sharpe 比的平方
\[
    \frac{z^{\top}\bar{r}+(1-z^{\top}\iota)r_{F}-r_{F}}{z^{\top}\Sigma z}=\frac{z^{\top}\bar{\eta}}{z^{\top}\Sigma z}
\]
其中 $z$ 是风险组合的权重，最大化的一阶条件是
\[
    \bar{\eta}\,z^{\top}\Sigma z=z\,\bar{\eta}^{\top}\Sigma z
\]
可以验证 $z_{T}=\dfrac{1}{\iota^{\top}\Sigma^{-1}\bar{\eta}}\Sigma^{-1}\bar{\eta}$ 即为其解，因而切点组合的 Sharpe 比最大。
\end{solution}



% !TeX root = ../../金融经济学_王江_讲义及习题.tex
% 知识点：references/金融经济学/金融经济学-王江_261013.pdf，印刷页 197–207（PDF 页 212–222）
% 习题：references/金融经济学/史上最全最完整的金融经济学答案(王江版).pdf，页 68–75（印刷页 67–74）

\lecture{资本资产定价模型}

\ifdefined\WJHandoutKnowledge

第 12 章说明：若参与者有均值--方差偏好，风险证券的最优组合都沿切点组合方向。本章再加\textbf{市场出清}，把切点组合识别为所有风险证券的总供给组合，由此得到资本资产定价模型（CAPM）。

\begin{tcolorbox}[knowledge,title=13.1 市场组合与均衡]
证券 $n$ 的流通份额为 $\bar\theta_n$、价格为 $P_n$，市场价值为 $v_n=P_n\bar\theta_n$。全部证券的市场价值 $v_M=\sum_nv_n$，\textbf{市场组合}的价值权重为
\[
 z_{n,M}=\frac{P_n\bar\theta_n}{v_M},\qquad
 z_M=\begin{bmatrix}z_{1,M}&\cdots&z_{N,M}\end{bmatrix}^{\mathsf T}.
 \tag{13.2}
\]
这是所有可交易风险证券按市值组成的总供给组合。它不是随意选取的股票指数；要用某指数代理市场组合，还需说明覆盖范围和权重。

\textbf{定理 13.1}\quad 有无风险证券且所有参与者对风险资产的需求均为切点组合 $z_T$ 时，市场出清意味着 $z_M=z_T$。证明：参与者 $k$ 对风险组合投入 $a_k$，总风险资产需求为 $(\sum_ka_k)z_T$；债券净供给为零时，$\sum_ka_k=v_M$，故风险资产总需求正好是 $v_Mz_T$，与风险资产总供给 $v_Mz_M$ 相等。于是二者权重相同。

由第 12 章切点组合的贝塔关系得到\textbf{定理 13.2（CAPM）}：对任意证券 $n$，
\[
 \bar r_n-r_F=\beta_{n,M}(\bar r_M-r_F),\qquad
 \beta_{n,M}=\frac{\operatorname{Cov}(\widetilde r_n,\widetilde r_M)}
 {\operatorname{Var}(\widetilde r_M)}.
 \tag{13.3}
\]
证券的风险溢价是市场风险溢价乘以它对市场的\textbf{贝塔}。定价关系要求均值--方差选择、可交易与市场出清等条件，不是对任意收益数据的恒等式。
\end{tcolorbox}

\begin{tcolorbox}[knowledge,title=13.3 系统风险与证券市场线]
把资产的超额收益投影到市场超额收益上：
\[
 \widetilde r_n-r_F
 =\alpha_n+\beta_{n,M}(\widetilde r_M-r_F)+\widetilde\varepsilon_n,
 \quad\mathrm E[\widetilde\varepsilon_n]=0,
 \quad\operatorname{Cov}(\widetilde\varepsilon_n,\widetilde r_M)=0.
 \tag{13.4}
\]
贝塔项是\textbf{系统风险}，残差项是与市场不相关的\textbf{剩余风险}。方差分解为
\[
 \sigma_n^2=\beta_{n,M}^2\sigma_M^2+\sigma_{\varepsilon_n}^2.
\]
CAPM 要求 $\alpha_n=0$；剩余方差即使很大，也不会仅因自身波动大而获得风险溢价。证券市场线以贝塔为横轴、期望收益率为纵轴，截距是 $r_F$，斜率是 $\bar r_M-r_F$。这与以\emph{总标准差}为横轴、只适用于有效组合的资本市场线不同。

\begin{tcolorbox}[example]
\textbf{例 13.1}\quad 若 $r_F=4\%$、$\bar r_M=12\%$，市场风险溢价为 $8\%$。贝塔分别为 $0,1.5,-1.5$ 的资产，CAPM 预期风险溢价分别为 $0,12\%,-12\%$，对应预期收益率 $4\%,16\%,-8\%$。负贝塔资产在市场下跌时倾向提供保险，因而可接受低于无风险利率的预期收益；高总波动率并不能改变这一贝塔结论。
\end{tcolorbox}
\end{tcolorbox}

\begin{tcolorbox}[knowledge,title=13.4 无风险资产缺席时的 CAPM]
没有可交易无风险资产时，第 12 章仍保证参与者各自的最优组合位于均值--方差前沿。市场出清使市场组合 $M$ 也位于该前沿。取与 $M$ 收益零协方差的前沿组合 $zcm$，\textbf{定理 13.3}给出
\[
 \bar r_q-\bar r_{zcm}
 =\beta_{q,M}(\bar r_M-\bar r_{zcm})
 \tag{13.5}
\]
对任意组合 $q$ 成立。这里的 $\bar r_{zcm}$ 是\textbf{零贝塔组合收益率}，不是不存在的无风险利率；$zcm$ 自身仍有风险。这是无风险资产版本的推广，不能把 (13.5) 的截距直接写作 $r_F$。
\end{tcolorbox}

\begin{tcolorbox}[knowledge,title=13.5 CAPM 作为均衡结果]
均值--方差偏好可由不同条件得到。\textbf{二次效用}使期望效用取决于均值和方差，但要限制财富域以保持边际效用为正。若资产\textbf{联合正态}且参与者具有 CARA 效用，则组合终值正态，期望效用等价于最大化
\[
 \mathrm E[\widetilde w_k]-\frac{a_k}{2}\operatorname{Var}(\widetilde w_k).
\]
令风险证券期末支付向量均值为 $\bar x$、协方差矩阵为 $\Sigma$、0 期价格为 $P$，参与者 $k$ 的最优证券持仓为
\[
 \theta_k=\frac1{a_k}\Sigma^{-1}[\bar x-(1+r_F)P].
\]
对所有人加总并令总需求等于总供给 $\bar\theta$，记 $1/a=\sum_k1/a_k$，则
\[
 P=\frac{\bar x-a\Sigma\bar\theta}{1+r_F}.
\]
将价格和市场组合的价值定义代入，可再得到 (13.3)。这一推导说明在正态支付、CARA 偏好、共同信念和完整可交易等\textbf{充分条件}下，CAPM 可以由一般均衡产生；并不表示这些条件在现实中自动成立。
\end{tcolorbox}

\begin{tcolorbox}[summary,title=本章小结]
均值--方差偏好使每个人持有相同的有效风险组合；市场出清把该组合识别为市值加权的市场组合。CAPM 于是用市场贝塔定价系统风险，剩余风险不单独获得溢价。无风险资产缺席时，用零贝塔前沿组合的收益率取代无风险利率。二次效用与正态支付加 CARA 偏好提供两类可推导 CAPM 的具体条件。
\end{tcolorbox}

\fi

% !TeX root = ../../金融经济学_王江_讲义及习题.tex
% 13_CAPM 课后习题（由章节文件拆分）

\begin{problem}
\problemno{13.1*} 如下的经济有两个时期，0 和 1：
\begin{itemize}
    \item 有 $N$ 只风险证券（股票）和一只无风险证券。
    \item 无风险利率为 0。
    \item 风险证券的 1 期支付向量为 $\tilde{v}$，$\tilde{v}$ 是均值为 $\bar{v}$、协方差阵为 $\Sigma$ 的 $N$ 维正态随机变量。假设 $\Sigma$ 是正定的。
    \item 股票市场组合的 1 期支付为 $\tilde{m}=\iota^{\mathrm{T}}\tilde{v}$，其中 $\iota$ 是元素全为 1 的 $(N\times 1)$ 列向量。
    \item 经济中有 $K$ 个投资者，所有投资者都具有相同的 CARA 偏好：$u(w)=-e^{-w}$。
    \item 每一投资者的禀赋是 $\bar{\theta}_{k}$ 股股票，$k=1,\cdots,K$。（注意 $\bar{\theta}_{k}$ 是由股票份额组成的列向量）且 $\sum_{k}\bar{\theta}_{k}=\iota$。
    \item 每一投资者在 1 期也有劳动收入 $y_{k}(\tilde{m}+\tilde{\varepsilon})$，其中 $\tilde{\varepsilon}$ 服从均值为 0、方差为 $\sigma^{2}$ 的正态分布，而 $y_{k}$ 是给定的。
\end{itemize}
假设 $\tilde{\varepsilon}$ 和 $\tilde{v}$ 不相关。
\begin{enumerate}[label=(\alph*)]
    \item 求解投资者（在给定价格向量下）的股票需求函数。
    \item 求解均衡股票价格。
    \item 证明 CAPM 成立。
    \item 均衡配置是 Pareto 最优吗？为什么？证明答案。如果 $\tilde{\varepsilon}=0$，则答案如何变化？[提示：在 Pareto 最优配置当中，投资者的相对边际效用在任何状态下都是相等的。]
\end{enumerate}
\end{problem}
\begin{solution}
\begin{enumerate}[label=(\alph*)]
    \item 假设股票价格向量为 $S$，参与者持有的股票头寸为 $\theta_{k}$，那么参与者 $k$ 的优化问题是
    \[
        \begin{aligned}
            \max_{\theta_{k}}\;&\mathrm{E}[-e^{-\tilde{w}_{k}}]\\
            \text{s.t.}\;&\tilde{w}_{k}=w_{k}+\bar{\theta}_{k}^{\mathrm{T}}S-\theta_{k}^{\mathrm{T}}S+\theta_{k}^{\mathrm{T}}\tilde{v}+y_{k}(\iota^{\mathrm{T}}\tilde{v}+\tilde{\varepsilon})
        \end{aligned}
    \]
    CARA 且 $\tilde{v}$、$\tilde{\varepsilon}$ 正态，最大化 $E[-e^{-\tilde{w}_{k}}]$ 等价于最大化 $E[\tilde{w}_{k}]-\tfrac12\operatorname{Var}(\tilde{w}_{k})$。记常数项 $A_{k}=w_{k}+\bar{\theta}_{k}^{\mathrm{T}}S$，则
    \[
        \begin{aligned}
            E[\tilde{w}_{k}]&=A_{k}+\theta_{k}^{\mathrm{T}}(\bar{v}-S)+y_{k}\iota^{\mathrm{T}}\bar{v}\\
            \operatorname{Var}(\tilde{w}_{k})&=(\theta_{k}+y_{k}\iota)^{\mathrm{T}}\Sigma(\theta_{k}+y_{k}\iota)+y_{k}^{2}\sigma^{2}
        \end{aligned}
    \]
    对 $\theta_{k}$ 求导：$-(\bar{v}-S)+\Sigma(\theta_{k}+y_{k}\iota)=0$，解得
    \[
        \theta_{k}^{*}=\Sigma^{-1}(\bar{v}-S)-y_{k}\iota
    \]

    \item 市场出清 $\sum_{k}\theta_{k}^{*}=\iota$ 给出
    \[
        K\Sigma^{-1}(\bar{v}-S)-\Bigl(\sum_{k}y_{k}\Bigr)\iota=\iota
    \]
    即 $\bar{v}-S=\dfrac{1}{K}\bigl(\sum_{k}y_{k}+1\bigr)\Sigma\iota$，因而股价为
    \[
        S=\bar{v}-\frac{1}{K}\Bigl(\sum y_{k}+1\Bigr)\Sigma\iota
    \]

    \item 这里无风险利率为 $0$，CAPM 是 $\bar{r}_{i}=\beta_{i}\bar{r}_{m}$，其中 $\beta_{i}=\operatorname{Cov}(\tilde{r}_{i},\tilde{r}_{m})/\operatorname{Var}(\tilde{r}_{m})$，$\tilde{r}_{m}$ 是市场组合收益率。下面验证这个等式。因为
    \[
        \begin{aligned}
            \frac{\bar{r}_{i}}{\beta_{i}\bar{r}_{m}}
            &=\Bigl(\frac{\bar{v}_{i}}{S_{i}}-1\Bigr)\Big/\Biggl[\frac{\operatorname{Cov}(\tilde{v}_{i},\,\iota^{\mathrm{T}}\tilde{v})(\iota^{\mathrm{T}}S)}{\operatorname{Var}(\iota^{\mathrm{T}}\tilde{v})S_{i}}\Bigl(\frac{\iota^{\mathrm{T}}\bar{v}}{\iota^{\mathrm{T}}S}-1\Bigr)\Biggr]\\
            &=(\bar{v}_{i}-S_{i})\big/\bigl[\operatorname{Cov}(\tilde{v}_{i},\,\iota^{\mathrm{T}}\tilde{v})(\iota^{\mathrm{T}}\bar{v}-\iota^{\mathrm{T}}S)\bigr]\\
            &=\frac{1}{K}\Bigl(\sum y_{k}+1\Bigr)(\Sigma\iota)_{i}\Big/\frac{(\Sigma\iota)_{i}\bigl(\tfrac{1}{K}(\sum y_{k}+1)(\iota^{\mathrm{T}}\Sigma\iota)\bigr)}{\iota^{\mathrm{T}}\Sigma\iota}\\
            &=1
        \end{aligned}
    \]
    即 $\bar{r}_{i}=\beta_{i}\bar{r}_{m}$，因而 CAPM 成立；

    \item 对于参与者 $k$ 来说，1 期边际效用为：
    \[
        e^{-\tilde{w}_{k}}=e^{-w_{k}+\bar{\theta}_{k}^{\mathrm{T}}S+(\Sigma^{-1}(\bar{v}-S))^{\mathrm{T}}(\bar{v}-S)+y_{k}\tilde{\varepsilon}}
    \]
    除非 $y_{k}$ 全部相等，否则边际效用之比（1 期边际效用除以 0 期边际效用）是不会相等的，因而均衡配置不是 Pareto 最优的；如果 $\tilde{\varepsilon}=0$，那么参与者的边际效用之比最多相差一个常数，但是在均衡条件下，他们对股票的定价是相同的，因而边际效用之比相等，均衡配置是 Pareto 最优的。
\end{enumerate}
\end{solution}

\begin{problem}
\problemno{13.2**} 在练习 13.1 定义的经济中，现在假设 $\tilde{\varepsilon}$ 与 $\tilde{v}$ 相关。令 $\operatorname{Cov}[\tilde{v},\tilde{\varepsilon}]=c$，这里 $c$ 是一 $(N\times 1)$ 维向量。
\begin{enumerate}[label=(\alph*)]
    \item 求解投资者的股票需求函数。是否得到两基金分离？证明答案。如果两基金分离不成立，能否得到 $K$ 基金分离？$K$ 应为多少？
    \item 求解均衡股票价格。
    \item CAPM 成立吗？解释它为什么成立或不成立。列出导致不同股票的期望收益不同的原因并做出解释。
    \item 讨论均衡配置的 Pareto 最优性。
\end{enumerate}
\end{problem}
\begin{solution}
\begin{enumerate}[label=(\alph*)]
    \item 用和第 1 题类似的方法可以求得参与者 $k$ 的股票需求函数为
    \[
        \theta_{k}=\Sigma^{-1}(\bar{v}-S)-y_{k}(\iota+c)
    \]
    此时两基金分离不成立，但是三基金分离成立：一是组合头寸向量为 $\Sigma^{-1}(\bar{v}-S)$ 的股票组合，另一是组合头寸向量为 $\iota+c$ 的股票组合，最后是无风险债券；

    \item 类似地，可以求得均衡股价为
    \[
        S=\bar{v}-\frac{1}{K}\Biggl[\Bigl(\sum y_{k}+1\Bigr)\Sigma\iota+\sum y_{k}\Sigma c\Biggr]
    \]

    \item 考虑
    \[
        \begin{aligned}
            \frac{\bar{r}_{i}}{\beta_{i}\bar{r}_{m}}
            &=\Bigl(\frac{\bar{v}_{i}}{S_{i}}-1\Bigr)\Big/\Biggl[\frac{\operatorname{Cov}(\tilde{v}_{i},\,\iota^{\mathrm{T}}\tilde{v})(\iota^{\mathrm{T}}S)}{\operatorname{Var}(\iota^{\mathrm{T}}\tilde{v})S_{i}}\Bigl(\frac{\iota^{\mathrm{T}}\bar{v}}{\iota^{\mathrm{T}}S}-1\Bigr)\Biggr]\\
            &=(\bar{v}_{i}-S_{i})\big/\bigl[\operatorname{Cov}(\tilde{v}_{i},\,\iota^{\mathrm{T}}\tilde{v})(\iota^{\mathrm{T}}\bar{v}-\iota^{\mathrm{T}}S)\bigr]\\
            &=\frac{1}{K}\Biggl(\Bigl(\sum y_{k}+1\Bigr)(\Sigma\iota)+\sum y_{k}\Sigma c\Biggr)_{i}\Big/\frac{(\Sigma\iota)_{i}\bigl(\tfrac{1}{K}\bigl\{(\sum y_{k}+1)(\iota^{\mathrm{T}}\Sigma\iota)+\sum y_{k}\Sigma c\bigr\}\bigr)}{\iota^{\mathrm{T}}\Sigma\iota}\\
            &\ne 1
        \end{aligned}
    \]
    因而 CAPM 不成立：上式一般不等于 $1$，期望收益不能只写成市场组合 beta 的倍数。这是因为需要用股票来对冲与股票相关的其它风险。股票 $i$ 的期望收益率为
    \[
        \bar{r}_{i}=\frac{\bar{v}_{i}}{S_{i}}-1=\frac{1}{1-\frac{1}{K\bar{v}_{i}}\bigl((\sum y_{k}+1)(\Sigma\iota)+\sum y_{k}\Sigma c\bigr)_{i}}-1
    \]
    因而导致股票期望收益率不同的因素有：
    \begin{itemize}
        \item 期望支付；
        \item 与其它股票的相关性以及自身的风险（方差）；
        \item 与收入中的风险 $\tilde{\varepsilon}$ 的相关性。
    \end{itemize}

    \item 考虑参与者 $k$ 的 1 期边际效用：
    \[
        e^{-\tilde{w}_{k}}=e^{-w_{k}+\bar{\theta}_{k}^{\mathrm{T}}S+(\Sigma^{-1}(\bar{v}-S))^{\mathrm{T}}(\bar{v}-S)+y_{k}(\tilde{\varepsilon}-c^{\mathrm{T}}\tilde{v})}
    \]
    除非 $\tilde{\varepsilon}=c^{\mathrm{T}}\tilde{v}$ 或者 $y_{k}$ 全部相等，否则边际效用之比（1 期边际效用除以 0 期边际效用）是不会相等的，均衡配置也不是 Pareto 最优的。
\end{enumerate}
\end{solution}

\begin{problem}
\problemno{13.3**} 考虑一个由两个国家（$a$ 和 $b$）组成的、两个时期（0 和 1）的经济。这两个国家使用相同的货币。
\begin{itemize}
    \item 两个国家中的无风险资产是一样的，无风险利率为 $r_{F}$。
    \item 国家 $a$ 中有 $N$ 只风险证券，国家 $b$ 中也有 $N$ 只风险证券。$\tilde{v}_{a}=\bar{v}+\tilde{z}_{a}$、$\tilde{v}_{b}=\bar{v}+\tilde{z}_{b}$ 分别是 $a$-资产、$b$-资产的 1 期支付。假设 $\bar{v}\sim N(0,\sigma_{w}^{2}I)$，这里 $I$ 是 $N\times N$ 单位阵，$\tilde{v}_{k}\sim N(\bar{v},\Sigma_{k})$，$k=a,b$，且 $\bar{v}$、$\tilde{z}_{a}$ 和 $\tilde{z}_{b}$ 是独立的。
    \item 每一风险资产的总供给都为 1。
    \item 国家 $a$ 有 $K_{a}$ 个投资者而国家 $b$ 有 $K_{b}$ 个投资者。所有投资都具有风险厌恶系数为 $\alpha$、依赖于 1 期财富的 CARA 效用函数。
\end{itemize}
令 $p_{a}$ 和 $p_{b}$ 为国家 $a$ 和国家 $b$ 的资产价格向量。
\begin{enumerate}[label=(\alph*)]
    \item 假设投资者可以持有两个国家的资产。计算国家 $a$ 和国家 $b$ 的投资者的资产需求。
    \item 计算均衡价格。
    \item 计算所有资产的均衡风险溢价。
    \item 证明国际化的 CAPM 成立，即
    \[
        \bar{r}-r_{F}=\beta(\bar{r}_{wm}-r_{F}),\quad(\bar{r}=\mathrm{E}[\tilde{r}]),
    \]
    其中 $\beta$ 是相对于世界市场组合的 beta 值，世界市场组合是由国家 $a$ 和国家 $b$ 的市场组合构成的等权重组合（也即两国所有资产的组合）。
\end{enumerate}
\end{problem}
\begin{solution}
\begin{enumerate}[label=(\alph*)]
    \item 考虑国家 $a$ 中的参与者 $k$，他的优化问题是
    \[
        \begin{aligned}
            \max_{\theta_{k,a},\,\theta_{k,b}}\;&\mathrm{E}[-e^{-\alpha\tilde{w}_{k}}]\\
            \text{s.t.}\;&\tilde{w}_{k}=(w_{k}-\theta_{k,a}P_{a}-\theta_{k,b}P_{b})(1+r_{F})+\theta_{k,a}^{\mathrm{T}}\tilde{v}_{a}+\theta_{k,b}^{\mathrm{T}}\tilde{v}_{b}
        \end{aligned}
    \]
    求得最优的股票需求量为
    \[
        \begin{aligned}
            \theta_{k,a}^{*}&=\frac{1}{\alpha}\Sigma_{a}^{-1}(\bar{v}-P_{a}(1+r_{F}))-\frac{\sigma_{w}^{2}}{\alpha}\Sigma_{a}^{-1}(\sigma_{w}^{2}(\Sigma_{a}^{-1}+\Sigma_{b}^{-1})+I)^{-1}\\
            &\quad\times\bigl(\Sigma_{a}^{-1}(\bar{v}-P_{a}(1+r_{F}))+\Sigma_{b}^{-1}(\bar{v}-P_{b}(1+r_{F}))\bigr)\\
            \theta_{k,b}^{*}&=\frac{1}{\alpha}\Sigma_{b}^{-1}(\bar{v}-P_{b}(1+r_{F}))-\frac{\sigma_{w}^{2}}{\alpha}\Sigma_{b}^{-1}(\sigma_{w}^{2}(\Sigma_{a}^{-1}+\Sigma_{b}^{-1})+I)^{-1}\\
            &\quad\times\bigl(\Sigma_{a}^{-1}(\bar{v}-P_{a}(1+r_{F}))+\Sigma_{b}^{-1}(\bar{v}-P_{b}(1+r_{F}))\bigr)
        \end{aligned}
    \]
    类似地可以求得国家 $b$ 中参与者 $k$ 的需求函数有
    \[
        \theta_{k,a}'^{*}=\theta_{k,a}^{*},\qquad\theta_{k,b}'^{*}=\theta_{k,b}^{*}
    \]

    \item 市场出清条件是
    \[
        \begin{aligned}
            \sum\theta_{k,a}^{*}+\sum\theta_{k,a}'^{*}&=\iota\\
            \sum\theta_{k,b}^{*}+\sum\theta_{k,b}'^{*}&=\iota
        \end{aligned}
    \]
    因而可以求得均衡股价为
    \[
        \begin{aligned}
            P_{a}(1+r_{F})&=\bar{v}-\frac{\alpha}{K}\sigma_{w}^{2}\iota-\frac{\alpha}{2K}\Sigma_{a}\iota\\
            P_{b}(1+r_{F})&=\bar{v}-\frac{\alpha}{K}\sigma_{w}^{2}\iota-\frac{\alpha}{2K}\Sigma_{b}\iota
        \end{aligned}
    \]

    \item 国家 $a$ 中资产 $k$ 的风险溢价为
    \[
        \bar{r}_{k}-r_{F}=\frac{\bar{v}_{k}}{P_{a,k}}-(1+r_{F})=\frac{\frac{\alpha}{2K}(2\sigma_{w}^{2}\iota+\Sigma_{a}\iota)_{k}(1+r_{F})}{\bar{v}_{k}-\frac{\alpha}{2K}(2\sigma_{w}^{2}\iota+\Sigma_{a}\iota)_{k}}
    \]
    同样地，对于国家 $b$ 中资产 $i$ 有
    \[
        \bar{r}_{i}-r_{F}=\frac{\bar{v}_{i}}{P_{b,i}}-(1+r_{F})=\frac{\frac{\alpha}{2K}(2\sigma_{w}^{2}\iota+\Sigma_{a}\iota)_{i}(1+r_{F})}{\bar{v}_{i}-\frac{\alpha}{2K}(2\sigma_{w}^{2}\iota+\Sigma_{a}\iota)_{i}}
    \]

    \item 对于国家 $a$ 中的资产 $k$，
    \[
        \begin{aligned}
            \beta_{k}&=\operatorname{Cov}\Bigl(\frac{\tilde{v}_{k}}{P_{a,k}}-1,\,\frac{\iota^{\mathrm{T}}(\tilde{v}_{a}+\tilde{v}_{b})}{\iota^{\mathrm{T}}(P_{a}+P_{b})}-1\Bigr)\Big/\operatorname{Var}\Bigl(\frac{\iota^{\mathrm{T}}(\tilde{v}_{a}+\tilde{v}_{b})}{\iota^{\mathrm{T}}(P_{a}+P_{b})}\Bigr)\\
            &=\frac{(2\sigma_{w}^{2}\iota+\Sigma_{a}\iota)_{k}(\iota^{\mathrm{T}}(P_{a}+P_{b}))}{\iota^{\mathrm{T}}(4\sigma_{w}^{2}I+\Sigma_{a}+\Sigma_{b})\iota}\\
            \bar{r}_{wm}&=\frac{\tfrac12\iota^{\mathrm{T}}(\tilde{v}_{a}+\tilde{v}_{b})}{\iota^{\mathrm{T}}(P_{a}+P_{b})}\\
            \bar{r}_{wm}-r_{F}&=\frac{\frac{\alpha}{2K}\iota^{\mathrm{T}}(4\sigma_{w}^{2}I+\Sigma_{a}+\Sigma_{b})\iota}{\iota^{\mathrm{T}}(P_{a}+P_{b})}
        \end{aligned}
    \]
    因此，$\bar{r}_{k}-r_{F}=\beta_{k}(\bar{r}_{wm}-r_{F})$，可以验证对于国家 $b$ 中的资产此式也成立，因而国际化的 CAPM 成立。
\end{enumerate}
\end{solution}

\begin{problem}
\problemno{13.4**} 继续考虑练习 13.3 所描述的经济。假设国家 $a$ 中的投资者可以持有国家 $b$ 的资产而国家 $b$ 中的投资者不能持有国家 $a$ 的资产。
\begin{enumerate}[label=(\alph*)]
    \item 重复 (a)--(c) 中的计算。
    \item 现在证明
    \[
        \bar{r}-r_{F}=\gamma_{a}(\bar{r}_{a}-r_{F})+\gamma_{b}(\bar{r}_{b}-r_{F}).
    \]
    其中 $\bar{r}_{a}$ 和 $\bar{r}_{b}$ 分别是国家 $a$ 和 $b$ 的市场组合的期望收益。计算 $\gamma_{a}$ 和 $\gamma_{b}$ 并讨论所得到的结论。
\end{enumerate}
\end{problem}
\begin{solution}
略
\end{solution}

\begin{problem}
\problemno{13.5*} 考虑如下的具有两个时期 0 和 1 的 Arrow--Debreu 经济：
\begin{itemize}
    \item 单一的消费品（也作为计量单位）。
    \item 经济在 $t=1$ 期的状态为 $\omega\in\Omega=\{1,\cdots,\Omega\}$。
    \item $\pi_{\omega}>0$ 是状态 $\omega$，$\omega\in\Omega$ 的概率。
    \item 证券市场上存在 Arrow--Debreu 状态或有证券的完全集合。令 $\phi_{\omega}$ 为状态 $\omega$ 或有要求权（只在状态 $\omega$ 下支付 1 单位消费品）的价格。
    \item 总禀赋是 $\{e_{\omega}:\omega\in\Omega\}$，假设 $e_{\omega}>0$，$\forall\,\omega\in\Omega$。
\end{itemize}
\begin{enumerate}[label=(\alph*)]
    \item 写出 $\phi_{\omega}$，$\omega\in\Omega$ 的无套利条件。
    \item 一个状态或有要求权的组合是 $w=\begin{bmatrix}w_{1} \\ \cdots \\ w_{\omega} \\ \cdots \\ w_{\Omega}\end{bmatrix}$，其中 $w_{\omega}$ 是投资于状态 $\omega$ 状态或有要求权的份额，$\sum_{\omega}w_{\omega}=1$。
    \begin{itemize}
        \item 计算组合收益率的均值和方差。
        \item 计算两个组合 $w$ 和 $w'$ 之间的协方差。
    \end{itemize}
    \item 考虑如下的两个组合 $w^{a}$ 和 $w^{b}$
    \[
        \forall\,\omega\in\Omega:\quad w_{\omega}^{a}=\frac{\phi_{\omega}}{\sum_{\omega}\phi_{\omega}},\quad w_{\omega}^{b}=\frac{\phi_{\omega}q_{\omega}}{\sum_{\omega}\phi_{\omega}q_{\omega}}
    \]
    其中 $q_{\omega}=\phi_{\omega}/\pi_{\omega}$。
    \begin{itemize}
        \item 验证 $w^{a}$ 是 $t=0$ 期价格为 1 的无风险组合。
        \item 证明所有的前沿组合都是组合 $w^{a}$ 和 $w^{b}$ 的凸组合。
    \end{itemize}
\end{enumerate}
\end{problem}
\begin{solution}
\begin{enumerate}[label=(\alph*)]
    \item 无套利的条件是 $\phi_{\omega}>0$，$\sum\phi_{\omega}\le 1$；
    \item
    \begin{itemize}
        \item 收益率 $\tilde{r}(\omega)=\frac{w_{\omega}}{\phi_{\omega}}-1$，因而收益率的期望值是 $\bar{r}=\sum\pi_{\omega}\frac{w_{\omega}}{\phi_{\omega}}-1$，方差是 $\operatorname{Var}(\tilde{r})=\sum\pi_{\omega}\bigl(\frac{w_{\omega}}{\phi_{\omega}}\bigr)^{2}-\bigl(\sum\pi_{\omega}\frac{w_{\omega}}{\phi_{\omega}}\bigr)^{2}$
        \item
        \[
            \operatorname{Cov}(\tilde{r},\,\tilde{r}')=\sum\pi_{\omega}\frac{w_{\omega}}{\phi_{\omega}}\frac{w_{\omega}'}{\phi_{\omega}}-\sum\pi_{\omega}\frac{w_{\omega}}{\phi_{\omega}}\sum\pi_{\omega}\frac{w_{\omega}'}{\phi_{\omega}}
        \]
    \end{itemize}
    \item
    \begin{itemize}
        \item $w^{a}$ 的 1 期支付为 $\frac{w_{\omega}^{a}}{\phi_{\omega}}=\frac{1}{\phi_{0}}$ 是一常数；而它的价格为 1，即它是 0 期价格为 1 的无风险组合。
        \item 考虑下面的优化问题
        \[
            \begin{aligned}
                \max_{w_{\omega}}\;&\frac12\sum\pi_{\omega}\Bigl(\frac{w_{\omega}}{\phi_{\omega}}\Bigr)^{2}\\
                \text{s.t.}\;&\sum\pi_{\omega}\frac{w_{\omega}}{\phi_{\omega}}-1=\bar{r}_{p}\\
                &\sum w_{\omega}=1
            \end{aligned}
        \]
        构造 Lagrangian 函数
        \[
            \mathcal{L}=\frac12\sum\pi_{\omega}\Bigl(\frac{w_{\omega}}{\phi_{\omega}}\Bigr)^{2}+\lambda_{r}\Bigl(\bar{r}_{p}+1-\pi_{\omega}\frac{w_{\omega}}{\phi_{\omega}}\Bigr)+\lambda_{1}\Bigl(1-\sum w_{\omega}\Bigr)
        \]
        一阶条件为
        \[
            \pi_{\omega}\frac{w_{\omega}}{\phi_{\omega}^{2}}=\lambda_{r}\frac{w_{\omega}}{\phi_{\omega}}+\lambda_{1}
        \]
        \[
            \begin{aligned}
                \Rightarrow w_{\omega}&=\lambda_{r}\phi_{\omega}+\lambda_{1}\phi_{\omega}q_{\omega}\\
                &=\frac{(1+\bar{r}_{p})\sum\phi_{\omega}-1}{\sum\phi_{\omega}q_{\omega}-(\sum\phi_{\omega})^{2}}\phi_{\omega}q_{\omega}+\frac{\sum\phi_{\omega}-(1+\bar{r}_{p})\sum\phi_{\omega}q_{\omega}}{\sum\phi_{\omega}q_{\omega}-(\sum\phi_{\omega})^{2}}\phi_{\omega}\\
                &=\delta\frac{\phi_{\omega}q_{\omega}}{\sum\phi_{\omega}q_{\omega}}+(1-\delta)\frac{\phi_{\omega}}{\sum\phi_{\omega}}
            \end{aligned}
        \]
        其中，
        \[
            \delta=\frac{((1+\bar{r}_{p})\sum\phi_{\omega}-1)\sum\phi_{\omega}q_{\omega}}{\sum\phi_{\omega}q_{\omega}-(\sum\phi_{\omega})^{2}}
        \]
        因为约束都起限制作用，因而 $\lambda_{r}>0$，$\lambda_{1}>0$，因而任意前沿组合都是 $w^{a}$，$w^{b}$ 的凸组合。
    \end{itemize}
\end{enumerate}
\end{solution}

\begin{problem}
\problemno{13.6**} 继续考虑练习 13.5 中的经济。
\begin{enumerate}[label=(\alph*)]
    \item 令 $u(\cdot)$ 为代表性参与者的效用函数。证明为使 CAPM 成立 $u(\cdot)$ 必须具有平方函数的形式。
    \item 假设所有的资产都可由 CAPM 定价。证明一份在状态 $\omega$（$\omega\in\Omega$）下支付 $x_{\omega}$ 单位消费品的证券价格由下式给出
    \[
        p_{x}=\alpha\mathrm{E}[\tilde{x}]-\beta\operatorname{Cov}[\tilde{x},\tilde{e}].
    \]
    其中 $\alpha$ 和 $\beta$ 是两个和 $x$ 无关的常数。
\end{enumerate}
\end{problem}
\begin{solution}
\begin{enumerate}[label=(\alph*)]
    \item 对于支付为 $\tilde{x}$ 的资产而言，其价格为 $P_{x}=\mathrm{E}\bigl[\frac{u'(\tilde{e})}{u'(e_{0})}\tilde{x}\bigr]=\operatorname{Cov}\bigl(\frac{u'(\tilde{e})}{u'(e_{0})},\,\tilde{x}\bigr)+\frac{\bar{x}}{1+r_{F}}$，如果 CAPM 成立，那么有
    \[
        \begin{aligned}
            \frac{\bar{x}}{P_{x}}-(1+r_{F})&=\frac{\operatorname{Cov}(\tilde{x},\,\tilde{e})P_{e}}{P_{x}\operatorname{Var}(\tilde{e})}\Bigl(\frac{\bar{e}}{P_{e}}-(1+r_{F})\Bigr)\\
            \bar{x}-P_{x}(1+r_{F})&=\frac{\operatorname{Cov}(\tilde{x},\,\tilde{e})}{\operatorname{Var}(\tilde{e})}\bigl(\bar{e}-P_{e}(1+r_{F})\bigr)
        \end{aligned}
    \]
    因而对于任意的 $\tilde{x}$ 都有
    \[
        \operatorname{Cov}\Biggl(\tilde{x},\,\frac{\bar{e}-P_{e}(1+r_{F})}{\operatorname{Var}(\tilde{e})}\tilde{e}+\frac{u'(\tilde{e})(1+r_{F})}{u'(e_{0})}\Biggr)=0
    \]
    由此可得 $u'(\tilde{e})$ 是 $\tilde{e}$ 的线性函数，$u$ 是二次函数；

    \item 由上面的
    \[
        \bar{x}-P_{x}(1+r_{F})=\frac{\operatorname{Cov}(\tilde{x},\,\tilde{e})}{\operatorname{Var}(\tilde{e})}\bigl(\bar{e}-P_{e}(1+r_{F})\bigr)
    \]
    可得，
    \[
        \begin{aligned}
            P_{x}&=\frac{1}{1+r_{F}}\bar{x}-\frac{\bar{e}-P_{e}(1+r_{F})}{(1+r_{F})\operatorname{Var}(\tilde{e})}\operatorname{Cov}(\tilde{x},\,\tilde{e})\\
            &=\alpha\bar{x}-\beta\operatorname{Cov}(\tilde{x},\,\tilde{e})
        \end{aligned}
    \]
    其中，$\alpha=\frac{1}{1+r_{F}}$，$\beta=\frac{\bar{e}-P_{e}(1+r_{F})}{(1+r_{F})\operatorname{Var}(\tilde{e})}$ 是两个与 $\tilde{x}$ 无关的常数。
\end{enumerate}
\end{solution}

\begin{problem}
\problemno{13.7*} 考虑练习 13.5 中的经济。假设只有两个可能状态 $a$ 和 $b$，$e_{b}>e_{a}$。称支付总禀赋的证券为股票。考虑一份以股票为标的资产、执行价格为 $K\in(e_{a},e_{b})$ 的看涨期权。
\begin{enumerate}[label=(\alph*)]
    \item 能用 CAPM（代表性参与者具有一般形式的 $u(\cdot)$）为期权定价吗？解释所得的结论。
    \item 能把上面的结论推广到多余两个状态的情形吗？给出解释。
\end{enumerate}
\end{problem}
\begin{solution}
\begin{enumerate}[label=(\alph*)]
    \item 如果 1 期只有两个状态，并且 $e_{b}>e_{a}$，那么可由
    \[
        \bar{e}=\phi_{a}e_{a}+\phi_{b}e_{b},\quad\sigma_{e}^{2}=\phi_{a}(e_{a}-\bar{e})^{2}+\phi_{b}(e_{b}-\bar{e})^{2}
    \]
    求得 $e_{a}=e_{a}(\bar{e},\,\sigma_{e}^{2})$，$e_{b}=e_{b}(\bar{e},\,\sigma_{e}^{2})$，因此，
    \[
        \mathrm{E}[u(\tilde{e})]=\phi_{a}u(e_{a})+\phi_{b}u(e_{b})=v(\bar{e},\,\sigma_{e}^{2})
    \]
    即代表性参与者具有均值一方差偏好，按照正文中的结论，我们可以知道 CAPM 成立，因此可以用 CAPM 为期权定价；
    \item 如果状态数大于 2，我们不能用均值、方差表示出各个状态的禀赋，代表性参与者不一定具有均值方差偏好（如果效用函数是二次函数则具有均值一方差偏好），CAPM 不一定成立。
\end{enumerate}
\end{solution}



% !TeX root = ../../金融经济学_王江_讲义及习题.tex
% 知识点：references/金融经济学/金融经济学-王江_261013.pdf，印刷页 213–223（PDF 页 228–238）
% 习题：references/金融经济学/史上最全最完整的金融经济学答案(王江版).pdf，页 76–79（印刷页 75–78）

\lecture{套利定价理论}

\ifdefined\WJHandoutKnowledge

CAPM 用市场组合的贝塔衡量系统风险。本章先不指定风险因子的经济含义，而假设收益可以由少数共同因子与剩余风险分解；再研究无套利给预期收益施加的限制。\textbf{精确复制}与\textbf{大市场中的近似复制}会得到不同强度的结论。

\begin{tcolorbox}[knowledge,title=14.1 资产收益的因子结构]
设 $N$ 只资产有 $F$ 个共同因子，收益满足线性因子模型
\[
 \widetilde r_n=\bar r_n+\sum_{k=1}^{F}b_{n,k}\widetilde f_k
 +\widetilde\varepsilon_n,\qquad n=1,\ldots,N.
 \tag{14.1}
\]
$\mathrm E[\widetilde f_k]=\mathrm E[\widetilde\varepsilon_n]=0$；$b_{n,k}$ 是资产 $n$ 对因子 $k$ 的\textbf{载荷}，$\widetilde\varepsilon_n$ 是因子无法解释的剩余风险。把因子标准化为 $\operatorname{Cov}(\widetilde f)=I$，并令因子与剩余风险不相关，则向量式为
\[
 \widetilde r=\bar r+B\widetilde f+\widetilde\varepsilon,
 \qquad \operatorname{Cov}(\widetilde r)=BB^{\mathsf T}
 +\operatorname{Cov}(\widetilde\varepsilon).
 \tag{14.2}
\]
共同因子解释跨资产共同波动；剩余风险能否充分分散，取决于它在不同资产间的协方差结构。因子可经可逆线性变换重新表示，所以单个因子的名称和载荷并不天然唯一；定价关系在相应变换后保持一致。
\end{tcolorbox}

\begin{tcolorbox}[knowledge,title=14.2 精确因子模型与无套利]
先假设无剩余风险，只有一个因子：$\widetilde r_n=\bar r_n+b_n\widetilde f$。若有两只资产的载荷不同，适当的多空组合可以使因子载荷为 0，从而得到无风险组合；也可构造载荷为 1 的\textbf{因子组合}。若无套利，所有无风险组合的收益须相同，且拥有同样因子载荷的组合须同价。于是\textbf{定理 14.1（单因子 APT）}：
\[
 \bar r_n=r_F+b_n\lambda,\qquad\forall n,
 \tag{14.5}
\]
其中 $\lambda$ 是单位因子载荷的风险溢价。若某资产不满足上式，买入便宜的一边、卖出由无风险组合和因子组合复制的另一边，可形成零成本、无随机支付差且严格正收益的套利。

精确 $F$ 因子模型写成 $\widetilde r=\bar r+B\widetilde f$。对任何零成本且零因子载荷的组合 $z$，无套利要求 $z^{\mathsf T}\bar r=0$。线性代数因此给出\textbf{定理 14.2}：
\[
 \bar r=\lambda_0\mathbf1+B\lambda,
 \quad\text{即 }\bar r_n=\lambda_0+\sum_{k=1}^{F}b_{n,k}\lambda_k.
 \tag{14.6}
\]
若存在收益线性无关的无风险组合，$\lambda_0=r_F$；因子组合可交易时，$\lambda_k$ 是第 $k$ 个标准化因子组合超出无风险利率的期望收益。因子协方差矩阵为单位阵只是标准化选择，不是经济上必须相互独立。
\end{tcolorbox}

\begin{tcolorbox}[knowledge,title=14.3 剩余风险、极限套利与近似 APT]
若每只资产还带有 $\widetilde\varepsilon_n$，有限资产的因子组合通常\textbf{不能精确复制}单只资产：消去共同因子后仍留下特殊风险。因此 (14.6) 不再由有限市场的简单复制严格推出。

考察越来越大的市场。\textbf{定义 14.1}\quad 若存在零成本组合序列，使其预期支付持续增大，而支付方差趋于 0，就称存在\textbf{极限套利}。这一概念强于只排除每个有限市场中的常规无套利机会。

\textbf{定理 14.3}\quad 若收益由固定数目 $F$ 的共同因子与有界、可充分分散的剩余风险描述，且不存在极限套利，则预期收益向量相对于某个 $F$ 因子线性定价向量的\textbf{平方偏差总和有界}：
\[
 \sum_{i=1}^{n}
 (\bar r_i^{(n)}-\lambda_0^{(n)}-b_i^{(n)\mathsf T}\lambda^{(n)})^2
 \le A,
\]
其中 $A$ 与资产数 $n$ 无关。证明的要点是把预期收益投影到由常数与因子载荷张成的空间；若正交残差的平方范数无界，可放大一个零成本、零因子载荷的组合，使其期望收益趋于无穷而剩余风险方差趋于零，违反无极限套利。

\textbf{定理 14.4}\quad 上述平方和界意味着：对任意固定 $\epsilon>0$，定价误差超过 $\epsilon$ 的资产数量至多为 $A/\epsilon^2$。故随着资产数增加，定价式对\textbf{大多数资产}近似成立；不能说每一只资产都被精确定价。

\textbf{边界例子。}即使每个有限市场没有常规套利，若越来越多独立特殊风险资产的预期超额收益保持正值，组合可分散其风险却保留收益，从而形成极限套利。无极限套利不是普通无套利的自动推论，也不是市场均衡的完整条件。
\end{tcolorbox}

\begin{tcolorbox}[knowledge,title=14.5 APT 作为均衡结果]
前述 APT 主要由支付结构和无套利出发。要说明该关系是\textbf{均衡}价格，还需偏好、禀赋、可交易性和出清等条件。\textbf{定理 14.5}给出一组充分条件：证券支付具有因子结构；无风险支付、各共同因子及参与者的禀赋可由证券支付线性表出；参与者严格风险厌恶；均衡配置为内点。在这些条件下，均衡中的个体消费可在因子支付空间内表示，个体边际效用定价条件给每个因子组合赋予溢价 $\lambda_k$，从而推出 (14.6)。

这一结果说明精确 APT 可与均衡相容，但其条件比“因子模型加无套利”更强。与 CAPM 相比，APT 容许多个系统因子，却不直接告诉我们因子的经济身份或风险溢价数值；这些仍需模型设定或数据识别。
\end{tcolorbox}

\begin{tcolorbox}[summary,title=本章小结]
精确因子模型中，因子支付可复制，常规无套利给出精确的线性预期收益关系。含剩余风险的大市场中，无极限套利只保证多数资产的近似定价，可能存在少数偏差较大的资产。若还要把 APT 解释成均衡结果，必须补充偏好、可交易支付与市场出清的条件。
\end{tcolorbox}

\fi

% !TeX root = ../../金融经济学_王江_讲义及习题.tex
% 14_APT 课后习题（由章节文件拆分）

\begin{problem}
\problemno{14.1} 市场上有无穷多的证券。它们的收益率由单因子模型给定：
\[
    \tilde{r}_{n}=\bar{r}_{n}+b_{n}\tilde{f}+\tilde{u}_{n},\quad n=1,2,\ldots
\]
这里 $\mathrm{E}[\tilde{f}]=0$；$\forall\,n,n'$，$\mathrm{E}[\tilde{u}_{n}]=0$，$\mathrm{E}[\tilde{f}\tilde{u}_{n}]=0$；$\operatorname{Var}[\tilde{u}_{n}]\le\sigma^{2}$；以及 $\mathrm{E}[\tilde{u}_{n}\tilde{u}_{n'}]=0$。一个组合由 $\{w_{1},w_{2},\ldots\}$ 给定。如当 $N$ 很大时，
\[
    \sum_{n}^{N}w_{n}^{2}\simeq O(N^{-1})
\]
则称它是充分分散的，其中 $N$ 代表组合中绝对权重最大的股票数目而 $O(\cdot)$ 表示同阶或更高阶。
\begin{enumerate}[label=(\alph*)]
    \item 计算充分分散组合的方差。
    \item 证明它不受单个证券的特殊风险的影响。
    \item 更一般地讨论什么因素决定了充分分散组合的风险，这里风险的度量是收益率的方差。
\end{enumerate}
\end{problem}
\begin{solution}
\begin{enumerate}[label=(\alph*)]
    \item 组合收益率为
    \[
        \tilde{r}_{p}=\sum w_{n}\tilde{r}_{n}=\sum w_{n}\bar{r}_{n}+\Bigl(\sum w_{n}b_{n}\Bigr)\tilde{f}+\sum w_{n}\tilde{u}_{n}
    \]
    组合风险为
    \[
        \operatorname{Var}(\tilde{r}_{p})=\Bigl(\sum w_{n}b_{n}\Bigr)^{2}\sigma_{f}^{2}+\sum w_{n}^{2}\operatorname{Var}(\tilde{u}_{n})
    \]
    因为组合是充分分散的，故
    \[
        \sum w_{n}^{2}\operatorname{Var}(\tilde{u}_{n})\le\sigma^{2}\sum w_{n}^{2}\to 0,\quad N\to\infty
    \]
    因此
    \[
        \operatorname{Var}(\tilde{r}_{p})=\Bigl(\sum w_{n}b_{n}\Bigr)^{2}\sigma_{f}^{2}
    \]
    \item $\operatorname{Var}(\tilde{r}_{p})$ 中没有 $\operatorname{Var}(\tilde{u}_{n})$ 项，即组合风险不受特殊风险的影响；
    \item 影响充分分散组合的风险有两个：一是因子的风险，另一个是组合的载荷。
\end{enumerate}
\end{solution}

\begin{problem}
\problemno{14.2} 假设资产 $A$、$B$、$C$ 的收益服从精确两因子模型：
\begin{align}
    r_{A}&=a_{A}+b_{A,1}\tilde{f}_{1}+b_{A,2}\tilde{f}_{2}\tag{14.1}\\
    r_{B}&=a_{B}+b_{B,1}\tilde{f}_{1}+b_{B,2}\tilde{f}_{2}\tag{14.2}\\
    r_{C}&=a_{C}+b_{C,1}\tilde{f}_{1}+b_{C,2}\tilde{f}_{2}\tag{14.3}
\end{align}
其中 $\tilde{f}_{1}$ 和 $\tilde{f}_{2}$ 是两个均值为 0、方差为 1 的因子。$r_{F}$ 是无风险利率。为消除套利机会，$a$、$b$ 以及 $r_{F}$ 应该满足什么条件？
\end{problem}
\begin{solution}
构造一个组合，组合权重为 $\begin{bmatrix}x_{1} \\ x_{2} \\ 1-x_{1}-x_{2}\end{bmatrix}$，那么组合的收益率为
\[
    \begin{aligned}
        \tilde{r}_{p}&=\bigl[(a_{A}-a_{C})x_{1}+(a_{B}-a_{C})x_{2}\bigr]+\bigl[(b_{A,1}-b_{C,1})x_{1}+(b_{B,1}-b_{C,1})x_{2}\bigr]\tilde{f}_{1}\\
        &\quad+\bigl[(b_{A,2}-b_{C,2})x_{1}+(b_{B,2}-b_{C,2})x_{2}\bigr]\tilde{f}_{2}
    \end{aligned}
\]
如果不存在套利机会，那么
\[
    \begin{aligned}
        (b_{A,1}-b_{C,1})x_{1}+(b_{B,1}-b_{C,1})x_{2}&=0\\
        (b_{A,2}-b_{C,2})x_{1}+(b_{B,2}-b_{C,2})x_{2}&=0\\
        (a_{A}-a_{C})x_{1}+(a_{B}-a_{C})x_{2}&=r_{F}
    \end{aligned}
\]
也就是说，若 $x\in\{y\in\mathrm{R}^{2}:y\perp\begin{bmatrix}b_{A,1}-b_{C,1} \\ b_{B,1}-b_{C,1}\end{bmatrix},\,y\perp\begin{bmatrix}b_{A,2}-b_{C,2} \\ b_{B,2}-b_{C,2}\end{bmatrix}\}$，那么 $x^{\mathrm{T}}\begin{bmatrix}a_{A}-a_{C} \\ a_{B}-a_{C}\end{bmatrix}=r_{F}$。
\end{solution}

\begin{problem}
\problemno{14.3} 考虑一个有无穷多资产和 $K$ 个因子的因子模型：
\[
    \tilde{r}_{n}=\bar{r}_{n}+b_{n}\tilde{f}+\tilde{u}_{n},\quad n=1,2,\ldots
\]
其中 $\tilde{f}=\begin{bmatrix}\tilde{f}_{1} \\ \ldots \\ \tilde{f}_{K}\end{bmatrix}$ 是因子向量，$b_{n}=\begin{bmatrix}b_{n,1} \\ \ldots \\ b_{n,K}\end{bmatrix}$ 资产 $n$ 的因子载荷向量，而 $\forall\,n,n'$，$\mathrm{E}[\tilde{f}]=0$，$\mathrm{E}[\tilde{u}_{n}]=0$，$\mathrm{E}[\tilde{f}\tilde{f}^{\mathrm{T}}]=I$，$\mathrm{E}[\tilde{f}\tilde{u}_{n}]=0$，$\mathrm{E}[\tilde{u}_{n}\tilde{u}_{n'}]=0$，其中 $I$ 是单位阵。考虑另外一组因子 $\tilde{f}'$，它们是 $\tilde{f}$ 的线性变换：
\[
    \tilde{f}'=\mathrm{A}\tilde{f}
\]
这里 $\mathrm{A}$ 是一个 $(K\times K)$ 的矩阵。
\begin{enumerate}[label=(\alph*)]
    \item 证明存在 $\tilde{f}'$ 使得
    \[
        \tilde{r}_{n}=\bar{r}_{n}+b_{n}\tilde{f}'+\tilde{u}_{n},\quad n=1,2,\ldots
    \]
    其中 $\mathrm{E}[\tilde{f}']=0$，$\mathrm{E}[\tilde{f}'\tilde{f}'^{\mathrm{T}}]=I$，$\mathrm{E}[\tilde{f}'\tilde{u}_{n}]=0$，$\forall\,n$。因此，因子模型不是唯一的。
    \item 证明因子模型有许多因子不同但却等价的表述形式。
\end{enumerate}
\end{problem}
\begin{solution}
\begin{enumerate}[label=(\alph*)]
    \item 若 $A$ 为正交阵（即 $A^{\mathrm{T}}A=I$），那么
    \[
        \begin{aligned}
            \tilde{r}_{n}&=\bar{r}_{n}+b_{n}^{\mathrm{T}}\tilde{f}+\tilde{u}_{n}\\
            &=\bar{r}_{n}+b_{n}^{\mathrm{T}}A^{\mathrm{T}}A\tilde{f}+\tilde{u}_{n}\\
            &=\bar{r}_{n}+(Ab_{n})^{\mathrm{T}}(A\tilde{f})+\tilde{u}_{n}\\
            &=\bar{r}_{n}+b_{n}'^{\mathrm{T}}\tilde{f}'+\tilde{u}_{n}
        \end{aligned}
    \]
    可以验证：
    \[
        \mathrm{E}[\tilde{f}']=0,\quad\mathrm{E}[\tilde{f}'\tilde{f}'^{\mathrm{T}}]=I,\quad\mathrm{E}[\tilde{f}'\tilde{u}_{n}]=0
    \]
    \item 因为 $K$ 阶正交阵有很多，因而等价的表述形式也有很多。
\end{enumerate}
\end{solution}

\begin{problem}
\problemno{14.4*} 给定练习 14.3 描述的因子模型，APT 要求
\[
    \bar{r}_{n}-r_{F}=\lambda^{\mathrm{T}}b_{n}
\]
其中 $\lambda=\begin{bmatrix}\lambda_{1} \\ \lambda_{2} \\ \ldots \\ \lambda_{K}\end{bmatrix}$ 是因子风险溢价向量。证明存在一个等价的因子模型，它的因子向量为 $\tilde{f}^{*}$ 而对应的因子风险溢价向量为 $\lambda^{*}=\begin{bmatrix}\lambda_{1}^{*} \\ 0 \\ \ldots \\ 0\end{bmatrix}$，即只有第一个因子有风险溢价而其它因子虽然代表了不可分散风险，却没有风险溢价。
\end{problem}
\begin{solution}
考虑这样的正交阵：$A=\bigl[\frac{1}{\sqrt{\lambda^{\mathrm{T}}\lambda}}\lambda,\,U\bigr]$，其中
\[
    U=[U_{1},\,U_{2},\,\cdots,\,U_{K-1}],\quad U_{i}^{\mathrm{T}}U_{j}=\delta_{ij},\quad U_{i}^{\mathrm{T}}\lambda=0
\]
\[
    \delta_{ij}=\begin{cases}
        1&i=j\\
        0&i\neq j
    \end{cases}
\]
可以验证：$AA^{\mathrm{T}}=I$，即 $A$ 是正交阵，且 $A^{\mathrm{T}}\lambda=\begin{bmatrix}\sqrt{\lambda^{\mathrm{T}}\lambda} \\ \,0 \\ \,0 \\ \,\cdots \\ \,0\end{bmatrix}$，那么
\[
    \begin{aligned}
        \tilde{r}_{n}&=\bar{r}_{n}+b_{n}^{\mathrm{T}}AA^{\mathrm{T}}\tilde{f}+\tilde{u}_{n}\\
        &=\bar{r}_{n}+(A^{\mathrm{T}}b_{n})^{\mathrm{T}}(A^{\mathrm{T}}\tilde{f})+\tilde{u}_{n}\\
        &=\bar{r}_{n}+b_{n}^{*\,\mathrm{T}}\tilde{f}^{*}+\tilde{u}_{n}
    \end{aligned}
\]
是一个等价表述。并且 $\bar{r}_{n}-r_{F}=\lambda^{*\,\mathrm{T}}b_{n}^{*}$，其中，$\lambda^{*}=A^{\mathrm{T}}\lambda=\begin{bmatrix}\sqrt{\lambda^{\mathrm{T}}\lambda} \\ \,0 \\ \,0 \\ \,\cdots \\ \,0\end{bmatrix}$。
\end{solution}



% !TeX root = ../../金融经济学_王江_讲义及习题.tex
% 源 PDF：references/金融经济学/史上最全最完整的金融经济学答案(王江版).pdf，页 80–86（印刷页 79–85）

\lecture{完全市场下的公司财务}

\ifdefined\WJHandoutKnowledge
\fi

% !TeX root = ../../金融经济学_王江_讲义及习题.tex
% 15_完全市场下的公司财务 课后习题（由章节文件拆分）

\begin{problem}
\problemno{15.1*} 考虑一个具有两个时期 0 和 1 的经济。在 1 期有两个可能状态 $a$ 和 $b$，发生概率分别为 $\pi_{a}$ 和 $\pi_{b}$。存在生产技术如下：如果在 0 期投入 $I$，那么在 1 期的产出将是
\[
    Y=\begin{cases}
        \lambda_{a}I^{\alpha},&\text{状态 }a\\
        \lambda_{b}I^{\alpha},&\text{状态 }b
    \end{cases}
\]
其中 $0<\lambda_{b}<\lambda_{a}<1$ 且 $0<\alpha<1$。假设参与者 1 拥有这项技术但却没有任何其它资产。他具有如下形式的效用函数：
\[
    \log(c_{0})+\pi_{a}\log c_{1a}+\pi_{b}\log c_{1b}.
\]
证券市场上只存在无风险债券，它的利率为 $r_{F}$。
\begin{enumerate}[label=(\alph*)]
    \item 写出参与者 1 的优化问题。
    \item 求解他的最优生产、消费和投资选择。
    \item 说明他的生产决策如何依赖于利率。并给出解释。
    \item 计算在最优选择下他的期望效用。
\end{enumerate}
\end{problem}
\begin{solution}
\begin{enumerate}[label=(\alph*)]
    \item 假设参与者的 0 期消费为 $c_{0}$、投资为 $I$，那么参与者的优化问题是
    \[
        \begin{aligned}
            \max_{c_{0},\,I}\;&\log c_{0}+\pi_{a}\log(c_{a})+\pi_{b}\log(c_{b})\qquad(*)\\
            \text{s.t.}\;&c_{a}=\lambda_{a}I^{\alpha}-(c_{0}+I)(1+r_{F})\\
            &c_{b}=\lambda_{b}I^{\alpha}-(c_{0}+I)(1+r_{F})
        \end{aligned}
    \]
    \item 上述优化问题的一阶条件是
    \[
        \begin{aligned}
            \frac{1}{c_{0}}+\pi_{a}\frac{-(1+r_{F})}{\lambda_{a}I^{\alpha}-(c_{0}+I)(1+r_{F})}+\pi_{b}\frac{-(1+r_{F})}{\lambda_{b}I^{\alpha}-(c_{0}+I)(1+r_{F})}&=0\\
            \pi_{a}\frac{\lambda_{a}I^{\alpha-1}\alpha-(1+r_{F})}{\lambda_{a}I^{\alpha}-(c_{0}+I)(1+r_{F})}+\pi_{b}\frac{\lambda_{b}I^{\alpha-1}\alpha-(1+r_{F})}{\lambda_{b}I^{\alpha}-(c_{0}+I)(1+r_{F})}&=0
        \end{aligned}
    \]
    第一式是消费 Euler：$1/c_{0}=(1+r_{F})E[1/\tilde{c}_{1}]$。第二式说：用 1 期边际效用加权后，边际产出 $\alpha\lambda_{\omega}I^{\alpha-1}$ 的期望等于 $1+r_{F}$。两式联立决定最优 $c_{0}^{*}$、$I^{*}$（一般无更简的显式解）。另外参与者须在 0 期卖空 $(c_{0}^{*}+I^{*})$ 单位的无风险债券；
    \item 由生产一阶条件，边际效用加权的 $E[\mathrm{MPK}]$ 必须等于 $1+r_{F}$。$\alpha<1$ 时 MPK 随 $I$ 递减，故 $r_{F}$ 上升迫使 $I^{*}$ 下降：借款为生产融资的成本提高。
    \item 把求得的 $c_{0}^{*}$，$I^{*}$ 代入 $(*)$ 式即可求得期望效用，记作 $EU_{1}$。
\end{enumerate}
\end{solution}

\begin{problem}
\problemno{15.2*} 在练习 15.1 描述的经济中，现在假设证券市场还包括一只证券，它在状态 $a$ 支付 $\lambda_{a}$、在状态 $b$ 支付 $\lambda_{b}$ 而它的价格为 $S$。
\begin{enumerate}[label=(\alph*)]
    \item 证明现在证券市场是完全的。
    \item 用 $r_{F}$ 和 $S$ 表示状态价格。
    \item 描述参与者 1 的优化问题。
    \item 计算他的总财富。证明它总是正的。
    \item 求解他的最优生产、消费和投资选择。
    \item 用 NPV 法则解释他的投资选择。
    \item 他的生产技术的 NPV 是多少？这个如何随着状态价格的改变而改变？解释所得到的结论。
    \item 计算引入风险证券参与者 1 得到的期望效用的增加值，用 0 期财富的确定性等价表示。
\end{enumerate}
\end{problem}
\begin{solution}
\begin{enumerate}[label=(\alph*)]
    \item 两个状态，两只支付独立的证券，因而证券市场是完全的；
    \item 由
    \[
        \begin{aligned}
            \frac{1}{1+r_{F}}&=\phi_{a}+\phi_{b}\\
            S&=\phi_{a}\lambda_{a}+\phi_{b}\lambda_{b}
        \end{aligned}
    \]
    可以求得状态价格为
    \[
        \begin{aligned}
            \phi_{a}&=\frac{S(1+r_{F})-\lambda_{b}}{(\lambda_{a}-\lambda_{b})(1+r_{F})}\\
            \phi_{b}&=\frac{\lambda_{a}-S(1+r_{F})}{(\lambda_{a}-\lambda_{b})(1+r_{F})}
        \end{aligned}
    \]
    \item 假设参与者的 0 期消费为 $c_{0}$、股票头寸为 $\theta_{S}$、生产投入为 $I$，那么他的优化问题是
    \[
        \begin{aligned}
            \max_{c_{0},\,\theta_{S},\,I}\;&\log c_{0}+\pi_{a}\log(c_{a})+\pi_{b}\log(c_{b})\\
            \text{s.t.}\;&c_{a}=\lambda_{a}I^{\alpha}+\theta_{S}\lambda_{a}-(\theta_{S}S+c_{0}+I)(1+r_{F})\\
            &c_{b}=\lambda_{b}I^{\alpha}+\theta_{S}\lambda_{b}-(\theta_{S}S+c_{0}+I)(1+r_{F})
        \end{aligned}
    \]
    \item 参与者的总财富为
    \[
        \begin{aligned}
            w&=c_{0}+\phi_{a}\bigl(\lambda_{a}I^{\alpha}+\theta_{S}\lambda_{a}-(\theta_{S}S+c_{0}+I)(1+r_{F})\bigr)\\
            &\quad+\phi_{b}\bigl(\lambda_{b}I^{\alpha}+\theta_{S}\lambda_{b}-(\theta_{S}S+c_{0}+I)(1+r_{F})\bigr)\\
            &=SI^{\alpha}-I
        \end{aligned}
    \]
    因为当 $I=0$ 时，$w=0$，因而在参与者选择最优生产投入时，$w>0$；
    \item     完全市场下生产与消费分离。把财富约束写成 $c_{0}+\phi_{a}c_{a}+\phi_{b}c_{b}=SI^{\alpha}-I$。拉格朗日函数为
    \[
        \mathcal{L}=\log c_{0}+\pi_{a}\log c_{a}+\pi_{b}\log c_{b}+\lambda\bigl(SI^{\alpha}-I-c_{0}-\phi_{a}c_{a}-\phi_{b}c_{b}\bigr)
    \]
    对 $I$ 的一阶条件即 NPV 法则
    \[
        \alpha SI^{\alpha-1}=1\implies I^{*}=(\alpha S)^{\frac{1}{1-\alpha}}
    \]
    消费侧一阶条件：$1/c_{0}=\lambda$，$\pi_{\omega}/c_{\omega}=\lambda\phi_{\omega}$。代入预算并注意 $\pi_{a}+\pi_{b}=1$，得 $c_{0}^{*}=w^{*}/2$，$c_{\omega}^{*}=\pi_{\omega}w^{*}/(2\phi_{\omega})$，其中 $w^{*}=SI^{*\alpha}-I^{*}$。再用 $c_{a}$、$c_{b}$ 反解股票头寸，得到
    \[
        I^{*}=(\alpha S)^{\frac{1}{1-\alpha}},\quad w^{*}=SI^{*\alpha}-I^{*},\quad c_{0}^{*}=\tfrac12 w^{*},\quad\theta_{S}^{*}=\frac{\frac{\pi_{a}}{2\phi_{a}}-\frac{\pi_{b}}{2\phi_{b}}}{\lambda_{a}-\lambda_{b}}w^{*}-I^{*\alpha}
    \]
    另外，参与者须在 0 期卖空 $\theta_{S}^{*}S+c_{0}^{*}+I^{*}$ 单位的无风险债券；
    \item
    \[
        NPV=(\phi_{a}\lambda_{a}I^{\alpha}+\phi_{b}\lambda_{b}I^{\alpha})-I=SI^{\alpha}-I
    \]
    最优的生产投入应是 $(NPV)'=0$，即 $I=(\alpha S)^{\frac{1}{1-\alpha}}$，恰是上面所求得的最优生产投入；
    \item 代入计算得
    \[
        \begin{aligned}
            NPV^{*}&=S(\alpha S)^{\frac{\alpha}{1-\alpha}}-(\alpha S)^{\frac{1}{1-\alpha}}\\
            &=(\alpha^{\frac{\alpha}{1-\alpha}}-\alpha^{\frac{1}{1-\alpha}})(\phi_{a}\lambda_{a}+\phi_{b}\lambda_{b})^{\frac{1}{1-\alpha}}
        \end{aligned}
    \]
    由此可以看到，NPV 会随着状态价格的上升而增加：一方面是在相同生产投入条件下产出的净现值增加；另外，根据 NPV 法则，参与者会增加生产投入，这也会使得 NPV 增加；
    \item 在最优选择下，我们可计算得期望效用为 $EU_{2}$。在第 1 题中，0 期禀赋为 0 时计算得到的期望效用是 $EU_{1}$，我们可以通过 $EU_{1}(\Delta w)=EU_{2}$ 求效用增加值的确定性等价 $\Delta w$。
\end{enumerate}
\end{solution}

\begin{problem}
\problemno{15.3**} 考虑练习 15.1 描述的经济。现在假设经济中除了那里描述的参与者 1 以外还有一个参与者 2。参与者 2 没有生产技术，但经济赋予了他 $e_{0}$ 单位的 0 期消费。他的效用函数和参与者 1 的一样。证券市场和练习 15.2 中的一样。
\begin{enumerate}[label=(\alph*)]
    \item 求解两个参与者在给定状态价格下的最优生产、消费和投资选择。
    \item 求解均衡状态价格、总消费和消费配置。
    \item 证明均衡配置是 Pareto 最优的。
\end{enumerate}
\end{problem}
\begin{solution}
\begin{enumerate}[label=(\alph*)]
    \item 在第 2 题中已经求解了参与者 1 的最优化问题，现在考虑参与者 2 的最优化问题（假设他持有的无风险债券、股票的头寸分别为 $\theta_{B}$，$\theta_{S}$）
    \[
        \begin{aligned}
            \max_{\theta_{B},\,\theta_{S}}\;&\log c_{0}+\pi_{a}\log(c_{a})+\pi_{b}\log(c_{b})\\
            \text{s.t.}\;&c_{0}=e_{0}-\theta_{B}-S\theta_{S}\\
            &c_{a}=\theta_{B}(1+r_{F})+\lambda_{a}\theta_{S}\\
            &c_{b}=\theta_{B}(1+r_{F})+\lambda_{b}\theta_{S}
        \end{aligned}
    \]
    完全市场下与练习 15.2 同一拉格朗日，Euler 给出 $c_{0}^{*}=e_{0}/2$，$c_{\omega}^{*}=\pi_{\omega}e_{0}/(2\phi_{\omega})$。由 1 期支付反解头寸：
    \[
        \theta_{S}^{*}=\frac{c_{a}^{*}-c_{b}^{*}}{\lambda_{a}-\lambda_{b}}=\frac{1}{\lambda_{a}-\lambda_{b}}\Bigl(\frac{\pi_{a}}{2\phi_{a}}-\frac{\pi_{b}}{2\phi_{b}}\Bigr)e_{0},\qquad
        \theta_{B}^{*}=\frac{c_{a}^{*}-\lambda_{a}\theta_{S}^{*}}{1+r_{F}}
    \]
    \item 市场出清条件是
    \[
        \begin{aligned}
            \theta_{1,S}^{*}S+c_{1,0}^{*}+I^{*}&=\theta_{2,b}^{*}\\
            \theta_{1,S}^{*}+\theta_{2,S}^{*}&=0
        \end{aligned}
    \]
    再加上 $\frac{1}{1+r_{F}}=\phi_{a}+\phi_{b}$，$S=\phi_{a}\lambda_{a}+\phi_{b}\lambda_{b}$ 可以求得均衡状态价格为
    \[
        \begin{aligned}
            \phi_{a}&=\frac{\pi_{a}e_{0}^{1-\alpha}}{\lambda_{a}\alpha^{\alpha}(1+\alpha)^{1-\alpha}}\\
            \phi_{b}&=\frac{\pi_{b}e_{0}^{1-\alpha}}{\lambda_{b}\alpha^{\alpha}(1+\alpha)^{1-\alpha}}
        \end{aligned}
    \]
    最优的生产投入为 $I^{*}=\frac{\alpha e_{0}}{1+\alpha}$，最优消费为
    \[
        \begin{aligned}
            c_{1,0}^{*}&=\frac{1-\alpha}{2(1+\alpha)}e_{0},\quad c_{2,0}^{*}=\tfrac12 e_{0},\quad C_{0}^{*}=\frac{1}{1+\alpha}e_{0}\\
            c_{1,a}^{*}&=\frac{\pi_{a}}{2\phi_{a}}\frac{1-\alpha}{1+\alpha}e_{0},\quad c_{2,a}^{*}=\frac{\pi_{a}}{2\phi_{a}}e_{0},\quad C_{a}^{*}=\lambda_{a}\Bigl(\frac{\alpha e_{0}}{1+\alpha}\Bigr)^{\alpha}\\
            c_{1,b}^{*}&=\frac{\pi_{b}}{2\phi_{b}}\frac{1-\alpha}{1+\alpha}e_{0},\quad c_{2,b}^{*}=\frac{\pi_{b}}{2\phi_{b}}e_{0},\quad C_{b}^{*}=\lambda_{b}\Bigl(\frac{\alpha e_{0}}{1+\alpha}\Bigr)^{\alpha}
        \end{aligned}
    \]
    \item 考虑两个参与者的边际效用之比：
    \[
        \begin{aligned}
            \frac{\pi_{a}c_{1,0}^{*}}{c_{1,a}^{*}}&=\phi_{a}=\frac{\pi_{a}c_{2,0}^{*}}{c_{2,a}^{*}}\\
            \frac{\pi_{b}c_{1,0}^{*}}{c_{1,b}^{*}}&=\phi_{b}=\frac{\pi_{b}c_{2,0}^{*}}{c_{2,b}^{*}}
        \end{aligned}
    \]
    因而均衡配置是 Pareto 最优的。
\end{enumerate}
\end{solution}

\begin{problem}
\problemno{15.4**} 考虑练习 15.3 中的经济。
\begin{enumerate}[label=(\alph*)]
    \item 构建代表性参与者。
    \item 通过求解代表性参与者的优化问题而得到均衡配置。
    \item 用代表性参与者的边际效用计算状态价格。
    \item 讨论 $\lambda_{b}$、$\lambda_{a}$ 以及 $\alpha$ 如何影响状态价格和均衡配置。
\end{enumerate}
\end{problem}
\begin{solution}
\begin{enumerate}[label=(\alph*)]
    \item 我们知道，对于对数形式的效用函数，代表性参与者的效用函数也是对数形式的，且与财富在参与者之间的分布无关；
    \item 对于代表性参与者，他的优化问题是
    \[
        \max_{I}\ \log(e_{0}-I)+\pi_{a}\log(\lambda_{a}I^{\alpha})+\pi_{b}\log(\lambda_{b}I^{\alpha})
    \]
    目标函数等于 $\log(e_{0}-I)+\alpha\log I+$ 常数。一阶条件
    \[
        -\frac{1}{e_{0}-I}+\frac{\alpha}{I}=0
    \]
    即 $I=\alpha(e_{0}-I)$，故 $I^{*}=\dfrac{\alpha e_{0}}{1+\alpha}$，$c_{0}^{*}=\dfrac{e_{0}}{1+\alpha}$。1 期总消费为 $C_{\omega}^{*}=\lambda_{\omega}(I^{*})^{\alpha}$；
    \item 状态价格为
    \[
        \begin{aligned}
            \phi_{a}&=\frac{\pi_{a}C_{0}}{C_{a}}=\pi_{a}\frac{\frac{e_{0}}{1+\alpha}}{\lambda_{a}\bigl(\frac{\alpha e_{0}}{1+\alpha}\bigr)^{\alpha}}=\frac{\pi_{a}e_{0}^{1-\alpha}}{\lambda_{a}\alpha^{\alpha}(1+\alpha)^{1-\alpha}}\\
            \phi_{b}&=\frac{\pi_{b}e_{0}^{1-\alpha}}{\lambda_{b}\alpha^{\alpha}(1+\alpha)^{1-\alpha}}
        \end{aligned}
    \]
    \item 均衡配置只受到 $\alpha$ 的影响，$\alpha$ 增大，相同生产投入的产出增加，因而参与者会减少当前消费，增加生产投入；当 $\alpha$ 增大时，状态价格均变小，而当 $\lambda_{a}$，$\lambda_{b}$ 增加时，相应状态的状态价格变小，因而 $\alpha$，$\lambda_{a}$，$\lambda_{b}$ 增加会增加产出，因而状态价格变小。
\end{enumerate}
\end{solution}

\begin{problem}
\problemno{15.5**} 与练习 15.2 中的经济一样。在这里我们不假设有一组给定的可交易证券，而是假设拥有生产技术的参与者 1 可以成立一个公司、而公司可以到市场上发行债券和股票，但没有其它可交易的证券。一单位（无风险）公司债在 1 期有 1 单位支付，而一股股票在 1 期有 $\hat{\lambda}I^{\alpha}$ 单位支付，在状态 $a$ 和 $b$ 下 $\hat{\lambda}=\lambda_{a}$，$\lambda_{b}$，$I$ 是公司的生产投资。
\begin{enumerate}[label=(\alph*)]
    \item 给定利率 $r_{F}$ 和股价 $S$，求解参与者 1 的优化问题，即最优的投资额、在市场发行的债券和股票的数量以及他的消费。
    \item 在这种情况下，有无最优负债水平？这和 Modigliani--Miller 定理是否一致？解释原因。
    \item 求解市场均衡，特别地，求解均衡利率、股价、生产水平、发行的债券和股票的数量以及两个参与者的消费配置。
\end{enumerate}
\end{problem}
\begin{solution}
\begin{enumerate}[label=(\alph*)]
    \item 假设参与者的投资额为 $I$，在市场上发行的债券和股票分别为 $\theta_{B}$，$\theta_{S}$，那么他的优化问题是
    \[
        \begin{aligned}
            \max_{\theta_{B},\,\theta_{S},\,I}\;&\log c_{0}+\pi_{a}\log(c_{a})+\pi_{b}\log(c_{b})\\
            \text{s.t.}\;&c_{0}=\theta_{B}+S\theta_{S}-I\\
            &c_{a}=\lambda_{a}I^{\alpha}(1-\theta_{S})-\theta_{B}(1+r_{F})\\
            &c_{b}=\lambda_{b}I^{\alpha}(1-\theta_{S})-\theta_{B}(1+r_{F})
        \end{aligned}
    \]
    一阶条件为
    \[
        \begin{aligned}
            \frac{1}{c_{0}}+\pi_{a}\frac{-(1+r_{F})}{c_{a}}+\pi_{b}\frac{-(1+r_{F})}{c_{b}}&=0\\
            \frac{S}{c_{0}}-\frac{\pi_{a}\lambda_{a}I^{\alpha}}{c_{a}}-\frac{\pi_{b}\lambda_{b}I^{\alpha}}{c_{b}}&=0\\
            \frac{1}{c_{0}}-\frac{\pi_{a}\lambda_{a}\alpha I^{\alpha-1}(1-\theta_{S})}{c_{a}}+\frac{\pi_{b}\lambda_{b}\alpha I^{\alpha-1}(1-\theta_{S})}{c_{b}}&=0
        \end{aligned}
    \]
    由上述方程组可以解得最优的 $\theta_{B}^{*}$，$\theta_{S}^{*}$，$I^{*}$。
    \item 给定 $r_{F}$，$S$，参与者的债券和股票的最优发行量是定值，与 MM 定理不一致，这是因为此时的证券市场是不完全的。
    \item 对于参与者 2 来说，他的最优化问题是
    \[
        \begin{aligned}
            \max_{\theta_{B}',\,\theta_{S}'}\;&\log c_{0}+\pi_{a}\log(c_{a})+\pi_{b}\log(c_{b})\\
            \text{s.t.}\;&c_{0}=e_{0}-\theta_{B}'-S\theta_{S}'\\
            &c_{a}=\lambda_{a}I^{\alpha}\theta_{S}'+\theta_{B}'(1+r_{F})\\
            &c_{b}=\lambda_{b}I^{\alpha}\theta_{S}'+\theta_{B}'(1+r_{F})
        \end{aligned}
    \]
    求解出最优的 $\theta_{B}'^{*}$，$\theta_{S}'^{*}$ 以后，再加上市场出清条件
    \[
        \theta_{B}^{*}=\theta_{B}'^{*},\quad\theta_{S}^{*}=\theta_{S}'^{*}
    \]
    就可以求解出均衡利率、均衡股价，从而可以得到均衡的生产水平、债券和股票的发行量、以及消费配置。
\end{enumerate}
\end{solution}

\begin{problem}
\problemno{15.6*} 在 15.5 现在假设市场上存在一只无风险证券可供参与者交易。Modigliani--Miller 定理成立，即公司的融资决策不影响它的投资决策和均衡配置。解释这种情形和上述情形之间的差别。
\end{problem}
\begin{solution}
如果市场上有无风险证券可以交易，假设参与者 1 购买的债券数量为 $\theta_{B}'$，那么他的优化问题就变成了
\[
    \begin{aligned}
        \max_{\theta_{B},\,\theta_{B}',\,\theta_{S},\,I}\;&\log c_{0}+\pi_{a}\log(c_{a})+\pi_{b}\log(c_{b})\\
        \text{s.t.}\;&c_{0}=\theta_{B}+\theta_{B}'+S\theta_{S}-I\\
        &c_{a}=\lambda_{a}I^{\alpha}(1-\theta_{S})-(\theta_{B}+\theta_{B}')(1+r_{F})\\
        &c_{b}=\lambda_{b}I^{\alpha}(1-\theta_{S})-(\theta_{B}+\theta_{B}')(1+r_{F})
    \end{aligned}
\]
这时我们只能得到 $\theta_{B}+\theta_{B}'$ 的值，而不能确定债券的发行量，因而公司债券发行量的大小（融资决策）不影响公司的投资决策和均衡配置。这是因为公司除了发行债券外，还可以到市场上买卖无风险债券，可以通过调整无风险债券的头寸来补充公司债券发行量的增减，因此债券的发行量不再影响到公司的投资决策以及消费配置，
\end{solution}




\end{document}
